% Generated by build_chapter.py from the GY1965 files. Do not edit by hand; edit the files and rebuild.
\chapter{The academic year 1965--66}\label{ch:1965}
\noindent\textbf{Scope and editorial position.}
This chapter covers all 37 catalogue records, 53--89, from September 1965 to
March 1966, in the order of the handwriting. Its 298 mathematical
occurrences include repeated explanations and short interventions. The plan
in record 57, the year heading in record 78 and Kushnirenko's title-only
record 79 are preserved without invented mathematical content.
The supplied 52-page bundle begins with two carry-over pages,
\GSpage{P000088--P000089}, already represented by 1964-50.11--24;
Kohn's new-year talk starts on \GSpage{P000090}. Actual page boundaries
therefore replace the catalogue's overlapping attachment ranges.

\paragraph{The main mathematical lines.}
The most sustained autumn sequence develops the Lefschetz formula for
elliptic complexes: symbols, Hodge decomposition, positive-eigenspace
cancellation, regularized traces and the local fixed-point contribution.
The 1964 Harvard notes of Atiyah and Bott already contained this programme;
the later journal publication does not make Gelfand's exposition a new
1965 theorem \citep{atiyahbott1964notes,atiyahbott1967lefschetz}.
Kazhdan's parallel sequence supplies affine algebraic geometry, coherent
sheaves and derived functors. Gelfand's characteristic-form discussions
interact with Arnold's accounts of singularities, normal forms, canonical
transformations and ideal-fluid motion.

The spring records bring several more direct publication matches.
Gelfand--Kirillov's 1966 paper contains the Weyl-field and growth-dimension
results of record 80. The affine-plane transform constructions of record 82
match the 1967 Gelfand--Graev--Shapiro paper. Lidskii's all-zero zeta
continuation and Gindikin's multiple-time parabolic construction likewise
match precise 1967 results \citep{gelfandkirillov1966corps,gelfand1967planes,
lidskii1967zeros,gindikin1967multitime}. Agranovich's calculus had already
appeared in his 1965 survey \citep{agranovich1965elliptic}. These different
publication histories are kept separate from literal correctness and from
claims of historical priority.

\begin{center}
\begin{tabular}{@{}clrr@{}}
\toprule Class & Meaning & Occurrences & Share\\\midrule
A & Published soon after by the speaker/Gelfand & 22 & 7.4\%\\
B & Developed by others & 7 & 2.3\%\\
C & Partly developed & 1 & 0.3\%\\
D & Not developed & 0 & 0.0\%\\
E & Already known before the talk & 246 & 82.6\%\\
F & Not identifiable from the notes & 22 & 7.4\%\\
\bottomrule
\end{tabular}
\end{center}

\paragraph{What the critical reading changes.}
The apparatus separates a faithful transcription of a mathematical mistake
from a mistaken transcription. Examples of the former include the omitted
logarithm in growth dimension and the two extraneous factors in the top's
cyclic momenta. A direct counterexample shows why the torus intersection
bound needs transversality: a Hamiltonian graph on \(T^2\) can have three
intersections, so the unqualified \(2^2\) claim fails. An ellipsoid likewise
disproves the fixed-direction interpretation of the slice-volume formula.
An actual reading correction is decisive for Gindikin: the scan says
\ru{гиперб. оператор}, a hyperbolic operator, where the supplied text says
pseudo-operator. The original paper then provides an exact restricted
match. Restored hypotheses are labelled as reconstructions throughout.

\paragraph{Questions and a bounded development.}
The problem index follows 9 definite questions, including resolved ones.
There is one class-C occurrence: the request to find all automorphisms of
the Weyl division ring, GP-1965-08. Its unrestricted form is not
classified in the verified literature; the search did not locate a later
full solution. The development fixes \(D=Q(A_1(\mathbb C))\) and the
rational subfield \(\mathbb C(x)\). It gives complete proofs of its known
stabilizer normal form, an exact pole-and-residue criterion for inner
rational translations, the resulting translation subgroup in the outer
automorphism group, and finite-group cocycle removal. These are explicitly
restricted answers with no novelty claim. Two parameter runs and an
independent written review check the proofs and computations. The remaining
obstruction is moving the chosen rational subfield; no conclusion about
the whole automorphism group follows. There are no class-D items: the
unidentified fragments lack enough information to define research problems.

\paragraph{Continuities with 1964--65.}
The Weyl-rank question links directly to 1964-38.1, and the final
lattice-gas and thermodynamic-limit lines repeat 1964-3.1--2. The elliptic
sequence develops the language of 1964-28.2--5; the derived-functor
construction continues 1964-49.8. Smirnov's positive-definite-function
fragment is related to 1964-13.1, but its missing centrality and topology
prevent identification of the questions. These are links of different
strength, not assertions that every later talk settles an earlier problem.
The supplied previous index contains no problem identifiers, and its
cumulative bibliography is empty. New problem identifiers and checked
bibliography keys are therefore necessary. All present-status assessments
are dated 9 November 2026.

\GStalk{53}{1965-09-20}{J. J. Kohn}{P000090-P000091}{The d-bar Neumann problem and the Levi problem}
Kohn discusses solving \(\bar\partial u=\alpha\), the Levi form, and the passage from a local supporting holomorphic polynomial to a global nonextendible holomorphic function. The second page introduces the Dolbeault complex, the Neumann quadratic form, a boundary estimate, and the Kohn--Rossi extension paper. These are accounts of work already available before this talk, not announcements inferred solely from later publication dates \citep{kohn1963potential,kohn1963harmonic,kohn1964harmonic,kohnrossi1965extension}. The dated heading begins on P000090; P000088--P000089 belong to the preceding invariant-subspace discussion. A short reading programme at the end names algebraic groups and algebraic geometry.

\begin{GSitem}{1965-53.1}{E}{Compatibility and global solvability of the d-bar equation}
\GSfield{Kind; source}{question; \GSpage{P000090}}
\GSfield{The notes}{Given the coefficients \(\alpha_j\), the notes ask when \(u_{\bar z_j}=\alpha_j\) is soluble and how \(u\) depends on them; they write the cross-derivative compatibility conditions.}
\GSfield{Statement}{Let \(\Omega\Subset\mathbb C^n\) have \(C^\infty\) strongly pseudoconvex boundary and let \(\alpha=\sum_j\alpha_j\,d\bar z_j\in C^\infty(\overline\Omega;\Lambda^{0,1})\). If \(\partial\alpha_j/\partial\bar z_k=\partial\alpha_k/\partial\bar z_j\), there is \(u\in C^\infty(\overline\Omega)\) with \(\bar\partial u=\alpha\); the canonical solution is \(u=\bar\partial^*N_1\alpha\). Without a global condition on \(\Omega\), compatibility alone need not give global solvability.}
\GSfield{Reconstruction}{The global strong pseudoconvexity and smoothness assumptions are supplied from the ensuing argument. The notes begin with only a bounded domain; they do not state a general necessary-and-sufficient theorem for arbitrary domains.}
\GSfield{Fate}{Already part of the expository framework in \citet{kohn1963potential}.}
\GSfield{Present status}{As of 9 November 2026, the precise statement is a standard established result or definition.}
\GSfield{Assessment}{literal statement: correct with missing hypotheses; publication match: partial; verification: read the result; development: annotated; confidence: high; problem \hyperref[prob:GP-1965-01]{GP-1965-01}}
\end{GSitem}
\begin{GSitem}{1965-53.2}{E}{The complex tangent space and Levi form}
\GSfield{Kind; source}{definition; \GSpage{P000090}}
\GSfield{The notes}{A defining function \(r\) is negative inside \(M\); the complex tangent vectors satisfy \(\sum r_{z_j}a^j=0\), and the displayed Hermitian form is \(\sum r_{z_i\bar z_j}a^i\bar a^j\).}
\GSfield{Statement}{For a \(C^2\) real defining function with \(dr\ne0\) on \(\partial\Omega\), the complex tangent space at \(p\) is \(\ker\partial r(p)\). Its Levi form is \(\mathcal L_r(p;a)=\sum_{j,k}r_{z_j\bar z_k}(p)a_j\bar a_k\); strong pseudoconvexity means positivity for every nonzero complex tangent vector. Multiplying \(r\) by a positive smooth factor multiplies this restricted form by that factor.}
\GSfield{Reconstruction}{The sign convention depends on \(\Omega=\{r<0\}\); positivity on the full ambient complex tangent space is not required.}
\GSfield{Fate}{Already part of the expository framework in \citet{kohn1963potential}.}
\GSfield{Present status}{As of 9 November 2026, the precise statement is a standard established result or definition.}
\GSfield{Assessment}{literal statement: correct; publication match: exact; verification: read the result; development: annotated; confidence: high}
\end{GSitem}
\begin{GSitem}{1965-53.3}{E}{A local holomorphic support function}
\GSfield{Kind; source}{construction; \GSpage{P000090}}
\GSfield{The notes}{Positivity of the Levi form is used to construct a holomorphic polynomial \(W\) whose zero set touches the boundary at \(p\) and otherwise lies outside the domain near \(p\).}
\GSfield{Statement}{At a strongly pseudoconvex point \(p\) of a \(C^2\) boundary, a holomorphic polynomial \(W\) of degree at most two can be chosen with \(W(p)=0\), \(dW(p)\ne0\), and \(W(z)\ne0\) on \((\overline\Omega\cap U)\setminus\{p\}\), after shrinking \(U\). One obtains \(W\) from the holomorphic part of the second-order Taylor expansion of a suitably modified defining function.}
\GSfield{Reconstruction}{A quadratic change in the defining function can make the real estimate strict in all relevant directions. The conclusion is local; the sketch does not prove global pseudoconvexity.}
\GSfield{Fate}{Already part of the expository framework in \citet{kohn1963potential}.}
\GSfield{Present status}{As of 9 November 2026, the precise statement is a standard established result or definition.}
\GSfield{Assessment}{literal statement: correct with missing hypotheses; publication match: exact; verification: read the result; development: annotated; confidence: high}
\end{GSitem}
\begin{GSitem}{1965-53.4}{E}{Globalizing the local pole}
\GSfield{Kind; source}{method; \GSpage{P000090}}
\GSfield{The notes}{With a cut-off \(g=1\) near \(p\), the notes solve \(\bar\partial u=\bar\partial(g/W)\) and form \(h=u-g/W\). The final sentence attributes noncontinuation to \(u\).}
\GSfield{Statement}{Under the preceding smooth strongly pseudoconvex domain assumptions, choose \(g\) supported in the support-function neighbourhood. The form \(\alpha=\bar\partial(g/W)\), extended by zero, is smooth on \(\overline\Omega\) and \(\bar\partial\)-closed. A smooth solution \(u\) gives a holomorphic \(h=u-g/W\) on \(\Omega\) with a singularity at \(p\), so \(h\) does not extend holomorphically across \(p\).}
\GSfield{Reconstruction}{It is \(h\), not the globally smooth correction \(u\), whose boundary singularity is forced. Smoothness of \(\alpha\) follows because \(\bar\partial g=0\) near the only local zero of \(W\).}
\GSfield{Fate}{Already part of the expository framework in \citet{kohn1963potential}.}
\GSfield{Present status}{As of 9 November 2026, the precise statement is a standard established result or definition.}
\GSfield{Assessment}{literal statement: false as written; publication match: exact; verification: read the result; development: annotated; confidence: high; problem \hyperref[prob:GP-1965-01]{GP-1965-01}}
\end{GSitem}
\begin{GSitem}{1965-53.5}{E}{The Dolbeault complex}
\GSfield{Kind; source}{construction; \GSpage{P000091}}
\GSfield{The notes}{The spaces \(a^{p,q}\) of forms and the arrows \(a^{p,q-1}\to a^{p,q}\to a^{p,q+1}\) are written with differential \(\bar\partial\).}
\GSfield{Statement}{On a complex manifold, \(C^\infty\) differential forms split by bidegree and \(\bar\partial:A^{p,q}\to A^{p,q+1}\) satisfies \(\bar\partial^2=0\). On a domain with boundary the maximal \(L^2\) realizations form a Hilbert complex; domains must be imposed before defining Hilbert adjoints.}
\GSfield{Reconstruction}{The notes alternate \(a\) and \(\alpha\) as space letters. The Hilbert completion and operator domains are made explicit here.}
\GSfield{Fate}{Already part of the expository framework in \citet{kohn1963potential}.}
\GSfield{Present status}{As of 9 November 2026, the precise statement is a standard established result or definition.}
\GSfield{Assessment}{literal statement: correct; publication match: exact; verification: read the result; development: annotated; confidence: high}
\end{GSitem}
\begin{GSitem}{1965-53.6}{E}{The Neumann quadratic form and canonical solution}
\GSfield{Kind; source}{method; \GSpage{P000091}}
\GSfield{The notes}{The quadratic form is \(Q(\varphi,\psi)=(\bar\partial\varphi,\bar\partial\psi)+(\bar\partial^*\varphi,\bar\partial^*\psi)\); a positive higher-order equation is proposed to solve the first-order compatible system.}
\GSfield{Statement}{On \(L^2\) forms, put \(\mathcal D(Q)=\mathcal D(\bar\partial)\cap\mathcal D(\bar\partial^*)\). The closed nonnegative form defines the self-adjoint \(\bar\partial\)-Neumann Laplacian \(\Box=\bar\partial\bar\partial^*+\bar\partial^*\bar\partial\). If its inverse \(N_q\) exists on the orthogonal complement of its harmonic kernel, a closed \((p,q)\)-form orthogonal to that kernel is solved by \(\bar\partial^*N_q\alpha\).}
\GSfield{Reconstruction}{The first-order equation follows using the complex identity and commutation with \(N\), not merely by solving an arbitrary positive second-order equation. The displayed domain in the notebook must also include \(\mathcal D(\bar\partial)\).}
\GSfield{Fate}{Already part of the expository framework in \citet{kohn1963potential}.}
\GSfield{Present status}{As of 9 November 2026, the precise statement is a standard established result or definition.}
\GSfield{Assessment}{literal statement: correct with missing hypotheses; publication match: exact; verification: read the result; development: annotated; confidence: high}
\end{GSitem}
\begin{GSitem}{1965-53.7}{E}{The boundary estimate and its limitations}
\GSfield{Kind; source}{claim; \GSpage{P000091}}
\GSfield{The notes}{The estimate \(\int_{\partial M}|\varphi|^2\,dS\le C Q(\varphi,\varphi)\) is displayed, followed by the warning that it does not control the sum of squares of all first derivatives.}
\GSfield{Statement}{For a bounded smooth strongly pseudoconvex domain in \(\mathbb C^n\), \(q\ge1\), and smooth \((p,q)\)-forms satisfying the \(\bar\partial^*\) boundary condition, the boundary \(L^2\) norm is bounded by \(C Q\). More generally on a strongly pseudoconvex manifold an added \(L^2\) term or projection away from harmonic forms is needed. For \(0<q<n\), this is not a full \(H^1\) coercive estimate; at top antiholomorphic degree the boundary condition is of Dirichlet type.}
\GSfield{Reconstruction}{\ours\ The degree and geometric hypotheses are indispensable: at \(q=0\), nonzero holomorphic functions can have \(Q=0\) and a nonzero boundary norm. The notebook's one-dimensional comparison refers to the top-degree situation.}
\GSfield{Fate}{Already part of the expository framework in \citet{kohn1963potential}.}
\GSfield{Present status}{As of 9 November 2026, the precise statement is a standard established result or definition.}
\GSfield{Assessment}{literal statement: correct with missing hypotheses; publication match: exact; verification: read the result; development: annotated; confidence: high}
\end{GSitem}
\begin{GSitem}{1965-53.8}{F}{A reference to general overdetermined systems}
\GSfield{Kind; source}{remark; \GSpage{P000091}}
\GSfield{The notes}{The notes refer to work of Spencer on general overdetermined systems but give no operator class or statement.}
\GSfield{Statement}{No unique theorem about general overdetermined systems can be extracted. The identifiable idea is to exploit a differential complex and its compatibility operator when reducing a first-order system to a second-order variational problem.}
\GSfield{Reconstruction}{This is a programme-level reference, not a statement that every overdetermined system is solved by the Neumann construction.}
\GSfield{Fate}{\citet{kohn1963potential} discusses the related complex Neumann work and cites Kohn--Spencer, but this does not identify which general Spencer result the speaker meant.}
\GSfield{Present status}{As of 9 November 2026, the cited survey and its references establish the context; the notebook reference remains underdetermined.}
\GSfield{Assessment}{literal statement: not determinable; publication match: thematic only; verification: unresolved lead; development: annotated; confidence: low}
\end{GSitem}
\begin{GSitem}{1965-53.9}{E}{The Kohn-Rossi extension reference}
\GSfield{Kind; source}{remark; \GSpage{P000091}}
\GSfield{The notes}{The notes name Kohn and Rossi, Annals of Mathematics 1965, and emphasize holomorphic continuation of forms into the domain.}
\GSfield{Statement}{A precise special case is: if \(\Omega\Subset\mathbb C^n\), \(n\ge2\), has connected smooth strongly pseudoconvex boundary, every smooth CR function on \(\partial\Omega\) is the boundary value of a holomorphic function smooth on \(\overline\Omega\). Extension for higher-degree forms requires the appropriate tangential compatibility and cohomological conditions.}
\GSfield{Reconstruction}{The function statement is a safe modern specialization, not a transcription of the unspecified assertion about all forms. No unrestricted forms-extension theorem is attributed to this one-line entry.}
\GSfield{Fate}{The named article is \citet{kohnrossi1965extension}; its bibliographic identity is confirmed in Kohn's author bibliography. The original theorem text was not obtained in this search.}
\GSfield{Present status}{As of 9 November 2026, the function-extension theorem is established; the exact higher-form assertion intended here remains unspecified.}
\GSfield{Assessment}{literal statement: correct with missing hypotheses; publication match: partial; verification: indirect; development: annotated; confidence: medium}
\end{GSitem}

\GStalk{54}{1965-09-27}{R. Finn}{P000091-P000092}{Total curvature of complete open surfaces}
The brief entry recalls Gauss--Bonnet and Cohn--Vossen, illustrates the difference between closed and complete surfaces, and states the aim of replacing an inequality by an equality. Its final sentence is at the top of P000092, before Gelfand's new heading. Finn's 1965 article is a plausible publication context, but the surviving entry contains no end-defect formula with which to establish an exact theorem match \citep{cohnvossen1935curvature,finn1965conformal}.

\begin{GSitem}{1965-54.1}{E}{Complete surfaces need not satisfy the closed Gauss-Bonnet equality}
\GSfield{Kind; source}{remark; \GSpage{P000091}}
\GSfield{The notes}{A paraboloid is given as a complete, nonclosed surface with \(C=2\pi\) and \(\chi=1\); a second schematic surface is marked with a strict inequality.}
\GSfield{Statement}{The graph \(z=x^2+y^2\), with its induced metric, is complete, is homeomorphic to \(\mathbb R^2\), and has total Gaussian curvature \(2\pi\). Completeness alone does not force equality: the Euclidean plane has \(\chi=1\) and total curvature zero.}
\GSfield{Reconstruction}{\oursComputed\ The plane is a simple explicit replacement for the unnamed strict-inequality sketch, not an identification of that drawing. For the paraboloid, \(K=4/(1+4r^2)^2\), \(dA=r\sqrt{1+4r^2}\,dr\,d\theta\), whose integral is \(2\pi\).}
\GSfield{Fate}{Already part of the expository framework in \citet{cohnvossen1935curvature}.}
\GSfield{Present status}{As of 9 November 2026, the precise statement is a standard established result or definition.}
\GSfield{Assessment}{literal statement: correct; publication match: partial; verification: read the result; development: annotated; confidence: high}
\end{GSitem}
\begin{GSitem}{1965-54.2}{E}{Cohn-Vossen's inequality}
\GSfield{Kind; source}{theorem; \GSpage{P000091}}
\GSfield{The notes}{The statement is abbreviated to \(C\le2\pi\chi\) for every complete surface.}
\GSfield{Statement}{If a connected complete boundaryless Riemannian surface has finite topological type and integrable Gaussian curvature, then \(\int_M K\,dA\le2\pi\chi(M)\). The original formulation permits the corresponding well-defined total-curvature limit; absolute integrability is a convenient sufficient hypothesis here.}
\GSfield{Reconstruction}{Finite topological type and existence of total curvature are supplied. Without a meaningful integral and Euler characteristic the unqualified expression is not a theorem.}
\GSfield{Fate}{This is the pre-existing result in \citet[Theorem 6]{cohnvossen1935curvature}.}
\GSfield{Present status}{As of 9 November 2026, the precise statement is a standard established result or definition.}
\GSfield{Assessment}{literal statement: correct with missing hypotheses; publication match: exact; verification: read the result; development: annotated; confidence: high}
\end{GSitem}
\begin{GSitem}{1965-54.3}{F}{Finn's proposed equality with a defect term}
\GSfield{Kind; source}{question; \GSpage{P000092}}
\GSfield{The notes}{The sole statement about Finn's own aim is to replace the inequality by an equality.}
\GSfield{Statement}{The entry does not specify the class of metrics, the ends, or a proposed defect invariant. Consequently no unique equality theorem can be reconstructed from it.}
\GSfield{Reconstruction}{An identity obtained simply by defining a defect as \(2\pi\chi-\int K\,dA\) would be tautological; assigning Finn a particular geometric formula requires information absent here.}
\GSfield{Fate}{The title and date of \citet{finn1965conformal} fit the discussion, but the accessible publisher record establishes metadata, not the equality intended in this notebook.}
\GSfield{Present status}{As of 9 November 2026, the search identifies a likely contemporaneous publication context, not an unambiguously formulated unresolved problem.}
\GSfield{Assessment}{literal statement: not determinable; publication match: partial; verification: publisher metadata checked; development: annotated; confidence: low}
\end{GSitem}

\GStalk{55}{1965-09-27}{I. M. Gelfand}{P000092-P000093}{Atiyah-Bott: differential operators and elliptic complexes}
Gelfand begins an exposition attributed to Atiyah, introducing symbols, elliptic complexes, and the de Rham and Dolbeault examples. The speaker is Gelfand, not Atiyah. The reference underlying the sequence is the Harvard lecture material dated autumn 1964, preceding these meetings; its Russian translation appeared in 1966, and the journal article followed in 1967 \citep{atiyahbott1964notes,atiyahbott1967lefschetz}. The initial list of applications includes Brouwer, Lefschetz, eigenvectors, and Weyl characters. The symbol and finiteness discussions continue the topics indexed as 1964-28.2, 1964-28.4, and 1964-28.5; the underqualified earlier finiteness assertion 1964-28.3 should not be treated as a stronger theorem.

\begin{GSitem}{1965-55.1}{E}{Brouwer's theorem among the fixed-point applications}
\GSfield{Kind; source}{remark; \GSpage{P000092}}
\GSfield{The notes}{Brouwer's theorem is listed among the results contained in the proposed fixed-point theory.}
\GSfield{Statement}{Every continuous map from a closed finite-dimensional real ball to itself has a fixed point. Its relation to Lefschetz theory uses homotopy invariance and a suitable compact-space formulation; the smooth transverse elliptic-complex formula itself is not directly a theorem about arbitrary maps of a ball.}
\GSfield{Reconstruction}{The ball, continuity, and finite-dimensionality are supplied. The entry gives only the name and does not record the reduction.}
\GSfield{Fate}{Already part of the expository framework in \citet{atiyahbott1964notes}.}
\GSfield{Present status}{As of 9 November 2026, the precise statement is a standard established result or definition.}
\GSfield{Assessment}{literal statement: correct with missing hypotheses; publication match: thematic only; verification: read the result; development: annotated; confidence: medium}
\end{GSitem}
\begin{GSitem}{1965-55.2}{E}{Lefschetz's fixed-point theorem in the applications list}
\GSfield{Kind; source}{remark; \GSpage{P000092}}
\GSfield{The notes}{The ordinary Lefschetz formula is named before elliptic complexes are defined.}
\GSfield{Statement}{For a self-map \(f\) of a finite polyhedron, a nonzero Lefschetz number \(\sum_q(-1)^q\operatorname{tr}(f^*:H^q(X;\mathbb Q)\to H^q(X;\mathbb Q))\) forces a fixed point. For a smooth map of a closed manifold with nondegenerate fixed points, this number is \(\sum_p\operatorname{sgn}\det(I-df_p)\).}
\GSfield{Reconstruction}{These are precise familiar versions of the name in the notes; their derivation is discussed in record 65.}
\GSfield{Fate}{Already part of the expository framework in \citet{atiyahbott1964notes}.}
\GSfield{Present status}{As of 9 November 2026, the precise statement is a standard established result or definition.}
\GSfield{Assessment}{literal statement: correct with missing hypotheses; publication match: partial; verification: read the result; development: annotated; confidence: high}
\end{GSitem}
\begin{GSitem}{1965-55.3}{E}{The eigenvector application}
\GSfield{Kind; source}{remark; \GSpage{P000092}}
\GSfield{The notes}{The notes say that a linear space has an eigenvector.}
\GSfield{Statement}{Every endomorphism of a nonzero finite-dimensional complex vector space has an eigenvector. If it is invertible, its induced projective transformation has a fixed point, which is exactly an invariant line.}
\GSfield{Reconstruction}{The endomorphism and the finite-dimensional complex setting must be supplied. The literal wording is not a statement about every real operator or every infinite-dimensional operator.}
\GSfield{Fate}{Already part of the expository framework in \citet{atiyahbott1964notes}.}
\GSfield{Present status}{As of 9 November 2026, the precise statement is a standard established result or definition.}
\GSfield{Assessment}{literal statement: correct with missing hypotheses; publication match: thematic only; verification: read the result; development: annotated; confidence: medium}
\end{GSitem}
\begin{GSitem}{1965-55.4}{E}{Weyl's character formula in the applications list}
\GSfield{Kind; source}{remark; \GSpage{P000092}}
\GSfield{The notes}{Weyl's theorem on traces of representations is named without a group, weights, or formula.}
\GSfield{Statement}{For a compact connected simply connected semisimple group, choose positive roots, let \(W\) be its Weyl group, and let \(\rho\) be half the sum of positive roots. The irreducible representation of dominant integral highest weight \(\lambda\) has character \(\chi_\lambda(t)=\frac{\sum_{w\in W}\epsilon(w)e^{w(\lambda+\rho)}(t)}{\sum_{w\in W}\epsilon(w)e^{w\rho}(t)}\) on regular elements of a maximal torus, where \(\epsilon(w)=\det w\). Extension across the apparent singular set is by continuity.}
\GSfield{Reconstruction}{This supplies the standard meaning of the named formula. The relevant flag manifold, bundle, and fixed-point derivation are not present in this entry.}
\GSfield{Fate}{Already part of the expository framework in \citet{atiyahbott1964notes}.}
\GSfield{Present status}{As of 9 November 2026, the precise statement is a standard established result or definition.}
\GSfield{Assessment}{literal statement: correct with missing hypotheses; publication match: thematic only; verification: read the result; development: annotated; confidence: medium}
\end{GSitem}
\begin{GSitem}{1965-55.5}{E}{Differential complexes of vector bundles}
\GSfield{Kind; source}{definition; \GSpage{P000092}}
\GSfield{The notes}{Bundles \(E_1,\ldots,E_n\), spaces of sections, and a sequence of differential maps satisfying \(d^2=0\) are introduced.}
\GSfield{Statement}{A differential complex on a smooth manifold is a finite sequence of smooth finite-rank vector bundles and differential operators \(d_i:\Gamma(E_i)\to\Gamma(E_{i+1})\) such that \(d_{i+1}d_i=0\). Its \(i\)-th cohomology is \(\ker d_i/\operatorname{im}d_{i-1}\).}
\GSfield{Reconstruction}{The trivial line-bundle example identifies sections with smooth scalar functions. Compactness is unnecessary for the definition.}
\GSfield{Fate}{Already part of the expository framework in \citet{atiyahbott1964notes}.}
\GSfield{Present status}{As of 9 November 2026, the precise statement is a standard established result or definition.}
\GSfield{Assessment}{literal statement: correct; publication match: exact; verification: read the result; development: annotated; confidence: high}
\end{GSitem}
\begin{GSitem}{1965-55.6}{E}{Coordinate-free principal symbols}
\GSfield{Kind; source}{definition; \GSpage{P000092}}
\GSfield{The notes}{The operator acts on \(f^m s\), where \(f(x)=0\), and its value at \(x\) is used to define the symbol in terms of \(df_x\) and \(s(x)\).}
\GSfield{Statement}{For an order-\(m\) differential operator \(P\), \(P(f^ms)(x)=m!\,\sigma_{\mathrm{alg}}(P)(x,df_x)s(x)\), where \(\sigma_{\mathrm{alg}}(P)(x,\xi)=\sum_{|\alpha|=m}a_\alpha(x)\xi^\alpha\) uses ordinary derivatives. Equivalently the notebook's symbol is \(m!\sigma_{\mathrm{alg}}\). The convention with \(D_j=(1/i)\partial_{x_j}\) introduces the corresponding powers of \(i\).}
\GSfield{Reconstruction}{The formula is stated for an already local differential operator. It is not silently promoted to a characterization of arbitrary linear maps on sections; the notebook compresses the locality and order discussion.}
\GSfield{Fate}{Already part of the expository framework in \citet{atiyahbott1964notes}.}
\GSfield{Present status}{As of 9 November 2026, the precise statement is a standard established result or definition.}
\GSfield{Assessment}{literal statement: correct with missing hypotheses; publication match: partial; verification: read the result; development: annotated; confidence: high}
\end{GSitem}
\begin{GSitem}{1965-55.7}{E}{The factor two in the Laplace symbol}
\GSfield{Kind; source}{claim; \GSpage{P000092}}
\GSfield{The notes}{For the Laplacian, the notes obtain \(2(\eta_1^2+\cdots+\eta_n^2)\).}
\GSfield{Statement}{If \(\Delta=\sum_j\partial_{x_j}^2\) and \(f(x)=0\), then \(\Delta(f^2s)(x)=2|df_x|^2s(x)\). Thus the factor two is correct for the preceding unnormalized symbol definition.}
\GSfield{Reconstruction}{The evaluation at \(x\), where \(f(x)=0\), is essential. This is neither the globally valid product rule nor the symbol convention of the positive Laplacian \(-\sum\partial_j^2\).}
\GSfield{Fate}{Already part of the expository framework in \citet{atiyahbott1964notes}.}
\GSfield{Present status}{As of 9 November 2026, the precise statement is a standard established result or definition.}
\GSfield{Assessment}{literal statement: correct with missing hypotheses; publication match: partial; verification: read the result; development: annotated; confidence: high}
\end{GSitem}
\begin{GSitem}{1965-55.8}{E}{Ellipticity of a complex}
\GSfield{Kind; source}{definition; \GSpage{P000092}}
\GSfield{The notes}{The sequence of fibre maps given by the symbols is required to be exact for every \(\eta\ne0\).}
\GSfield{Statement}{A differential complex is elliptic if for each nonzero real cotangent vector \(\xi\in T_x^*X\), the symbol sequence \(0\to E_{0,x}\to\cdots\to E_{N,x}\to0\) is exact. Each map uses its own operator order; rescaling an individual symbol by a nonzero conventional constant leaves exactness unchanged.}
\GSfield{Reconstruction}{Nonzero real cotangent vectors are meant, even for complex vector bundles. Finite-dimensional cohomology additionally requires global hypotheses supplied later.}
\GSfield{Fate}{Already part of the expository framework in \citet{atiyahbott1964notes}.}
\GSfield{Present status}{As of 9 November 2026, the precise statement is a standard established result or definition.}
\GSfield{Assessment}{literal statement: correct; publication match: exact; verification: read the result; development: annotated; confidence: high}
\end{GSitem}
\begin{GSitem}{1965-55.9}{E}{The Laplace operator as a two-term elliptic complex}
\GSfield{Kind; source}{construction; \GSpage{P000093}}
\GSfield{The notes}{The first example is the trivial line bundle and the Laplace operator between its spaces of sections.}
\GSfield{Statement}{On a Riemannian manifold, the Laplace operator on functions has invertible scalar principal symbol at each nonzero real covector; hence \(0\to\Gamma(1)\xrightarrow{\Delta}\Gamma(1)\to0\) is an elliptic complex.}
\GSfield{Reconstruction}{No compactness or boundary condition is required for symbol ellipticity. Fredholm or cohomological finiteness conclusions do not follow for an arbitrary open domain.}
\GSfield{Fate}{Already part of the expository framework in \citet{atiyahbott1964notes}.}
\GSfield{Present status}{As of 9 November 2026, the precise statement is a standard established result or definition.}
\GSfield{Assessment}{literal statement: correct; publication match: exact; verification: read the result; development: annotated; confidence: high}
\end{GSitem}
\begin{GSitem}{1965-55.10}{E}{The de Rham symbol is an exact Koszul complex}
\GSfield{Kind; source}{construction; \GSpage{P000093}}
\GSfield{The notes}{The notes identify the symbols of \(d\) with exterior multiplication by \(\eta\) and state that \(e\wedge\eta=0\) implies \(e=a\wedge\eta\).}
\GSfield{Statement}{The de Rham differential has symbol \(\xi\wedge-\), up to the chosen factor and graded sign. For \(\xi\ne0\), choose \(v\) with \(\xi(v)=1\). Then \(\iota_v(\xi\wedge-)+\xi\wedge\iota_v=\operatorname{id}\), so the exterior-algebra symbol complex is exact in every degree.}
\GSfield{Reconstruction}{\ours\ The contraction identity supplies the brief algebraic justification. The notebook's product-rule display misses a sign for odd-degree \(s\): at \(f(x)=0\), \(d(fs)(x)=df_x\wedge s(x)=(-1)^q s(x)\wedge df_x\) for a \(q\)-form. These degree-dependent signs do not affect symbol exactness.}
\GSfield{Fate}{Already part of the expository framework in \citet{atiyahbott1964notes}.}
\GSfield{Present status}{As of 9 November 2026, the precise statement is a standard established result or definition.}
\GSfield{Assessment}{literal statement: false as written; publication match: exact; verification: read the result; development: annotated; confidence: high}
\end{GSitem}
\begin{GSitem}{1965-55.11}{E}{The Dolbeault elliptic complex}
\GSfield{Kind; source}{construction; \GSpage{P000093}}
\GSfield{The notes}{A third example is written simply as \(\bar\partial\), and so on.}
\GSfield{Statement}{For fixed \(p\) on a complex manifold, \((A^{p,\bullet},\bar\partial)\) is elliptic: at a nonzero real covector \(\xi\), the symbol is exterior multiplication by its nonzero \((0,1)\) part. The same contraction argument proves exactness.}
\GSfield{Reconstruction}{The use of real covectors matters; a nonzero complexified covector may have zero \((0,1)\) part.}
\GSfield{Fate}{Already part of the expository framework in \citet{atiyahbott1964notes}.}
\GSfield{Present status}{As of 9 November 2026, the precise statement is a standard established result or definition.}
\GSfield{Assessment}{literal statement: correct; publication match: exact; verification: read the result; development: annotated; confidence: high}
\end{GSitem}

\GStalk{56}{1965-10-04}{I. M. Gelfand}{P000093-P000095}{Atiyah-Bott: Lefschetz numbers, local terms, and Hodge theory}
The exposition resumes below the dated continuation heading on P000093. Gelfand defines endomorphisms of complexes and their Lefschetz numbers, describes a map with bundle lifts and simple fixed points, states the elliptic fixed-point formula, and starts its spectral proof. The main formulas require a consistent pullback convention, compactness, and the distinction between nonnegative spectrum and strictly positive eigenspaces \citep{atiyahbott1964notes}.

\begin{GSitem}{1965-56.1}{E}{The symbol formula recalled}
\GSfield{Kind; source}{definition; \GSpage{P000093}}
\GSfield{The notes}{The continuation begins by repeating \(d(f^m s)(x_0)=\sigma(d,\eta)s(x_0)\), with \(\eta=df_{x_0}\).}
\GSfield{Statement}{For an order-\(m\) differential operator, with \(f(x_0)=0\), the value of \(d(f^m s)\) at \(x_0\) is \(m!\) times the algebraically normalized principal symbol at \(df_{x_0}\), acting on \(s(x_0)\).}
\GSfield{Reconstruction}{This is a repeat occurrence of the unnormalized symbol convention. The same pointwise-vanishing assumption and factorial convention as record 55 are retained.}
\GSfield{Fate}{Already part of the expository framework in \citet{atiyahbott1964notes}.}
\GSfield{Present status}{As of 9 November 2026, the precise statement is a standard established result or definition.}
\GSfield{Assessment}{literal statement: correct with missing hypotheses; publication match: exact; verification: read the result; development: annotated; confidence: high}
\end{GSitem}
\begin{GSitem}{1965-56.2}{E}{Elliptic operators beyond differential operators}
\GSfield{Kind; source}{remark; \GSpage{P000093}}
\GSfield{The notes}{A brief line notes that elliptic operators need not be differential operators.}
\GSfield{Statement}{Ellipticity also makes sense for classical pseudodifferential operators: an order-\(m\) operator is elliptic when its homogeneous principal symbol is invertible away from the zero section. For example the identity plus a nonzero smoothing integral operator is elliptic of order zero and need not be a differential operator.}
\GSfield{Reconstruction}{\ours\ The example specifies one standard extension of the term. The notebook does not yet define the broader class; its oscillatory definition appears in record 61.}
\GSfield{Fate}{Already part of the expository framework in \citet{atiyahbott1964notes}.}
\GSfield{Present status}{As of 9 November 2026, the precise statement is a standard established result or definition.}
\GSfield{Assessment}{literal statement: correct with missing hypotheses; publication match: partial; verification: read the result; development: annotated; confidence: high}
\end{GSitem}
\begin{GSitem}{1965-56.3}{E}{Complex endomorphisms and induced cohomology maps}
\GSfield{Kind; source}{definition; \GSpage{P000093}}
\GSfield{The notes}{The maps \(T_i\) commute with the differentials and induce \(T_i^*\) on the cohomology.}
\GSfield{Statement}{If \(d_iT_i=T_{i+1}d_i\), then \(T_i\) preserves \(\ker d_i\) and carries \(\operatorname{im}d_{i-1}\) to itself, so it induces a well-defined linear map \(H^i(T)\) on \(\ker d_i/\operatorname{im}d_{i-1}\).}
\GSfield{Reconstruction}{The star on \(T_i^*\) denotes the induced cohomology map here, not the metric adjoint.}
\GSfield{Fate}{Already part of the expository framework in \citet{atiyahbott1964notes}.}
\GSfield{Present status}{As of 9 November 2026, the precise statement is a standard established result or definition.}
\GSfield{Assessment}{literal statement: correct; publication match: exact; verification: read the result; development: annotated; confidence: high}
\end{GSitem}
\begin{GSitem}{1965-56.4}{E}{Finiteness of elliptic-complex cohomology}
\GSfield{Kind; source}{theorem; \GSpage{P000093}}
\GSfield{The notes}{All cohomology groups of an elliptic complex are said to be finite-dimensional.}
\GSfield{Statement}{For a finite elliptic complex of smooth finite-rank bundles on a compact smooth manifold without boundary, every smooth cohomology group is finite-dimensional. On a manifold with boundary, ellipticity of the interior symbol alone is insufficient.}
\GSfield{Reconstruction}{\ours\ Compactness and absence of boundary are supplied. A noncompact complex can have infinite-dimensional cohomology, for example holomorphic functions in degree zero of a Dolbeault complex.}
\GSfield{Fate}{Already part of the expository framework in \citet{atiyahbott1964notes}.}
\GSfield{Present status}{As of 9 November 2026, the precise statement is a standard established result or definition.}
\GSfield{Assessment}{literal statement: correct with missing hypotheses; publication match: exact; verification: read the result; development: annotated; confidence: high}
\end{GSitem}
\begin{GSitem}{1965-56.5}{E}{The Lefschetz number of an elliptic complex}
\GSfield{Kind; source}{definition; \GSpage{P000093}}
\GSfield{The notes}{The alternating trace of the induced cohomology endomorphisms is denoted \(L(T,E)\).}
\GSfield{Statement}{For a finite cohomologically finite complex, \(L(T,E)=\sum_i(-1)^i\operatorname{tr}H^i(T)\). A consistent grading is part of this definition: shifting every degree by one reverses the sign.}
\GSfield{Reconstruction}{The notebook labels bundles starting at 1 while also referring to the usual de Rham convention. We use the actual cohomological degree, beginning at zero for forms.}
\GSfield{Fate}{Already part of the expository framework in \citet{atiyahbott1964notes}.}
\GSfield{Present status}{As of 9 November 2026, the precise statement is a standard established result or definition.}
\GSfield{Assessment}{literal statement: correct with missing hypotheses; publication match: exact; verification: read the result; development: annotated; confidence: high}
\end{GSitem}
\begin{GSitem}{1965-56.6}{E}{Geometric lifts of a map to the complex}
\GSfield{Kind; source}{construction; \GSpage{P000094}}
\GSfield{The notes}{A map \(f:X\to X\) and maps of the pulled-back bundles define the operators \(T_i\), but the fibre display uses \(f^{-1}x\).}
\GSfield{Statement}{For a smooth map \(f\) and a bundle morphism \(\phi_i:f^*E_i\to E_i\), define \(T_is(x)=\phi_i(x)s(f(x))\), with \(\phi_i(x):E_{i,f(x)}\to E_{i,x}\). Assume these operators form a chain map. An inverse-map convention is possible when \(f\) is a diffeomorphism, but the kernel and derivative conventions must then change together.}
\GSfield{Reconstruction}{The source's \(f^{-1}x\) cannot define a fibre for an arbitrary noninvertible map. The displayed pullback and inverse conventions are not silently conflated.}
\GSfield{Fate}{Already part of the expository framework in \citet{atiyahbott1964notes}.}
\GSfield{Present status}{As of 9 November 2026, the precise statement is a standard established result or definition.}
\GSfield{Assessment}{literal statement: false as written; publication match: partial; verification: read the result; development: annotated; confidence: high}
\end{GSitem}
\begin{GSitem}{1965-56.7}{E}{Simple fixed points and transversality}
\GSfield{Kind; source}{definition; \GSpage{P000094}}
\GSfield{The notes}{The notes call a fixed point simple when \(\det(I-df_p)\ne0\) and describe transverse intersection of the graph and diagonal.}
\GSfield{Statement}{A fixed point \(p\) of a smooth self-map is nondegenerate if \(I-df_p:T_pX\to T_pX\) is invertible. This is equivalent to transversality of its graph and the diagonal at \((p,p)\), and implies that \(p\) is isolated. On compact \(X\), if every fixed point is nondegenerate, there are finitely many.}
\GSfield{Fate}{Already part of the expository framework in \citet{atiyahbott1964notes}.}
\GSfield{Present status}{As of 9 November 2026, the precise statement is a standard established result or definition.}
\GSfield{Assessment}{literal statement: correct; publication match: exact; verification: read the result; development: annotated; confidence: high}
\end{GSitem}
\begin{GSitem}{1965-56.8}{E}{The local Lefschetz contribution}
\GSfield{Kind; source}{definition; \GSpage{P000094}}
\GSfield{The notes}{The local number is an alternating fibre trace divided by the absolute determinant of \(1-df\).}
\GSfield{Statement}{With the preceding pullback convention and a nondegenerate fixed point, the local contribution is \(\nu_p(T)=\frac{\sum_i(-1)^i\operatorname{tr}\phi_i(p)}{|\det_{\mathbb R}(I-df_p)|}\). The numerator is the supertrace on the graded bundle fibre, not the cohomological trace.}
\GSfield{Reconstruction}{The real determinant in the denominator is distinguished from the complex determinant used after simplifying a Dolbeault application.}
\GSfield{Fate}{Already part of the expository framework in \citet{atiyahbott1964notes}.}
\GSfield{Present status}{As of 9 November 2026, the precise statement is a standard established result or definition.}
\GSfield{Assessment}{literal statement: correct; publication match: exact; verification: read the result; development: annotated; confidence: high}
\end{GSitem}
\begin{GSitem}{1965-56.9}{E}{The Atiyah-Bott fixed-point formula}
\GSfield{Kind; source}{theorem; \GSpage{P000094}}
\GSfield{The notes}{The sum of the local numbers over simple fixed points is asserted to equal \(L(T,E)\).}
\GSfield{Statement}{Let \(X\) be compact and without boundary, let \(E^\bullet\) be an elliptic complex, and let a smooth map \(f:X\to X\) have only nondegenerate fixed points. If smooth zero-order bundle lifts define a chain map \(T\), then \(L(T,E)=\sum_{p=f(p)}\nu_p(T)\).}
\GSfield{Reconstruction}{All global and chain-map hypotheses are gathered here. The theorem concerns a geometric operator, not an arbitrary endomorphism of the section complex.}
\GSfield{Fate}{The theorem was already in the autumn 1964 Harvard notes \citep{atiyahbott1964notes}; the later article \citep{atiyahbott1967lefschetz} does not turn Gelfand's exposition into a new announcement.}
\GSfield{Present status}{As of 9 November 2026, the precise statement is a standard established result or definition.}
\GSfield{Assessment}{literal statement: correct with missing hypotheses; publication match: exact; verification: read the result; development: annotated; confidence: high}
\end{GSitem}
\begin{GSitem}{1965-56.10}{E}{Formal adjoints and the Hodge Laplacian}
\GSfield{Kind; source}{construction; \GSpage{P000095}}
\GSfield{The notes}{After choosing metrics, the adjoint and \(\Box=dd^*+d^*d\) are introduced.}
\GSfield{Statement}{For a first-order elliptic complex on a closed manifold with Hermitian bundle metrics and a positive density, put \(\Box_i=d_{i-1}d_{i-1}^*+d_i^*d_i\). It is elliptic, formally self-adjoint, and nonnegative: \((\Box_i u,u)=\|d_{i-1}^*u\|^2+\|d_i u\|^2\).}
\GSfield{Reconstruction}{For a complex whose differentials have unequal positive orders, the unweighted sum need not have the required elliptic principal symbol; one needs an equal-order reduction or suitable powers. The simple formula is used here in its first-order setting.}
\GSfield{Fate}{Already part of the expository framework in \citet{atiyahbott1964notes}.}
\GSfield{Present status}{As of 9 November 2026, the precise statement is a standard established result or definition.}
\GSfield{Assessment}{literal statement: correct with missing hypotheses; publication match: exact; verification: read the result; development: annotated; confidence: high}
\end{GSitem}
\begin{GSitem}{1965-56.11}{E}{Discrete spectrum and harmonic representatives}
\GSfield{Kind; source}{theorem; \GSpage{P000095}}
\GSfield{The notes}{The notes describe positive eigenvalues of finite multiplicity, followed by a separate zero-eigenvalue space representing cohomology.}
\GSfield{Statement}{The self-adjoint elliptic Hodge Laplacian on a closed manifold has nonnegative discrete spectrum, finite-dimensional smooth eigenspaces, and a complete \(L^2\) orthonormal eigenbasis. The harmonic space \(\ker\Box_i=\ker d_i\cap\ker d_{i-1}^*\) maps isomorphically to \(H^i(E^\bullet)\).}
\GSfield{Reconstruction}{Nonnegative replaces the earlier unrestricted positive. Finite multiplicity means finitely many independent eigenvectors for each eigenvalue, not finitely many eigenvalues overall.}
\GSfield{Fate}{Already part of the expository framework in \citet{atiyahbott1964notes}.}
\GSfield{Present status}{As of 9 November 2026, the precise statement is a standard established result or definition.}
\GSfield{Assessment}{literal statement: false as written; publication match: exact; verification: read the result; development: annotated; confidence: high}
\end{GSitem}
\begin{GSitem}{1965-56.12}{E}{Positive-eigenvalue complexes are acyclic}
\GSfield{Kind; source}{claim; \GSpage{P000095}}
\GSfield{The notes}{The differentials preserve the eigenspaces of \(\Box\), and the notes separate the positive-eigenvalue complexes from the harmonic one.}
\GSfield{Statement}{The identities \(d\Box=\Box d\) and \(d^*\Box=\Box d^*\) give complexes on each eigenspace. If \(\lambda>0\), \(h=d^*/\lambda\) satisfies \(dh+hd=\operatorname{id}\), so that eigenspace complex has zero cohomology.}
\GSfield{Reconstruction}{The contracting homotopy is the explicit proof later supplied in Bernstein's intervention. No commutation of a separate chain map \(T\) with \(\Box\) is assumed.}
\GSfield{Fate}{Already part of the expository framework in \citet{atiyahbott1964notes}.}
\GSfield{Present status}{As of 9 November 2026, the precise statement is a standard established result or definition.}
\GSfield{Assessment}{literal statement: correct; publication match: exact; verification: read the result; development: annotated; confidence: high}
\end{GSitem}
\begin{GSitem}{1965-56.13}{E}{Formal trace cancellation}
\GSfield{Kind; source}{method; \GSpage{P000095}}
\GSfield{The notes}{The notes compare the sum of diagonal matrix entries, the sum of eigenvalues, and integration of a kernel on the diagonal, anticipating cancellation in an alternating trace.}
\GSfield{Statement}{For a trace-class operator these are valid trace descriptions under their usual convergence hypotheses; a smooth kernel on a closed manifold gives a trace by diagonal integration. A geometric pullback is generally not trace class. The cohomological supertrace must therefore be obtained using a smoothing regularization before applying finite-dimensional cancellation.}
\GSfield{Reconstruction}{These are formal motivations, not permission to identify divergent ordinary traces. Later records construct the required regularization.}
\GSfield{Fate}{Already part of the expository framework in \citet{atiyahbott1964notes}.}
\GSfield{Present status}{As of 9 November 2026, the precise statement is a standard established result or definition.}
\GSfield{Assessment}{literal statement: correct with missing hypotheses; publication match: exact; verification: read the result; development: annotated; confidence: high}
\end{GSitem}

\GStalk{57}{1965-10-11}{I. M. Gelfand}{P000096}{Seminar programme for 1965-66}
The seven agenda entries propose Agranovich on integro-differential operators; Kirillov, Gelfand, or Dynin on Atiyah--Singer and the Cartan--Schwartz and Vishik seminars; Minlos on statistical physics; Graev on adeles; Kirillov and Gelfand on algebraic groups; Gelfand on Bott--Chern and zero distribution of sections; and Kazhdan on Riemann--Roch--Grothendieck. These are assignments and reading topics, not theorem statements. The Bott--Chern reference can be identified as the 1965 Acta Mathematica paper \citep{bottchern1965hermitian}. Kazhdan's talk begins separately in the lower part of the same page.


\GStalk{58}{1965-10-11}{D. A. Kazhdan}{P000096-P000100}{Algebraic varieties, rational equivalence, and the approach to Riemann-Roch}
The announced Riemann--Roch discussion starts with the foundations of classical algebraic geometry: affine sets, regular functions, ringed spaces, gluing, products, and separation. It then introduces cycles and their equivalence via \(\mathbb P^1\), ending with an unfinished proposed map \(K(X)\to A(X)\). The framework belongs to the earlier theory of coherent sheaves and Grothendieck's Riemann--Roch theorem \citep{serre1955coherent,borelserre1958riemannroch}. The last lines do not state a Riemann--Roch formula or a complete axiom system for Chern classes.

\begin{GSitem}{1965-58.1}{E}{The announced characteristic classes of a vector bundle}
\GSfield{Kind; source}{construction; \GSpage{P000096}}
\GSfield{The notes}{The first line under Kazhdan's heading associates a vector sheaf with a cohomology class, called a Chern class.}
\GSfield{Statement}{A complex vector bundle \(E\) of rank \(r\) on a paracompact Hausdorff space has characteristic classes \(c_i(E)\in H^{2i}(X;\mathbb Z)\), \(0\le i\le r\), with \(c_0=1\). In the smooth algebraic setting, locally free sheaves have algebraic classes \(c_i(E)\in CH^i(X)\); over \(\mathbb C\), the cycle-class comparison gives their topological counterparts.}
\GSfield{Reconstruction}{Finite-rank locally free sheaves, rather than arbitrary coherent sheaves on an unspecified singular space, are used for this introductory statement. No particular class or axiom is supplied in the source line.}
\GSfield{Fate}{Already part of the expository framework in \citet{borelserre1958riemannroch}.}
\GSfield{Present status}{As of 9 November 2026, the precise statement is a standard established result or definition.}
\GSfield{Assessment}{literal statement: correct with missing hypotheses; publication match: partial; verification: read the result; development: annotated; confidence: high}
\end{GSitem}
\begin{GSitem}{1965-58.2}{E}{Affine algebraic sets and their coordinate rings}
\GSfield{Kind; source}{definition; \GSpage{P000096}}
\GSfield{The notes}{Over an algebraically closed field, the common zeros of given polynomials form an affine variety; polynomial functions are identified with the quotient by the ideal spanned by the equations.}
\GSfield{Statement}{Over an algebraically closed field \(k\), for \(I\subset k[x_1,\ldots,x_n]\), the classical algebraic set \(X=V(I)\) has ring of polynomial functions \(k[X]=k[x_1,\ldots,x_n]/I(X)=k[x]/\sqrt I\). One obtains \(k[x]/I\) itself if \(I\) is radical, or if one deliberately uses the possibly nonreduced affine scheme \(\operatorname{Spec}(k[x]/I)\).}
\GSfield{Reconstruction}{\ours\ The notebook works with point sets and restrictions of functions, so passing to the radical is required. For example \(x^2=0\) and \(x=0\) define the same classical point set but different quotient rings.}
\GSfield{Fate}{Already part of the expository framework in \citet{serre1955coherent}.}
\GSfield{Present status}{As of 9 November 2026, the precise statement is a standard established result or definition.}
\GSfield{Assessment}{literal statement: false as written; publication match: exact; verification: read the result; development: annotated; confidence: high}
\end{GSitem}
\begin{GSitem}{1965-58.3}{E}{Regular functions on principal open sets}
\GSfield{Kind; source}{definition; \GSpage{P000097-P000098}}
\GSfield{The notes}{Additional nonvanishing conditions are allowed; on the punctured line the functions are \(P(x)/x^m\), and principal opens provide a basis.}
\GSfield{Statement}{For \(X=\operatorname{Spec}A\) of finite type over an algebraically closed field and \(f\in A\), \(D(f)\) is affine with coordinate ring \(A_f\). In the reduced classical interpretation these are locally rational functions with no poles on \(D(f)\). On \(\mathbb A^1\setminus\{0\}\), the ring is \(k[t,t^{-1}]\).}
\GSfield{Reconstruction}{For reducible varieties, a single rational-function-field description is inadequate; the localization formulation handles zero divisors correctly. Finitely many nonvanishing conditions combine as a product.}
\GSfield{Fate}{Already part of the expository framework in \citet{serre1955coherent}.}
\GSfield{Present status}{As of 9 November 2026, the precise statement is a standard established result or definition.}
\GSfield{Assessment}{literal statement: correct with missing hypotheses; publication match: exact; verification: read the result; development: annotated; confidence: high}
\end{GSitem}
\begin{GSitem}{1965-58.4}{E}{Prevarieties as locally affine ringed spaces}
\GSfield{Kind; source}{definition; \GSpage{P000097}}
\GSfield{The notes}{A prevariety is locally isomorphic to an affine variety; the notes emphasize that its functions and ringed-space structure accompany its points.}
\GSfield{Statement}{A classical algebraic prevariety is a space with a sheaf of \(k\)-algebras locally isomorphic to affine algebraic varieties with their regular-function sheaves, subject to the customary finite-type cover condition. The local isomorphisms are isomorphisms of ringed spaces, not just homeomorphisms.}
\GSfield{Reconstruction}{The finite-type and compatibility conventions distinguish this from an arbitrary locally ringed space or an arbitrary gluing of sets.}
\GSfield{Fate}{Already part of the expository framework in \citet{serre1955coherent}.}
\GSfield{Present status}{As of 9 November 2026, the precise statement is a standard established result or definition.}
\GSfield{Assessment}{literal statement: correct with missing hypotheses; publication match: exact; verification: read the result; development: annotated; confidence: high}
\end{GSitem}
\begin{GSitem}{1965-58.5}{E}{The structure sheaf and its restriction maps}
\GSfield{Kind; source}{construction; \GSpage{P000097-P000098}}
\GSfield{The notes}{The sheaf of rings is specified on a basis of sets defined by a nonvanishing polynomial, with sections identified by their values.}
\GSfield{Statement}{The presheaf \(D(f)\mapsto A_f\), with localization restriction maps, extends uniquely to the structure sheaf \(\mathcal O_X\) on an affine algebraic variety. Local regular functions that agree on overlaps glue uniquely.}
\GSfield{Reconstruction}{An arbitrary assignment of rings to open sets would not suffice: the restriction maps and gluing axiom are part of the construction.}
\GSfield{Fate}{Already part of the expository framework in \citet{serre1955coherent}.}
\GSfield{Present status}{As of 9 November 2026, the precise statement is a standard established result or definition.}
\GSfield{Assessment}{literal statement: correct; publication match: exact; verification: read the result; development: annotated; confidence: high}
\end{GSitem}
\begin{GSitem}{1965-58.6}{E}{Morphisms respect functions contravariantly}
\GSfield{Kind; source}{definition; \GSpage{P000098}}
\GSfield{The notes}{The diagram for an isomorphism of ringed spaces includes a map between the bases and an induced map on sheaves of functions in the opposite direction.}
\GSfield{Statement}{A morphism of locally ringed spaces \(f:X\to Y\) consists of a continuous map and \(f^\#:\mathcal O_Y\to f_*\mathcal O_X\) inducing local homomorphisms of stalks. For affine varieties this is contravariantly equivalent to a homomorphism of their coordinate \(k\)-algebras.}
\GSfield{Reconstruction}{The notes' emphasis on maps induced by the base is made precise by pullback of regular functions and the local condition. Abstract bundle isomorphisms of the underlying sheaf spaces are not enough.}
\GSfield{Fate}{Already part of the expository framework in \citet{serre1955coherent}.}
\GSfield{Present status}{As of 9 November 2026, the precise statement is a standard established result or definition.}
\GSfield{Assessment}{literal statement: correct with missing hypotheses; publication match: exact; verification: read the result; development: annotated; confidence: high}
\end{GSitem}
\begin{GSitem}{1965-58.7}{E}{The punctured affine plane is not affine}
\GSfield{Kind; source}{claim; \GSpage{P000097-P000098}}
\GSfield{The notes}{The example \(\mathbb C^2\setminus\{0\}\) is called nonaffine but is covered by the two affine opens \(x\ne0\) and \(y\ne0\).}
\GSfield{Statement}{For any field \(k\), \(U=\mathbb A^2_k\setminus\{0\}\) is covered by \(D(x)\) and \(D(y)\), and \(\Gamma(U,\mathcal O_U)=k[x,y]_x\cap k[x,y]_y=k[x,y]\) inside \(k(x,y)\). If \(U\) were affine, the canonical morphism \(U\to\operatorname{Spec}\Gamma(U,\mathcal O_U)=\mathbb A^2_k\) would be an isomorphism, contrary to its missing origin.}
\GSfield{Reconstruction}{\ours\ The intersection computation supplies an algebraic proof and does not depend on analytic Hartogs extension. The classical point-set formulation in the notebook uses \(k=\mathbb C\).}
\GSfield{Fate}{Already part of the expository framework in \citet{serre1955coherent}.}
\GSfield{Present status}{As of 9 November 2026, the precise statement is a standard established result or definition.}
\GSfield{Assessment}{literal statement: correct; publication match: partial; verification: read the result; development: annotated; confidence: high}
\end{GSitem}
\begin{GSitem}{1965-58.8}{E}{Realizing a principal open as a closed affine variety}
\GSfield{Kind; source}{construction; \GSpage{P000098-P000099}}
\GSfield{The notes}{The condition \(Q(x)\ne0\) is replaced by the extra equation \(Q(x)y-1=0\).}
\GSfield{Statement}{If \(X\subset\mathbb A^n\) has coordinate ring \(A\), the maps \(x\mapsto(x,Q(x)^{-1})\) and projection identify \(D(Q)\) with the closed subvariety of \(X\times\mathbb A^1\) defined by \(Qy-1=0\). Algebraically, \(A_Q\cong A[y]/(Qy-1)\).}
\GSfield{Fate}{Already part of the expository framework in \citet{serre1955coherent}.}
\GSfield{Present status}{As of 9 November 2026, the precise statement is a standard established result or definition.}
\GSfield{Assessment}{literal statement: correct; publication match: exact; verification: read the result; development: annotated; confidence: high}
\end{GSitem}
\begin{GSitem}{1965-58.9}{E}{Products and tensor-product coordinate rings}
\GSfield{Kind; source}{construction; \GSpage{P000099}}
\GSfield{The notes}{For affine factors the notebook adds the two extended ideals in the polynomial ring of both variable sets and obtains a tensor product of coordinate rings; general products are glued from affine products.}
\GSfield{Statement}{If \(X=\operatorname{Spec}A\) and \(Y=\operatorname{Spec}B\) are affine \(k\)-schemes, \(X\times_kY=\operatorname{Spec}(A\otimes_kB)\). For \(A=k[x]/I\), \(B=k[y]/J\), the tensor product is \(k[x,y]/(I\,k[x,y]+J\,k[x,y])\). In the classical reduced finite-type setting over algebraically closed \(k\), these products are again reduced; compatible affine charts glue the general product.}
\GSfield{Reconstruction}{The operation on ideals is a sum inside the ambient polynomial ring, not an assertion of direct-sum independence.}
\GSfield{Fate}{Already part of the expository framework in \citet{serre1955coherent}.}
\GSfield{Present status}{As of 9 November 2026, the precise statement is a standard established result or definition.}
\GSfield{Assessment}{literal statement: correct; publication match: exact; verification: read the result; development: annotated; confidence: high}
\end{GSitem}
\begin{GSitem}{1965-58.10}{E}{The Zariski topology on a product is not the product topology}
\GSfield{Kind; source}{claim; \GSpage{P000099}}
\GSfield{The notes}{The notes explicitly warn that the product of varieties is not the ordinary topological product.}
\GSfield{Statement}{For an infinite algebraically closed field, \(\mathbb A^1(k)\) has the cofinite topology. The diagonal in \(\mathbb A^2(k)\) is Zariski closed, but it is dense in the product of the two cofinite topologies. Thus the algebraic product topology is strictly finer in this example.}
\GSfield{Reconstruction}{\ours\ This distinguishes the product in varieties from the product in topological spaces and also clarifies the subsequent diagonal criterion.}
\GSfield{Fate}{Already part of the expository framework in \citet{serre1955coherent}.}
\GSfield{Present status}{As of 9 November 2026, the precise statement is a standard established result or definition.}
\GSfield{Assessment}{literal statement: correct; publication match: exact; verification: read the result; development: annotated; confidence: high}
\end{GSitem}
\begin{GSitem}{1965-58.11}{E}{Separation and the closed diagonal}
\GSfield{Kind; source}{definition; \GSpage{P000099}}
\GSfield{The notes}{A prevariety whose diagonal in \(X\times X\) is closed is called an algebraic variety; an aside compares this with Hausdorff separation and then incorrectly with the sheaf axiom.}
\GSfield{Statement}{A classical prevariety over \(k\) is separated if its diagonal is closed in the algebraic product. In scheme language the diagonal must be a closed immersion. The Hausdorff criterion uses the ordinary topological product; the sheaf condition instead consists of unique gluing of compatible local sections.}
\GSfield{Reconstruction}{The handwritten sentence that a presheaf is a sheaf when its diagonal is closed has no valid meaning in the displayed framework. It is not a criterion for sheaves.}
\GSfield{Fate}{Already part of the expository framework in \citet{serre1955coherent}.}
\GSfield{Present status}{As of 9 November 2026, the precise statement is a standard established result or definition.}
\GSfield{Assessment}{literal statement: false as written; publication match: exact; verification: read the result; development: annotated; confidence: high}
\end{GSitem}
\begin{GSitem}{1965-58.12}{E}{The affine line with doubled origin}
\GSfield{Kind; source}{construction; \GSpage{P000099}}
\GSfield{The notes}{Two affine lines are glued along the punctured line and the reader is asked to check that the diagonal is not closed.}
\GSfield{Statement}{Glue two copies of \(\mathbb A^1_k\) by the identity on \(D(t)\). The resulting prevariety has distinct origins \(0_1,0_2\). In the chart \(U_1\times U_2\cong\mathbb A^2\), its diagonal meets the chart in \(\{(t,t):t\ne0\}\), whose closure contains \((0_1,0_2)\); hence the diagonal is not closed.}
\GSfield{Reconstruction}{\ours\ The diagram is interpreted as two retained origins, not an infinitesimal double point or a nonreduced scheme.}
\GSfield{Fate}{Already part of the expository framework in \citet{serre1955coherent}.}
\GSfield{Present status}{As of 9 November 2026, the precise statement is a standard established result or definition.}
\GSfield{Assessment}{literal statement: correct; publication match: partial; verification: read the result; development: annotated; confidence: high}
\end{GSitem}
\begin{GSitem}{1965-58.13}{E}{Algebraic cycles}
\GSfield{Kind; source}{definition; \GSpage{P000100}}
\GSfield{The notes}{Formal integer combinations \(\sum a_iV_i\) of subvarieties are called cycles, with cycle classes described as more precise than cohomology.}
\GSfield{Statement}{The group \(Z_r(X)\) is the free abelian group on integral closed \(r\)-dimensional subvarieties of \(X\). One can take the direct sum over \(r\). Over \(\mathbb C\), cycle-class maps connect suitable Chow groups to cohomology or Borel--Moore homology, but rational equivalence is generally finer than homological equivalence.}
\GSfield{Reconstruction}{Irreducibility, reduced structure, and the dimension grading are supplied. The vague comparison with cohomology is a direction of information, not an assertion that the cycle-class map is injective.}
\GSfield{Fate}{Already part of the expository framework in \citet{borelserre1958riemannroch}.}
\GSfield{Present status}{As of 9 November 2026, the precise statement is a standard established result or definition.}
\GSfield{Assessment}{literal statement: correct with missing hypotheses; publication match: partial; verification: read the result; development: annotated; confidence: high}
\end{GSitem}
\begin{GSitem}{1965-58.14}{E}{Equivalence through a family over the projective line}
\GSfield{Kind; source}{definition; \GSpage{P000100}}
\GSfield{The notes}{Two cycles are declared equivalent when they are fibres at \(a,b\in\mathbb P^1\) of a cycle \(W\subset X\times\mathbb P^1\).}
\GSfield{Statement}{Rational equivalence on \(Z_r(X)\) is generated by differences of the intersection cycles \(W\cdot(X\times\{a\})-W\cdot(X\times\{b\})\), where \(W\) is an \((r+1)\)-cycle with the relevant intersections proper and \(a\ne b\). The quotient is \(CH_r(X)\). Equivalently one uses divisors of rational functions on integral \((r+1)\)-dimensional subvarieties.}
\GSfield{Reconstruction}{Intersection multiplicities and generation of an equivalence relation are supplied. Replacing \(\mathbb P^1\) by an arbitrary connected parameter curve would concern algebraic equivalence, a different relation.}
\GSfield{Fate}{Already part of the expository framework in \citet{borelserre1958riemannroch}.}
\GSfield{Present status}{As of 9 November 2026, the precise statement is a standard established result or definition.}
\GSfield{Assessment}{literal statement: correct with missing hypotheses; publication match: exact; verification: read the result; development: annotated; confidence: high}
\end{GSitem}
\begin{GSitem}{1965-58.15}{E}{Linear equivalence of divisors and the graph of a rational function}
\GSfield{Kind; source}{construction; \GSpage{P000100}}
\GSfield{The notes}{For curves \(U,V\) on a surface, the notes use a function with zeros on \(U\), poles on \(V\), and its graph in \(X\times\mathbb P^1\).}
\GSfield{Statement}{On a normal integral variety, Weil divisors \(D_0,D_\infty\) are linearly equivalent if \(D_0-D_\infty=\operatorname{div}(f)\) for a nonzero rational function \(f\). The closure of the graph of \(f:X\dashrightarrow\mathbb P^1\) realizes the corresponding rational-equivalence relation with multiplicities.}
\GSfield{Reconstruction}{One must prescribe orders of zeros and poles, not merely their supports. At indeterminacy points, the graph closure replaces the naive set \((x,f(x))\).}
\GSfield{Fate}{Already part of the expository framework in \citet{borelserre1958riemannroch}.}
\GSfield{Present status}{As of 9 November 2026, the precise statement is a standard established result or definition.}
\GSfield{Assessment}{literal statement: correct with missing hypotheses; publication match: exact; verification: read the result; development: annotated; confidence: high}
\end{GSitem}
\begin{GSitem}{1965-58.16}{F}{The unfinished additive map from K-theory to cycle classes}
\GSfield{Kind; source}{question; \GSpage{P000100}}
\GSfield{The notes}{Under the heading Chern classes, the notes write a linear map \(K(X)\to A(X)\), begin axiom 1, and cancel its continuation.}
\GSfield{Statement}{No unique map is defined. The Chern character is additive \(K_0(X)\to CH^*(X)\otimes\mathbb Q\) for smooth \(X\); the total Chern class satisfies a multiplicative Whitney formula and is not an additive map to the integral Chow group. These are distinct candidates.}
\GSfield{Reconstruction}{Neither rational coefficients nor the decisive axioms survive. It would be an unsupported repair to relabel the entry as a complete definition of the Chern character, or to supply Grothendieck--Riemann--Roch as its statement.}
\GSfield{Fate}{\citet{borelserre1958riemannroch} supplies the pre-existing distinction between Chern classes and the Chern character; it does not resolve the cancelled wording.}
\GSfield{Present status}{As of 9 November 2026, both surrounding theories are established, while the specific transformation intended in these last lines remains unidentified.}
\GSfield{Assessment}{literal statement: not determinable; publication match: partial; verification: read the result; development: annotated; confidence: low}
\end{GSitem}

\GStalk{59}{1965-10-18}{I. M. Gelfand}{P000101-P000102}{Elliptic complexes and the Hodge decomposition}
Gelfand repeats the symbol and geometric-endomorphism setup, chooses a density and metrics, and introduces the formal adjoint and Laplacian. The eigenspace cohomology assertion reaches the top of P000102. Dynin and Bernstein then intervene before the heading indicating that Gelfand resumes; these three headings are kept as separate catalogue records \citep{atiyahbott1964notes}.

\begin{GSitem}{1965-59.1}{E}{Repeated symbol and ellipticity setup}
\GSfield{Kind; source}{definition; \GSpage{P000101}}
\GSfield{The notes}{The formula \(d(f^ms)(x_0)=\sigma(d,df_{x_0})s(x_0)\) and the exact symbol sequence are written again.}
\GSfield{Statement}{For a finite differential complex, the principal symbol of an order-\(m\) operator may be normalized by the displayed \(f^m\) formula, incorporating the factor \(m!\). The complex is elliptic when the fibre symbol sequence is exact at every nonzero real covector.}
\GSfield{Reconstruction}{The condition \(f(x_0)=0\), explicit in the earlier exposition, is still needed. This is a repeat occurrence of the definitions in record 55, not a new theorem.}
\GSfield{Fate}{Already part of the expository framework in \citet{atiyahbott1964notes}.}
\GSfield{Present status}{As of 9 November 2026, the precise statement is a standard established result or definition.}
\GSfield{Assessment}{literal statement: correct with missing hypotheses; publication match: exact; verification: read the result; development: annotated; confidence: high}
\end{GSitem}
\begin{GSitem}{1965-59.2}{E}{Repeated geometric endomorphisms}
\GSfield{Kind; source}{construction; \GSpage{P000101}}
\GSfield{The notes}{A map \(f:X\to X\) and bundle maps give \(T\); the simplest example is differential forms, where \(f\) supplies the fibre map.}
\GSfield{Statement}{A smooth map \(f\) acts on differential forms by pullback \(f^*\), satisfying \(d f^*=f^*d\). More generally, bundle lifts \(\phi_i:f^*E_i\to E_i\) give a complex endomorphism only if the resulting operators commute with the differentials.}
\GSfield{Reconstruction}{The bundle diagram is reconstructed with the pullback convention. A generic bundle map covering \(f\) does not automatically produce a chain map in the required direction.}
\GSfield{Fate}{Already part of the expository framework in \citet{atiyahbott1964notes}.}
\GSfield{Present status}{As of 9 November 2026, the precise statement is a standard established result or definition.}
\GSfield{Assessment}{literal statement: correct with missing hypotheses; publication match: exact; verification: read the result; development: annotated; confidence: high}
\end{GSitem}
\begin{GSitem}{1965-59.3}{E}{Repeated finite-dimensionality theorem}
\GSfield{Kind; source}{theorem; \GSpage{P000101}}
\GSfield{The notes}{Theorem 1 says that the elliptic complex has finite-dimensional homology.}
\GSfield{Statement}{A finite elliptic complex on a closed compact smooth manifold has finite-dimensional smooth cohomology in every degree.}
\GSfield{Reconstruction}{Compactness and absence of boundary are implicit in this continuation, not consequences of the symbol definition. The notebook's homology is the cohomology of its degree-raising complex.}
\GSfield{Fate}{Already part of the expository framework in \citet{atiyahbott1964notes}.}
\GSfield{Present status}{As of 9 November 2026, the precise statement is a standard established result or definition.}
\GSfield{Assessment}{literal statement: correct with missing hypotheses; publication match: exact; verification: read the result; development: annotated; confidence: high}
\end{GSitem}
\begin{GSitem}{1965-59.4}{E}{Densities, duality, and formal adjoints}
\GSfield{Kind; source}{construction; \GSpage{P000101}}
\GSfield{The notes}{A nowhere vanishing top form and Hermitian structures define a duality map \(\alpha\), and \(d^*=\alpha^{-1}d'\alpha\) is written.}
\GSfield{Statement}{A positive smooth density on \(X\) and Hermitian metrics on the bundles define \(L^2\) pairings. The formal adjoint \(d_i^*\) is the unique differential operator satisfying \((d_i u,v)=(u,d_i^*v)\) for compactly supported smooth sections; equivalently it is obtained from the formal transpose by the metric-density identifications.}
\GSfield{Reconstruction}{A density removes an unnecessary orientability assumption. The notation suppresses different \(\alpha_i\) in adjacent degrees and the anti-linearity required for complex Hermitian duality. On a manifold with boundary a Hilbert-space adjoint additionally involves domains.}
\GSfield{Fate}{Already part of the expository framework in \citet{atiyahbott1964notes}.}
\GSfield{Present status}{As of 9 November 2026, the precise statement is a standard established result or definition.}
\GSfield{Assessment}{literal statement: correct with missing hypotheses; publication match: partial; verification: read the result; development: annotated; confidence: high}
\end{GSitem}
\begin{GSitem}{1965-59.5}{E}{The repeated Laplacian and finite eigenspaces}
\GSfield{Kind; source}{construction; \GSpage{P000101}}
\GSfield{The notes}{The notes form \(\Box=d^*d+dd^*\), define \(\Gamma_\lambda(E_i)\) by \(\Box s=\lambda s\), and infer finite-dimensional eigenspaces from ellipticity.}
\GSfield{Statement}{For a first-order elliptic complex on a closed compact manifold, \(\Box_i=d_i^*d_i+d_{i-1}d_{i-1}^*\) is a nonnegative elliptic self-adjoint operator with compact resolvent. Thus each eigenspace is finite-dimensional and consists of smooth sections.}
\GSfield{Reconstruction}{Ellipticity on a noncompact space alone would not imply this spectral conclusion. The first-order or compatible equal-order setup is retained.}
\GSfield{Fate}{Already part of the expository framework in \citet{atiyahbott1964notes}.}
\GSfield{Present status}{As of 9 November 2026, the precise statement is a standard established result or definition.}
\GSfield{Assessment}{literal statement: correct with missing hypotheses; publication match: exact; verification: read the result; development: annotated; confidence: high}
\end{GSitem}
\begin{GSitem}{1965-59.6}{E}{Commutation with the Laplacian}
\GSfield{Kind; source}{claim; \GSpage{P000101}}
\GSfield{The notes}{The identity \(\Box d=d\Box\) is displayed and used to construct the complex on each eigenspace.}
\GSfield{Statement}{On smooth sections of a differential complex, \(d_{i+1}d_i=0\) gives \(\Box_{i+1}d_i=d_i\Box_i\). Hence \(d_i\) maps the \(\lambda\)-eigenspace in degree \(i\) into that in degree \(i+1\). The analogous adjoint identity also holds.}
\GSfield{Reconstruction}{The equality is degreewise. It does not imply that an arbitrary chain endomorphism \(T\) commutes with \(\Box\).}
\GSfield{Fate}{Already part of the expository framework in \citet{atiyahbott1964notes}.}
\GSfield{Present status}{As of 9 November 2026, the precise statement is a standard established result or definition.}
\GSfield{Assessment}{literal statement: correct; publication match: exact; verification: read the result; development: annotated; confidence: high}
\end{GSitem}
\begin{GSitem}{1965-59.7}{E}{Cohomology of the eigenspace complexes}
\GSfield{Kind; source}{theorem; \GSpage{P000102}}
\GSfield{The notes}{The Hodge assertion is quoted without proof: the eigenspace complex has the original cohomology at \(\lambda=0\) and zero cohomology at \(\lambda\ne0\).}
\GSfield{Statement}{Under the preceding closed compact elliptic-complex assumptions, \(H^i(\Gamma_\lambda(E^\bullet))=0\) for \(\lambda>0\), and inclusion identifies \(\Gamma_0(E_i)\) with \(H^i(E^\bullet)\).}
\GSfield{Reconstruction}{The two conclusions have different proofs: an algebraic contracting homotopy for positive \(\lambda\), and the global Hodge decomposition for the harmonic representatives.}
\GSfield{Fate}{Already part of the expository framework in \citet{atiyahbott1964notes}.}
\GSfield{Present status}{As of 9 November 2026, the precise statement is a standard established result or definition.}
\GSfield{Assessment}{literal statement: correct; publication match: exact; verification: read the result; development: annotated; confidence: high}
\end{GSitem}

\GStalk{60}{1965-10-18}{A. S. Dynin}{P000102}{Intervention on positive-eigenvalue cohomology}
The short intervention insists that nonzero eigenspaces give zero cohomology. The written justification then identifies cohomology with the kernel of a shifted Laplacian and asks about reapplying Hodge theory. That equation is not a valid proof: the shifted operator is not the nonnegative Hodge Laplacian of the original complex. Bernstein's immediately following intervention gives the direct contracting-homotopy argument.

\begin{GSitem}{1965-60.1}{E}{Acyclicity and the invalid shifted-Laplacian shortcut}
\GSfield{Kind; source}{claim; \GSpage{P000102}}
\GSfield{The notes}{Dynin says the nonzero-\(\lambda\) cohomology is zero, but the boxed formula reads \(H^i(\Gamma_\lambda(E))=\ker(\Box-\lambda)\), followed by a proposed Hodge argument.}
\GSfield{Statement}{For the Hodge Laplacian, the positive-\(\lambda\) eigenspace complex is acyclic. It is not generally isomorphic in cohomology to \(\ker(\Box-\lambda)\): this kernel is the entire, often nonzero, \(\lambda\)-eigenspace. The correct contraction is \(d^*/\lambda\).}
\GSfield{Reconstruction}{\ours\ One cannot replace \(\Box\) by \(\Box-\lambda\) in the harmonic-representative theorem while leaving the differential and inner product unchanged. The source's intended conclusion is correct, while the displayed identification is false.}
\GSfield{Fate}{Already part of the expository framework in \citet{atiyahbott1964notes}.}
\GSfield{Present status}{As of 9 November 2026, the precise statement is a standard established result or definition.}
\GSfield{Assessment}{literal statement: false as written; publication match: partial; verification: read the result; development: annotated; confidence: high}
\end{GSitem}

\GStalk{62}{1965-10-18}{I. N. Bernstein}{P000102}{Intervention giving the contracting homotopy}
Bernstein's named intervention follows Dynin and precedes the words indicating Gelfand's return. It derives the vanishing of positive-eigenvalue cohomology directly from \(\lambda s=dd^*s\). Catalogue record 62 is therefore placed before record 61 in actual source order.

\begin{GSitem}{1965-62.1}{E}{Bernstein's direct proof of acyclicity}
\GSfield{Kind; source}{method; \GSpage{P000102}}
\GSfield{The notes}{For \(s\in\Gamma_\lambda(E_i)\), \(ds=0\), and \(\lambda\ne0\), the notes write \(s=d(d^*s/\lambda)\); they also state that \(d^*s\) is a smooth eigenfunction.}
\GSfield{Statement}{If \(\Box s=\lambda s\), \(\lambda>0\), and \(ds=0\), then \(\lambda s=(dd^*+d^*d)s=dd^*s\). Thus \(s=d(d^*s/\lambda)\); since \(d^*\Box=\Box d^*\), the primitive remains in the \(\lambda\)-eigenspace. Elliptic regularity ensures smoothness of the eigensections.}
\GSfield{Reconstruction}{This is a complete elementary proof inside the already established spectral setup. It does not establish the existence or completeness of that setup by itself.}
\GSfield{Fate}{The argument is an exposition of the standard Hodge-complex contraction used in \citet{atiyahbott1964notes}; no independent publication novelty is inferred from the intervention.}
\GSfield{Present status}{As of 9 November 2026, the precise statement is a standard established result or definition.}
\GSfield{Assessment}{literal statement: correct; publication match: exact; verification: read the result; development: annotated; confidence: high}
\end{GSitem}

\GStalk{61}{1965-10-18}{I. M. Gelfand}{P000102-P000104}{Harmonic compression, trace regularization, and pseudodifferential operators}
The heading indicating Gelfand's return comes after both interventions. He compresses the chain map to Laplacian eigenspaces, obtains alternating trace cancellation, introduces the weighted trace series, and begins an oscillatory definition of pseudodifferential operators. That definition continues into the top of P000104, above the next dated heading. The argument follows the 1964 Harvard notes, especially their spectral and pseudodifferential sections \citep{atiyahbott1964notes,kohnnirenberg1965algebra}.

\begin{GSitem}{1965-61.1}{E}{The harmonic subcomplex and spectral expansion}
\GSfield{Kind; source}{theorem; \GSpage{P000102}}
\GSfield{The notes}{The zero-eigenvalue case is attributed to the Hodge proof, \(\Gamma(E)\) is written as a sum of eigenspaces, and \(\Box s=0\) implies \(ds=0\).}
\GSfield{Statement}{For a closed compact first-order elliptic complex, \(L^2(E_i)\) is the Hilbert direct sum of its finite-dimensional Laplacian eigenspaces. A harmonic section satisfies both \(d_is=0\) and \(d_{i-1}^*s=0\), and represents a unique cohomology class. A smooth section has a convergent spectral expansion in every Sobolev norm.}
\GSfield{Reconstruction}{The displayed sum is not an algebraic direct sum: general smooth sections need infinitely many eigencomponents. Smooth spectral convergence uses elliptic estimates.}
\GSfield{Fate}{Already part of the expository framework in \citet{atiyahbott1964notes}.}
\GSfield{Present status}{As of 9 November 2026, the precise statement is a standard established result or definition.}
\GSfield{Assessment}{literal statement: correct with missing hypotheses; publication match: exact; verification: read the result; development: annotated; confidence: high}
\end{GSitem}
\begin{GSitem}{1965-61.2}{E}{Compression of a chain map to each eigenspace}
\GSfield{Kind; source}{construction; \GSpage{P000102}}
\GSfield{The notes}{The definition \(T_\lambda=\pi_\lambda T i_\lambda\) uses inclusion and orthogonal projection; the zero block is identified with the induced cohomology map.}
\GSfield{Statement}{Let \(T\) preserve the smooth complex but not necessarily the metric. Since the spectral projections commute with \(d\), \(T_\lambda=\pi_\lambda T i_\lambda\) is a finite-dimensional chain endomorphism of the \(\lambda\)-complex. On the harmonic complex, \(T_0\) represents \(H^\bullet(T)\).}
\GSfield{Reconstruction}{Compression does not require \(T\Gamma_\lambda\subset\Gamma_\lambda\). The equality with the cohomological action is after the Hodge identification, not literal preservation of each harmonic section by \(T\).}
\GSfield{Fate}{Already part of the expository framework in \citet{atiyahbott1964notes}.}
\GSfield{Present status}{As of 9 November 2026, the precise statement is a standard established result or definition.}
\GSfield{Assessment}{literal statement: correct; publication match: exact; verification: read the result; development: annotated; confidence: high}
\end{GSitem}
\begin{GSitem}{1965-61.3}{E}{Hopf trace cancellation on eigenspaces}
\GSfield{Kind; source}{theorem; \GSpage{P000102-P000103}}
\GSfield{The notes}{The alternating chain trace equals the alternating cohomology trace. The zero block contributes \(L(T,E)\), and every nonzero block contributes zero.}
\GSfield{Statement}{For a finite-dimensional bounded complex and a chain endomorphism \(S\), \(\sum_i(-1)^i\operatorname{tr}(S|C^i)=\sum_i(-1)^i\operatorname{tr}(H^i(S))\). Applying this to \(T_\lambda\) gives supertrace \(0\) for \(\lambda>0\), and \(L(T,E)\) for \(\lambda=0\).}
\GSfield{Reconstruction}{Finiteness is essential to this elementary trace identity. It is applied before the infinite sum is formed.}
\GSfield{Fate}{Already part of the expository framework in \citet{atiyahbott1964notes}.}
\GSfield{Present status}{As of 9 November 2026, the precise statement is a standard established result or definition.}
\GSfield{Assessment}{literal statement: correct; publication match: exact; verification: read the result; development: annotated; confidence: high}
\end{GSitem}
\begin{GSitem}{1965-61.4}{E}{The weighted trace series and its half-plane of convergence}
\GSfield{Kind; source}{construction; \GSpage{P000103}}
\GSfield{The notes}{The regularized trace is \(\operatorname{tr}^{z}(E_i)=\sum_\lambda(1+\lambda)^{-z}\operatorname{tr}(T_\lambda|E_i)\), asserted to converge for \(\operatorname{Re}z>\dim X/2\).}
\GSfield{Statement}{If \(X\) is closed of real dimension \(n\), \(\Box_i\) is a nonnegative elliptic operator of order two, and \(T_i\) is bounded on \(L^2(E_i)\), then \(Z_i(z)=\operatorname{tr}(T_i(1+\Box_i)^{-z})\) is holomorphic for \(\operatorname{Re}z>n/2\). Its spectral series is absolutely convergent there.}
\GSfield{Reconstruction}{\ours\ The bound follows from \(|\operatorname{tr}T_\lambda|\le\|T\|\dim\Gamma_\lambda\) and the order-two eigenvalue count. Smooth bundle-weighted pullback by a diffeomorphism is \(L^2\)-bounded; arbitrary smooth maps need a different continuity argument and cannot simply inherit this estimate.}
\GSfield{Fate}{Already part of the expository framework in \citet{atiyahbott1964notes}.}
\GSfield{Present status}{As of 9 November 2026, the precise statement is a standard established result or definition.}
\GSfield{Assessment}{literal statement: correct with missing hypotheses; publication match: exact; verification: read the result; development: annotated; confidence: high}
\end{GSitem}
\begin{GSitem}{1965-61.5}{E}{Constancy of the regularized supertrace}
\GSfield{Kind; source}{claim; \GSpage{P000103}}
\GSfield{The notes}{The equality \(L(T,E)=\sum_i(-1)^i\operatorname{tr}^{z}(E_i)\) is written for all \(z\), before continuation is proved.}
\GSfield{Statement}{In the common convergence half-plane, absolute convergence permits summation over eigenspaces and gives \(\sum_i(-1)^iZ_i(z)=L(T,E)\). Any analytic continuation of this supertrace is the same constant on its connected domain.}
\GSfield{Reconstruction}{The words for all \(z\) initially mean all \(z\) in the proven half-plane. Entire continuation of the individual geometric traces requires the transversality theorem in record 63; it does not follow from the cancellation identity alone.}
\GSfield{Fate}{Already part of the expository framework in \citet{atiyahbott1964notes}.}
\GSfield{Present status}{As of 9 November 2026, the precise statement is a standard established result or definition.}
\GSfield{Assessment}{literal statement: correct with missing hypotheses; publication match: exact; verification: read the result; development: annotated; confidence: high}
\end{GSitem}
\begin{GSitem}{1965-61.6}{F}{An unassigned warning about a pole at zero}
\GSfield{Kind; source}{remark; \GSpage{P000103}}
\GSfield{The notes}{A parenthetical line mentions analytic continuation and a pole at \(z=0\), without specifying which trace function it concerns.}
\GSfield{Statement}{The entry is insufficient to identify a pole assertion. An infinite ordinary trace at a formal parameter value does not imply a pole of its analytic continuation. Under the later transverse geometric hypotheses, the individual regularized traces are entire and therefore have no pole at zero.}
\GSfield{Reconstruction}{The warning is retained as ambiguous, rather than silently read as a universal claim about all spectral zeta functions or as a claim about the already constant supertrace.}
\GSfield{Fate}{The continuation result in \citet[Section 9]{atiyahbott1964notes} determines the transverse case; it does not identify the referent of this parenthesis.}
\GSfield{Present status}{As of 9 November 2026, the parenthetical referent remains unidentified; it is not a distinct open trace-regularization problem.}
\GSfield{Assessment}{literal statement: not determinable; publication match: partial; verification: read the result; development: annotated; confidence: low}
\end{GSitem}
\begin{GSitem}{1965-61.7}{E}{The circle model and the Laplacian sign}
\GSfield{Kind; source}{construction; \GSpage{P000103}}
\GSfield{The notes}{On \(S^1\), \(d=d/dt\) is followed by \(\Box=d^2/dt^2\); a weighted composition operator is also drawn.}
\GSfield{Statement}{For the standard periodic \(L^2\) metric, \((d/dt)^*=-d/dt\), so the nonnegative Laplacian on the circle complex is \(-d^2/dt^2\), with eigenvalues \(m^2\) for \(e^{imt}\) when the period is \(2\pi\). A geometric action on both functions and one-forms must satisfy the chain-map identity.}
\GSfield{Reconstruction}{\oursComputed\ The missing minus sign is a notebook sign error relative to its own definition \(dd^*+d^*d\), not a harmless choice of the signed scalar Laplacian.}
\GSfield{Fate}{Already part of the expository framework in \citet{atiyahbott1964notes}.}
\GSfield{Present status}{As of 9 November 2026, the precise statement is a standard established result or definition.}
\GSfield{Assessment}{literal statement: false as written; publication match: partial; verification: read the result; development: annotated; confidence: high}
\end{GSitem}
\begin{GSitem}{1965-61.8}{E}{Oscillatory asymptotics for pseudodifferential operators}
\GSfield{Kind; source}{definition; \GSpage{P000103-P000104}}
\GSfield{The notes}{A map \(P:\Gamma_c(E)\to\Gamma(F)\) is called pseudodifferential when \(e^{-i\lambda g}P(e^{i\lambda g}s)\) has an asymptotic expansion with descending orders tending to minus infinity and \(dg\ne0\) on the support of \(s\).}
\GSfield{Statement}{For a properly supported classical pseudodifferential operator of order \(m\), nonstationary phases give an expansion \(e^{-i\lambda g}P(e^{i\lambda g}s)\sim\sum_{j\ge0}\lambda^{m-j}a_j(g,s)\) as \(\lambda\to+\infty\), with local \(C^\infty\) remainder estimates; the leading term is \(\sigma_m(P)(x,dg_x)s(x)\). Polyhomogeneous classes can allow a more general discrete decreasing degree sequence.}
\GSfield{Reconstruction}{Continuity, uniform remainder estimates, and localization are needed for a characterization; mere pointwise existence of formal coefficients is insufficient. The notes do not state these analytic quantifiers.}
\GSfield{Fate}{The oscillatory formulation is in the Harvard exposition \citep{atiyahbott1964notes}; the contemporary operator calculus is \citet{kohnnirenberg1965algebra}.}
\GSfield{Present status}{As of 9 November 2026, the precise statement is a standard established result or definition.}
\GSfield{Assessment}{literal statement: correct with missing hypotheses; publication match: partial; verification: read the result; development: annotated; confidence: high}
\end{GSitem}

\GStalk{63}{1965-10-26}{I. M. Gelfand}{P000104-P000105}{Regularized traces and the local fixed-point trace}
The next dated continuation begins below the end of the pseudodifferential definition on P000104. Gelfand revisits index, finite-dimensional trace cancellation, and the spectral series; he then states the entire continuation theorem for a transverse geometric map and its equality at zero with the distributional flat trace. The final computation initially calls the delta integral the number of fixed points, omitting its Jacobian weight. The weights are restored explicitly in record 65 \citep{atiyahbott1964notes}.

\begin{GSitem}{1965-63.1}{E}{The repeated positive-eigenvalue decomposition}
\GSfield{Kind; source}{construction; \GSpage{P000104}}
\GSfield{The notes}{The Laplacian, its eigenspaces, and zero versus nonzero eigenvalues are recalled before introducing the finite-dimensional operators \(T_\lambda\).}
\GSfield{Statement}{In the closed compact first-order elliptic setting, the smooth eigenspace complexes are acyclic for \(\lambda>0\), and the zero complex is the harmonic realization of cohomology. The spectrum is nonnegative, not strictly positive when cohomology is present.}
\GSfield{Reconstruction}{This is another occurrence of the Hodge decomposition; the repeated positive wording must retain the separate harmonic case.}
\GSfield{Fate}{Already part of the expository framework in \citet{atiyahbott1964notes}.}
\GSfield{Present status}{As of 9 November 2026, the precise statement is a standard established result or definition.}
\GSfield{Assessment}{literal statement: correct with missing hypotheses; publication match: exact; verification: read the result; development: annotated; confidence: high}
\end{GSitem}
\begin{GSitem}{1965-63.2}{E}{The index as the Euler characteristic of a two-term complex}
\GSfield{Kind; source}{definition; \GSpage{P000104}}
\GSfield{The notes}{The two-term elliptic complex is associated with the difference of kernel and cokernel dimensions.}
\GSfield{Statement}{For an elliptic operator \(P:\Gamma(E)\to\Gamma(F)\) on a closed compact manifold, its Fredholm index is \(\dim\ker P-\dim\operatorname{coker}P\), equal to the Euler characteristic of the two-term complex in degrees zero and one. More generally \(L(\operatorname{id},E)=\sum_i(-1)^i\dim H^i(E)\).}
\GSfield{Reconstruction}{The Sobolev completions giving a Fredholm operator do not alter its smooth kernel or cokernel. This identity does not yet compute the topological index.}
\GSfield{Fate}{Already part of the expository framework in \citet{atiyahbott1964notes}.}
\GSfield{Present status}{As of 9 November 2026, the precise statement is a standard established result or definition.}
\GSfield{Assessment}{literal statement: correct with missing hypotheses; publication match: exact; verification: read the result; development: annotated; confidence: high}
\end{GSitem}
\begin{GSitem}{1965-63.3}{E}{Compression does not presume invariance under T}
\GSfield{Kind; source}{construction; \GSpage{P000104}}
\GSfield{The notes}{The notes stress that \(T_\lambda=\pi_\lambda T i_\lambda\) can be used even when \(T\) does not commute with the Laplacian.}
\GSfield{Statement}{Orthogonal compression of a chain map to the \(\lambda\)-eigenspace complex is itself a chain map because \(d\pi_\lambda=\pi_\lambda d\). The off-diagonal spectral blocks of \(T\) need not vanish. Its compressed supertrace is zero for \(\lambda>0\) and is the cohomological Lefschetz number for \(\lambda=0\).}
\GSfield{Reconstruction}{The second displayed cancellation on P000104 omits a summation sign: cancellation is over degrees and is not a separate zero trace in each degree.}
\GSfield{Fate}{Already part of the expository framework in \citet{atiyahbott1964notes}.}
\GSfield{Present status}{As of 9 November 2026, the precise statement is a standard established result or definition.}
\GSfield{Assessment}{literal statement: false as written; publication match: exact; verification: read the result; development: annotated; confidence: high}
\end{GSitem}
\begin{GSitem}{1965-63.4}{E}{The finite-dimensional trace lemma by splitting cycles and boundaries}
\GSfield{Kind; source}{method; \GSpage{P000105}}
\GSfield{The notes}{The proof decomposes \(A_i=Z_i+C_{Z_i}\) and \(Z_i=B_i+\text{complement}\), then compares the traces.}
\GSfield{Statement}{For a finite-dimensional complex choose \(C^i=B^i\oplus H^i\oplus L^i\), with \(Z^i=B^i\oplus H^i\). The differential identifies \(L^i\) with \(B^{i+1}\). A chain map preserves the filtration by boundaries and cycles; its diagonal blocks on these two identified spaces have equal traces and cancel in the alternating sum.}
\GSfield{Reconstruction}{\ours\ The complements need not be invariant under the chain map. Only its filtration and the induced diagonal blocks are used.}
\GSfield{Fate}{Already part of the expository framework in \citet{atiyahbott1964notes}.}
\GSfield{Present status}{As of 9 November 2026, the precise statement is a standard established result or definition.}
\GSfield{Assessment}{literal statement: correct; publication match: exact; verification: read the result; development: annotated; confidence: high}
\end{GSitem}
\begin{GSitem}{1965-63.5}{E}{Eigenvalue growth and trace regularization}
\GSfield{Kind; source}{claim; \GSpage{P000105}}
\GSfield{The notes}{The series with weights \((1+\lambda)^{-z}\) is repeated, alongside \(\lambda_n\sim n^{K/\dim X}\). A star on \(T_\lambda\) ambiguously suggests the cohomology action.}
\GSfield{Statement}{For a positive self-adjoint elliptic operator of order \(m>0\) on a closed \(n\)-manifold, its positive eigenvalues counted with multiplicity satisfy \(\lambda_j\sim c j^{m/n}\), with \(c>0\). Thus a bounded multiplier \(T\) gives a trace-class weighted product for \(\operatorname{Re}z>n/m\); for the Hodge Laplacian, \(m=2\).}
\GSfield{Reconstruction}{The exponent is an order parameter, not a universal constant. The series must use the trace of the chain-space compression \(T_\lambda\), not the trace on zero cohomology, if it is to regularize each individual chain operator. The latter replacement is valid only after taking the alternating sum.}
\GSfield{Fate}{Already part of the expository framework in \citet{atiyahbott1964notes}.}
\GSfield{Present status}{As of 9 November 2026, the precise statement is a standard established result or definition.}
\GSfield{Assessment}{literal statement: correct with missing hypotheses; publication match: partial; verification: read the result; development: annotated; confidence: high}
\end{GSitem}
\begin{GSitem}{1965-63.6}{E}{Entire continuation for transverse geometric operators}
\GSfield{Kind; source}{theorem; \GSpage{P000105}}
\GSfield{The notes}{The continuation theorem says that each regularized trace is analytic on the whole plane and that its value at zero equals the flat trace; the identity operator is explicitly excluded from this argument.}
\GSfield{Statement}{Let a smooth zero-order geometric bundle lift \(T\) cover a self-map of a closed manifold whose graph is transverse to the diagonal. In the Atiyah--Bott elliptic regularization, \(\operatorname{tr}(T(1+\Box)^{-z})\), initially defined in a sufficiently far right half-plane, extends to an entire function of \(z\), with value at zero equal to the distributional flat trace of \(T\).}
\GSfield{Reconstruction}{This is a theorem special to transverse geometric traces, not a claim that ordinary spectral zeta functions are entire. The identity on a positive-dimensional manifold is not transverse. The general-map convergence half-plane is not justified merely by an \(L^2\) bound from the preceding simplified argument.}
\GSfield{Fate}{This is the continuation theorem in the spectral proof of \citet[Section 9]{atiyahbott1964notes}. The apparently strong word entire is therefore not itself a notebook error.}
\GSfield{Present status}{As of 9 November 2026, the precise statement is a standard established result or definition.}
\GSfield{Assessment}{literal statement: correct with missing hypotheses; publication match: exact; verification: read the result; development: annotated; confidence: high}
\end{GSitem}
\begin{GSitem}{1965-63.7}{E}{The distribution kernel and flat trace}
\GSfield{Kind; source}{construction; \GSpage{P000105}}
\GSfield{The notes}{For \(Ts(x)=\varphi(x)s(f(x))\), the notes use the delta kernel and integrate it on the diagonal.}
\GSfield{Statement}{Relative to a smooth density, the geometric kernel is \(K_T(x,y)=\varphi(x)\delta_{f(x)}(y)\). If the fixed points are nondegenerate, its restriction to the diagonal is defined distributionally and gives \(\operatorname{tr}^{\,b}T=\sum_{p=f(p)}\operatorname{tr}\varphi(p)/|\det_{\mathbb R}(I-df_p)|\).}
\GSfield{Reconstruction}{Transversality is what legitimizes the delta-function substitution. Arbitrary distributions cannot be restricted to the diagonal merely by writing \(K(x,x)\).}
\GSfield{Fate}{Already part of the expository framework in \citet{atiyahbott1964notes}.}
\GSfield{Present status}{As of 9 November 2026, the precise statement is a standard established result or definition.}
\GSfield{Assessment}{literal statement: correct with missing hypotheses; publication match: exact; verification: read the result; development: annotated; confidence: high}
\end{GSitem}
\begin{GSitem}{1965-63.8}{E}{The delta trace is a weighted fixed-point sum}
\GSfield{Kind; source}{claim; \GSpage{P000105}}
\GSfield{The notes}{The scalar expression \(\int\delta(x,f(x))\,dx\) is called the number of fixed points.}
\GSfield{Statement}{At isolated nondegenerate fixed points, \(\int\delta(x-f(x))\,dx=\sum_{p=f(p)}|\det(I-df_p)|^{-1}\), with the coordinate-density convention consistently applied. Unit local weights suffice for equality with an unweighted count; they are not necessary for equality of the total sums.}
\GSfield{Reconstruction}{\oursComputed\ The missing absolute Jacobian determinant is essential. It reappears on P000106, so the local source trajectory itself supports this correction. \ours\ On \(\mathbb{RP}^1\), the map \(x\mapsto3x\) has fixed points zero and infinity, with derivatives \(3\) and \(1/3\) in their respective local coordinates. Its weights \(1/2\) and \(3/2\) sum to the unweighted count \(2\), even though neither weight is one. This independently checked example prevents treating a sufficient normalization as necessary.}
\GSfield{Fate}{Already part of the expository framework in \citet{atiyahbott1964notes}.}
\GSfield{Present status}{As of 9 November 2026, the precise statement is a standard established result or definition.}
\GSfield{Assessment}{literal statement: false as written; publication match: exact; verification: read the result; development: annotated; confidence: high}
\end{GSitem}

\GStalk{64}{1965-10-26}{D. A. Kazhdan}{P000105-P000106}{Finite-type and coherent sheaves}
The five-minute intervention begins at the bottom of P000105 and continues into P000106. Kazhdan defines finite generation and coherence through relation sheaves, warns that one chosen presentation is not enough over a general ring, and states the closure and local-generation properties. Serre's coherent-sheaf paper is explicitly named \citep{serre1955coherent}.

\begin{GSitem}{1965-64.1}{E}{Sheaves of modules of finite type}
\GSfield{Kind; source}{definition; \GSpage{P000105}}
\GSfield{The notes}{Near each point finitely many sections generate every stalk; they define a map \(\mathcal O_U^n\to M_U\).}
\GSfield{Statement}{An \(\mathcal O_X\)-module \(M\) is of finite type if every point has a neighbourhood \(U\) on which there is a surjection \(\mathcal O_U^n\twoheadrightarrow M|_U\) for some finite \(n\). Surjectivity is equivalent to generation on every stalk in \(U\).}
\GSfield{Reconstruction}{Stalks, not geometric fibres after tensoring with residue fields, are the modules meant by the notebook's fibre terminology.}
\GSfield{Fate}{Already part of the expository framework in \citet{serre1955coherent}.}
\GSfield{Present status}{As of 9 November 2026, the precise statement is a standard established result or definition.}
\GSfield{Assessment}{literal statement: correct; publication match: exact; verification: read the result; development: annotated; confidence: high}
\end{GSitem}
\begin{GSitem}{1965-64.2}{E}{Coherence requires finite relations for every finite local tuple}
\GSfield{Kind; source}{definition; \GSpage{P000105-P000106}}
\GSfield{The notes}{The kernel of \(\mathcal O_U^n\to M_U\) is the relation sheaf and must be of finite type for every finite set of sections. The notes say that one fixed set is insufficient, using a one-point non-Noetherian-ring example.}
\GSfield{Statement}{A sheaf \(M\) is coherent if it is of finite type and, for every finite local tuple of sections, the kernel of the corresponding map \(\mathcal O^n\to M\) is of finite type. One finite presentation alone implies coherence when the structure sheaf is coherent, but not over an arbitrary ring.}
\GSfield{Reconstruction}{\ours\ Non-Noetherian alone is not enough for the proposed counterexample: some non-Noetherian rings are coherent. An explicit noncoherent ring is \(A=k[x_1,x_2,\ldots]/(x_1x_j:j\ge2)\); \(A\) has its trivial one-generator presentation, but the annihilator of \(x_1\) is the non-finitely-generated ideal \((x_2,x_3,\ldots)\).}
\GSfield{Fate}{Already part of the expository framework in \citet{serre1955coherent}.}
\GSfield{Present status}{As of 9 November 2026, the precise statement is a standard established result or definition.}
\GSfield{Assessment}{literal statement: correct with missing hypotheses; publication match: exact; verification: read the result; development: annotated; confidence: high}
\end{GSitem}
\begin{GSitem}{1965-64.3}{E}{Finite-type submodules of coherent sheaves}
\GSfield{Kind; source}{theorem; \GSpage{P000106}}
\GSfield{The notes}{If \(N\subset M\) is of finite type, it is said to be coherent.}
\GSfield{Statement}{If \(M\) is a coherent \(\mathcal O_X\)-module and \(N\subset M\) is a finite-type submodule sheaf, then \(N\) is coherent: every relation among sections of \(N\) is the same relation when viewed inside \(M\).}
\GSfield{Reconstruction}{Coherence of the ambient \(M\) is required. Without it, taking \(N=M\) would incorrectly assert that every finite-type sheaf is coherent.}
\GSfield{Fate}{Already part of the expository framework in \citet{serre1955coherent}.}
\GSfield{Present status}{As of 9 November 2026, the precise statement is a standard established result or definition.}
\GSfield{Assessment}{literal statement: correct with missing hypotheses; publication match: exact; verification: read the result; development: annotated; confidence: high}
\end{GSitem}
\begin{GSitem}{1965-64.4}{E}{The two-out-of-three property for coherent sheaves}
\GSfield{Kind; source}{theorem; \GSpage{P000106}}
\GSfield{The notes}{For \(0\to M'\to M\to M''\to0\), coherence of two terms is asserted to imply coherence of the third.}
\GSfield{Statement}{For an exact sequence of sheaves of modules over a sheaf of rings, if two of \(M',M,M''\) are coherent, then so is the third. Kernels, cokernels, and extensions of coherent modules remain coherent.}
\GSfield{Reconstruction}{This is the coherent-module property, not the corresponding assertion for arbitrary finite-type modules over arbitrary rings.}
\GSfield{Fate}{Already part of the expository framework in \citet{serre1955coherent}.}
\GSfield{Present status}{As of 9 November 2026, the precise statement is a standard established result or definition.}
\GSfield{Assessment}{literal statement: correct; publication match: exact; verification: read the result; development: annotated; confidence: high}
\end{GSitem}
\begin{GSitem}{1965-64.5}{E}{Generators at one stalk generate nearby}
\GSfield{Kind; source}{claim; \GSpage{P000106}}
\GSfield{The notes}{Sections generating the module at \(x\) are said to generate it in a neighbourhood of \(x\).}
\GSfield{Statement}{If \(M\) is a finite-type sheaf near \(x\), and finitely many sections defined near \(x\) generate \(M_x\), they generate \(M\) on a possibly smaller neighbourhood. Indeed the cokernel of the induced finite free map is of finite type and has zero stalk at \(x\), so it vanishes nearby.}
\GSfield{Reconstruction}{\ours\ Finite type of \(M\) must be retained. An arbitrary sheaf can have zero stalk at a point without vanishing on any neighbourhood.}
\GSfield{Fate}{Already part of the expository framework in \citet{serre1955coherent}.}
\GSfield{Present status}{As of 9 November 2026, the precise statement is a standard established result or definition.}
\GSfield{Assessment}{literal statement: correct with missing hypotheses; publication match: exact; verification: read the result; development: annotated; confidence: high}
\end{GSitem}

\GStalk{66}{1965-10-26}{S. I. Gelfand}{P000106}{The algebraic definition of K-theory}
S. Gelfand introduces the Grothendieck group of projective modules, the splitting of its exact-sequence relations, localization over local rings, and the connection with vector bundles. He also asks what changes if all finite-type modules are used over a ring denoted \(E(X)\), whose exact identity is not explained. This heading precedes the new date 1 November and the Smirnov entry; it is therefore earlier in source order than catalogue record 65.

\begin{GSitem}{1965-66.1}{E}{K0 from finitely generated projective modules}
\GSfield{Kind; source}{definition; \GSpage{P000106}}
\GSfield{The notes}{The free abelian group on projective modules over a commutative ring \(A\) is divided by the exact-sequence relations \([P]-[P_1]-[P_2]\).}
\GSfield{Statement}{For a unital ring \(A\), \(K_0(A)\) is generated by isomorphism classes of finitely generated projective \(A\)-modules, subject to \([P]=[P_1]+[P_2]\) for exact sequences \(0\to P_1\to P\to P_2\to0\). Commutativity is needed for the usual vector-bundle interpretation, but not for this additive definition.}
\GSfield{Reconstruction}{\ours\ Finite generation and isomorphism classes must be supplied. Allowing all projectives with countable sums permits an Eilenberg swindle: \(Q=\bigoplus_{n\ge0}P\cong P\oplus Q\) would force every \([P]\) to be zero.}
\GSfield{Fate}{Already part of the expository framework in \citet{borelserre1958riemannroch}.}
\GSfield{Present status}{As of 9 November 2026, the precise statement is a standard established result or definition.}
\GSfield{Assessment}{literal statement: correct with missing hypotheses; publication match: partial; verification: read the result; development: annotated; confidence: high}
\end{GSitem}
\begin{GSitem}{1965-66.2}{E}{The exact sequences of projectives split}
\GSfield{Kind; source}{claim; \GSpage{P000106}}
\GSfield{The notes}{A parenthesis writes \(P=P_1+P_2\) after the defining relation.}
\GSfield{Statement}{If \(0\to P_1\to P\to P_2\to0\) is exact and \(P_2\) is projective, the quotient map admits a section, hence \(P\cong P_1\oplus P_2\). Thus \(K_0(A)\) is the group completion of the direct-sum monoid of finitely generated projectives.}
\GSfield{Reconstruction}{\ours\ The plus sign means an isomorphism to a direct sum; no canonical splitting is asserted.}
\GSfield{Fate}{Already part of the expository framework in \citet{borelserre1958riemannroch}.}
\GSfield{Present status}{As of 9 November 2026, the precise statement is a standard established result or definition.}
\GSfield{Assessment}{literal statement: correct; publication match: partial; verification: read the result; development: annotated; confidence: high}
\end{GSitem}
\begin{GSitem}{1965-66.3}{F}{Replacing projectives by all finite-type modules}
\GSfield{Kind; source}{question; \GSpage{P000106}}
\GSfield{The notes}{The notes ask what happens if one takes \(E(X)\) and all modules of finite type; the ring is not defined.}
\GSfield{Statement}{A standard comparison, once a ring and exact category are specified, is between \(K_0(A)\) and the Grothendieck group of finitely generated modules, denoted \(G_0(A)\) in the Noetherian setting. For a regular Noetherian ring of finite global dimension, finite projective resolutions identify these groups. This does not identify the unnamed \(E(X)\) or assert equality for it.}
\GSfield{Reconstruction}{\ours\ The sheaf-versus-global-ring distinction and the exact class of modules matter. For example \(A=k[\varepsilon]/(\varepsilon^2)\) has \(K_0(A)\cong\mathbb Z\) and \(G_0(A)\cong\mathbb Z\), but the natural map sends \([A]\) to \(2[k]\) and is not an isomorphism.}
\GSfield{Fate}{\citet{borelserre1958riemannroch} distinguishes locally free and coherent-sheaf Grothendieck groups and explains the role of resolutions. It does not settle which ring the notebook abbreviation denotes.}
\GSfield{Present status}{As of 9 November 2026, the standard K0/G0 comparison is established, but the particular question involving E(X) cannot be assigned a unique answer from this entry.}
\GSfield{Assessment}{literal statement: not determinable; publication match: partial; verification: read the result; development: annotated; confidence: medium}
\end{GSitem}
\begin{GSitem}{1965-66.4}{E}{Vector bundles and projective modules of continuous functions}
\GSfield{Kind; source}{theorem; \GSpage{P000106}}
\GSfield{The notes}{The topological interpretation refers to locally trivial bundles and the sheaf of continuous functions.}
\GSfield{Statement}{For a compact Hausdorff space \(X\), global sections give an equivalence between finite-rank complex vector bundles on \(X\) and finitely generated projective modules over \(C(X,\mathbb C)\). The real version uses real bundles and real continuous functions.}
\GSfield{Reconstruction}{Compactness and the global section ring are supplied for the Serre--Swan formulation. The notes themselves mention a sheaf of rings and do not assert this equivalence for every noncompact space and every finiteness convention.}
\GSfield{Fate}{This is the established correspondence of \citet{swan1962vector}, whose bibliographic identity was checked against Swan's own bibliography.}
\GSfield{Present status}{As of 9 November 2026, the precise statement is a standard established result or definition.}
\GSfield{Assessment}{literal statement: correct with missing hypotheses; publication match: partial; verification: indirect; development: annotated; confidence: high}
\end{GSitem}
\begin{GSitem}{1965-66.5}{E}{Projective modules are locally free}
\GSfield{Kind; source}{theorem; \GSpage{P000106}}
\GSfield{The notes}{Localization at a prime is said to make a projective module free, and a projective module over a local ring is said to be free.}
\GSfield{Statement}{If \(A\) is a commutative ring and \(P\) is finitely generated projective, then \(P_{\mathfrak p}\) is free over the local ring \(A_{\mathfrak p}\) for every prime \(\mathfrak p\). Its associated sheaf is locally free of finite rank on \(\operatorname{Spec}A\); conversely every finite-rank locally free sheaf on an affine scheme corresponds to a finitely generated projective module.}
\GSfield{Reconstruction}{The finite generation required by K0 is retained throughout. Local freeness means a neighbourhood statement, stronger than merely knowing each fibre is a vector space.}
\GSfield{Fate}{Already part of the expository framework in \citet{borelserre1958riemannroch}.}
\GSfield{Present status}{As of 9 November 2026, the precise statement is a standard established result or definition.}
\GSfield{Assessment}{literal statement: correct with missing hypotheses; publication match: partial; verification: read the result; development: annotated; confidence: high}
\end{GSitem}

\GStalk{67}{1965-11-01}{S. V. Smirnov}{P000106}{Extreme positive-definite functions on nilpotent groups}
The three-line entry asks for extreme points among positive-definite functions on groups obtained from nilpotent algebras by the Campbell--Hausdorff formula. It does not state the field, topology, normalization, or centrality condition. Smirnov's 1966 paper is a strong contemporaneous lead, received on 3 December 1965, but it treats normalized central positive-definite functions over a discrete field of characteristic zero. That restriction prevents an exact identification of the unrestricted notebook question \citep{smirnov1966positive}. The earlier normalized-central-function topic 1964-13.1 is thematically related; the missing centrality here prevents identifying the two questions.

\begin{GSitem}{1965-67.1}{E}{The Campbell-Hausdorff nilpotent group}
\GSfield{Kind; source}{construction; \GSpage{P000106}}
\GSfield{The notes}{The groups are obtained from a nilpotent algebra by introducing multiplication with the Campbell--Hausdorff formula.}
\GSfield{Statement}{For a finite-dimensional nilpotent Lie algebra \(\mathfrak n\) over a field of characteristic zero, the Baker--Campbell--Hausdorff series terminates and defines a polynomial associative group law on the underlying affine space, with identity zero and inverse \(X\mapsto-X\). Over \(\mathbb R\), this is the simply connected nilpotent Lie group in exponential coordinates.}
\GSfield{Reconstruction}{The algebra must be a Lie algebra and the characteristic condition ensures the rational coefficients are defined. A topology must be chosen separately if continuity of representations is intended.}
\GSfield{Fate}{This is pre-existing nilpotent-group structure, already underlying \citet{kirillov1962unitary}; \citet{smirnov1966positive} uses the algebraic discrete-field version.}
\GSfield{Present status}{As of 9 November 2026, the precise statement is a standard established result or definition.}
\GSfield{Assessment}{literal statement: correct with missing hypotheses; publication match: partial; verification: indirect; development: annotated; confidence: high}
\end{GSitem}
\begin{GSitem}{1965-67.2}{F}{Which convex set of positive-definite functions is intended?}
\GSfield{Kind; source}{question; \GSpage{P000106}}
\GSfield{The notes}{The notes ask to describe extreme points of positive-definite functions; the parenthesis is the abbreviation \ru{непр. предст.}. The English text expands this as continuous representations.}
\GSfield{Statement}{A positive-definite function satisfies \(\sum_{j,k}c_j\overline{c_k}\varphi(g_k^{-1}g_j)\ge0\). To ask for nontrivial extreme points one normally fixes \(\varphi(e)=1\). Extremes of all normalized positive-definite functions concern pure states and irreducible cyclic representations; extremes of the normalized central subset instead concern factor characters. The notebook does not determine which set, field, or topology is intended.}
\GSfield{Reconstruction}{\ours\ The Russian abbreviation can mean continuous or irreducible and must remain unresolved. Centrality cannot be supplied solely because the later paper needs it. Without normalization, the positive-definite cone has no nonzero extreme points as a convex set, since \(\varphi=\tfrac12(0+2\varphi)\).}
\GSfield{Fate}{\citet{smirnov1966positive} classifies normalized central positive-definite functions for nilpotent algebraic groups over discrete characteristic-zero fields. Its extremals are obtained from central characters of algebraic quotients, extended by zero outside the preimage of the quotient centre, with the paper's nontriviality condition on algebraic one-parameter central subgroups. This is a precise partial match, not a solution of the broader unspecified question.}
\GSfield{Present status}{As of 9 November 2026, the 1966 original theorem has been read and its central-class restriction verified. The notebook question remains underdetermined; it is not classified as an unresolved modern research problem.}
\GSfield{Assessment}{literal statement: not determinable; publication match: partial; verification: read the result; development: annotated; confidence: medium}
\end{GSitem}

\GStalk{65}{1965-11-01}{I. M. Gelfand}{P000106-P000107}{The fixed-point trace and its real and complex applications}
After Smirnov's entry, Gelfand returns to Atiyah. He computes the delta-kernel trace with a Jacobian denominator, repeats the finite-complex trace lemma, and applies the formula to de Rham theory, complex structures, and projective transformations. The real de Rham calculation is identifiable; the complex-case notation and factors are too compressed to recover a unique asserted formula. The local inverse-map and ordinary-map conventions also conflict \citep{atiyahbott1964notes}.

\begin{GSitem}{1965-65.1}{E}{The weighted delta-kernel fixed-point trace}
\GSfield{Kind; source}{construction; \GSpage{P000106}}
\GSfield{The notes}{The operator first uses \(s(f^{-1}(x))\), but its kernel then uses \(\delta(y-f(x))\); the trace is a sum with denominator \(|1-J_\rho(f)|\).}
\GSfield{Statement}{With \(Ts(x)=\phi(x)s(f(x))\), nondegenerate fixed points give \(\operatorname{tr}^{\,b}T=\sum_{p=f(p)}\operatorname{tr}\phi(p)/|\det_{\mathbb R}(I-df_p)|\). If instead \(Ts(x)=\phi(x)s(f^{-1}(x))\), replace \(f\) by \(f^{-1}\) consistently and assume it is a diffeomorphism.}
\GSfield{Reconstruction}{The denominator bars stand for an absolute determinant in dimensions greater than one. Keeping the initial inverse while retaining the later ordinary-map kernel changes the operator and generally its local weight.}
\GSfield{Fate}{Already part of the expository framework in \citet{atiyahbott1964notes}.}
\GSfield{Present status}{As of 9 November 2026, the precise statement is a standard established result or definition.}
\GSfield{Assessment}{literal statement: false as written; publication match: partial; verification: read the result; development: annotated; confidence: high}
\end{GSitem}
\begin{GSitem}{1965-65.2}{E}{The finite-complex trace lemma repeated}
\GSfield{Kind; source}{theorem; \GSpage{P000106-P000107}}
\GSfield{The notes}{The notes again define \(H^i=\ker d/\operatorname{im}d\), require that \(T\) commute with \(d\), and equate the two alternating traces.}
\GSfield{Statement}{For a bounded complex of finite-dimensional vector spaces and a chain endomorphism \(T\), \(\sum_i(-1)^i\operatorname{tr}T_i=\sum_i(-1)^i\operatorname{tr}H^i(T)\). The Euler-characteristic identity is the specialization \(T=\operatorname{id}\).}
\GSfield{Reconstruction}{The words finite analogue on P000107 supply the intended scope; the formula is not an ordinary-trace identity for the unsmoothed infinite-dimensional section complex.}
\GSfield{Fate}{Already part of the expository framework in \citet{atiyahbott1964notes}.}
\GSfield{Present status}{As of 9 November 2026, the precise statement is a standard established result or definition.}
\GSfield{Assessment}{literal statement: correct; publication match: exact; verification: read the result; development: annotated; confidence: high}
\end{GSitem}
\begin{GSitem}{1965-65.3}{E}{Recovery of the ordinary Lefschetz formula}
\GSfield{Kind; source}{theorem; \GSpage{P000107}}
\GSfield{The notes}{The first application gives the usual Lefschetz number as the sum of fixed-point multiplicities, written as a sum of signs.}
\GSfield{Statement}{For a smooth self-map of a closed compact real manifold with only nondegenerate fixed points, the de Rham pullback has \(L(f)=\sum_q(-1)^q\operatorname{tr}(f^*|H^q_{\mathrm{dR}}(X))=\sum_{p=f(p)}\operatorname{sgn}\det(I-df_p)\).}
\GSfield{Reconstruction}{The numerator from exterior powers cancels the absolute determinant in the denominator, leaving the local index. At degenerate isolated fixed points the index is still defined, but is not given by this determinant-sign expression.}
\GSfield{Fate}{Already part of the expository framework in \citet{atiyahbott1964notes}.}
\GSfield{Present status}{As of 9 November 2026, the precise statement is a standard established result or definition.}
\GSfield{Assessment}{literal statement: correct with missing hypotheses; publication match: exact; verification: read the result; development: annotated; confidence: high}
\end{GSitem}
\begin{GSitem}{1965-65.4}{E}{Exterior-power traces give the de Rham numerator}
\GSfield{Kind; source}{claim; \GSpage{P000107}}
\GSfield{The notes}{The local trace calculation is abbreviated as \(\operatorname{tr}\varphi_i=\operatorname{det}-1\), followed by eigenvalue products with alternating signs.}
\GSfield{Statement}{For any endomorphism \(A\) of a finite-dimensional vector space, \[\sum_{q=0}^n(-1)^q\operatorname{tr}(\Lambda^q A)=\det(I-A)\]. With \(A=(df_p)^*\), this is the de Rham fibre supertrace.}
\GSfield{Reconstruction}{\ours\ An individual degree trace is not the determinant. The sum over degrees and the determinant argument must be restored; the identity holds without diagonalizability, by polynomial continuity or triangularization.}
\GSfield{Fate}{Already part of the expository framework in \citet{atiyahbott1964notes}.}
\GSfield{Present status}{As of 9 November 2026, the precise statement is a standard established result or definition.}
\GSfield{Assessment}{literal statement: false as written; publication match: partial; verification: read the result; development: annotated; confidence: high}
\end{GSitem}
\begin{GSitem}{1965-65.5}{F}{The unresolved complex-case specialization}
\GSfield{Kind; source}{claim; \GSpage{P000107}}
\GSfield{The notes}{The complex case mentions holomorphic forms and \(H^{i,0}\), then displays \(|1-J(f)|^2\) in a denominator, \(2\operatorname{Re}(1/(1-J_f))\), and an uncertain conjugate average.}
\GSfield{Statement}{No single equality is determined by these lines. For comparison, the standard holomorphic Lefschetz formula for a holomorphic bundle \(E\) on a compact complex manifold and a nondegenerate holomorphic map with a bundle lift is \(\sum_q(-1)^q\operatorname{tr}(H^q(T):H^q(X,E)\to H^q(X,E))=\sum_p\operatorname{tr}\phi_p/\det_{\mathbb C}(I-df_p)\).}
\GSfield{Reconstruction}{The Dolbeault degree, the bundle being resolved, the meaning of the realification, and the factor in the final average are missing or inconsistent. The displayed comparison formula is not attributed verbatim to the notebook.}
\GSfield{Fate}{The relevant holomorphic specialization is part of \citet{atiyahbott1964notes}, but this does not uniquely repair the source's bidegrees and factors.}
\GSfield{Present status}{As of 9 November 2026, the holomorphic fixed-point formula is established; the precise complex-case formula intended on this page is not recoverable with confidence.}
\GSfield{Assessment}{literal statement: not determinable; publication match: partial; verification: read the result; development: annotated; confidence: low}
\end{GSitem}
\begin{GSitem}{1965-65.6}{E}{Real and complex traces and determinants}
\GSfield{Kind; source}{claim; \GSpage{P000107}}
\GSfield{The notes}{The last sentence of the complex calculation says that the real trace is the sum of the complex trace and its conjugate; the drawing records \(|a|^2\) for multiplication by \(a\).}
\GSfield{Statement}{For a complex-linear map \(A\), its underlying real map satisfies \(\operatorname{tr}_{\mathbb R}A=2\operatorname{Re}\operatorname{tr}_{\mathbb C}A\) and \(\det_{\mathbb R}A=|\det_{\mathbb C}A|^2\). Thus \(\operatorname{Re}z=(z+\bar z)/2\), whereas the realified trace is the sum, without the factor one-half.}
\GSfield{Reconstruction}{These exact linear-algebra identities clarify the different factors, but do not choose which real or complex Lefschetz number was intended in the preceding unfinished calculation.}
\GSfield{Fate}{Already part of the expository framework in \citet{atiyahbott1964notes}.}
\GSfield{Present status}{As of 9 November 2026, the precise statement is a standard established result or definition.}
\GSfield{Assessment}{literal statement: correct; publication match: partial; verification: read the result; development: annotated; confidence: high}
\end{GSitem}
\begin{GSitem}{1965-65.7}{E}{Projective transformations and eigenlines}
\GSfield{Kind; source}{remark; \GSpage{P000107}}
\GSfield{The notes}{Fractional-linear transformations of projective space are named as an application.}
\GSfield{Statement}{If \(A\in GL_{n+1}(\mathbb C)\), fixed points of its projective action on \(\mathbb P^n\) are its invariant lines. For distinct eigenvalues these are \(n+1\) nondegenerate fixed points. The projective transformation acts trivially on \(H^{2q}(\mathbb P^n;\mathbb C)\), so its ordinary Lefschetz number is \(n+1\). Existence of an eigenline holds without the distinct-eigenvalue assumption.}
\GSfield{Reconstruction}{The generic nondegeneracy condition allows direct use of the local formula; repeated eigenvalues require a limiting or ordinary topological fixed-point argument. The notes do not record a formula for a specified transformation.}
\GSfield{Fate}{Already part of the expository framework in \citet{atiyahbott1964notes}.}
\GSfield{Present status}{As of 9 November 2026, the precise statement is a standard established result or definition.}
\GSfield{Assessment}{literal statement: correct with missing hypotheses; publication match: partial; verification: read the result; development: annotated; confidence: high}
\end{GSitem}

\GStalk{68}{1965-11-01}{D. A. Kazhdan}{P000107-P000108}{Injectives and right derived functors}
Kazhdan defines injective objects and enough injectives, then constructs right derived functors of a left exact covariant functor through their long exact sequence, vanishing on injectives, and naturality. The final square and existence statement occupy the top of P000108 before Matsaev's new dated heading. This resumes the homological-algebra material already indexed as 1964-49.8; it is not a new open problem \citep{grothendieck1957homological}.

\begin{GSitem}{1965-68.1}{E}{Injective objects and the extension property}
\GSfield{Kind; source}{definition; \GSpage{P000107}}
\GSfield{The notes}{An object \(Q\) is called injective when \(\operatorname{Hom}(-,Q)\) is exact; a diagram expresses extension of a morphism from a subobject.}
\GSfield{Statement}{In an abelian category \(\mathcal A\), \(Q\) is injective if every map \(A'\to Q\) extends across every monomorphism \(A'\hookrightarrow A\). Equivalently the contravariant functor \(\operatorname{Hom}_{\mathcal A}(-,Q)\) sends each short exact sequence to an exact sequence \(0\to\operatorname{Hom}(A'',Q)\to\operatorname{Hom}(A,Q)\to\operatorname{Hom}(A',Q)\to0\).}
\GSfield{Reconstruction}{The variance reverses the order of the sequence. The left exactness of Hom is automatic; surjectivity of the last map is the injectivity condition.}
\GSfield{Fate}{Already part of the expository framework in \citet{grothendieck1957homological}.}
\GSfield{Present status}{As of 9 November 2026, the precise statement is a standard established result or definition.}
\GSfield{Assessment}{literal statement: correct; publication match: exact; verification: publisher metadata checked; development: annotated; confidence: high}
\end{GSitem}
\begin{GSitem}{1965-68.2}{E}{Enough injectives}
\GSfield{Kind; source}{definition; \GSpage{P000107}}
\GSfield{The notes}{The category is said to have enough injectives if every object embeds into an injective; sheaves are given as the example.}
\GSfield{Statement}{An abelian category has enough injectives if for each \(A\) there is a monomorphism \(A\hookrightarrow Q\) with \(Q\) injective. The categories of sheaves of abelian groups, and of modules over a sheaf of unital rings on a space, have enough injectives.}
\GSfield{Reconstruction}{The sheaf example concerns the abelian module category, not the category of coherent sheaves by itself, which usually does not have enough injective objects within that subcategory.}
\GSfield{Fate}{Already part of the expository framework in \citet{grothendieck1957homological}.}
\GSfield{Present status}{As of 9 November 2026, the precise statement is a standard established result or definition.}
\GSfield{Assessment}{literal statement: correct with missing hypotheses; publication match: exact; verification: publisher metadata checked; development: annotated; confidence: high}
\end{GSitem}
\begin{GSitem}{1965-68.3}{E}{Right derived functors of a left exact functor}
\GSfield{Kind; source}{construction; \GSpage{P000107}}
\GSfield{The notes}{A covariant left exact functor is introduced and its derived functors are required to start with \(F^0=F\).}
\GSfield{Statement}{If \(\mathcal A,\mathcal B\) are abelian categories, \(\mathcal A\) has enough injectives, and \(F:\mathcal A\to\mathcal B\) is additive and left exact, choose an injective resolution \(A\to I^\bullet\) and define \(R^iF(A)=H^i(F(I^\bullet))\). Comparison maps make the result functorial and independent up to canonical isomorphism of the resolution.}
\GSfield{Reconstruction}{Additivity is made explicit. The superscript denotes right derived functors of a covariant left exact functor, despite the notebook heading's generic use of homology.}
\GSfield{Fate}{Already part of the expository framework in \citet{grothendieck1957homological}.}
\GSfield{Present status}{As of 9 November 2026, the precise statement is a standard established result or definition.}
\GSfield{Assessment}{literal statement: correct with missing hypotheses; publication match: exact; verification: publisher metadata checked; development: annotated; confidence: high}
\end{GSitem}
\begin{GSitem}{1965-68.4}{E}{The long exact sequence of derived functors}
\GSfield{Kind; source}{theorem; \GSpage{P000107}}
\GSfield{The notes}{For \(0\to A\to B\to C\to0\), the notes display the beginning of the long exact sequence; the last visible arrow after \(F^1(B)\) is labelled \(\delta\).}
\GSfield{Statement}{Each short exact sequence gives \(0\to F(A)\to F(B)\to F(C)\xrightarrow{\delta^0}R^1F(A)\to R^1F(B)\to R^1F(C)\xrightarrow{\delta^1}R^2F(A)\to\cdots\). Connecting maps go from \(R^iF(C)\) to \(R^{i+1}F(A)\).}
\GSfield{Reconstruction}{The arrow \(R^1F(B)\to R^1F(C)\) is induced by \(B\to C\), not a connecting map. The missing term beyond the frame is not treated as legible source text.}
\GSfield{Fate}{Already part of the expository framework in \citet{grothendieck1957homological}.}
\GSfield{Present status}{As of 9 November 2026, the precise statement is a standard established result or definition.}
\GSfield{Assessment}{literal statement: false as written; publication match: exact; verification: publisher metadata checked; development: annotated; confidence: high}
\end{GSitem}
\begin{GSitem}{1965-68.5}{E}{Naturality of the connecting maps}
\GSfield{Kind; source}{claim; \GSpage{P000108}}
\GSfield{The notes}{The two short exact sequences are related by a map, and a square is drawn with \(F^0(A)\) at its upper left but \(F^0(C')\) at its lower left; the vertical arrows point upward.}
\GSfield{Statement}{For a morphism of short exact sequences from \(0\to A\to B\to C\to0\) to \(0\to A'\to B'\to C'\to0\), naturality says \(R^{i+1}F(A\to A')\,\delta^i=\delta'^i\,R^iF(C\to C')\). The square has \(R^iF(C)\) and \(R^{i+1}F(A)\) in its top row and the primed counterparts below, with downward induced arrows.}
\GSfield{Reconstruction}{The original square's upper-left letter and the direction of its arrows disagree with the stated morphism of sequences. The modern square is a correction, not a claim that the handwritten A was C.}
\GSfield{Fate}{Already part of the expository framework in \citet{grothendieck1957homological}.}
\GSfield{Present status}{As of 9 November 2026, the precise statement is a standard established result or definition.}
\GSfield{Assessment}{literal statement: false as written; publication match: exact; verification: publisher metadata checked; development: annotated; confidence: high}
\end{GSitem}
\begin{GSitem}{1965-68.6}{E}{Existence and uniqueness through effaceable delta-functors}
\GSfield{Kind; source}{theorem; \GSpage{P000107-P000108}}
\GSfield{The notes}{The conditions include vanishing \(F^i(Q)=0\) for injectives and a final theorem that the derived functors exist and are unique when there are enough injectives.}
\GSfield{Statement}{The right derived functors form a cohomological \(\delta\)-functor with \(R^0F=F\), natural connecting maps, and \(R^iF(Q)=0\) for injective \(Q\), \(i>0\). Because every object embeds into an injective, the positive-degree functors are effaceable; hence this \(\delta\)-functor is universal and unique up to the unique isomorphism extending the identity of \(F\).}
\GSfield{Reconstruction}{Uniqueness concerns the compatible system of functors and connecting maps. A list of groups with unrelated long exact sequences would not supply the required naturality or universal property.}
\GSfield{Fate}{Already part of the expository framework in \citet{grothendieck1957homological}.}
\GSfield{Present status}{As of 9 November 2026, the precise statement is a standard established result or definition.}
\GSfield{Assessment}{literal statement: correct with missing hypotheses; publication match: exact; verification: publisher metadata checked; development: annotated; confidence: high}
\end{GSitem}

\GStalk{69}{1965-11-15}{V. I. Matsaev}{P000108-P000109}{The Atiyah--Bott fixed-point formula}
The record begins at the dated Matsaev heading on P000108, after the preceding discussion of derived functors, and ends before Kazhdan on P000109. It gives the cohomological Lefschetz number of a geometric endomorphism of an elliptic complex, complex-power regularization, and the transverse-distribution interpretation of the local fixed-point contributions. The independent near-contemporary match is Bott's November 1965 Bourbaki exposition, which itself cites the 1964 Harvard notes. \citep{bott1966report}

\begin{GSitem}{1965-69.1}{E}{Lefschetz number of a geometric endomorphism}
\GSfield{Kind; source}{definition; \GSpage{P000108}}
\GSfield{The notes}{For \(\cdots\to\Gamma(E_i)\xrightarrow d\Gamma(E_{i+1})\to\cdots\), the maps \(T_i\) commute with \(d\); the notes define \(L(T,E)=\sum_i(-1)^i\operatorname{tr}T_i^*\).}
\GSfield{Statement}{Let \(X\) be a compact smooth manifold without boundary, \((E_\bullet,d)\) a finite elliptic complex, and \(T_i\) continuous endomorphisms of its section spaces satisfying \(d_iT_i=T_{i+1}d_i\). Ellipticity makes \(H^i(E_\bullet)\) finite dimensional. The Lefschetz number is the alternating trace of the induced maps on these cohomology spaces.}
\GSfield{Reconstruction}{Compactness, absence of boundary, finiteness and ellipticity are the analytic setting required by the proof, rather than words repeated at this heading.}
\GSfield{Fate}{Established before this talk. \citep{bott1966report}}
\GSfield{Present status}{As of 9 November 2026, this is an established result in the stated setting; the notebook occurrence does not leave a separate open problem.}
\GSfield{Assessment}{literal statement: correct with missing hypotheses; publication match: exact; verification: read the result; development: none; confidence: high}
\end{GSitem}
\begin{GSitem}{1965-69.2}{E}{Transverse fixed-point formula}
\GSfield{Kind; source}{theorem; \GSpage{P000108}}
\GSfield{The notes}{The local term is \(\sigma(\rho)=\sum_i(-1)^i\operatorname{Sp}\varphi_i(\rho)/|\det(1-I_f(\rho))|\); the total is a sum over fixed points.}
\GSfield{Statement}{Assume additionally that \(T_i s(x)=\varphi_i(x)s(f(x))\), where \(f:X\to X\) is smooth and \(\varphi_i:f^*E_i\to E_i\) intertwines the complex. If every fixed point \(\rho\) is transverse, equivalently \(\det(I-df_\rho)\ne0\), then \[L(T,E)=\sum_{f(\rho)=\rho}\frac{\sum_i(-1)^i\operatorname{tr}\varphi_i(\rho)}{|\det(I-df_\rho)|}.\] There are finitely many such points on compact \(X\).}
\GSfield{Reconstruction}{The geometric form of \(T\), transversality and the fixed-point index of the outer sum are supplied explicitly. The denominator belongs to the real smooth formula; holomorphic variants package numerator and denominator differently.}
\GSfield{Fate}{Bott's Expos\'e 295 states this formula and dates the lecture to November 1965; its bibliography identifies the Atiyah--Bott Harvard notes of fall 1964. Thus later printing does not make the theorem new in this seminar. \citep{bott1966report}}
\GSfield{Present status}{As of 9 November 2026, this is an established result in the stated setting; the notebook occurrence does not leave a separate open problem.}
\GSfield{Assessment}{literal statement: correct with missing hypotheses; publication match: exact; verification: read the result; development: none; confidence: high}
\end{GSitem}
\begin{GSitem}{1965-69.3}{E}{Kernel proof and its trace-class obstruction}
\GSfield{Kind; source}{method; \GSpage{P000108}}
\GSfield{The notes}{The notes call the formal proof trivial: \(L=\sum_i(-1)^i\operatorname{tr}T_i\), with kernel \(\varphi\delta(y-f(x))\); its diagonal gives the fixed-point sum.}
\GSfield{Statement}{The finite-dimensional Euler trace identity does not by itself justify traces on infinite-dimensional section spaces. For a transverse geometric endomorphism the distribution kernel restricts to the diagonal and defines a flat trace; its value is the displayed fixed-point sum. Ordinary operator trace need not exist. A smoothing or complex-power argument is required to identify the alternating flat trace with the cohomological trace.}
\GSfield{Reconstruction}{The distinction between ordinary trace and diagonal distributional (flat) trace makes explicit the obstruction acknowledged by the subsequent regularization.}
\GSfield{Fate}{Established before this talk. \citep{bott1966report}}
\GSfield{Present status}{As of 9 November 2026, this is an established result in the stated setting; the notebook occurrence does not leave a separate open problem.}
\GSfield{Assessment}{literal statement: correct with missing hypotheses; publication match: exact; verification: read the result; development: annotated; confidence: high}
\end{GSitem}
\begin{GSitem}{1965-69.4}{E}{Complex-power regularization and eigenspace cancellation}
\GSfield{Kind; source}{construction; \GSpage{P000108}}
\GSfield{The notes}{\(T_{i,z}=(1+\Box_i)^{-z}T_i\), \(\Box=d\delta+\delta d\). A displayed eigenseries lacks a trace on its finite-dimensional summands; the alternating trace is said to continue to all \(z\).}
\GSfield{Statement}{Choose metrics and the formal adjoints \(\delta\). For \(\operatorname{Re}z\) sufficiently large, \((1+\Box_i)^{-z}T_i\) is trace class, with trace \[\sum_\lambda(1+\lambda)^{-z}\operatorname{tr}(\pi_{i,\lambda}T_i\iota_{i,\lambda}).\] The eigenspace complexes at positive \(\lambda\) are acyclic, and their projected endomorphisms have vanishing alternating trace. The remaining harmonic-space trace equals \(L(T,E)\), so the supertrace is constant on its initial domain and extends as that constant.}
\GSfield{Reconstruction}{The finite-dimensional trace in the eigenseries and the projected chain map are explicit repairs. Commutation of \(T\) with \(d\) need not imply commutation with \(\Box\); the projection formulation avoids that unwarranted assumption.}
\GSfield{Fate}{Established before this talk. \citep{bott1966report}}
\GSfield{Present status}{As of 9 November 2026, this is an established result in the stated setting; the notebook occurrence does not leave a separate open problem.}
\GSfield{Assessment}{literal statement: correct with missing hypotheses; publication match: exact; verification: read the result; development: annotated; confidence: high}
\end{GSitem}
\begin{GSitem}{1965-69.5}{E}{Off-diagonal regularity and the value at zero}
\GSfield{Kind; source}{claim; \GSpage{P000108}}
\GSfield{The notes}{The complex-power kernel is smooth away from the diagonal. With separated cutoffs \(\mu_1,\mu_2\), \(\mu_1(1+\Box)^{-z}\mu_2\) has a smooth kernel. At \(z=0\) the continued alternating trace is an integral as a number.}
\GSfield{Statement}{Complex powers of a positive elliptic operator have kernels smooth off the diagonal; localization by smooth cutoffs with disjoint supports is smoothing. Under the fixed-point transversality assumptions the regularized trace has the appropriate holomorphic continuation at zero and there equals the diagonal distributional trace. An unrestricted integral of the unsmoothed kernel is not being asserted.}
\GSfield{Reconstruction}{The unreadable support condition is reconstructed as disjoint supports from both the kernel argument and Bott's regularization. The parenthetical name is identified bibliographically with Kotake--Narasimhan; the scan alone does not settle its letters.}
\GSfield{Fate}{Established before this talk. \citep{bott1966report}}
\GSfield{Present status}{As of 9 November 2026, this is an established result in the stated setting; the notebook occurrence does not leave a separate open problem.}
\GSfield{Assessment}{literal statement: correct with missing hypotheses; publication match: partial; verification: read the result; development: none; confidence: medium}
\end{GSitem}
\begin{GSitem}{1965-69.6}{B}{Multiplying distributions conormal to transverse submanifolds}
\GSfield{Kind; source}{question; \GSpage{P000109}}
\GSfield{The notes}{When can distributions singular on transverse submanifolds be multiplied? A sufficient example is \(P\delta_S\), with \(P\) pseudodifferential; the notes prescribe special smooth approximations.}
\GSfield{Statement}{For embedded submanifolds \(S,T\subset X\) meeting transversely and properly supported classical pseudodifferential operators \(P,Q\), the conormal distributions \(P\delta_S\) and \(Q\delta_T\) have a canonical product. Their wavefront sets lie in \(N^*S\setminus0\) and \(N^*T\setminus0\); transversality excludes opposing covectors over a common point. Singular support alone, without a conormal or wavefront restriction, is insufficient to state such a rule.}
\GSfield{Reconstruction}{The wavefront formulation is a modern reconstruction of the sufficient class, not terminology attributed to the notebook. Bott's exposition supplies the transverse-distribution regularization relevant to the fixed-point argument.}
\GSfield{Fate}{The fixed-point flat-trace case is already explicit in Bott's contemporary exposition. The general wavefront multiplication criterion is given and proved by H\"ormander, Theorem 2.5.10, printed pp. 126--127, and yields the conormal sufficient class in the question. Class B records this later precise microlocal development, not a claim that transverse delta multiplication was unknown in 1965. \citep{bott1966report,hormander1971fourier}}
\GSfield{Present status}{As of 9 November 2026, this is an established result in the stated setting; the notebook occurrence does not leave a separate open problem.}
\GSfield{Assessment}{literal statement: correct with missing hypotheses; publication match: exact; verification: read the result; development: annotated; confidence: high; problem \hyperref[prob:GP-1965-02]{GP-1965-02}}
\end{GSitem}
\begin{GSitem}{1965-69.7}{E}{The transverse delta product}
\GSfield{Kind; source}{claim; \GSpage{P000109}}
\GSfield{The notes}{\(\delta(x_1)\delta(x_2)=\delta(x_1,x_2)\), accompanied by crossing lines and a comment that the graph meets the diagonal transversely.}
\GSfield{Statement}{In coordinates \((x_1,x_2,z)\), the product of the integration distributions on \(x_1=0\) and \(x_2=0\) is the distribution \(\phi\mapsto\int\phi(0,0,z)\,dz\), with the chosen coordinate densities. Tensor-product mollifiers give this limit. More generally, a change of transverse defining equations contributes the usual Jacobian determinant.}
\GSfield{Reconstruction}{\ours{} }
\GSfield{Fate}{Established before this talk. \citep{bott1966report}}
\GSfield{Present status}{As of 9 November 2026, this is an established result in the stated setting; the notebook occurrence does not leave a separate open problem.}
\GSfield{Assessment}{literal statement: correct; publication match: exact; verification: read the result; development: annotated; confidence: high; problem \hyperref[prob:GP-1965-02]{GP-1965-02}}
\end{GSitem}

\GStalk{70}{1965-11-15}{D. A. Kazhdan}{P000109-P000111}{Coherent algebraic sheaves on affine varieties}
The Kazhdan heading lies on the lower part of P000109; his proof continues through the beginning of P000111. Principal affine opens, coherence of the structure sheaf, localization of sections, and affine acyclicity are presented in the tradition of Serre's \emph{Faisceaux alg\'ebriques coh\'erents}. An incomplete affineness-converse line is preserved as incomplete. \citep{serre1955coherent}

\begin{GSitem}{1965-70.1}{E}{Principal affine opens and localization}
\GSfield{Kind; source}{construction; \GSpage{P000109}}
\GSfield{The notes}{\(X_P=\{P\ne0\}\) forms a Zariski basis, with sections \(Q/P^n\). The notes say that \(X_P\) is again an affine space.}
\GSfield{Statement}{If \(X=\operatorname{Spec}A\) is an affine variety over a field and \(P\in A\), then \(D(P)=\operatorname{Spec}A_P\), and \(\Gamma(D(P),\mathcal O_X)=A_P\). It is affine, realized by \(Pt=1\) in \(X\times\mathbb A^1\), but need not be an affine space. For example \(D(x)\subset\mathbb A^1\) is \(\mathbb G_m\).}
\GSfield{Reconstruction}{Here and below \(A\) is the coordinate ring of \(X\), a quotient of a polynomial ring in general.}
\GSfield{Fate}{Established before this talk. \citep{serre1955coherent}}
\GSfield{Present status}{As of 9 November 2026, this is an established result in the stated setting; the notebook occurrence does not leave a separate open problem.}
\GSfield{Assessment}{literal statement: false as written; publication match: exact; verification: read the result; development: annotated; confidence: high}
\end{GSitem}
\begin{GSitem}{1965-70.2}{E}{Finite type and coherent algebraic sheaves}
\GSfield{Kind; source}{definition; \GSpage{P000109}}
\GSfield{The notes}{An algebraic sheaf is an \(\mathcal O_X\)-module. Finite type means locally a quotient of a finite direct sum; coherence additionally requires every relation sheaf for a finite set of sections to be of finite type.}
\GSfield{Statement}{An \(\mathcal O_X\)-module \(\mathcal F\) is coherent if it is locally finitely generated and, for every open \(U\) and finite family of local sections, the kernel of the resulting map \(\mathcal O_U^r\to\mathcal F|_U\) is locally finitely generated. Finite generation alone is not the definition on an arbitrary ringed space.}
\GSfield{Fate}{Established before this talk. \citep{serre1955coherent}}
\GSfield{Present status}{As of 9 November 2026, this is an established result in the stated setting; the notebook occurrence does not leave a separate open problem.}
\GSfield{Assessment}{literal statement: correct; publication match: exact; verification: read the result; development: none; confidence: high}
\end{GSitem}
\begin{GSitem}{1965-70.3}{E}{Coherence of structure sheaves by Noetherianity}
\GSfield{Kind; source}{theorem; \GSpage{P000109-P000110}}
\GSfield{The notes}{The relation module of \((Q_1/P^n,\ldots,Q_r/P^n)\) is reduced to a kernel over \(k[x_1,\ldots,x_n]\), which is finitely generated by Hilbert's theorem.}
\GSfield{Statement}{A finitely generated algebra over a field is Noetherian. Kernels of maps between finite free modules are finitely generated, and localization commutes with their kernels. Consequently \(\mathcal O_X\) is coherent on an affine finite-type variety and, by locality, on every locally Noetherian variety.}
\GSfield{Reconstruction}{For a general affine variety replace the polynomial ring by its defining quotient; the kernel generates the relation sheaf, rather than being a fibre without localization.}
\GSfield{Fate}{Established before this talk. \citep{serre1955coherent}}
\GSfield{Present status}{As of 9 November 2026, this is an established result in the stated setting; the notebook occurrence does not leave a separate open problem.}
\GSfield{Assessment}{literal statement: correct with missing hypotheses; publication match: exact; verification: read the result; development: none; confidence: high}
\end{GSitem}
\begin{GSitem}{1965-70.4}{E}{Closed subvarieties and quotient sheaves}
\GSfield{Kind; source}{construction; \GSpage{P000110}}
\GSfield{The notes}{For closed \(Y\subset X\), the ideal sheaf \(J_Y\) is coherent and \(\mathcal O_X/J_Y=\mathcal O_Y\).}
\GSfield{Statement}{For a closed immersion \(i:Y\hookrightarrow X\) of Noetherian varieties, the defining ideal sheaf \(\mathcal I_Y\) and the quotient \(\mathcal O_X/\mathcal I_Y\) are coherent; as sheaves on \(X\), the latter is \(i_*\mathcal O_Y\). The structure sheaf on \(Y\) itself is obtained by restriction to its own ringed space.}
\GSfield{Reconstruction}{The pushforward distinguishes the two underlying spaces suppressed in the notes.}
\GSfield{Fate}{Established before this talk. \citep{serre1955coherent}}
\GSfield{Present status}{As of 9 November 2026, this is an established result in the stated setting; the notebook occurrence does not leave a separate open problem.}
\GSfield{Assessment}{literal statement: correct with missing hypotheses; publication match: exact; verification: read the result; development: none; confidence: high}
\end{GSitem}
\begin{GSitem}{1965-70.5}{E}{Global generation and acyclicity on an affine variety}
\GSfield{Kind; source}{theorem; \GSpage{P000110-P000111}}
\GSfield{The notes}{First the notes state vanishing for coherent sheaves, and later say that any sheaf is generated by sections and has trivial higher cohomology.}
\GSfield{Statement}{For \(X=\operatorname{Spec}A\), a quasi-coherent sheaf \(\widetilde M\) is generated by its global sections and satisfies \(H^q(X,\widetilde M)=0\) for \(q>0\). On a Noetherian affine variety, coherent sheaves correspond to finite \(A\)-modules. The assertion is not true for arbitrary sheaves.}
\GSfield{Reconstruction}{The coherent or quasi-coherent hypothesis must be restored to the second formulation; it is already present in the preceding line.}
\GSfield{Fate}{Established before this talk. \citep{serre1955coherent}}
\GSfield{Present status}{As of 9 November 2026, this is an established result in the stated setting; the notebook occurrence does not leave a separate open problem.}
\GSfield{Assessment}{literal statement: false as written; publication match: exact; verification: read the result; development: annotated; confidence: high}
\end{GSitem}
\begin{GSitem}{1965-70.6}{E}{Computing cohomology with an affine cover}
\GSfield{Kind; source}{method; \GSpage{P000110-P000111}}
\GSfield{The notes}{Affine acyclicity permits calculation with an affine cover; the proof ends with the vanishing of cover cohomology for a principal-open cover \((X_{p_i})\).}
\GSfield{Statement}{For a separated variety and a finite affine cover, finite intersections are affine. The Cech complex of a quasi-coherent sheaf therefore computes its sheaf cohomology. On an affine scheme a finite principal-open cover has an exact augmented Cech complex for a quasi-coherent sheaf.}
\GSfield{Reconstruction}{Separatedness, quasi-coherence and the acyclicity of intersections are the hypotheses used by the Cech-to-derived comparison.}
\GSfield{Fate}{Established before this talk. \citep{serre1955coherent}}
\GSfield{Present status}{As of 9 November 2026, this is an established result in the stated setting; the notebook occurrence does not leave a separate open problem.}
\GSfield{Assessment}{literal statement: correct with missing hypotheses; publication match: exact; verification: read the result; development: none; confidence: high}
\end{GSitem}
\begin{GSitem}{1965-70.7}{F}{The unfinished converse to affine acyclicity}
\GSfield{Kind; source}{remark; \GSpage{P000110}}
\GSfield{The notes}{The line labelled Chevalley's theorem stops at \(H^i=0\Longrightarrow\), with no consequent.}
\GSfield{Statement}{The natural comparison is Serre's affineness criterion: for a Noetherian scheme, vanishing of \(H^q(X,\mathcal F)\) for every quasi-coherent \(\mathcal F\) and every \(q>0\) characterizes affineness. The notebook does not give the quantifiers or the conclusion; it cannot establish which variant or attribution the speaker intended.}
\GSfield{Reconstruction}{The affineness criterion is offered only as a historically and contextually natural reconstruction, not as a recovered missing phrase.}
\GSfield{Fate}{Serre's affine theory supplies a close comparison, but the unfinished sentence is insufficient for an exact identification of a separate theorem called Chevalley's. \citep{serre1955coherent}}
\GSfield{Present status}{As of 9 November 2026, the standard affineness criterion is settled; the exact assertion of this unfinished occurrence remains unidentifiable. No new open-problem entry is created.}
\GSfield{Assessment}{literal statement: not determinable; publication match: partial; verification: read the result; development: none; confidence: medium}
\end{GSitem}
\begin{GSitem}{1965-70.8}{E}{Hilbert Nullstellensatz and the coordinate ring}
\GSfield{Kind; source}{theorem; \GSpage{P000110}}
\GSfield{The notes}{If \(P\) vanishes on the variety of an ideal \(J\), then \(P^k\in J\); the notes then identify the ring of sections with \(A=k[x_1,\ldots,x_n]\).}
\GSfield{Statement}{Over an algebraically closed field \(k\), the strong Nullstellensatz gives \(I(V(J))=\sqrt J\) for \(J\subset k[x_1,\ldots,x_n]\). If \(X\) is a reduced affine variety, \(\Gamma(X,\mathcal O_X)=k[x_1,\ldots,x_n]/I(X)\), not the ambient polynomial ring unless \(X=\mathbb A^n\).}
\GSfield{Reconstruction}{The field condition and the quotient implicit in the symbol \(A\) are stated explicitly.}
\GSfield{Fate}{Established before this talk. \citep{serre1955coherent}}
\GSfield{Present status}{As of 9 November 2026, this is an established result in the stated setting; the notebook occurrence does not leave a separate open problem.}
\GSfield{Assessment}{literal statement: correct with missing hypotheses; publication match: exact; verification: read the result; development: annotated; confidence: high}
\end{GSitem}
\begin{GSitem}{1965-70.9}{E}{A section supported on a principal divisor is killed by a power}
\GSfield{Kind; source}{claim; \GSpage{P000110}}
\GSfield{The notes}{Lemma 1: if a section \(s\) of a coherent sheaf on affine \(X\) vanishes on \(X_P\), then \(P^m s=0\) for some \(m\).}
\GSfield{Statement}{For \(\mathcal F=\widetilde M\) on \(\operatorname{Spec}A\), a global section \(s\in M\) restricts to zero in \(M_P\) exactly when \(P^m s=0\) for some nonnegative integer \(m\). The localized-module criterion proves the statement, and in fact does not require finite generation of \(M\).}
\GSfield{Reconstruction}{\ours{} }
\GSfield{Fate}{Established before this talk. \citep{serre1955coherent}}
\GSfield{Present status}{As of 9 November 2026, this is an established result in the stated setting; the notebook occurrence does not leave a separate open problem.}
\GSfield{Assessment}{literal statement: correct; publication match: exact; verification: read the result; development: annotated; confidence: high}
\end{GSitem}
\begin{GSitem}{1965-70.10}{E}{The ideal of the origin on the affine line}
\GSfield{Kind; source}{remark; \GSpage{P000110}}
\GSfield{The notes}{The notes identify \(k[x]\) with germs of polynomials vanishing at zero and call the quotient one dimensional.}
\GSfield{Statement}{The ideal of global polynomials vanishing at zero is \((x)\subset k[x]\), with \(k[x]/(x)\simeq k\). At the stalk the corresponding ideal is \((x)k[x]_{(x)}\subset k[x]_{(x)}\). The ideal, ambient ring and local ring must be distinguished.}
\GSfield{Fate}{Established before this talk. \citep{serre1955coherent}}
\GSfield{Present status}{As of 9 November 2026, this is an established result in the stated setting; the notebook occurrence does not leave a separate open problem.}
\GSfield{Assessment}{literal statement: false as written; publication match: exact; verification: read the result; development: annotated; confidence: high}
\end{GSitem}
\begin{GSitem}{1965-70.11}{E}{Clearing denominators in a coherent subsheaf}
\GSfield{Kind; source}{claim; \GSpage{P000110-P000111}}
\GSfield{The notes}{Lemma 2: \(\mathcal F\subset\mathcal O_X^n\); a section over \(X_P\) extends to \(X\) after multiplication by a sufficiently high power of \(P\).}
\GSfield{Statement}{Let \(X=\operatorname{Spec}A\) be Noetherian and \(\mathcal F\subset\mathcal O_X^n\) coherent. For \(s\in\Gamma(D(P),\mathcal F)\), first clear its coordinates in \((A_P)^n\), obtaining \(t\in A^n\). The image of \(t\) in \(\Gamma(X,\mathcal O_X^n/\mathcal F)\) vanishes after localization and is killed by a further power of \(P\); hence that power times \(t\) lies in \(\Gamma(X,\mathcal F)\).}
\GSfield{Reconstruction}{\ours{} The sentence on P000111 saying that the cleared section may fail to belong to \(\mathcal O_X^n\) must mathematically refer to failure to belong to \(\mathcal F\). The ambient free module already contains the cleared vector.}
\GSfield{Fate}{Established before this talk. \citep{serre1955coherent}}
\GSfield{Present status}{As of 9 November 2026, this is an established result in the stated setting; the notebook occurrence does not leave a separate open problem.}
\GSfield{Assessment}{literal statement: correct with missing hypotheses; publication match: exact; verification: read the result; development: annotated; confidence: high}
\end{GSitem}

\GStalk{71}{1965-11-22}{I. M. Gelfand}{P000111-P000113}{Value distribution and Chern forms after Bott and Chern}
The dated heading on P000111 explicitly names Bott--Chern, \emph{Acta Mathematica} 114 (1965). Gelfand moves from Picard and Nevanlinna theory to spectral boundary conditions and an explicitly doubtful Banach-space extension, then introduces bundle-valued differential forms, connections, curvature and Chern forms. The record ends immediately before the Arnold heading on P000113. \citep{bottchern1965hermitian}

\begin{GSitem}{1965-71.1}{E}{Picard omission and the Nevanlinna defect relation}
\GSfield{Kind; source}{theorem; \GSpage{P000111}}
\GSfield{The notes}{The notes say that a meromorphic quotient \(a(z)/b(z)\) omits at most two values and that the sum of its defects is at most two.}
\GSfield{Statement}{For a nonconstant meromorphic function \(f:\mathbb C\to\mathbb P^1\), at most two points of \(\mathbb P^1\) are omitted. With the usual characteristic \(T(r,f)\) and counting functions, the Nevanlinna defects satisfy \(\sum_{a\in\mathbb P^1}\delta(a,f)\le2\). The ratio model uses a reduced pair of entire functions, with no common zero; the associated projective curve is the invariant object.}
\GSfield{Reconstruction}{Nonconstancy and the projective interpretation exclude the immediate constant-function counterexample.}
\GSfield{Fate}{Established before this talk. \citep{bottchern1965hermitian}}
\GSfield{Present status}{As of 9 November 2026, this is an established result in the stated setting; the notebook occurrence does not leave a separate open problem.}
\GSfield{Assessment}{literal statement: correct with missing hypotheses; publication match: exact; verification: read the result; development: none; confidence: high}
\end{GSitem}
\begin{GSitem}{1965-71.2}{E}{Higher-dimensional hyperplane omission requires general position}
\GSfield{Kind; source}{claim; \GSpage{P000111}}
\GSfield{The notes}{An entire curve \((a,b,c)\) in three-dimensional space is said to meet every plane through zero except possibly three; a finite number is asserted in dimension \(n\).}
\GSfield{Statement}{For a linearly nondegenerate holomorphic curve \(f:\mathbb C\to\mathbb P^{n-1}\), a collection of omitted hyperplanes in general position has at most \(n\) members. The general-position requirement concerns the collection: any \(n\) defining linear forms are independent. Without nondegeneracy even a constant projective curve avoids infinitely many hyperplanes, so the notebook's unqualified formulation is false.}
\GSfield{Reconstruction}{Linear nondegeneracy and general position are the precise classical hypotheses supplied by the cited value-distribution context; no unrestricted total-hyperplane bound is claimed here.}
\GSfield{Fate}{Established before this talk. \citep{bottchern1965hermitian}}
\GSfield{Present status}{As of 9 November 2026, this is an established result in the stated setting; the notebook occurrence does not leave a separate open problem.}
\GSfield{Assessment}{literal statement: correct with missing hypotheses; publication match: exact; verification: read the result; development: annotated; confidence: high}
\end{GSitem}
\begin{GSitem}{1965-71.3}{E}{The endpoint analytic curve for Sturm--Liouville problems}
\GSfield{Kind; source}{construction; \GSpage{P000111}}
\GSfield{The notes}{The equation is \(-y''+p(x)y=\lambda y\) on \([0,\pi]\). One margin has \(y(0)=y'(0)=0\), but the subsequent displayed solution is normalized by \(y(0,\lambda)=0,\ y'(0,\lambda)=1\).}
\GSfield{Statement}{For a fixed integrable potential \(p\), the solution with initial data \((0,1)\) is entire in \(\lambda\), as are \(y(\pi,\lambda)\) and \(y'(\pi,\lambda)\). For fixed \((a,b)\ne(0,0)\), the separated right condition \(a y'(\pi)+b y(\pi)=0\) gives the characteristic function \(a y'(\pi,\lambda)+b y(\pi,\lambda)\). Its zeros are precisely the eigenvalues for the left Dirichlet condition. Initial data \((0,0)\) instead produce only the zero solution.}
\GSfield{Reconstruction}{Integrability of the fixed potential and nonzero boundary coefficients are sufficient standard assumptions. The normalization repair is internal: the scan itself supplies \(y'(0,\lambda)=1\) below the erroneous margin.}
\GSfield{Fate}{The differential-equation observation is classical and is independently checked by the Volterra integral equation; the cited Bott--Chern paper supplies the surrounding value-distribution theme, not this spectral argument. \citep{bottchern1965hermitian}}
\GSfield{Present status}{As of 9 November 2026, this is an established result in the stated setting; the notebook occurrence does not leave a separate open problem.}
\GSfield{Assessment}{literal statement: false as written; publication match: thematic only; verification: indirect; development: annotated; confidence: high}
\end{GSitem}
\begin{GSitem}{1965-71.4}{F}{A value-distribution route to spectral completeness}
\GSfield{Kind; source}{question; \GSpage{P000111-P000112}}
\GSfield{The notes}{Gelfand asks how general considerations about an analytic curve in a function space could give completeness. A theorem for a Banach-space curve missing only finitely many hyperplanes is proposed with \ru{возм., неверно}.}
\GSfield{Statement}{The notebook specifies neither a Banach space, a topology of completeness, nor a class of hyperplanes or curves sufficient to identify a spectral-completeness theorem. Moreover, a literal finite-omission assertion fails even for a linearly nondegenerate entire curve in \(\ell^2\): \[F(z)=(1,e^z,e^{2z}/2!,e^{3z}/3!,\ldots).\] Every coordinate is nonzero, so every coordinate hyperplane is omitted. If a continuous linear functional annihilates \(F\), the entire power series in \(w=e^z\) vanishes on \(\mathbb C^*\), so all its coefficients vanish.}
\GSfield{Reconstruction}{\ours{} The displayed Hilbert-space counterexample and its linear nondegeneracy proof are this task's annotation of the broad proposed assertion. They do not purport to refute an unspecified stronger hypothesis that Gelfand might have intended.}
\GSfield{Fate}{Bott--Chern does not state this Banach-space completeness theorem. Its finite-dimensional setting cannot be promoted to one. The notes themselves signal doubt; the intended spectral question remains insufficiently specified. \citep{bottchern1965hermitian}}
\GSfield{Present status}{As of 9 November 2026, the raw finite-omission proposal has the explicit counterexample above. A precise additional-hypothesis completeness problem has not been recovered, so this is F rather than an allegedly neglected open problem.}
\GSfield{Assessment}{literal statement: not determinable; publication match: thematic only; verification: indirect; development: annotated; confidence: medium}
\end{GSitem}
\begin{GSitem}{1965-71.5}{F}{Distribution of simultaneous zeros in several complex variables}
\GSfield{Kind; source}{question; \GSpage{P000112}}
\GSfield{The notes}{The notes ask about the zeros of \(n\) functions in \(n\) variables and say that no satisfactory analogue of Nevanlinna theory exists.}
\GSfield{Statement}{This is a research programme, not a single quantified conjecture: the source does not specify the domain, exhaustion, growth hypotheses, target geometry or defect inequality. Bott--Chern treats zero distribution of holomorphic sections in Hermitian bundles, providing a concrete portion of the programme. No theorem covering every possible meaning of satisfactory theory can be formulated from the sentence.}
\GSfield{Reconstruction}{The exact subject of the named paper supplies context without supplying absent hypotheses.}
\GSfield{Fate}{The explicitly cited 1965 paper already develops one version of the several-variable zero-distribution theory. Its scope does not turn this broad remark into a settled or unsettled universal conjecture. \citep{bottchern1965hermitian}}
\GSfield{Present status}{As of 9 November 2026, substantial several-variable value-distribution theories exist, but the notebook sentence has no precise residual question. No problem-index entry is assigned.}
\GSfield{Assessment}{literal statement: not determinable; publication match: partial; verification: read the result; development: none; confidence: medium}
\end{GSitem}
\begin{GSitem}{1965-71.6}{E}{Bundle-valued forms and a connection}
\GSfield{Kind; source}{construction; \GSpage{P000112}}
\GSfield{The notes}{\(A(X,E)=A(X)\otimes\Gamma(E)\), followed by a tensor-base entry read as \(A^0(E)\); the connection satisfies \(D(fs)=df\,s+fDs\).}
\GSfield{Statement}{For a smooth finite-rank vector bundle \(E\to X\), \(A^\bullet(X,E)=\Omega^\bullet(X)\otimes_{C^\infty(X)}\Gamma(E)\). A connection is a real- or complex-linear map \(D:\Gamma(E)\to\Omega^1(X,E)\) with the displayed Leibniz rule. It extends to bundle-valued forms by \(D(\alpha\otimes s)=d\alpha\otimes s+(-1)^{\deg\alpha}\alpha\wedge Ds\).}
\GSfield{Reconstruction}{The tensor base is \(A^0(X)=C^\infty(X)\), as the very next tensor formula in the notes confirms. The graded extension is made explicit because its sign controls subsequent curvature calculations.}
\GSfield{Fate}{Established before this talk. \citep{bottchern1965hermitian}}
\GSfield{Present status}{As of 9 November 2026, this is an established result in the stated setting; the notebook occurrence does not leave a separate open problem.}
\GSfield{Assessment}{literal statement: correct with missing hypotheses; publication match: exact; verification: read the result; development: annotated; confidence: high}
\end{GSitem}
\begin{GSitem}{1965-71.7}{E}{Curvature in the notebook frame convention}
\GSfield{Kind; source}{definition; \GSpage{P000112}}
\GSfield{The notes}{For \(Ds_i=\sum_j\theta_{ij}s_j\), the curvature is written \(K=d\theta-\theta\wedge\theta\). Matrix multiplication explains why \(\theta\wedge\theta\) need not vanish.}
\GSfield{Statement}{Writing the frame as the column \(s=(s_i)\) with \(Ds=\theta s\) and extending \(D\) by the graded Leibniz rule gives \(D^2s=(d\theta-\theta\wedge\theta)s\). Thus \(K\) is a matrix of two-forms representing the curvature in this frame convention. The minus sign is consistent and should not be replaced by the plus sign used with another matrix convention.}
\GSfield{Fate}{Established before this talk. \citep{bottchern1965hermitian}}
\GSfield{Present status}{As of 9 November 2026, this is an established result in the stated setting; the notebook occurrence does not leave a separate open problem.}
\GSfield{Assessment}{literal statement: correct; publication match: exact; verification: read the result; development: annotated; confidence: high}
\end{GSitem}
\begin{GSitem}{1965-71.8}{E}{Chern forms, normalization and cohomology classes}
\GSfield{Kind; source}{construction; \GSpage{P000112-P000113}}
\GSfield{The notes}{The displayed determinant on P000112 has \(2\pi/i\), and the text calls the result closed two-forms; P000113 says that it is independent of the connection.}
\GSfield{Statement}{For a complex vector bundle with connection, the total Chern form in the preceding convention is \[c(E,D)=\det\left(I+\frac{i}{2\pi}K\right)=1+c_1(E,D)+\cdots+c_r(E,D),\qquad c_j\in\Omega^{2j}(X).\] These forms are closed. Their de Rham cohomology classes are independent of the connection; their actual differential-form representatives generally vary with it.}
\GSfield{Reconstruction}{The reciprocal coefficient on P000112 is a notebook mathematical error, not a transcription correction: the scan has it, whereas P000116 and P000119 give \(i/(2\pi)\). The distinction between a form and its class repairs the independence claim.}
\GSfield{Fate}{Established before this talk. \citep{bottchern1965hermitian}}
\GSfield{Present status}{As of 9 November 2026, this is an established result in the stated setting; the notebook occurrence does not leave a separate open problem.}
\GSfield{Assessment}{literal statement: false as written; publication match: exact; verification: read the result; development: annotated; confidence: high}
\end{GSitem}
\begin{GSitem}{1965-71.9}{E}{Change of frame and invariant determinant}
\GSfield{Kind; source}{theorem; \GSpage{P000113}}
\GSfield{The notes}{\(s'=As\) gives \(dA+A\theta=\theta'A\) and \(AK=K'A\). The intermediate differentiated line lacks a prime and has the wrong graded sign.}
\GSfield{Statement}{For an invertible matrix \(A\) of smooth functions, \[\theta'=dA\,A^{-1}+A\theta A^{-1},\qquad K'=AKA^{-1}.\] Differentiating the first frame identity requires \(dA\wedge\theta+A\,d\theta=d\theta' A-\theta'\wedge dA\). Conjugacy of \(K\), together with commutativity of even scalar forms, makes \(\det(I+tK)\) independent of frame.}
\GSfield{Reconstruction}{\ours{} The displayed conjugacy is correct. The defective intervening differentiated identity is corrected as a mathematical annotation after checking the scan; it is not silently replaced in the source transcription.}
\GSfield{Fate}{Established before this talk. \citep{bottchern1965hermitian}}
\GSfield{Present status}{As of 9 November 2026, this is an established result in the stated setting; the notebook occurrence does not leave a separate open problem.}
\GSfield{Assessment}{literal statement: false as written; publication match: exact; verification: read the result; development: annotated; confidence: high}
\end{GSitem}

\GStalk{72}{1965-11-22}{V. I. Arnold}{P000113-P000116}{Singularities and stability of differentiable mappings}
Arnold's heading begins below Gelfand's frame calculation on P000113; the final comment on a continuous smooth invariant extends onto the first lines of P000116. The record distinguishes Morse, fold and cusp singularities, discusses rank strata and Boardman symbols, and reports then-current claims of Thom and Mather. Arnold's 1968 lectures contain especially close matches, but their later publication is not evidence that the older Whitney and Morse results originated in this seminar. \citep{arnold1968singularities,boardman1967singularities,mather1971nice}

\begin{GSitem}{1965-72.1}{E}{Left--right equivalence and structural stability of map germs}
\GSfield{Kind; source}{definition; \GSpage{P000113}}
\GSfield{The notes}{The source studies \(f:M^m\to N^n\), or \((\mathbb R^m,0)\to(\mathbb R^n,0)\), up to changes of variables in source and target, and asks whether small perturbations remain equivalent.}
\GSfield{Statement}{Two germs are smoothly left--right equivalent when \(g=k\circ f\circ h^{-1}\) for local diffeomorphisms \(h,k\). Local stability allows the singular point to move under perturbation. For a global map, stability means that an open neighbourhood in the Whitney smooth topology lies in its orbit under source and target diffeomorphisms. A topological variant replaces the equivalences by homeomorphisms.}
\GSfield{Reconstruction}{Local stability with movable source/target points must be distinguished from equivalence of germs fixing the origin and from global stability.}
\GSfield{Fate}{Established before this talk. \citep{arnold1968singularities}}
\GSfield{Present status}{As of 9 November 2026, this is an established result in the stated setting; the notebook occurrence does not leave a separate open problem.}
\GSfield{Assessment}{literal statement: correct with missing hypotheses; publication match: exact; verification: read the result; development: none; confidence: high}
\end{GSitem}
\begin{GSitem}{1965-72.2}{E}{Morse normal form and genericity}
\GSfield{Kind; source}{theorem; \GSpage{P000113}}
\GSfield{The notes}{For \(n=1\), nondegenerate \(d^2f\) gives a quadratic form, classified by index and signature; generic critical points are isolated, and the relevant set is open and dense for compact source.}
\GSfield{Statement}{If \(df_p=0\) and the Hessian is nonsingular, smooth local coordinates give \(f=f(p)-x_1^2-\cdots-x_\lambda^2+x_{\lambda+1}^2+\cdots+x_m^2\). Morse functions form an open dense subset of \(C^\infty(M,\mathbb R)\) for compact \(M\). Global left--right stability additionally requires pairwise distinct critical values; nondegenerate critical points alone supply local stability.}
\GSfield{Fate}{Established before this talk. \citep{arnold1968singularities}}
\GSfield{Present status}{As of 9 November 2026, this is an established result in the stated setting; the notebook occurrence does not leave a separate open problem.}
\GSfield{Assessment}{literal statement: correct with missing hypotheses; publication match: exact; verification: read the result; development: annotated; confidence: high}
\end{GSitem}
\begin{GSitem}{1965-72.3}{E}{Whitney fold and cusp normal forms for surface maps}
\GSfield{Kind; source}{theorem; \GSpage{P000113}}
\GSfield{The notes}{The three displayed models are regular, fold \((x_1^2,x_2)\), and cusp \((x_1,x_2^3-x_1x_2)\); the source attributes approximation by these types to Whitney (1955).}
\GSfield{Statement}{A generic smooth map between surfaces has only regular points, folds and cusps; its critical set is a smooth curve, with isolated cusp points. The listed normal forms represent the local stable types. For global stability of a proper map, additional multisingularity conditions are needed: distinct fold branches meet transversely, only double fold intersections occur, and cusp images are separated from other critical images.}
\GSfield{Reconstruction}{The global transversality conditions are not supplied by the list of local models. They are stated to prevent reading a local classification as a complete global criterion.}
\GSfield{Fate}{Established before this talk. \citep{arnold1968singularities,malgrange1966ideals}}
\GSfield{Present status}{As of 9 November 2026, this is an established result in the stated setting; the notebook occurrence does not leave a separate open problem.}
\GSfield{Assessment}{literal statement: correct with missing hypotheses; publication match: exact; verification: read the result; development: none; confidence: high}
\end{GSitem}
\begin{GSitem}{1965-72.4}{E}{Cusp critical curve and its image}
\GSfield{Kind; source}{claim; \GSpage{P000114}}
\GSfield{The notes}{For \((y_1,y_2)=(x_1,x_2^3-x_1x_2)\), the determinant is \(3x_2^2-x_1\). The source then writes \(y_2^2=-y_1^3\), \(y_1>0\), and initially calls the origin a fold.}
\GSfield{Statement}{The critical curve is \((x_1,x_2)=(3t^2,t)\), mapped to \((y_1,y_2)=(3t^2,-2t^3)\). Thus \(27y_2^2=4y_1^3\). The origin is a cusp; every other nearby critical point is a fold. The family \(t\mapsto t^3-Ct\) is strictly increasing for \(C<0\), and has two critical points for \(C>0\), producing the familiar one-versus-three inverse-image regions.}
\GSfield{Reconstruction}{\oursComputed{} The signs and coefficients follow directly by substitution, independently of publication identification. Both erroneous statements are visible in the scan; the later line correctly says that points away from zero are folds.}
\GSfield{Fate}{Established before this talk. \citep{arnold1968singularities}}
\GSfield{Present status}{As of 9 November 2026, this is an established result in the stated setting; the notebook occurrence does not leave a separate open problem.}
\GSfield{Assessment}{literal statement: false as written; publication match: exact; verification: read the result; development: annotated; confidence: high}
\end{GSitem}
\begin{GSitem}{1965-72.5}{A}{The perturbation of a complex square by a conjugate-linear term}
\GSfield{Kind; source}{problem; \GSpage{P000114}}
\GSfield{The notes}{The explicit exercise is \(z\mapsto z^2+\varepsilon\overline z\), with \(\varepsilon\ne0\), and the predicted singularities are folds and cusps.}
\GSfield{Statement}{As a real map of the plane, \(F_\varepsilon\) has Jacobian \(4|z|^2-|\varepsilon|^2\). Its critical locus is the circle \(|z|=|\varepsilon|/2\). Rotations and positive rescalings reduce the caustic to \(w(t)=\tfrac14e^{2it}+\tfrac12e^{-it}\): its derivative vanishes exactly at \(e^{3it}=1\), giving three ordinary cusps and folds elsewhere. The critical image is a three-cusped hypocycloid.}
\GSfield{Reconstruction}{\oursComputed{} The determinant, parameterization and cusp test are independently computed in this task. A nonzero complex \(\varepsilon\) is reduced to the same model by source/target rotations.}
\GSfield{Fate}{Arnold's 1968 exposition, Example 6, printed pp. 4--5, gives exactly this perturbation and its three-cusp picture. No pre-November-1965 publication of this exact example was documented in this check; class A follows the verified later publication by the speaker, rather than inferring antiquity from its elementary character. \citep{arnold1968singularities}}
\GSfield{Present status}{As of 9 November 2026, this is an established result in the stated setting; the notebook occurrence does not leave a separate open problem.}
\GSfield{Assessment}{literal statement: correct; publication match: exact; verification: read the result; development: annotated; confidence: high; problem \hyperref[prob:GP-1965-03]{GP-1965-03}}
\end{GSitem}
\begin{GSitem}{1965-72.6}{E}{Rank strata and their expected codimension}
\GSfield{Kind; source}{claim; \GSpage{P000114}}
\GSfield{The notes}{The notes discuss the independent minors of a rectangular rank-\(r\) matrix and call the corresponding \(\Sigma_r\) smooth submanifolds of the source.}
\GSfield{Statement}{The rank-exactly-\(r\) locus in \(\operatorname{Hom}(\mathbb R^m,\mathbb R^n)\) is smooth of codimension \((m-r)(n-r)\). For a map \(f\), its inverse image under the first jet is a smooth source submanifold of this codimension when the jet section is transverse to that matrix-rank stratum. Arbitrary maps need not have smooth rank loci.}
\GSfield{Reconstruction}{The generic transversality condition separates a smooth universal jet stratum from its possibly singular inverse image.}
\GSfield{Fate}{Established before this talk. \citep{boardman1967singularities,arnold1968singularities}}
\GSfield{Present status}{As of 9 November 2026, this is an established result in the stated setting; the notebook occurrence does not leave a separate open problem.}
\GSfield{Assessment}{literal statement: correct with missing hypotheses; publication match: exact; verification: read the result; development: annotated; confidence: high}
\end{GSitem}
\begin{GSitem}{1965-72.7}{E}{Iterated coranks and the fold--cusp Boardman symbols}
\GSfield{Kind; source}{construction; \GSpage{P000114-P000115}}
\GSfield{The notes}{A decreasing rank sequence is introduced on P000114; P000115 changes to a corank convention, calls \(\Sigma_0\) regular, and labels a cusp \(\Sigma_{10}\).}
\GSfield{Statement}{With Boardman's convention, the first entry is \(\dim\ker df\), and successive entries are obtained by intrinsic derivative or Jacobian-ideal extensions. For a generic surface map, a fold has symbol \((1,0)\), while a cusp has \((1,1,0)\). Successively restricting to smooth rank strata is an explanatory model only where the relevant transversality holds; it is not a definition for every germ.}
\GSfield{Reconstruction}{The notation switch and incorrect cusp label are notebook errors. The corrected cusp symbol is also explicit in Arnold's later exposition, section 6.}
\GSfield{Fate}{Established before this talk. \citep{boardman1967singularities,arnold1968singularities}}
\GSfield{Present status}{As of 9 November 2026, this is an established result in the stated setting; the notebook occurrence does not leave a separate open problem.}
\GSfield{Assessment}{literal statement: false as written; publication match: exact; verification: read the result; development: annotated; confidence: high}
\end{GSitem}
\begin{GSitem}{1965-72.8}{B}{Boardman strata in jet space and generic transversality}
\GSfield{Kind; source}{theorem; \GSpage{P000115}}
\GSfield{The notes}{The notes name Boardman and replace the naive repeated-rank construction by a decomposition in spaces of truncated Taylor series.}
\GSfield{Statement}{For every admissible finite Boardman symbol \(I\), the universal stratum \(\Sigma^I\) is a smooth locally closed submanifold of a sufficiently high jet space, with a universal codimension. A jet section transverse to it has a smooth inverse-image stratum of that codimension. Smooth maps can be approximated in the fine smooth topology by maps transverse to all these strata.}
\GSfield{Reconstruction}{The finite-jet order, admissibility and transversality qualifications are supplied from the primary statement.}
\GSfield{Fate}{Boardman's Theorems 6.1--6.3 give the precise general construction and transversality statements in 1967. This is a substantive later publication match, not merely a shared title. \citep{boardman1967singularities}}
\GSfield{Present status}{As of 9 November 2026, this is an established result in the stated setting; the notebook occurrence does not leave a separate open problem.}
\GSfield{Assessment}{literal statement: correct with missing hypotheses; publication match: exact; verification: read the result; development: none; confidence: high}
\end{GSitem}
\begin{GSitem}{1965-72.9}{E}{Smooth, analytic and formal equivalence are different categories}
\GSfield{Kind; source}{method; \GSpage{P000115}}
\GSfield{The notes}{The record compares \(C^\infty\) germs, analytic germs and formal power series. Formal coordinate changes have invertible first derivative.}
\GSfield{Statement}{A formal map germ fixing zero is invertible under formal composition exactly when its linear part is invertible. Analytic changes also require convergence; smooth changes may contain flat terms invisible to every finite jet and to the full Taylor series. Consequently equality or equivalence of Taylor series alone does not establish smooth or analytic equivalence without a finite-determinacy or convergence theorem.}
\GSfield{Fate}{Established before this talk. \citep{arnold1968singularities,malgrange1966ideals}}
\GSfield{Present status}{As of 9 November 2026, this is an established result in the stated setting; the notebook occurrence does not leave a separate open problem.}
\GSfield{Assessment}{literal statement: correct with missing hypotheses; publication match: exact; verification: read the result; development: annotated; confidence: high}
\end{GSitem}
\begin{GSitem}{1965-72.10}{B}{Smooth moduli and the incompleteness of Boardman symbols}
\GSfield{Kind; source}{claim; \GSpage{P000115-P000116}}
\GSfield{The notes}{The notes mention a \(\mathbb R^9\to\mathbb R^9\) example, infinitely many or continuous invariants in larger dimensions, and a smooth invariant that is not topological.}
\GSfield{Statement}{Boardman symbols are coarse invariants and do not classify left--right orbits. Thom's quadratic-map examples show continuous moduli of smooth equivalence and failure of density of smooth stable maps beginning at the pair \((9,9)\). Such smooth moduli need not be invariants of topological equivalence. The notes do not display a particular germ or prove that one germ has an infinite list of independent moduli.}
\GSfield{Reconstruction}{Continuous moduli and infinitely many equivalence classes are not conflated with infinitely many independent invariants of a specified germ.}
\GSfield{Fate}{Arnold's 1968 section 3 supplies the quadratic-moduli argument and Mather's nice-dimension calculation gives the precise density boundary. The loose phrase about infinitely many invariants is matched only partially. \citep{arnold1968singularities,mather1971nice}}
\GSfield{Present status}{As of 9 November 2026, this is an established result in the stated setting; the notebook occurrence does not leave a separate open problem.}
\GSfield{Assessment}{literal statement: correct with missing hypotheses; publication match: partial; verification: read the result; development: none; confidence: medium}
\end{GSitem}
\begin{GSitem}{1965-72.11}{B}{Thom topological classification outside infinite codimension}
\GSfield{Kind; source}{claim; \GSpage{P000115}}
\GSfield{The notes}{Thom is said to remove an infinite-codimension part of the space of power series so that the remainder is topologically stable; an unpublished proof is discussed.}
\GSfield{Statement}{The identifiable historical formulation is a decomposition of germs into countably many finite-codimension pieces consisting of topologically equivalent germs, with an exceptional part of infinite jet-codimension. This is distinct from saying that every such germ is individually an open orbit. A related global theorem asserts density of topologically stable maps from a compact smooth manifold to any smooth manifold, in every dimension pair.}
\GSfield{Reconstruction}{The distinction between an equisingularity stratum and an open stable orbit is essential. No blanket theorem for differential equations is inferred from this map-germ statement.}
\GSfield{Fate}{Arnold's 1968 Theorem 4 repeats the germ formulation and still marks its proof unpublished. Mather's 1970 notes set up the stratification and isotopy machinery and identify the compact-source density theorem as their aim. These opened sources give a partial, not exact, match to the notebook's compressed assertion. \citep{arnold1968singularities,mather1970topological}}
\GSfield{Present status}{As of 9 November 2026, topological stability of smooth maps belongs to the established Thom--Mather theory. The exact formulation and proof represented by the 1965 oral report cannot be recovered solely from the notebook; this entry does not claim to have read an original unpublished proof.}
\GSfield{Assessment}{literal statement: correct with missing hypotheses; publication match: partial; verification: read the result; development: none; confidence: medium}
\end{GSitem}
\begin{GSitem}{1965-72.12}{B}{Infinitesimal stability and the source--target linearization}
\GSfield{Kind; source}{theorem; \GSpage{P000115}}
\GSfield{The notes}{For every variation \(\varphi\), the equation \(\varphi(x)=f'(x)h(x)+k(y)\) should be solvable; this is attributed to Mather, followed by the remark that the proof is incorrect.}
\GSfield{Statement}{For a proper smooth map \(f:M\to N\), infinitesimal stability is the surjectivity of \[\Gamma(TM)\oplus\Gamma(TN)\longrightarrow\Gamma(f^*TN),\qquad(h,k)\longmapsto df(h)+k\circ f.\] In Mather's smooth theory it is equivalent to stability. Properness, the topology of maps, and the global or multigerm setting matter; an arbitrary Frechet-space inverse-function theorem does not prove the assertion.}
\GSfield{Reconstruction}{The proper-map formulation is a modern precise version. Mather's original article was verified at publisher-metadata level; the later Arnold exposition was read, so the exact Mather proof is not claimed as checked.}
\GSfield{Fate}{Mather's 1969 paper establishes the sufficient direction; Arnold's 1968 exposition reports the subsequently successful smooth proof and develops an analytic local argument. The 1965 criticism cannot be transferred to the final theorem. \citep{mather1969stability,arnold1968singularities}}
\GSfield{Present status}{As of 9 November 2026, this is an established result in the stated setting; the notebook occurrence does not leave a separate open problem.}
\GSfield{Assessment}{literal statement: correct with missing hypotheses; publication match: exact; verification: publisher metadata checked; development: none; confidence: high}
\end{GSitem}
\begin{GSitem}{1965-72.13}{F}{The unreadable finite-degree test for stability}
\GSfield{Kind; source}{remark; \GSpage{P000115}}
\GSfield{The notes}{The parenthesis says it suffices to test polynomial \(\varphi\) of degree equal to a dimension, followed by an overwritten \(\ell\)- or \(N\)-like sign.}
\GSfield{Statement}{No numerical finite-jet threshold can be recovered. Finite determinacy and algebraic tests for stability do exist, but they depend on the relevant local algebra, the source and target dimensions, the equivalence group, and the distinction between germs and multigerms. None of those choices identifies the missing integer here.}
\GSfield{Fate}{The sources confirm the surrounding finite-determinacy programme, not the unreadable bound. \citep{mather1969stability,arnold1968singularities}}
\GSfield{Present status}{As of 9 November 2026, this occurrence remains an unresolved reading and identification problem, not an open finite-determinacy conjecture.}
\GSfield{Assessment}{literal statement: not determinable; publication match: thematic only; verification: unresolved lead; development: none; confidence: low}
\end{GSitem}
\begin{GSitem}{1965-72.14}{B}{The low-dimensional stability range and generic immersions}
\GSfield{Kind; source}{claim; \GSpage{P000115}}
\GSfield{The notes}{The source lists \(\mathbb R^2\to\mathbb R^2,\mathbb R^3\), calls the \(\mathbb R^4\) target case unknown, says \(\mathbb R^5\) has no singularities, and gives density through \(m=n=8\).}
\GSfield{Statement}{Mather's nice-dimension criterion gives density of stable proper maps for the pairs \((2,2),(2,3),(2,4),(2,5)\), and for equal dimensions \((m,m)\) exactly when \(m\le8\). Generic maps of a surface into a four-manifold are already immersions locally: the rank-drop stratum has codimension \(4-2+1=3>2\); target dimension five also eliminates generic double points. An immersion and an embedding must not be identified.}
\GSfield{Reconstruction}{\oursComputed{} The generic-immersion codimension calculation is independent. It shows that the notebook's \(\mathbb R^4\) uncertainty cannot safely be interpreted as uncertainty about absence of local rank singularities.}
\GSfield{Fate}{The exact density region is calculated in Mather's paper. The opened Russian translation reproduces the criterion and explicitly warns about nonproper maps. The 1965 unknown-status annotation is historical, not current. \citep{mather1971nice}}
\GSfield{Present status}{As of 9 November 2026, this is an established result in the stated setting; the notebook occurrence does not leave a separate open problem.}
\GSfield{Assessment}{literal statement: correct with missing hypotheses; publication match: exact; verification: read the result; development: annotated; confidence: high}
\end{GSitem}

\GStalk{73}{1965-11-29}{I. M. Gelfand}{P000116}{Chern forms, Bianchi identity and curvature jets}
After the residual Arnold sentence at the top of P000116, the dated Gelfand heading restarts the connection and Chern-form construction. The new mathematical content is the Bianchi/transgression argument and a precise question about relations among curvature and its covariant derivatives. The latter is compared with the metric-jet isomorphism theorem, not classified as an untouched problem. \citep{bottchern1965hermitian,jentsch2026jet}

\begin{GSitem}{1965-73.1}{E}{Repeated moving-frame construction of Chern forms}
\GSfield{Kind; source}{construction; \GSpage{P000116}}
\GSfield{The notes}{The connection, Cartan moving frame, \(K=d\theta-\theta\wedge\theta\), and \(\det(I+iK/2\pi)\) are repeated. This time the normalization is correct.}
\GSfield{Statement}{For a smooth complex bundle with connection in the notebook frame convention, the curvature matrix defines frame-invariant even-degree forms by \(c(E,D)=\det(I+iK/(2\pi))\). Each entry of \(K\) has degree two, so the determinant is taken in the commutative algebra of even scalar forms. The construction repeats items 1965-71.6--1965-71.9 and supplies an internal check on the reciprocal coefficient on P000112.}
\GSfield{Fate}{Established before this talk. \citep{bottchern1965hermitian}}
\GSfield{Present status}{As of 9 November 2026, this is an established result in the stated setting; the notebook occurrence does not leave a separate open problem.}
\GSfield{Assessment}{literal statement: correct; publication match: exact; verification: read the result; development: none; confidence: high}
\end{GSitem}
\begin{GSitem}{1965-73.2}{E}{Bianchi identity and closed characteristic forms}
\GSfield{Kind; source}{theorem; \GSpage{P000116}}
\GSfield{The notes}{The displayed Bianchi identity is \(dK=-[K,\theta]\), described as the derivative of the curvature tensor.}
\GSfield{Statement}{For \(K=d\theta-\theta\wedge\theta\), direct exterior differentiation gives \(dK=\theta\wedge K-K\wedge\theta=-[K,\theta]\). Equivalently the covariant exterior derivative of the curvature is zero. Applying invariant polynomials and cyclicity of trace makes each component of the total Chern form closed.}
\GSfield{Fate}{Established before this talk. \citep{bottchern1965hermitian}}
\GSfield{Present status}{As of 9 November 2026, this is an established result in the stated setting; the notebook occurrence does not leave a separate open problem.}
\GSfield{Assessment}{literal statement: correct; publication match: exact; verification: read the result; development: annotated; confidence: high}
\end{GSitem}
\begin{GSitem}{1965-73.3}{E}{Connection variation and transgression}
\GSfield{Kind; source}{method; \GSpage{P000116}}
\GSfield{The notes}{The notebook says that the form is independent of the connection and sketches expressing an integral by a boundary integral so that variation in the interior does not affect it.}
\GSfield{Statement}{For two connections \(D_0,D_1\) on a fixed complex bundle, the corresponding Chern forms differ by exact forms: \(c_j(E,D_1)-c_j(E,D_0)=dT_j(D_0,D_1)\). One derives this along the affine path \(D_t\), using Bianchi and invariant-polynomial differentiation. Consequently integrals over closed cycles are independent of the connection, but pointwise forms usually are not.}
\GSfield{Reconstruction}{The transgression statement is the precise content of the sketched variation argument. For a manifold with boundary, the boundary term vanishes only under an appropriate fixed-boundary or relative condition.}
\GSfield{Fate}{Established before this talk. \citep{bottchern1965hermitian}}
\GSfield{Present status}{As of 9 November 2026, this is an established result in the stated setting; the notebook occurrence does not leave a separate open problem.}
\GSfield{Assessment}{literal statement: false as written; publication match: exact; verification: read the result; development: annotated; confidence: high}
\end{GSitem}
\begin{GSitem}{1965-73.4}{E}{Natural metric tensors from curvature and its derivatives}
\GSfield{Kind; source}{claim; \GSpage{P000116}}
\GSfield{The notes}{Gelfand asks how to find all tensors formed from \(g_{\alpha\beta}\) and its derivatives, then says they are algebraic expressions in curvature and its covariant derivatives.}
\GSfield{Statement}{A finite-order natural tensor constructed from a Riemannian metric is determined in normal coordinates by the curvature tensor and finitely many covariant derivatives, equivariantly under changes of orthonormal frame. For polynomial natural tensors, invariant theory expresses them using tensor products, permutations and contractions with the metric and its inverse (and the volume tensor in an oriented setting). A general smooth natural invariant need only be a smooth function of such data; it need not be polynomial.}
\GSfield{Reconstruction}{Finite order, naturality and the distinction between polynomial and arbitrary smooth dependence are supplied explicitly. The jet-isomorphism part is checked in the opened primary source; no classification of unrestricted nonlocal constructions is asserted.}
\GSfield{Fate}{The normal-coordinate replacement principle is classical. Jentsch's revised 2026 treatment gives an explicit modern jet-isomorphism formulation and credits the older theory; it is not cited as its first discovery. \citep{jentsch2026jet}}
\GSfield{Present status}{As of 9 November 2026, this is an established result in the stated setting; the notebook occurrence does not leave a separate open problem.}
\GSfield{Assessment}{literal statement: correct with missing hypotheses; publication match: partial; verification: read the result; development: none; confidence: high}
\end{GSitem}
\begin{GSitem}{1965-73.5}{E}{Relations among covariant derivatives of curvature}
\GSfield{Kind; source}{question; \GSpage{P000116}}
\GSfield{The notes}{The definite question is \ru{Но какие алгебр. соотн.?}; the first Bianchi-type relations among covariant derivatives are mentioned.}
\GSfield{Statement}{Fix a Euclidean vector space \(V\) and an order \(k\). Which collections \((R,\nabla R,\ldots,\nabla^kR)\) at one point occur as the curvature jet of a smooth Riemannian metric with prescribed value at that point? The algebraic curvature symmetries, differentiated first and second Bianchi identities, and differentiated Ricci commutation identities are necessary; the metric-jet isomorphism theorem makes the corresponding algebraic-jet conditions sufficient. Nonlinear lower-order curvature terms in commutation identities cannot be discarded.}
\GSfield{Reconstruction}{A finite-order pointwise realization problem makes the notebook's all-relations question precise. It does not ask for a minimal presentation of scalar invariant rings in every dimension. The exact earliest publication of the full modern nonlinear all-relations formulation has not been independently established; mathematical verification of that formulation and historical attribution are separate checks.}
\GSfield{Fate}{The finite-order realization statement is read precisely in Jentsch's Definition 1 and Theorem 3, which the author labels a classical theorem and traces in its original form to Cartan. The normal-tensor replacement and curvature-reduction theory already appears in the pre-1965 Thomas/Veblen/Schouten tradition, explicitly identified in Jany\v{s}ka's historical introduction. Thomas's 1926 paper is verified at publisher-metadata level; its original proof was not obtained. Class E reflects this classical chronology, not the 2026 date of the checked explicit exposition. \citep{jentsch2026jet,janyska2011reduction,thomas1926identities}}
\GSfield{Present status}{As of 9 November 2026, the finite-order curvature-jet realization question in this formulation is settled. Explicit reconstruction formulae continue to be useful, but no precise unsolved remainder is asserted.}
\GSfield{Assessment}{literal statement: correct; publication match: exact; verification: read the result; development: none; confidence: medium; problem \hyperref[prob:GP-1965-04]{GP-1965-04}}
\end{GSitem}
\begin{GSitem}{1965-73.6}{E}{Normal-coordinate metric coefficients and their identities}
\GSfield{Kind; source}{method; \GSpage{P000116}}
\GSfield{The notes}{The source expands the metric without a linear term and says that all coefficients in a normal system are tensors; some Taylor indices are overwritten or repeated.}
\GSfield{Statement}{In geodesic normal coordinates about \(p\), \(g_{ij}(0)\) is the prescribed inner product and \(\partial_k g_{ij}(0)=0\). Changing the frame at \(p\) changes normal coordinates linearly, so their higher coefficient arrays are tensors for that restricted frame-change action. In one standard curvature-sign convention, \(g_{ij}(x)=g_{ij}(0)-\tfrac13R_{ikjl}(p)x^kx^l+O(|x|^3)\). General coordinate partial derivatives of the metric are not tensors. The normal-coordinate identity \(g_x(x,v)=g_0(x,v)\) provides the coefficient constraints.}
\GSfield{Reconstruction}{The repeated dummy indices in the scan are not treated as a meaningful tensor formula. The precise normal-coordinate transformation group and curvature-sign convention are supplied.}
\GSfield{Fate}{Established before this talk. \citep{jentsch2026jet}}
\GSfield{Present status}{As of 9 November 2026, this is an established result in the stated setting; the notebook occurrence does not leave a separate open problem.}
\GSfield{Assessment}{literal statement: correct with missing hypotheses; publication match: exact; verification: read the result; development: annotated; confidence: high}
\end{GSitem}

\GStalk{74}{1965-11-29}{V. I. Arnold}{P000117-P000118}{Stability, preparation and intersections of invariant tori}
P000117 opens with Arnold on structural stability of flows and parameters, continues with planar desingularization and the Weierstrass--Malgrange division and preparation theorem, and P000118 introduces exact canonical transformations of \(T^n\times B^n\). The torus intersection claim is recognized as an early Arnold problem, with the notebook's graph argument recovered. Its \(2^n\) bound requires transversality; the exact-period condition also must be separated from an assumed Hamiltonian isotopy. \citep{malgrange1966ideals,chaperon1983questions,gong2023cuplength}

\begin{GSitem}{1965-74.1}{E}{Structural stability of a flow and parameter-dependent families}
\GSfield{Kind; source}{definition; \GSpage{P000117}}
\GSfield{The notes}{The notes write \(\dot x=f(x,C)\), a flow \(g^t\), and a conjugacy square; nearby systems should be related by a small change of the phase space.}
\GSfield{Statement}{For a vector field on a compact manifold, orbital structural stability means that all sufficiently close vector fields (in a specified \(C^r\) topology) are topologically equivalent by a homeomorphism sending oriented trajectories to oriented trajectories. Time-preserving conjugacy is stronger. A parameter family can cross nonstable systems even if its nearby parameter regions are structurally stable; it therefore requires bifurcation theory, not only classification of open stable classes.}
\GSfield{Reconstruction}{The equivalence category, topology and treatment of time are absent and are explicitly supplied as a conventional interpretation. The cited mapping exposition is contextual, not an exact theorem for flows.}
\GSfield{Fate}{Established before this talk. \citep{arnold1968singularities}}
\GSfield{Present status}{As of 9 November 2026, this is an established result in the stated setting; the notebook occurrence does not leave a separate open problem.}
\GSfield{Assessment}{literal statement: correct with missing hypotheses; publication match: thematic only; verification: indirect; development: none; confidence: high}
\end{GSitem}
\begin{GSitem}{1965-74.2}{F}{A centre and the unspecified degree of instability}
\GSfield{Kind; source}{claim; \GSpage{P000117}}
\GSfield{The notes}{A number of \ru{когрубости} is proposed, and a centre is said to have infinite degree of instability; concentric orbits are drawn.}
\GSfield{Statement}{The elementary content is that a planar centre is not structurally stable in the unrestricted class of smooth planar vector fields: an arbitrarily small radial component changes the closed orbits to spirals. A numerical codimension or infinite-degree statement is not determined until the equivalence group, analytic or smooth category, and permitted family are fixed. In a Hamiltonian or reversible subclass the stability question changes.}
\GSfield{Fate}{Singularity and bifurcation theory develop such codimension questions, but the precise centre claim in this notation is not identified. \citep{arnold1968singularities}}
\GSfield{Present status}{As of 9 November 2026, the unrestricted instability example is elementary; the notebook's numerical degree remains insufficiently specified. This is not converted into a current open problem.}
\GSfield{Assessment}{literal statement: not determinable; publication match: thematic only; verification: indirect; development: annotated; confidence: medium}
\end{GSitem}
\begin{GSitem}{1965-74.3}{F}{Thom strata transferred from map germs to differential equations}
\GSfield{Kind; source}{claim; \GSpage{P000117}}
\GSfield{The notes}{The display \(\vartheta=\vartheta^\infty\cup\bigcup_\gamma\vartheta^\gamma\) is followed by partly contradictory codimension annotations, finite-parameter descriptions, and a claim that points inside a stratum are diffeomorphic.}
\GSfield{Statement}{The notebook does not identify whether \(\vartheta\) consists of map germs, vector-field germs, or parameter families, or whether equivalence is smooth, orbital or topological. The Thom germ-classification statement for maps in item 1965-72.11 does not automatically give the asserted classification for differential equations. The claim that all one-parameter families lie in an infinite-codimension exceptional part is also uninterpretable without the missing quantifiers.}
\GSfield{Fate}{The opened Thom--Mather sources clarify the map-theoretic model but do not justify the transfer in this damaged passage. \citep{arnold1968singularities,mather1970topological}}
\GSfield{Present status}{As of 9 November 2026, no precise residual theorem or conjecture is extracted from these conflicting annotations. F records source indeterminacy, not absence of research.}
\GSfield{Assessment}{literal statement: not determinable; publication match: thematic only; verification: read the result; development: none; confidence: low}
\end{GSitem}
\begin{GSitem}{1965-74.4}{E}{Planar desingularization by point blow-ups}
\GSfield{Kind; source}{method; \GSpage{P000117}}
\GSfield{The notes}{Bendixson, \emph{Acta Mathematica} 1901, is cited for finite \(\sigma\)-processes applied to \(dy/dx=P/Q\), with a chart substitution \(y/x\), and a model \(x^pdy/dx=\lambda y+P(x,y)\).}
\GSfield{Statement}{For a germ of an analytic planar foliation with an isolated singularity, a point blow-up replaces the origin by its space of tangent directions; one affine chart has \(y=xu\). After removing common exceptional factors, finite reduction leads to elementary singularities, including nondegenerate singularities and saddle-nodes. The displayed equation is a local schematic model, not a complete normal form without conditions on \(p,\lambda\) and the remainder. The complex projective formulation in the notebook goes beyond the directly checked real planar discussion in Bendixson.}
\GSfield{Reconstruction}{Isolated singularity, removal of common factors and elementary local types are supplied. No uniqueness or analytic linearization of all reduced singularities is claimed.}
\GSfield{Fate}{Bendixson's primary paper is correctly identified and its repeated substitutions and local cases are read. The full complex-foliation reduction statement is a modern comparison, not an exact theorem attributed to those pages. \citep{bendixson1901courbes}}
\GSfield{Present status}{As of 9 November 2026, this is an established result in the stated setting; the notebook occurrence does not leave a separate open problem.}
\GSfield{Assessment}{literal statement: correct with missing hypotheses; publication match: partial; verification: read the result; development: none; confidence: medium}
\end{GSitem}
\begin{GSitem}{1965-74.5}{F}{The unspecified three-variable reduction problem}
\GSfield{Kind; source}{question; \GSpage{P000117}}
\GSfield{The notes}{Immediately after the planar statement the notes say \ru{Для 3-х пер. неизв.} and refer to simple singularities on a birationally equivalent manifold.}
\GSfield{Statement}{To state a three-dimensional resolution problem one must specify the rank or codimension of the foliation, the real or complex category, the allowed centres and whether weighted blow-ups are permitted, and the terminal class of singularities. The notebook supplies none of these. A bare analogy with the planar equation does not identify a single modern conjecture.}
\GSfield{Fate}{The named 1901 source is planar. The historical three-variable remark cannot be equated without hypotheses to any one later resolution theorem or remaining higher-dimensional problem. \citep{bendixson1901courbes}}
\GSfield{Present status}{As of 9 November 2026, no current open-status claim is made for this underspecified sentence. It is not added to the problem index.}
\GSfield{Assessment}{literal statement: not determinable; publication match: thematic only; verification: unresolved lead; development: none; confidence: low}
\end{GSitem}
\begin{GSitem}{1965-74.6}{E}{Distinguished Weierstrass polynomials}
\GSfield{Kind; source}{definition; \GSpage{P000117}}
\GSfield{The notes}{The notes write \(P(y)=y^n+\sum_{i=1}^{n-1}y^ia_i(x,z)\), with \(a_i(0,0)=0\); the constant coefficient is absent.}
\GSfield{Statement}{For parameter variables \(u=(x,z)\), a distinguished polynomial is \(P(y,u)=y^n+\sum_{i=0}^{n-1}a_i(u)y^i\), where each \(a_i(0)=0\). The coefficients lie in the relevant analytic, formal or smooth germ ring. The term \(a_0\) is essential; omitting it forces divisibility by \(y\) and does not describe a general distinguished polynomial.}
\GSfield{Fate}{Established before this talk. \citep{malgrange1966ideals}}
\GSfield{Present status}{As of 9 November 2026, this is an established result in the stated setting; the notebook occurrence does not leave a separate open problem.}
\GSfield{Assessment}{literal statement: false as written; publication match: exact; verification: read the result; development: annotated; confidence: high}
\end{GSitem}
\begin{GSitem}{1965-74.7}{E}{Analytic and formal Weierstrass division}
\GSfield{Kind; source}{theorem; \GSpage{P000117}}
\GSfield{The notes}{The stated division is \(f=PQ+\sum_{i=0}^{n-1}b_i(u)y^i\), with coefficients and quotient of the same class; the quotient by \((P)\) is finite over the parameter ring.}
\GSfield{Statement}{If \(P\) is a distinguished polynomial of degree \(n\), every analytic germ \(f(y,u)\) has a unique division \(f=PQ+R\), where \(Q\) is analytic and \(R=\sum_{i=0}^{n-1}b_i(u)y^i\) has analytic coefficients. The analogous formal statement holds over characteristic-zero formal power-series rings. The analytic or formal quotient by \((P)\) is free of rank \(n\) over the parameter germ ring, with basis \(1,y,\ldots,y^{n-1}\). This is finite module rank, not finite dimension over the scalar field when parameters remain.}
\GSfield{Fate}{Established before this talk. \citep{malgrange1966ideals}}
\GSfield{Present status}{As of 9 November 2026, this is an established result in the stated setting; the notebook occurrence does not leave a separate open problem.}
\GSfield{Assessment}{literal statement: correct with missing hypotheses; publication match: exact; verification: read the result; development: annotated; confidence: high}
\end{GSitem}
\begin{GSitem}{1965-74.8}{E}{Smooth division and its loss of uniqueness}
\GSfield{Kind; source}{theorem; \GSpage{P000117}}
\GSfield{The notes}{Malgrange's \(C^\infty\) analogue is said to hold but without uniqueness.}
\GSfield{Statement}{For a smooth distinguished polynomial near the origin and a smooth germ \(f\), smooth division produces a smooth quotient and a remainder polynomial in \(y\) of degree below \(n\), with smooth parameter coefficients. Such a representation need not be unique. The quotient module is finitely generated over the parameter ring; it must not automatically be called free of rank \(n\), because smooth flat terms can create relations.}
\GSfield{Reconstruction}{Finite generation, rather than freeness, is the statement justified in the smooth category.}
\GSfield{Fate}{Malgrange's primary monograph contains the preparation/division theorem and discusses nonunique remainders; its lectures were delivered in 1964 and foreword dated 1965. The 1966 imprint is therefore not evidence of a post-talk discovery. \citep{malgrange1966ideals}}
\GSfield{Present status}{As of 9 November 2026, this is an established result in the stated setting; the notebook occurrence does not leave a separate open problem.}
\GSfield{Assessment}{literal statement: correct; publication match: exact; verification: read the result; development: none; confidence: high}
\end{GSitem}
\begin{GSitem}{1965-74.9}{E}{Interpolation when the roots are real}
\GSfield{Kind; source}{method; \GSpage{P000117}}
\GSfield{The notes}{The notes propose Lagrange interpolation when the polynomial has \(n\) real roots, with a warning about smooth dependence.}
\GSfield{Statement}{For distinct real roots, the remainder is determined by the values of \(f\) at the roots through Lagrange interpolation. At repeated real roots, Hermite interpolation prescribes derivatives to the corresponding multiplicities. Divided-difference formulae explain why the remainder has smooth symmetric dependence when roots coalesce; the naive formula with denominators of root differences requires this limiting argument. In the general non-real-root case smooth uniqueness fails.}
\GSfield{Reconstruction}{Multiplicity and the divided-difference passage, not merely distinct-root interpolation, supply the omitted smoothness issue.}
\GSfield{Fate}{Established before this talk. \citep{malgrange1966ideals}}
\GSfield{Present status}{As of 9 November 2026, this is an established result in the stated setting; the notebook occurrence does not leave a separate open problem.}
\GSfield{Assessment}{literal statement: correct with missing hypotheses; publication match: exact; verification: read the result; development: none; confidence: high}
\end{GSitem}
\begin{GSitem}{1965-74.10}{E}{The flat-function counterexample to unique division}
\GSfield{Kind; source}{claim; \GSpage{P000117}}
\GSfield{The notes}{For a smooth function \(f(x)\) flat at zero, the notes divide by \(y^2+x^2\) in two ways and attribute smoothness of the nontrivial quotient to a Lojasiewicz argument.}
\GSfield{Statement}{Let \(f\in C^\infty(\mathbb R)\) satisfy \(f^{(j)}(0)=0\) for every \(j\). Put \(q(x,y)=f(x)/(x^2+y^2)\) off the origin and \(q(0,0)=0\). Then \(q\) is smooth and flat: every derivative is bounded by a finite sum of \(|f^{(j)}(x)|\,(x^2+y^2)^{-N}\), and flatness absorbs each denominator. Thus \(f=(x^2+y^2)\cdot0+f\) and \(f=(x^2+y^2)q+0\) are distinct smooth divisions when \(f\) is a nonzero germ at zero; for example, take \(f(x)=e^{-1/x^2}\) for \(x\ne0\) and \(f(0)=0\).}
\GSfield{Reconstruction}{\ours{} The derivative estimate is this task's explicit justification of the source example; no quantitative Lojasiewicz theorem is needed for this denominator.}
\GSfield{Fate}{Established before this talk. \citep{malgrange1966ideals}}
\GSfield{Present status}{As of 9 November 2026, this is an established result in the stated setting; the notebook occurrence does not leave a separate open problem.}
\GSfield{Assessment}{literal statement: correct; publication match: exact; verification: read the result; development: annotated; confidence: high}
\end{GSitem}
\begin{GSitem}{1965-74.11}{E}{Preparation of a germ regular in one variable}
\GSfield{Kind; source}{theorem; \GSpage{P000118}}
\GSfield{The notes}{If \(f(0,0)=0\) and some \(\partial_y^n f(0,0)\ne0\), the notes write \(f=P Q\), with distinguished \(P\) and \(Q(0,0)\ne0\); nonflatness is said to allow this.}
\GSfield{Statement}{If \(f\) has exact finite order \(n\) on the chosen \(y\)-axis, meaning \(\partial_y^j f(0,0)=0\) for \(j<n\) and \(\partial_y^n f(0,0)\ne0\), then \(f=U(y,u)(y^n+\sum_{i=0}^{n-1}a_i(u)y^i)\), with \(U(0,0)\ne0\) and \(a_i(0)=0\), in the analytic, formal or smooth setting. Analytic/formal preparation is unique in its normalized form; smooth preparation need not be. A nonflat germ admits a regular direction after a suitable linear coordinate change, not necessarily along the originally chosen axis.}
\GSfield{Fate}{Established before this talk. \citep{malgrange1966ideals}}
\GSfield{Present status}{As of 9 November 2026, this is an established result in the stated setting; the notebook occurrence does not leave a separate open problem.}
\GSfield{Assessment}{literal statement: correct with missing hypotheses; publication match: exact; verification: read the result; development: annotated; confidence: high}
\end{GSitem}
\begin{GSitem}{1965-74.12}{E}{An essential circle meets its area-preserving image}
\GSfield{Kind; source}{claim; \GSpage{P000118}}
\GSfield{The notes}{An area-preserving annulus diffeomorphism homotopic to the identity is said to satisfy \(\#(A\gamma\cap\gamma)\ge2\) for every closed curve surrounding the hole.}
\GSfield{Statement}{Let \(A\) be an area-preserving diffeomorphism of a compact annulus, isotopic to the identity, and let \(\gamma\) be an embedded essential circle in its interior. If \(A\gamma\) and \(\gamma\) meet transversely, they have at least two intersections. Were they disjoint, one would lie strictly on one side of the other and the enclosed annular region would contradict equality of the area between each curve and the same boundary component. A transverse intersection count between isotopic essential circles is even.}
\GSfield{Reconstruction}{\ours{} Embeddedness, a finite-area annulus and transversality give a clean sufficient version of the elementary claim. Homotopy alone for an arbitrary closed, possibly self-intersecting curve does not supply this argument.}
\GSfield{Fate}{Established before this talk. \citep{chaperon1983questions}}
\GSfield{Present status}{As of 9 November 2026, this is an established result in the stated setting; the notebook occurrence does not leave a separate open problem.}
\GSfield{Assessment}{literal statement: correct with missing hypotheses; publication match: partial; verification: read the result; development: annotated; confidence: high}
\end{GSitem}
\begin{GSitem}{1965-74.13}{E}{Globally canonical transformations and exactness}
\GSfield{Kind; source}{definition; \GSpage{P000118}}
\GSfield{The notes}{On \(T^n\times B^n\), the notes impose \(\oint_{A\gamma}p\,dq=\oint_\gamma p\,dq\) for closed loops, calling this globally canonical.}
\GSfield{Statement}{With \(\lambda=\sum p_i\,dq_i\), equality of these integrals for every closed loop is equivalent to \(A^*\lambda-\lambda=dS\) for a global smooth function \(S\), assuming a connected domain. Differentiating gives \(A^*d\lambda=d\lambda\), so \(A\) is exact symplectic. Ordinary symplecticity only makes the difference closed. A compactly supported Hamiltonian isotopy has exact time-one map; an exact endpoint homotopic to the identity should not be silently substituted for a specified Hamiltonian isotopy without an additional argument.}
\GSfield{Reconstruction}{\ours{} The de Rham period argument makes the globally canonical condition precise. This entry deliberately does not assert that all endpoint-exact maps in the notebook's unspecified domain lie in the Hamiltonian component.}
\GSfield{Fate}{Established before this talk. \citep{chaperon1983questions}}
\GSfield{Present status}{As of 9 November 2026, this is an established result in the stated setting; the notebook occurrence does not leave a separate open problem.}
\GSfield{Assessment}{literal statement: correct; publication match: exact; verification: read the result; development: annotated; confidence: high}
\end{GSitem}
\begin{GSitem}{1965-74.14}{B}{Arnold torus-intersection bound and the missing transversality}
\GSfield{Kind; source}{conjecture; \GSpage{P000118}}
\GSfield{The notes}{For the zero torus \(T=\{p=0\}\), the notebook labels \(\#(AT\cap T)\ge2^n\) a theorem, but proves only the graphical case under an additional assumption. No transversality is written.}
\GSfield{Statement}{A precise established version is: if \(\phi\) is a compactly supported Hamiltonian diffeomorphism of \(T^*T^n\) and \(\phi(T^n)\) meets the zero section transversely, then \(\#(\phi(T^n)\cap T^n)\ge\sum_j b_j(T^n;\mathbb F_2)=2^n\). Without transversality the universal cup-length lower bound is \(n+1\). The literal \(2^n\) assertion is false already for \(n=2\): \(S(x,y)=\sin x+\sin y-\sin(x+y)\) has exactly three critical points, and a small graph of \(dS\) is a Hamiltonian image of the zero section.}
\GSfield{Reconstruction}{\ours{} Hamiltonian isotopy, compact support and transverse intersection are sufficient modern hypotheses. The counterexample and its cutoff Hamiltonian realization are independently proved in the accompanying report. Broader variants of the Arnold conjecture are not imported as new notebook problems.}
\GSfield{Fate}{Chaperon's 1983 exposition records the torus progress; the transverse Floer bound and the cotangent-bundle nontransverse cup-length theorem are checked in the opened later sources. This solves the reconstructed Hamiltonian versions, not an unverified equivalence of all endpoint-exact hypotheses. \citep{chaperon1983questions,auroux2026floer,gong2023cuplength}}
\GSfield{Present status}{As of 9 November 2026, the stated transverse Hamiltonian torus bound and nontransverse cup-length bound are established. The literal unqualified bound fails. This occurrence has no claimed new C/D remainder.}
\GSfield{Assessment}{literal statement: false as written; publication match: partial; verification: read the result; development: annotated; confidence: high; problem \hyperref[prob:GP-1965-05]{GP-1965-05}}
\end{GSitem}
\begin{GSitem}{1965-74.15}{E}{The graph argument reduces intersections to Morse critical points}
\GSfield{Kind; source}{method; \GSpage{P000118}}
\GSfield{The notes}{The notes assume \(AT\) can be written \(p=p(q)\) and finish by invoking Morse theory on the torus.}
\GSfield{Statement}{If an exact Lagrangian image of the zero section is the graph of a one-form \(\alpha\) on \(T^n\), then \(\alpha=dS\). Its intersections with the zero section are exactly \(\operatorname{Crit}(S)\), and they are transverse precisely when \(S\) is Morse. Morse inequalities imply at least \(2^n\) such points. For a general smooth \(S\), Lusternik--Schnirelmann cup length gives at least \(n+1\) distinct critical points; it does not imply \(2^n\).}
\GSfield{Reconstruction}{The graph assumption and exact periods recover a generating function. The Morse nondegeneracy condition is missing from the notebook and is indispensable to its numerical conclusion.}
\GSfield{Fate}{Established before this talk. \citep{chaperon1983questions,gong2023cuplength}}
\GSfield{Present status}{As of 9 November 2026, this is an established result in the stated setting; the notebook occurrence does not leave a separate open problem.}
\GSfield{Assessment}{literal statement: correct with missing hypotheses; publication match: exact; verification: read the result; development: annotated; confidence: high; problem \hyperref[prob:GP-1965-05]{GP-1965-05}}
\end{GSitem}

\GStalk{75}{1965-12-06}{I. M. Gelfand}{P000119-P000121}{Hermitian connections, cohomology from rings, and analytic canonical forms}
The dated Gelfand heading begins P000119. The discussion derives the Chern connection and its curvature, then reports Watts's October 1965 algebraic model for Alexander--Spanier cohomology. The Levinson and Whitney discussion on the upper part of P000121 has no intervening speaker heading and is assigned to this record by source sequence; Arnold's new heading begins the fluid discussion below it. The repeated Chern-form argument contains a genuine graded-sign error, and the Whitney example has five factors although the prose twice counts four surfaces. \citep{bottchern1965hermitian,watts1965alexander,lladser2007uniqueness,arnold1968singularities}

\begin{GSitem}{1965-75.1}{E}{Connections and bundle-valued differential forms}
\GSfield{Kind; source}{definition; \GSpage{P000119}}
\GSfield{The notes}{\(D:\Gamma(E)\to\Gamma(T^*X\otimes E)\), \(D(fs)=df\,s+fDs\), and \(A(X,E)=A(X)\otimes_{A^0(X)}\Gamma(E)\).}
\GSfield{Statement}{For a smooth finite-rank complex bundle \(E\to X\), a connection is a complex-linear differential operator satisfying the displayed Leibniz identity. It extends to \(\Omega^q(X,E)\) by \(D(\alpha\otimes s)=d\alpha\otimes s+(-1)^q\alpha\wedge Ds\). Local tensor expressions give the bundle-valued form sheaf; on a compact smooth manifold the global tensor description follows from finite projectivity of \(\Gamma(E)\).}
\GSfield{Reconstruction}{The graded extension and a sufficient compactness hypothesis for the global algebraic tensor presentation are explicit. No completed tensor product is intended.}
\GSfield{Fate}{Established before this talk. \citep{bottchern1965hermitian}}
\GSfield{Present status}{As of 9 November 2026, this is an established result in the stated setting; the notebook occurrence does not leave a separate open problem.}
\GSfield{Assessment}{literal statement: correct with missing hypotheses; publication match: exact; verification: read the result; development: none; confidence: high}
\end{GSitem}
\begin{GSitem}{1965-75.2}{E}{Curvature and the connection-independent Chern classes}
\GSfield{Kind; source}{construction; \GSpage{P000119}}
\GSfield{The notes}{In the row-frame convention \(Ds_i=\sum_j\theta_{ij}s_j\), \(K=d\theta-\theta\wedge\theta\) and \(C(E,D)=\det(1+iK/(2\pi))\); the next line says that this does not depend on the connection.}
\GSfield{Statement}{With this convention the curvature matrix is \(K=d\theta-\theta\wedge\theta\). The even-degree entries commute as scalar forms, so the determinant defines closed forms \(c_j(E,D)\). Their de Rham cohomology classes are independent of \(D\); the differential forms themselves generally vary by exact transgression terms. In a column-frame convention the associated signs and matrix placement change together.}
\GSfield{Reconstruction}{The notebook's independence claim is repaired at the level of cohomology classes, not by changing its faithfully transcribed curvature sign.}
\GSfield{Fate}{Established before this talk. \citep{bottchern1965hermitian}}
\GSfield{Present status}{As of 9 November 2026, this is an established result in the stated setting; the notebook occurrence does not leave a separate open problem.}
\GSfield{Assessment}{literal statement: false as written; publication match: exact; verification: read the result; development: annotated; confidence: high}
\end{GSitem}
\begin{GSitem}{1965-75.3}{E}{Type decomposition on a holomorphic bundle}
\GSfield{Kind; source}{definition; \GSpage{P000119}}
\GSfield{The notes}{\(d=d'+d''\), \(D=D'+D''\), with bidegrees \((1,0)\) and \((0,1)\); the notes call these two connections.}
\GSfield{Statement}{On a complex manifold, \(d=\partial+\bar\partial\) and \(\Omega^k=\bigoplus_{p+q=k}\Omega^{p,q}\). Any complex connection splits into operators \(D'\) and \(D''\) of the corresponding bidegrees. They are partial connections, with Leibniz terms \(\partial f\) and \(\bar\partial f\), rather than two separate full connections. The local generators of scalar \((p,q)\)-forms use holomorphic coordinate functions and their conjugates; the \(f_i\) in that display are not sections of an arbitrary vector bundle.}
\GSfield{Reconstruction}{The word connection in the split display is interpreted as partial connection. The reference to analytic sections in the form generators is repaired to local holomorphic scalar coordinates.}
\GSfield{Fate}{Established before this talk. \citep{bottchern1965hermitian}}
\GSfield{Present status}{As of 9 November 2026, this is an established result in the stated setting; the notebook occurrence does not leave a separate open problem.}
\GSfield{Assessment}{literal statement: correct with missing hypotheses; publication match: exact; verification: read the result; development: none; confidence: high}
\end{GSitem}
\begin{GSitem}{1965-75.4}{E}{Existence and uniqueness of the Chern connection}
\GSfield{Kind; source}{theorem; \GSpage{P000119}}
\GSfield{The notes}{The headline says that a Hermitian metric determines a unique compatible connection; the numbered conditions additionally require \(D''s=0\) for holomorphic \(s\).}
\GSfield{Statement}{On a holomorphic vector bundle \(E\) with a smooth positive-definite Hermitian metric \(h\), there is a unique connection preserving \(h\) and satisfying \(D''=\bar\partial_E\). The metric alone allows many compatible connections; compatibility with the specified holomorphic structure supplies the uniqueness condition. A local holomorphic frame then has connection matrix of type \((1,0)\).}
\GSfield{Reconstruction}{The two numbered conditions, read together, recover the correct theorem despite the abbreviated headline.}
\GSfield{Fate}{Established before this talk. \citep{bottchern1965hermitian}}
\GSfield{Present status}{As of 9 November 2026, this is an established result in the stated setting; the notebook occurrence does not leave a separate open problem.}
\GSfield{Assessment}{literal statement: correct with missing hypotheses; publication match: exact; verification: read the result; development: none; confidence: high}
\end{GSitem}
\begin{GSitem}{1965-75.5}{E}{Hermitian pairing of bundle-valued forms}
\GSfield{Kind; source}{definition; \GSpage{P000119}}
\GSfield{The notes}{The displayed formula is \(\langle s\otimes\alpha,s'\otimes\beta\rangle=\langle s,s'\rangle\alpha\wedge\beta\), with no visible conjugation over \(\beta\).}
\GSfield{Statement}{With the Hermitian metric linear in its first argument, the natural sesquilinear form-valued pairing is \(h(s\otimes\alpha,s'\otimes\beta)=h(s,s')\alpha\wedge\overline\beta\). The conjugation is required for compatibility with moving a complex scalar between the section and form factors. For homogeneous bundle-valued forms, the differential compatibility formula has the usual graded sign.}
\GSfield{Reconstruction}{The scan and Russian transcription both omit the bar. It is a mathematical correction in this annotation, not evidence that the transcription lost a visible stroke.}
\GSfield{Fate}{Established before this talk. \citep{bottchern1965hermitian}}
\GSfield{Present status}{As of 9 November 2026, this is an established result in the stated setting; the notebook occurrence does not leave a separate open problem.}
\GSfield{Assessment}{literal statement: false as written; publication match: exact; verification: read the result; development: annotated; confidence: high}
\end{GSitem}
\begin{GSitem}{1965-75.6}{E}{The metric matrix determines the Chern connection matrix}
\GSfield{Kind; source}{construction; \GSpage{P000119-P000120}}
\GSfield{The notes}{\(N_{ij}=\langle s_i,s_j\rangle\), \(dN=\theta N+N\overline\theta^t\), hence \(d'N=\theta N\) and \(\theta=d'N\,N^{-1}\).}
\GSfield{Statement}{For a local holomorphic column frame \(s=(s_i)\), with \(Ds=\theta s\), and a metric linear in the first argument, \(N\) is a positive Hermitian matrix. Metric compatibility gives \(dN=\theta N+N\overline\theta^t\), and the \((1,0)\)-part yields the uniquely determined matrix \(\theta=(\partial N)N^{-1}\). This matrix has the connection transformation law under holomorphic frame changes.}
\GSfield{Fate}{Established before this talk. \citep{bottchern1965hermitian}}
\GSfield{Present status}{As of 9 November 2026, this is an established result in the stated setting; the notebook occurrence does not leave a separate open problem.}
\GSfield{Assessment}{literal statement: correct; publication match: exact; verification: read the result; development: none; confidence: high}
\end{GSitem}
\begin{GSitem}{1965-75.7}{E}{The Maurer--Cartan identity and the mistaken minus sign}
\GSfield{Kind; source}{claim; \GSpage{P000120}}
\GSfield{The notes}{The notes first write \(d'\theta=\theta\wedge\theta\), then question it, and differentiate \(\theta=d'N\,N^{-1}\) with a minus sign on the resulting product.}
\GSfield{Statement}{The positive identity is correct: \[\partial\theta=-\partial N\wedge\partial(N^{-1})=(\partial N)N^{-1}\wedge(\partial N)N^{-1}=\theta\wedge\theta.\] Both the graded Leibniz sign from the one-form \(\partial N\) and the sign in \(\partial N^{-1}=-N^{-1}(\partial N)N^{-1}\) are needed; they cancel.}
\GSfield{Reconstruction}{\ours{} This graded calculation identifies exactly the missing sign. The minus in the handwritten differentiated line is genuine, so it is retained in the source apparatus.}
\GSfield{Fate}{Established before this talk. \citep{bottchern1965hermitian}}
\GSfield{Present status}{As of 9 November 2026, this is an established result in the stated setting; the notebook occurrence does not leave a separate open problem.}
\GSfield{Assessment}{literal statement: false as written; publication match: exact; verification: read the result; development: annotated; confidence: high}
\end{GSitem}
\begin{GSitem}{1965-75.8}{E}{Curvature of type (1,1), Bianchi, and Chern-form types}
\GSfield{Kind; source}{theorem; \GSpage{P000120}}
\GSfield{The notes}{\(K=d''\theta\), \(d''K=0\), \(d'K=-[K,\theta]\); the notes conclude that all characteristic classes have type \((p,p)\).}
\GSfield{Statement}{For the Chern connection, \(K=\bar\partial\theta\) has type \((1,1)\), \(\bar\partial K=0\), and \(\partial K=-[K,\theta]\), where the bracket is the graded matrix commutator. Consequently \(c_j(E,h)\) has type \((j,j)\), is closed, and represents the de Rham Chern class. On a general non-Kahler manifold this statement concerns a representative; it does not presume a Hodge decomposition of all de Rham cohomology.}
\GSfield{Reconstruction}{The last phrase is restricted to the Chern forms of the holomorphic bundle and is separated from a global Hodge-decomposition assertion.}
\GSfield{Fate}{Established before this talk. \citep{bottchern1965hermitian}}
\GSfield{Present status}{As of 9 November 2026, this is an established result in the stated setting; the notebook occurrence does not leave a separate open problem.}
\GSfield{Assessment}{literal statement: correct with missing hypotheses; publication match: exact; verification: read the result; development: none; confidence: high}
\end{GSitem}
\begin{GSitem}{1965-75.9}{E}{Curvature of a Hermitian holomorphic line bundle}
\GSfield{Kind; source}{construction; \GSpage{P000120}}
\GSfield{The notes}{In rank one, \(N(x)>0\), \(\theta=d'\log N\) and \(K=d''d'\log N\); the characteristic class is the same up to a factor.}
\GSfield{Statement}{In a local nonvanishing holomorphic section \(s\), put \(N=h(s,s)\). Then \(K=\bar\partial\partial\log N\), and \(c_1(E,h)=\frac{i}{2\pi}\bar\partial\partial\log N\). Under holomorphic rescaling of \(s\), the extra logarithm is pluriharmonic, so the curvature glues globally. The order \(\bar\partial\partial\) matches the notebook's row-frame convention.}
\GSfield{Fate}{Established before this talk. \citep{bottchern1965hermitian}}
\GSfield{Present status}{As of 9 November 2026, this is an established result in the stated setting; the notebook occurrence does not leave a separate open problem.}
\GSfield{Assessment}{literal statement: correct; publication match: exact; verification: read the result; development: none; confidence: high}
\end{GSitem}
\begin{GSitem}{1965-75.10}{E}{Watts tensor complex and its differential}
\GSfield{Kind; source}{construction; \GSpage{P000120}}
\GSfield{The notes}{\(F^0=\Lambda\), \(F^p=\Lambda\otimes_KF^{p-1}\), insertion maps \(d_i\) add a unit, and \(d^p=\sum(-1)^id_i\). The prose calls \(F^p\) functions of \(p\) variables.}
\GSfield{Statement}{For a commutative unital \(K\)-algebra \(\Lambda\), set \(F^p=\Lambda^{\otimes_K(p+1)}\). For \(0\le i\le p+1\), let \(d_i:F^p\to F^{p+1}\) insert \(1\) in position \(i\), and put \(d^p=\sum_{i=0}^{p+1}(-1)^id_i\). The cosimplicial insertion identities imply \(d^{p+1}d^p=0\). For a function algebra, a simple tensor is a function of \(p+1\) variables, not \(p\); no completion to all continuous functions on the product is part of this algebraic definition.}
\GSfield{Reconstruction}{Commutativity, a unit, the limits of the alternating sum, and the variable count are supplied or corrected from Watts's exact construction.}
\GSfield{Fate}{Established before this talk. \citep{watts1965alexander}}
\GSfield{Present status}{As of 9 November 2026, this is an established result in the stated setting; the notebook occurrence does not leave a separate open problem.}
\GSfield{Assessment}{literal statement: correct with missing hypotheses; publication match: exact; verification: read the result; development: none; confidence: high}
\end{GSitem}
\begin{GSitem}{1965-75.11}{E}{Algebraic germs along the diagonal}
\GSfield{Kind; source}{definition; \GSpage{P000120}}
\GSfield{The notes}{\(\mu(f_0\otimes\cdots\otimes f_p)=f_0\cdots f_p\); for a maximal ideal \(M\subset\Lambda\), an element vanishes near \(M\) if it dies upon localization at \(\mu^{-1}M\).}
\GSfield{Statement}{The multiplication map \(\mu_p:F^p\to\Lambda\) is a surjective algebra homomorphism, so \(\mu_p^{-1}M\) is maximal. An element \(f\in F^p\) maps to zero in \((F^p)_{\mu_p^{-1}M}\) exactly if there exists \(g\notin\mu_p^{-1}M\) with \(gf=0\). For \(\Lambda=C(X,\mathbb R)\) and compact Hausdorff \(X\), maximal ideals correspond to points and \(\mu\) is restriction to the diagonal.}
\GSfield{Reconstruction}{The point interpretation is only asserted for the specified continuous real-function algebra; arbitrary maximal ideals need not have this topological interpretation.}
\GSfield{Fate}{Established before this talk. \citep{watts1965alexander}}
\GSfield{Present status}{As of 9 November 2026, this is an established result in the stated setting; the notebook occurrence does not leave a separate open problem.}
\GSfield{Assessment}{literal statement: correct with missing hypotheses; publication match: exact; verification: read the result; development: none; confidence: high}
\end{GSitem}
\begin{GSitem}{1965-75.12}{E}{The quotient by functions locally zero at the diagonal}
\GSfield{Kind; source}{construction; \GSpage{P000120}}
\GSfield{The notes}{\(N^p\) consists of elements vanishing in a strip around the diagonal; \(d^pN^p\subset N^{p+1}\), and the notes call the resulting \(H^p\) homology.}
\GSfield{Statement}{Let \(N^p=\bigcap_{M\in\operatorname{Max}\Lambda}\ker(F^p\to(F^p)_{\mu_p^{-1}M})\). The insertion differential preserves this family, so \(F^\bullet/N^\bullet\) is a cochain complex. Its degree-\(p\) group is cohomology, \(H^p=\ker d^p/\operatorname{im}d^{p-1}\). For a compact function space, the local condition is vanishing on some neighborhood of the whole diagonal.}
\GSfield{Reconstruction}{The missing algebraic definition of the diagonal neighborhood is filled by the intersection of localization kernels; the grading distinguishes cohomology from the source's informal homology label.}
\GSfield{Fate}{Established before this talk. \citep{watts1965alexander}}
\GSfield{Present status}{As of 9 November 2026, this is an established result in the stated setting; the notebook occurrence does not leave a separate open problem.}
\GSfield{Assessment}{literal statement: correct with missing hypotheses; publication match: exact; verification: read the result; development: none; confidence: high}
\end{GSitem}
\begin{GSitem}{1965-75.13}{E}{Recovering Alexander--Spanier cohomology from C(X)}
\GSfield{Kind; source}{theorem; \GSpage{P000120}}
\GSfield{The notes}{The final parenthesis says that the groups agree with ordinary Alexander--Spanier cohomology; Watts, PNAS 54(4), October 1965, is named.}
\GSfield{Statement}{For compact Hausdorff \(X\), the preceding construction with \(K=\mathbb R\) and \(\Lambda=C(X,\mathbb R)\) computes Alexander--Spanier cohomology \(H^p(X;\mathbb R)\). Thus this cohomology is determined functorially by the unital real algebra of continuous real functions. The statement does not identify the construction for every abstract algebra with the cohomology of an unspecified spectrum.}
\GSfield{Fate}{The exact two-page source was published in PNAS 54(4), 1027--1028, in October 1965 and is read in full. The talk is reporting that result, not posing a new cohomology problem. \citep{watts1965alexander}}
\GSfield{Present status}{As of 9 November 2026, this is an established result in the stated setting; the notebook occurrence does not leave a separate open problem.}
\GSfield{Assessment}{literal statement: correct with missing hypotheses; publication match: exact; verification: read the result; development: none; confidence: high}
\end{GSitem}
\begin{GSitem}{1965-75.14}{E}{Levinson polynomial form in one distinguished variable}
\GSfield{Kind; source}{theorem; \GSpage{P000121}}
\GSfield{The notes}{The notes cite Duke 1961, pp. 345--353, and assume finite order \(n\) in \(y\). An analytic coordinate change makes \(f\) a polynomial in the new \(y\), with analytic coefficients in the other variables.}
\GSfield{Statement}{Let \(f(y,x)\) be an analytic germ with \(\partial_y^j f(0,0)=0\) for \(j<n\) and \(\partial_y^n f(0,0)\ne0\). There are an analytic local coordinate \(u=u(y,x)\), with \(\partial_yu(0,0)\ne0\), and analytic \(a_j(x)\) such that \(f(y,x)=\sum_{j=0}^{n}a_j(x)u(y,x)^j\), where \(a_n(0)\ne0\) and \(a_j(0)=0\) for \(j<n\). This is polynomialization after a coordinate change, not merely Weierstrass preparation with an extra unit factor.}
\GSfield{Reconstruction}{Analytic germs are understood over \(\mathbb C\), or in the real-analytic version with real data. The inverse-function condition on the new coordinate makes the notebook's word diffeomorphism precise.}
\GSfield{Fate}{Levinson's 1961 bibliographic identification agrees with the notebook. The exact canonical-form statement is read in Lladser's later primary research paper, equation (2), which cites the 1960 announcement; the original 1961 proof was not obtained. \citep{levinson1961canonical,lladser2007uniqueness}}
\GSfield{Present status}{As of 9 November 2026, this is an established result in the stated setting; the notebook occurrence does not leave a separate open problem.}
\GSfield{Assessment}{literal statement: correct; publication match: exact; verification: indirect; development: none; confidence: medium}
\end{GSitem}
\begin{GSitem}{1965-75.15}{E}{Polynomialization of analytic functions of two variables}
\GSfield{Kind; source}{claim; \GSpage{P000121}}
\GSfield{The notes}{The next line says that with two variables one can make an actual polynomial, without further qualifications.}
\GSfield{Statement}{A safe precise version is that an analytic function germ with an isolated critical point is analytically right-equivalent to a polynomial germ, by finite determinacy; this includes the corresponding two-variable case. The precise maximum generality of Levinson's cited two-variable formulation has not been checked in its original full text here. The unrestricted wording in this line is not used to infer that every analytic germ in any dimension is polynomializable, or to identify a new unresolved classification problem.}
\GSfield{Reconstruction}{The isolated-critical-point condition is a sufficient reconstruction, not a recovered handwritten hypothesis. The still broader two-variable formulation remains a limitation of this literature check, not a claim of current openness.}
\GSfield{Fate}{The pre-talk Levinson publication is explicitly identified by later primary authors. Arnold's exposition contrasts two variables with Whitney's nonisolated three-variable obstruction. The original Levinson proof was not accessed, and no claim is made to have verified its full maximum generality. \citep{levinson1960polynomial,lladser2007uniqueness,arnold1968singularities}}
\GSfield{Present status}{As of 9 November 2026, this is an established result in the stated setting; the notebook occurrence does not leave a separate open problem.}
\GSfield{Assessment}{literal statement: correct with missing hypotheses; publication match: partial; verification: indirect; development: none; confidence: medium}
\end{GSitem}
\begin{GSitem}{1965-75.16}{E}{Whitney example with five branches}
\GSfield{Kind; source}{construction; \GSpage{P000121}}
\GSfield{The notes}{The example is \(xy(y-x)(y-xz)(x\varphi(z)-y)\), with \(\varphi(0)=0\) and \(\varphi\) transcendental; the notes say four surfaces but later use the two quadruples \((1,2,3,4)\) and \((1,2,3,5)\).}
\GSfield{Statement}{Choose an analytic germ \(\varphi\) algebraically transcendental over \(\mathbb C(z)\), for example \(e^z-1\). The five factors define five smooth surface germs meeting along the \(z\)-axis, with generically distinct tangent planes away from exceptional axis points. This is the Whitney-type construction of an analytic three-variable germ whose zero-set geometry obstructs polynomialization by a smooth local diffeomorphism. Its critical locus is nonisolated, so it does not contradict finite determinacy at isolated critical points.}
\GSfield{Reconstruction}{\ours{} Five, rather than four, branches are forced by the product and by the two listed quadruples. Algebraic transcendence, rather than merely being a nonpolynomial analytic function, is the property needed by the cross-ratio argument.}
\GSfield{Fate}{Arnold's 1968 exposition, printed p. 39, reproduces Whitney's five-factor counterexample with an exponential factor and attributes it to Whitney's 1965 work. The notebook's parameterized variant has the same mechanism, but its overwritten factors prevent a letter-for-letter exact match. \citep{arnold1968singularities}}
\GSfield{Present status}{As of 9 November 2026, this is an established result in the stated setting; the notebook occurrence does not leave a separate open problem.}
\GSfield{Assessment}{literal statement: false as written; publication match: partial; verification: read the result; development: annotated; confidence: medium}
\end{GSitem}
\begin{GSitem}{1965-75.17}{E}{The pair of cross-ratios obstructs algebraic realization}
\GSfield{Kind; source}{method; \GSpage{P000121}}
\GSfield{The notes}{The anharmonic ratio of tangent planes through a common line is called a diffeomorphism invariant. A curve in the two cross-ratio coordinates is transcendental for the example but would be algebraic for a polynomial.}
\GSfield{Statement}{Modulo the common tangent line, five tangent planes give five points of a projective line. A local diffeomorphism acts by one projective linear transformation on this normal plane, hence preserves the two cross-ratios formed using the first three marked planes. After fixed rational changes of cross-ratio coordinates, the resulting parameterized curve is \(z\mapsto(z,\varphi(z))\). A polynomial configuration yields an algebraic cross-ratio relation along its algebraic multiple-branch locus; choosing \(\varphi\) algebraically transcendental contradicts that relation. Permuting the finitely many branches only applies finitely many rational coordinate changes and does not remove the obstruction.}
\GSfield{Reconstruction}{The argument is local on the open axis stratum with distinct planes; it does not assume that the coordinate change itself is analytic. The five-branch count and algebraic-transcendence requirement make the sketch precise.}
\GSfield{Fate}{Established before this talk. \citep{arnold1968singularities}}
\GSfield{Present status}{As of 9 November 2026, this is an established result in the stated setting; the notebook occurrence does not leave a separate open problem.}
\GSfield{Assessment}{literal statement: correct with missing hypotheses; publication match: partial; verification: read the result; development: annotated; confidence: high}
\end{GSitem}

\GStalk{76}{1965-12-06}{V. I. Arnold}{P000121-P000122}{The topology of analytic stationary fluid flows}
Arnold's underlined heading on the lower part of P000121 begins this record, which ends immediately before the 13 December heading on P000122. It presents the commuting velocity and vorticity fields and the decomposition of an analytic steady flow by regular Bernoulli levels. Arnold's 1966 full paper gives an exact expanded match, explicitly referring to his earlier 1965 announcement. The notebook's separate compact-orbit lemma is too broad and is disproved here; individual dense streamlines must also be distinguished from the compact joint two-dimensional orbit. \citep{arnold1966geometrie}

\begin{GSitem}{1965-76.1}{E}{Vorticity and its transport in an ideal incompressible flow}
\GSfield{Kind; source}{definition; \GSpage{P000121}}
\GSfield{The notes}{The velocity is tangent to the boundary, its vorticity is \(\tau=\operatorname{rot}v\), and a partly diagrammatic evolution equation has an unspecified sign before a Poisson bracket.}
\GSfield{Statement}{For a smooth inviscid incompressible velocity \(v\) in an oriented Euclidean three-dimensional domain, put \(\omega=\nabla\times v\). With \([v,w]=(v\cdot\nabla)w-(w\cdot\nabla)v\), the unforced Euler equation implies \(\partial_t\omega+[v,\omega]=0\). The boundary condition for an impermeable wall is \(v\cdot n=0\). Reversing the vector-field bracket convention reverses the displayed bracket sign.}
\GSfield{Reconstruction}{A definite modern bracket convention resolves the handwritten \(\pm\) without guessing which convention was intended in that line.}
\GSfield{Fate}{Established before this talk. \citep{arnold1966geometrie}}
\GSfield{Present status}{As of 9 November 2026, this is an established result in the stated setting; the notebook occurrence does not leave a separate open problem.}
\GSfield{Assessment}{literal statement: correct with missing hypotheses; publication match: exact; verification: read the result; development: none; confidence: high}
\end{GSitem}
\begin{GSitem}{1965-76.2}{E}{The curl of a cross product and the bracket sign}
\GSfield{Kind; source}{claim; \GSpage{P000121}}
\GSfield{The notes}{\(\operatorname{rot}(a\ast b)=[a\ b]+a\operatorname{div}b-b\operatorname{div}a\), with \(\ast\) identified as vector product.}
\GSfield{Statement}{For the standard vector-field bracket used in item 1965-76.1, \[\nabla\times(a\times b)=-[a,b]+a\,\nabla\cdot b-b\,\nabla\cdot a.\] Thus the notebook's plus sign is consistent with its opposite bracket convention. In particular, commuting divergence-free fields have curl-free cross product regardless of the convention.}
\GSfield{Reconstruction}{\ours{} Component expansion checks the sign. This is a convention reconciliation, not a transcription correction.}
\GSfield{Fate}{Established before this talk. \citep{arnold1966geometrie}}
\GSfield{Present status}{As of 9 November 2026, this is an established result in the stated setting; the notebook occurrence does not leave a separate open problem.}
\GSfield{Assessment}{literal statement: correct with missing hypotheses; publication match: exact; verification: read the result; development: annotated; confidence: high}
\end{GSitem}
\begin{GSitem}{1965-76.3}{E}{Steady Euler flow, commuting vorticity, and global pressure}
\GSfield{Kind; source}{theorem; \GSpage{P000121-P000122}}
\GSfield{The notes}{The three conditions are \(\operatorname{div}v=0\), \([v,\operatorname{rot}v]=0\), and tangency to the boundary. A margin note says that this is the right definition in a simply connected domain.}
\GSfield{Statement}{A smooth steady unforced Euler solution satisfies \((v\cdot\nabla)v=-\nabla p\), so \(v\times\omega=\nabla B\) with \(B=p+|v|^2/2\), and therefore \([v,\omega]=0\). Conversely, for divergence-free \(v\), commutation implies that \((v\times\omega)^\flat\) is closed. A global Bernoulli function, hence pressure, follows when \(H^1(D;\mathbb R)=0\), or more generally when all periods of that closed one-form vanish. Tangency supplies the impermeable boundary condition but does not alone remove a cohomological obstruction.}
\GSfield{Reconstruction}{The handwritten simply-connected qualification is mathematically substantive and is retained; its sufficient cohomological generalization is made explicit.}
\GSfield{Fate}{Established before this talk. \citep{arnold1966geometrie}}
\GSfield{Present status}{As of 9 November 2026, this is an established result in the stated setting; the notebook occurrence does not leave a separate open problem.}
\GSfield{Assessment}{literal statement: correct with missing hypotheses; publication match: exact; verification: read the result; development: none; confidence: high}
\end{GSitem}
\begin{GSitem}{1965-76.4}{E}{Arnold analytic decomposition of a stationary flow}
\GSfield{Kind; source}{theorem; \GSpage{P000121-P000122}}
\GSfield{The notes}{For an analytic stationary velocity in a bounded domain, not everywhere collinear with its curl, finitely many analytic pieces separate cells with torus windings or cylindrical rotations.}
\GSfield{Statement}{Let \(D\) be a connected compact three-dimensional real-analytic Riemannian domain with real-analytic boundary, and let \(v\) be a real-analytic stationary incompressible Euler field tangent to the boundary, analytic up to it. Assume \(v\times\operatorname{curl}v\) is not identically zero. Outside a proper compact analytic exceptional set, regular Bernoulli level components are embedded tori or annuli. On a torus the streamlines are conjugate to a constant linear flow and hence are closed or dense; on an annulus they are closed. The critical values of the Bernoulli function, including its boundary restriction, determine finitely many regular cells.}
\GSfield{Reconstruction}{Compactness, analyticity through the boundary, the global Bernoulli function, and the exceptional critical-level set are explicit. Boundedness and interior analyticity alone do not give the stated finiteness.}
\GSfield{Fate}{Theorem 7 in Arnold's full 1966 paper, printed pp. 346--348, is an exact expanded match and cites his 1965 C. R. Acad. Sci. note, volume 261, p. 17. The result was therefore already announced before this December talk. \citep{arnold1966geometrie}}
\GSfield{Present status}{As of 9 November 2026, this is an established result in the stated setting; the notebook occurrence does not leave a separate open problem.}
\GSfield{Assessment}{literal statement: correct with missing hypotheses; publication match: exact; verification: read the result; development: none; confidence: high}
\end{GSitem}
\begin{GSitem}{1965-76.5}{E}{A common first integral for two commuting divergence-free fields}
\GSfield{Kind; source}{construction; \GSpage{P000122}}
\GSfield{The notes}{The compact-orbit proof writes \(\operatorname{rot}(a\ast b)=0\), then calls \(a\times b\) a differential and identifies orbits with level surfaces.}
\GSfield{Statement}{On an oriented volume three-manifold \((M,\mu)\), if \(a,b\) preserve \(\mu\) and commute, the one-form \(\alpha=\iota_b\iota_a\mu\) is closed. If \(H^1(M;\mathbb R)=0\), then \(\alpha=dF\). Both fields are tangent to the levels of \(F\). Where they are independent, \(dF\ne0\), and their joint orbit is open in the connected regular level component. This last step does not identify a singular-level orbit with an entire closed level set.}
\GSfield{Reconstruction}{\ours{} Cartan's identity gives \(d\alpha=\mathcal L_b\iota_a\mu-\iota_b\mathcal L_a\mu=0\). Regularity and rank two are the missing steps in the notebook's orbit-to-level inference.}
\GSfield{Fate}{Established before this talk. \citep{arnold1966geometrie}}
\GSfield{Present status}{As of 9 November 2026, this is an established result in the stated setting; the notebook occurrence does not leave a separate open problem.}
\GSfield{Assessment}{literal statement: correct with missing hypotheses; publication match: exact; verification: read the result; development: annotated; confidence: high}
\end{GSitem}
\begin{GSitem}{1965-76.6}{E}{A compact regular joint orbit is a torus}
\GSfield{Kind; source}{theorem; \GSpage{P000122}}
\GSfield{The notes}{The notes say that compact orbits without boundary are tori, with the one-parameter velocity action giving a winding.}
\GSfield{Statement}{Let \(C\) be a compact connected regular level surface without boundary on which the two commuting fields \(a,b\) are everywhere independent and tangent. Their complete flows define a locally free \(\mathbb R^2\)-action. Every orbit is open; connectedness makes the action transitive, so \(C\simeq\mathbb R^2/\Lambda\) for a discrete cocompact subgroup \(\Lambda\), necessarily a rank-two lattice. Thus \(C\simeq T^2\), and either field is a constant linear flow in these action coordinates. Its individual streamline may be dense and nonclosed although the joint \(\mathbb R^2\)-orbit is compact.}
\GSfield{Reconstruction}{\ours{} The action-and-lattice argument supplies the rank and completeness qualifications. Compact joint orbits are not confused with compact or closed individual streamlines.}
\GSfield{Fate}{Established before this talk. \citep{arnold1966geometrie}}
\GSfield{Present status}{As of 9 November 2026, this is an established result in the stated setting; the notebook occurrence does not leave a separate open problem.}
\GSfield{Assessment}{literal statement: correct with missing hypotheses; publication match: exact; verification: read the result; development: annotated; confidence: high}
\end{GSitem}
\begin{GSitem}{1965-76.7}{E}{The assertion that every joint orbit is compact is false}
\GSfield{Kind; source}{claim; \GSpage{P000122}}
\GSfield{The notes}{The separately labelled mathematical theorem assumes \(b_1=0\), a finite or compact manifold, and two commuting divergence-free fields, and concludes that their orbits are compact.}
\GSfield{Statement}{These hypotheses do not force compact joint orbits. On the round \(S^3\subset\mathbb C^2\), let \(a(z_1,z_2)=(iz_1,i\alpha z_2)\), with \(\alpha\notin\mathbb Q\), and let \(b=0\). Both fields are smooth, complete and divergence-free, they commute, and \(H^1(S^3;\mathbb R)=0\). Through a point with \(z_1z_2\ne0\), the joint \(\mathbb R^2\)-orbit equals the irrational \(a\)-orbit and is dense but nonclosed in a Clifford torus, hence noncompact. The common first integral here is constant; its level set is not the orbit.}
\GSfield{Reconstruction}{\ours{} The explicit counterexample isolates the missing rank/regular-level hypothesis. The term joint orbit is used deliberately: the failure is in the asserted \(\mathbb R^2\)-orbit theorem itself.}
\GSfield{Fate}{Arnold's published proof only uses compact regular Bernoulli components with two independent tangent fields, as in item 1965-76.6; it does not state this unrestricted compact-orbit lemma. \citep{arnold1966geometrie}}
\GSfield{Present status}{As of 9 November 2026, the literal standalone assertion is disproved by the displayed example; the regular-level torus theorem is established. No open compact-orbit problem is extracted.}
\GSfield{Assessment}{literal statement: false as written; publication match: partial; verification: read the result; development: annotated; confidence: high}
\end{GSitem}
\begin{GSitem}{1965-76.8}{E}{Boundary components give annuli with closed streamlines}
\GSfield{Kind; source}{claim; \GSpage{P000122}}
\GSfield{The notes}{The notes say that noncompact orbits are cylinders and illustrate fluid between rotating coaxial cylinders, with neighboring rings turning in opposite directions.}
\GSfield{Statement}{For a compact regular Bernoulli level component \(C\) meeting the boundary transversely in Arnold's setting, each boundary component is a closed velocity orbit. Commutation and independence propagate its period through \(C\), and \(C\) is an annulus with closed streamlines. It is not a theorem that every noncompact two-dimensional orbit of an \(\mathbb R^2\)-action is a cylinder: a trivial stabilizer gives a plane, while a rank-one stabilizer gives a cylinder. The rotating-cylinder sketch is a qualitative illustration and does not by itself establish an inviscid or viscous Taylor--Couette instability theorem.}
\GSfield{Reconstruction}{The boundary annulus is distinguished from a general noncompact homogeneous orbit. No viscosity, rotation rate, onset threshold or boundary-value problem is recoverable from the sketch.}
\GSfield{Fate}{Established before this talk. \citep{arnold1966geometrie}}
\GSfield{Present status}{As of 9 November 2026, this is an established result in the stated setting; the notebook occurrence does not leave a separate open problem.}
\GSfield{Assessment}{literal statement: false as written; publication match: partial; verification: read the result; development: annotated; confidence: high}
\end{GSitem}
\begin{GSitem}{1965-76.9}{F}{The everywhere-collinear or Beltrami exception}
\GSfield{Kind; source}{remark; \GSpage{P000121}}
\GSfield{The notes}{If velocity and vorticity are everywhere collinear, the notebook asserts that examples need not have the stated simple orbit structure, without giving an explicit example.}
\GSfield{Statement}{If \(v\times\omega\equiv0\), the Bernoulli function is constant on a connected domain, and the nonconstant-first-integral proof of the torus/annulus structure is unavailable. The notebook does not identify a particular Beltrami field or a specific dynamical obstruction. Arnold's 1966 discussion lists the trigonometric ABC-type field as an example of this exceptional class, but describes complicated streamline topology there as probable, not as proved in that paper.}
\GSfield{Fate}{The exceptional case is explicitly discussed in Arnold's paper, printed p. 347 and its footnote. The particular counterexample claimed orally is not identifiable from the notebook, and the paper's speculative language is not upgraded to a proof. \citep{arnold1966geometrie}}
\GSfield{Present status}{As of 9 November 2026, this occurrence identifies the limitation of the regular Bernoulli-level theorem but not a precise open problem or a verified historical example. No general Beltrami-flow conjecture is imported.}
\GSfield{Assessment}{literal statement: not determinable; publication match: partial; verification: read the result; development: none; confidence: medium}
\end{GSitem}

\GStalk{77}{1965-12-13}{V. I. Arnold}{P000122-P000123}{Ideal fluids and the geometry of invariant metrics}
The dated Arnold heading on P000122 starts this record. Newton's law and the incompressibility constraint lead to the Euler/Navier--Stokes equations; the discussion then sketches viscosity-dependent stability and gives the Lie-group interpretation of rigid-body and fluid motion. P000123 ends with the inertia operator, the bilinear adjoint \(B\), and a tangent/cotangent translation diagram. The actual Euler--Arnold evolution formula begins on the next page and belongs to the following record, not this one. Yudovich's named two-dimensional result is identified with his 1963 original paper. \citep{yudovich1963nonstationary,arnold1966geometrie}

\begin{GSitem}{1965-77.1}{E}{Material acceleration and the incompressible equations}
\GSfield{Kind; source}{definition; \GSpage{P000122}}
\GSfield{The notes}{\(DV/Dt=\partial_tv+v\partial_xv\), force consists of pressure, an external term and \(\nu\Delta v\), and \(\operatorname{div}v=0\). The notes call the system a dynamical system on fields satisfying boundary conditions.}
\GSfield{Statement}{For unit constant density in a smooth Euclidean domain, the material acceleration is \(\partial_tv+(v\cdot\nabla)v\). The incompressible Navier--Stokes equations are \(\partial_tv+(v\cdot\nabla)v=-\nabla p+f+\nu\Delta v\), \(\nabla\cdot v=0\); Euler is the inviscid case \(\nu=0\). Pressure enforces the constraint and is determined only up to a function of time. The boundary condition must be specified: impermeability \(v\cdot n=0\) is natural for Euler, while no-slip is a separate viscous condition. On a Riemannian manifold the advective term is \(\nabla_vv\), and a viscous vector Laplacian requires its own geometric convention.}
\GSfield{Reconstruction}{The repeated force symbol in the source is separated into total force and prescribed external force. Potential external forces may be absorbed into pressure. No boundary condition or global solution theory follows merely from the dynamical-system terminology.}
\GSfield{Fate}{Established before this talk. \citep{yudovich1963nonstationary,arnold1966geometrie}}
\GSfield{Present status}{As of 9 November 2026, this is an established result in the stated setting; the notebook occurrence does not leave a separate open problem.}
\GSfield{Assessment}{literal statement: correct with missing hypotheses; publication match: exact; verification: read the result; development: none; confidence: high}
\end{GSitem}
\begin{GSitem}{1965-77.2}{E}{Yudovich global two-dimensional existence and uniqueness}
\GSfield{Kind; source}{theorem; \GSpage{P000122}}
\GSfield{The notes}{The notes mention short-time existence and, in two dimensions, existence and uniqueness, with the explicit parenthesis Yudovich.}
\GSfield{Statement}{A precise unforced specialization is: on a bounded simply connected planar domain with \(C^2\) boundary, an initial divergence-free velocity tangent to the boundary with bounded vorticity determines a unique global Euler solution in the class of velocities with bounded vorticity. Smooth initial data yield the corresponding classical solution. The historical local classical three-dimensional result and the two-dimensional global bounded-vorticity result are different theorems; the Yudovich citation concerns inviscid Euler, not an assertion that viscosity is present.}
\GSfield{Reconstruction}{The function class, domain regularity and inviscid specialization are supplied. The short-time three-dimensional statement is acknowledged in the original introduction; its older proofs are not represented as read here.}
\GSfield{Fate}{Yudovich's 1963 original is opened and read: Theorem 5.1 gives uniqueness/existence and Theorem 6.1 summarizes the velocity/pressure solution with bounded initial vorticity. The simply connected unforced version avoids the additional circulation data treated in his multiply connected discussion. \citep{yudovich1963nonstationary}}
\GSfield{Present status}{As of 9 November 2026, this is an established result in the stated setting; the notebook occurrence does not leave a separate open problem.}
\GSfield{Assessment}{literal statement: correct with missing hypotheses; publication match: exact; verification: read the result; development: none; confidence: high}
\end{GSitem}
\begin{GSitem}{1965-77.3}{F}{Large-viscosity convergence to one equilibrium}
\GSfield{Kind; source}{claim; \GSpage{P000122}}
\GSfield{The notes}{Large viscosity is said to give a unique equilibrium, with all trajectories tending to the solution of the associated elliptic equation.}
\GSfield{Statement}{The statement needs an autonomous forcing, a specified bounded-domain boundary problem, a solution class, and a quantitative smallness condition relative to viscosity. One sufficient conditional criterion is the following: for a steady incompressible Navier--Stokes solution \(u_*\) in a bounded domain with viscosity \(\nu>0\), homogeneous no-slip boundary conditions and Poincare constant \(\lambda_1>0\), if \(\|\nabla u_*\|_{L^\infty}<\nu\lambda_1\), every sufficiently regular solution with the same forcing and boundary conditions obeys, on its interval of existence, \[\|v(t)-u_*\|_2^2\le e^{-2(\nu\lambda_1-\|\nabla u_*\|_\infty)t}\|v(0)-u_*\|_2^2.\] This is a stability criterion for an existing equilibrium; it does not by itself prove the equilibrium's existence as viscosity varies or global existence of the evolving classical solution. Here \(\lambda_1\) means \(\|\nabla w\|_2^2\ge \lambda_1\|w\|_2^2\), and the gradient bound uses the pointwise matrix operator norm (the Frobenius norm also suffices).}
\GSfield{Reconstruction}{\ours{} If \(w=v-u_*\), cancellation of the divergence-free transport terms yields \(\frac12\frac d{dt}\|w\|_2^2+\nu\|\nabla w\|_2^2=-\int(w\cdot\nabla)u_*\cdot w\); the estimate follows by Poincare and Gronwall. This explicitly bounded reconstruction is not silently substituted for the source's missing existence theorem.}
\GSfield{Fate}{No exact large-viscosity theorem is identified from this sentence. The cited Yudovich article is the independently identified inviscid result nearby, not a source for the viscous assertion. The displayed conditional estimate is proved here from the difference energy identity.}
\GSfield{Present status}{As of 9 November 2026, the supplied sufficient criterion is proved, while the exact theorem intended by the unqualified large-viscosity sentence remains unidentified. No generic Navier--Stokes open problem is assigned to it.}
\GSfield{Assessment}{literal statement: not determinable; publication match: thematic only; verification: indirect; development: annotated; confidence: medium}
\end{GSitem}
\begin{GSitem}{1965-77.4}{F}{What happens as viscosity decreases?}
\GSfield{Kind; source}{question; \GSpage{P000122}}
\GSfield{The notes}{The question is followed by a spectral sketch and the statement that eigenvalues of the variational equation move to the right when viscosity decreases.}
\GSfield{Statement}{For a specified branch \(u_\nu\) of steady solutions, one can study the spectrum of the linearized operator \(L_\nu w=-P((u_\nu\cdot\nabla)w+(w\cdot\nabla)u_\nu)+\nu P\Delta w\) on a chosen divergence-free space, with boundary conditions and projection \(P\). Decreasing \(\nu\) does not by itself give a general monotonicity theorem for each eigenvalue of this nonselfadjoint family; both the base flow and its operator domain matter. The notebook gives neither a base-flow family nor a threshold question.}
\GSfield{Reconstruction}{The operator formula is a conventional precise setup, not a recovered source formula. A generic instability diagram cannot be promoted to a definite new research problem.}
\GSfield{Fate}{Arnold's paper supplies the inviscid variational framework but does not identify the unspecified viscous spectral family in this sketch. \citep{arnold1966geometrie}}
\GSfield{Present status}{As of 9 November 2026, no well-posed viscosity-dependent spectral problem is recoverable from the passage. Its class F records that limitation and does not assert absence of later stability theory.}
\GSfield{Assessment}{literal statement: not determinable; publication match: thematic only; verification: indirect; development: none; confidence: low}
\end{GSitem}
\begin{GSitem}{1965-77.5}{E}{Hamiltonian spectral symmetry in the inviscid limit}
\GSfield{Kind; source}{claim; \GSpage{P000123}}
\GSfield{The notes}{At \(\nu=0\) the spectrum is said to be symmetric about both axes, with a partly unclear Hamiltonian parenthesis.}
\GSfield{Statement}{For a real finite-dimensional Hamiltonian linearization \(A=JH\), where \(J\) is invertible skew-symmetric and \(H\) symmetric, eigenvalues occur with their algebraic multiplicities in the set \(\lambda,-\lambda,\overline\lambda,-\overline\lambda\). Reality gives conjugation and the Hamiltonian identity gives sign symmetry. For Euler's infinite-dimensional Hamiltonian/Poisson operator, an analogous assertion requires a specified symplectic leaf, operator domains and a distinction between point and continuous spectrum; the finite-dimensional matrix statement alone is not a complete theorem about the PDE spectrum.}
\GSfield{Reconstruction}{\ours{} The matrix assertion follows because \(A^T=-HJ\) is similar to \(-A\), while a real matrix has conjugate-symmetric spectrum. The notebook's broad PDE statement is restricted to a justified conditional analogue.}
\GSfield{Fate}{Established before this talk. \citep{arnold1966geometrie}}
\GSfield{Present status}{As of 9 November 2026, this is an established result in the stated setting; the notebook occurrence does not leave a separate open problem.}
\GSfield{Assessment}{literal statement: correct with missing hypotheses; publication match: partial; verification: read the result; development: annotated; confidence: high}
\end{GSitem}
\begin{GSitem}{1965-77.6}{F}{The unreadable conditional spectral-stability theorem}
\GSfield{Kind; source}{claim; \GSpage{P000123}}
\GSfield{The notes}{A conditional theorem says that if an inviscid stable flow exists, its spectrum looks like the adjacent unlabelled sketch. A nearby Poiseuille remark calls a flow unstable without specifying geometry or parameters.}
\GSfield{Statement}{The sketch does not determine a precise spectral theorem. In a finite-dimensional Hamiltonian linearization, Lyapunov stability rules out eigenvalues with positive real part, and sign symmetry then puts eigenvalues on the imaginary axis; this does not by itself imply semisimplicity of a zero eigenvalue or nonlinear stability from spectral information. A Poiseuille-flow assertion would require planar or pipe geometry, boundary conditions, Reynolds number, perturbation space and the intended notion of instability.}
\GSfield{Reconstruction}{Only the limited matrix-theoretic implication is supplied as context; neither the shape of the unreadable diagram nor a particular Poiseuille theorem is invented.}
\GSfield{Fate}{The published energy-based stability criterion is more precise than this sketch and cannot be assumed to be its exact missing content. \citep{arnold1966geometrie}}
\GSfield{Present status}{As of 9 November 2026, the passage is too incomplete to identify a definite spectral or flow-instability problem. No current open-status claim follows.}
\GSfield{Assessment}{literal statement: not determinable; publication match: thematic only; verification: read the result; development: none; confidence: low}
\end{GSitem}
\begin{GSitem}{1965-77.7}{E}{Inertia operator and the left-invariant Riemannian metric}
\GSfield{Kind; source}{construction; \GSpage{P000123}}
\GSfield{The notes}{\(L:\mathfrak g\to\mathfrak g^*\) is symmetric positive definite and \(\langle a,b\rangle_g=\langle dL_{g^{-1}}a,dL_{g^{-1}}b\rangle_e\).}
\GSfield{Statement}{On a finite-dimensional Lie group \(G\), a symmetric positive-definite linear isomorphism \(\mathcal I:\mathfrak g\to\mathfrak g^*\) defines \(\langle u,v\rangle_e=(\mathcal Iu)(v)\). Left translation gives a Riemannian metric, with inertia map \(\mathcal I_g=(dL_{g^{-1}})^*\mathcal I\,dL_{g^{-1}}:T_gG\to T_g^*G\). In infinite dimensions a weak positive metric need not identify the tangent space with its full continuous dual; that extra functional-analytic conclusion is not automatic.}
\GSfield{Reconstruction}{The symbol \(\mathcal I\) distinguishes inertia from left translation. The exact tangent/cotangent placement repairs the ambiguity in the lower arrows of the handwritten diagram.}
\GSfield{Fate}{Established before this talk. \citep{arnold1966geometrie}}
\GSfield{Present status}{As of 9 November 2026, this is an established result in the stated setting; the notebook occurrence does not leave a separate open problem.}
\GSfield{Assessment}{literal statement: correct with missing hypotheses; publication match: exact; verification: read the result; development: none; confidence: high}
\end{GSitem}
\begin{GSitem}{1965-77.8}{E}{Free rigid-body motion as a geodesic on SO(3)}
\GSfield{Kind; source}{theorem; \GSpage{P000123}}
\GSfield{The notes}{For \(G=SO(3)\), geodesics describe rotation of a rigid body about a fixed point; the metric is kinetic energy.}
\GSfield{Statement}{For a rigid body with a fixed point and no applied torque, the positive inertia tensor defines a left-invariant kinetic-energy metric on \(SO(3)\). Its geodesics are precisely the free rotations, with the usual factor \(1/2\) in kinetic energy. An asymmetric top generally has a left-invariant metric that is not bi-invariant; no external potential or torque is included.}
\GSfield{Fate}{Arnold presents this as the classical Euler rigid-body model motivating the general group construction. \citep{arnold1966geometrie}}
\GSfield{Present status}{As of 9 November 2026, this is an established result in the stated setting; the notebook occurrence does not leave a separate open problem.}
\GSfield{Assessment}{literal statement: correct with missing hypotheses; publication match: exact; verification: read the result; development: none; confidence: high}
\end{GSitem}
\begin{GSitem}{1965-77.9}{A}{Euler flow as a geodesic of the right-invariant L2 metric}
\GSfield{Kind; source}{theorem; \GSpage{P000123}}
\GSfield{The notes}{\(G=\operatorname{SDiff}D\), \(\langle a,b\rangle=\int_D(a,b)\,dx\), and ideal-fluid flows are geodesics; this follows the preceding left-invariant construction without a stated change of convention.}
\GSfield{Statement}{For an oriented compact Riemannian domain, smooth fluid motions that preserve volume \(\eta_t\) with impermeable boundary are critical points of kinetic action for the right-invariant \(L^2\) metric on volume-preserving diffeomorphisms. The Eulerian velocity \(v=\dot\eta_t\circ\eta_t^{-1}\) satisfies the incompressible Euler equation. Inversion converts right-invariant to left-invariant notation and changes the associated velocity/bracket convention. A rigorous infinite-dimensional geodesic-flow existence theorem needs an appropriate Sobolev manifold and regularity analysis; the smooth formal identification alone does not supply it.}
\GSfield{Reconstruction}{The physical particle map uses right invariance. The explicit convention change prevents importing the left-invariant sign into Eulerian velocity formulas. The functional-analytic existence question is separated from this variational identity.}
\GSfield{Fate}{Arnold's 1966 paper, section 10, pp. 344--345, gives the exact right-invariant fluid interpretation and explains the sign change from the earlier left-invariant convention. This is the full post-talk source checked here; its introduction also mentions 1965 announcements, without a priority claim being made for this occurrence. \citep{arnold1966geometrie}}
\GSfield{Present status}{As of 9 November 2026, this is an established result in the stated setting; the notebook occurrence does not leave a separate open problem.}
\GSfield{Assessment}{literal statement: correct with missing hypotheses; publication match: exact; verification: read the result; development: none; confidence: high}
\end{GSitem}
\begin{GSitem}{1965-77.10}{E}{The affine group realizes the hyperbolic plane}
\GSfield{Kind; source}{construction; \GSpage{P000123}}
\GSfield{The notes}{For \(G=\{x\mapsto ax+b:a>0\}\), the notes identify the resulting dynamics with the geodesic flow on the Lobachevsky plane and ask to inspect the invariant metric.}
\GSfield{Statement}{Write group coordinates as \((b,a)\) with \(a>0\) and multiplication \((b,a)(b',a')=(b+ab',aa')\). Declaring the dilation and translation generators orthonormal at the identity gives the left-invariant metric \(ds^2=(db^2+da^2)/a^2\). This is the standard upper-half-plane hyperbolic metric of curvature \(-1\), so its geodesic flow realizes the stated example. Different overall metric scales change the curvature normalization.}
\GSfield{Reconstruction}{\ours{} The explicit group law, metric and curvature normalization are computed here. The cited Arnold source supplies the invariant-metric framework; this affine example is not asserted to have been located verbatim there.}
\GSfield{Fate}{Established before this talk. \citep{arnold1966geometrie}}
\GSfield{Present status}{As of 9 November 2026, this is an established result in the stated setting; the notebook occurrence does not leave a separate open problem.}
\GSfield{Assessment}{literal statement: correct; publication match: thematic only; verification: read the result; development: annotated; confidence: high}
\end{GSitem}
\begin{GSitem}{1965-77.11}{E}{The bilinear metric-adjoint operation B}
\GSfield{Kind; source}{definition; \GSpage{P000123}}
\GSfield{The notes}{\(\langle[a,b],c\rangle=\langle b,B(c,a)\rangle\).}
\GSfield{Statement}{On a finite-dimensional metric Lie algebra, this identity uniquely defines \(B(c,a)=\operatorname{ad}_a^{\dagger}c\), the metric adjoint of \(b\mapsto[a,b]\). It is bilinear and satisfies \(\langle B(c,a),b\rangle+\langle B(c,b),a\rangle=0\). For a weak infinite-dimensional metric, existence of such an adjoint in the chosen smooth Lie algebra must be established separately. This record defines \(B\); its evolution equation appears in the following record.}
\GSfield{Reconstruction}{The placement of \(c,a\) in \(B\) is checked against both scan and original publication; it is not replaced by a differently ordered adjoint convention.}
\GSfield{Fate}{Arnold's equation (10), printed p. 322, is an exact textual match, with the immediate antisymmetry identity (11). \citep{arnold1966geometrie}}
\GSfield{Present status}{As of 9 November 2026, this is an established result in the stated setting; the notebook occurrence does not leave a separate open problem.}
\GSfield{Assessment}{literal statement: correct with missing hypotheses; publication match: exact; verification: read the result; development: none; confidence: high}
\end{GSitem}
\begin{GSitem}{1965-77.12}{E}{Left and right trivializations of tangent and cotangent vectors}
\GSfield{Kind; source}{construction; \GSpage{P000123}}
\GSfield{The notes}{The lower diagram places \(T_gG\) between two copies of the Lie algebra, with downward inertia maps and dual arrows in the bottom row; the last star and one lower arrow are overwritten.}
\GSfield{Statement}{For a tangent velocity \(\dot g\), the left and right trivializations are \(u=dL_{g^{-1}}\dot g\) and \(w=dR_{g^{-1}}\dot g\). Their dual maps send Lie-algebra covectors to covectors at \(g\), with arrows contravariant to tangent maps. The left-invariant inertia square commutes by \(\mathcal I_g=(dL_{g^{-1}})^*\mathcal I\,dL_{g^{-1}}\). The same fixed \(\mathcal I\) in a right-hand square would require a right-invariant metric, unless an appropriate transformed inertia is inserted. The overwritten terminal node is read mathematically as a dual algebra, not as evidence for a bidual identification.}
\GSfield{Reconstruction}{Only the square justified by the stated left-invariant metric is asserted to commute. The diagram is not treated as evidence that the metric is automatically bi-invariant.}
\GSfield{Fate}{Established before this talk. \citep{arnold1966geometrie}}
\GSfield{Present status}{As of 9 November 2026, this is an established result in the stated setting; the notebook occurrence does not leave a separate open problem.}
\GSfield{Assessment}{literal statement: correct with missing hypotheses; publication match: partial; verification: read the result; development: none; confidence: medium}
\end{GSitem}

\GStalk{78}{}{-}{P000124}{The year 1966}
The boxed heading \ru{1966 г.} marks the resumption of the notebook. It contains no mathematical assertion. The attached catalogue gives a wider overlapping page range; the heading itself is on \GSpage{P000124}.


\GStalk{79}{1966-02-07}{A. G. Kushnirenko}{P000124}{On invariant dynamical systems: title only}
Only Kushnirenko\textquotesingle s name, date and the title \ru{Об инвариантных динамических систем} survive. No definition of the invariant, theorem or proof follows before Gelfand\textquotesingle s heading. This is not sufficient evidence for identifying the talk with Kushnirenko\textquotesingle s entropy estimates from the preceding year (1964-10.2--3). Accordingly it has no invented mathematical item.


\GStalk{80}{1966-02-07}{I. M. Gelfand}{P000124}{Ore fractions, Weyl fields, automorphisms and dimension}
The page begins with the Ore property and division rings of fractions, then distinguishes differential-operator rings \(R_{n,k}\) and their quotient division rings \(D_{n,k}\), introduces multiplicative leading symbols, and asks about determinants and all automorphisms. It ends with the growth definition of dimension and \(\dim R_{n,k}=2n+k\). The dimension formula in the handwriting lacks a logarithm in the numerator. Much of the page matches \citet{gelfandkirillov1966corps}; the paper was received on 19 January 1966, before this talk, but published afterwards. The general Gelfand--Kirillov conjecture for enveloping division rings is in that paper, not stated on this notebook page, and is not imported as a seminar item.

\begin{GSitem}{1965-80.1}{E}{A Noetherian domain satisfies the Ore condition}
\GSfield{Kind; source}{theorem; \GSpage{P000124}}
\GSfield{The notes}{\ru{Всякое нёт. кольцо (без дел. 0) --- кольцо Оре}; the displayed condition is \(a,b\ne0\Rightarrow \exists x,y\ne0:xa=yb\).}
\GSfield{Statement}{If \(R\) is an associative unital left-Noetherian domain, then \(Ra\cap Rb\ne0\) for every \(a,b\in R\setminus\{0\}\): \(R\) satisfies the left Ore condition. The right-sided assertion requires right Noetherianity. Enveloping algebras of finite-dimensional Lie algebras over a characteristic-zero field have both properties by the PBW filtration.}
\GSfield{Reconstruction}{The side of Noetherianity and the unit convention are supplied. PBW is contextual explanation, not an extra assertion attributed to this page.}
\GSfield{Fate}{Background ring theory, reproduced as Lemma 2 of \citet[\S1]{gelfandkirillov1966corps}; Ore localization predates the seminar \citep{ore1931linear}.}
\GSfield{Present status}{As of 9 November 2026: The stated result is established with the hypotheses supplied here.}
\GSfield{Assessment}{literal statement: correct with missing hypotheses; publication match: exact; verification: read the result; development: none; confidence: high}
\end{GSitem}
\begin{GSitem}{1965-80.2}{E}{The ideal-chain proof and its missing first nonzero coefficient}
\GSfield{Kind; source}{method; \GSpage{P000124}}
\GSfield{The notes}{The notes form the left ideals generated by \(a,ab,\ldots,ab^n\), write \(ab^{n+1}=x_0a+\cdots+x_nab^n\), and leave \(x_0a=(\quad)b\).}
\GSfield{Statement}{After the chain stabilizes, choose the least \(j\) with \(x_j\ne0\). Cancel the common right factor \(b^j\) in the domain. Then \[x_j a=\bigl(ab^{n-j}-x_{j+1}a-\cdots-x_nab^{n-j-1}\bigr)b\ne0.\] If \(j=n\), the parenthesis is simply \(a\). This produces a nonzero common left multiple without assuming \(x_0\ne0\).}
\GSfield{Reconstruction}{The notebook leaves the parenthesis empty and does not justify that its displayed \(x_0\) is nonzero. Choosing \(j\), rather than assuming \(x_0\ne0\), is the necessary completion.}
\GSfield{Fate}{The completed minimal-coefficient step is explicit in the proof of Lemma 2 in \citet{gelfandkirillov1966corps}.}
\GSfield{Present status}{As of 9 November 2026: The stated result is established with the hypotheses supplied here.}
\GSfield{Assessment}{literal statement: correct with missing hypotheses; publication match: exact; verification: read the result; development: annotated; confidence: high}
\end{GSitem}
\begin{GSitem}{1965-80.3}{E}{The Ore division ring of fractions}
\GSfield{Kind; source}{construction; \GSpage{P000124}}
\GSfield{The notes}{\ru{Кольцо Оре вклад. в тело}; elements are written \(ab^{-1}\), described as equivalence classes.}
\GSfield{Statement}{A unital domain satisfying the right Ore condition admits a division ring \(Q(R)\) in which every element is \(ab^{-1}\), \(a,b\in R\), \(b\ne0\). It is unique up to the isomorphism fixing \(R\). Left Ore gives the analogous construction with \(b^{-1}a\); a two-sided Ore domain has both descriptions.}
\GSfield{Reconstruction}{The preceding displayed condition is left-sided, whereas the written fractions are right-sided. Both sides hold for the Noetherian rings subsequently used; a one-sided hypothesis alone should not silently be swapped.}
\GSfield{Fate}{This is Ore\textquotesingle s fraction construction \citep{ore1931linear}, reviewed in \citet[\S1]{gelfandkirillov1966corps}.}
\GSfield{Present status}{As of 9 November 2026: The stated result is established with the hypotheses supplied here.}
\GSfield{Assessment}{literal statement: correct with missing hypotheses; publication match: exact; verification: read the result; development: none; confidence: high}
\end{GSitem}
\begin{GSitem}{1965-80.4}{F}{The unresolved Green-operator remark}
\GSfield{Kind; source}{claim; \GSpage{P000124}}
\GSfield{The notes}{An underlined line reads \ru{Дифф. опер. от [зачёркнуто] Грина[?], это тело!!}, followed by \ru{(и наоборот)}.}
\GSfield{Statement}{No definite theorem is recoverable: a differential-operator ring itself is generally not a division ring, whereas adjoining formal inverses in an Ore localization is. The uncertain word may refer to Green operators, but no domain, boundary conditions or inverse convention is specified.}
\GSfield{Reconstruction}{The two possible readings are kept apart. No analytic inverse, Green kernel or claim that every differential operator is invertible is supplied.}
\GSfield{Fate}{The algebraic context is \citet{gelfandkirillov1966corps}; no exact publication match for this damaged sentence was identified.}
\GSfield{Present status}{As of 9 November 2026: The fragment remains underdetermined; it is not evidence for an undeveloped mathematical problem.}
\GSfield{Assessment}{literal statement: not determinable; publication match: thematic only; verification: unresolved lead; development: none; confidence: low}
\end{GSitem}
\begin{GSitem}{1965-80.5}{E}{Weyl algebras with central parameters and their fraction fields}
\GSfield{Kind; source}{definition; \GSpage{P000124}}
\GSfield{The notes}{\(R_{n,k}\) is the ring of differential operators in \(n\) variables with \(k\) commuting parameters; \(D_{n,k}\) is its division ring of ratios.}
\GSfield{Statement}{Over a fixed characteristic-zero field \(L\), let \(R_{n,k}=A_n(L)\otimes_L L[z_1,\ldots,z_k]\), where \([p_i,q_j]=\delta_{ij}\), \([p_i,p_j]=[q_i,q_j]=0\), and all \(z_\ell\) are central. Set \(D_{n,k}=Q(R_{n,k})\). Here \(n,k\ge0\); \(D_{0,k}=L(z_1,\ldots,z_k)\).}
\GSfield{Reconstruction}{The coefficient field, characteristic, commutator sign and meaning of independent central parameters are made explicit. Polynomial coefficients, rather than arbitrary smooth coefficients, are essential.}
\GSfield{Fate}{Classical differential-operator and Ore constructions; this notation is formalized in \citet[\S\S2--3]{gelfandkirillov1966corps}.}
\GSfield{Present status}{As of 9 November 2026: The stated result is established with the hypotheses supplied here.}
\GSfield{Assessment}{literal statement: correct with missing hypotheses; publication match: exact; verification: read the result; development: none; confidence: high}
\end{GSitem}
\begin{GSitem}{1965-80.6}{A}{Different Weyl ranks and parameter counts give nonisomorphic division rings}
\GSfield{Kind; source}{theorem; \GSpage{P000124}}
\GSfield{The notes}{For \(D_{n,k}\), \ru{При разных n и k неизом.}}
\GSfield{Statement}{For a fixed characteristic-zero field \(L\), an \(L\)-algebra isomorphism \(D_{n,k}\simeq D_{n^{\prime},k^{\prime}}\) exists if and only if \((n,k)=(n^{\prime},k^{\prime})\). The same distinction holds for the rings \(R_{n,k}\). The division-ring statement needs the fraction-field invariant, not only the growth dimension of the polynomial subring.}
\GSfield{Reconstruction}{Isomorphisms are required to fix \(L\). This extends the earlier Weyl-rank item 1964-38.1; it is not the endomorphism conjecture for a fixed Weyl algebra.}
\GSfield{Fate}{Exactly the nonisomorphism theorem of \citet[\S4]{gelfandkirillov1966corps}; also announced by \citet{gelfandkirillov1966fields}.}
\GSfield{Present status}{As of 9 November 2026: The stated result is established with the hypotheses supplied here.}
\GSfield{Assessment}{literal statement: correct with missing hypotheses; publication match: exact; verification: read the result; development: none; confidence: high; problem \hyperref[prob:GP-1965-06]{GP-1965-06}}
\end{GSitem}
\begin{GSitem}{1965-80.7}{E}{Commutative leading symbols}
\GSfield{Kind; source}{method; \GSpage{P000124}}
\GSfield{The notes}{For an operator \(a\), \([a]\) denotes its leading part; \ru{В старш. ч. p и q коммут.}, and \([ab]=[a][b]\).}
\GSfield{Statement}{Give each Weyl generator degree one and each central parameter degree one (or give the parameters degree zero over their coefficient ring). Then \(\operatorname{gr}R_{n,k}\) is the corresponding commutative polynomial domain. For nonzero \(a,b\), degree is additive and \(\sigma(ab)=\sigma(a)\sigma(b)\). The principal-symbol map is multiplicative on nonzero elements, not an additive map from the filtered ring to its graded ring.}
\GSfield{Reconstruction}{The filtration must be specified. Its two natural versions are distinguished; no characteristic-zero operator ring is identified with its commutative associated graded ring.}
\GSfield{Fate}{The standard filtered-operator argument is used explicitly in \citet[\S3]{gelfandkirillov1966corps}.}
\GSfield{Present status}{As of 9 November 2026: The stated result is established with the hypotheses supplied here.}
\GSfield{Assessment}{literal statement: correct; publication match: exact; verification: read the result; development: none; confidence: high}
\end{GSitem}
\begin{GSitem}{1965-80.8}{A}{A well-defined multiplicative symbol of an Ore fraction}
\GSfield{Kind; source}{construction; \GSpage{P000124}}
\GSfield{The notes}{For \(b=a_1a_2^{-1}\), the notes take \([a_1]/[a_2]\) and say that it depends only on \(b\), not its representation. The short continuation at the right edge is uncertain.}
\GSfield{Statement}{The multiplicative symbol extends to \(D_{n,k}^{\times}\) by \(\sigma(a_1a_2^{-1})=\sigma(a_1)/\sigma(a_2)\). This is independent of the right-fraction expression, takes values in the multiplicative group of nonzero homogeneous rational functions of the chosen graded variables, and is invariant under inner automorphisms. It is not a division-ring homomorphism.}
\GSfield{Reconstruction}{Inner invariance is the publication\textquotesingle s corollary, not a confident decipherment of the damaged right-edge fragment. Multiplicativity follows by taking common Ore denominators; commutativity of graded symbols removes conjugation.}
\GSfield{Fate}{The independence statement and inner-invariance corollary are Lemma 4 of \citet[\S3]{gelfandkirillov1966corps}.}
\GSfield{Present status}{As of 9 November 2026: The stated result is established with the hypotheses supplied here.}
\GSfield{Assessment}{literal statement: correct; publication match: exact; verification: read the result; development: none; confidence: high}
\end{GSitem}
\begin{GSitem}{1965-80.9}{E}{A universal determinant over a division ring}
\GSfield{Kind; source}{question; \GSpage{P000124}}
\GSfield{The notes}{\ru{Универс. det --- отобр. в фактор по коммутанту. То же в теле, универс. --- ?}}
\GSfield{Statement}{For a division ring \(D\), what is the abelianization of the stable general linear group? The answer is \[K_1(D)=\mathrm{GL}_{\infty}(D)_{\mathrm{ab}}\simeq D^{\times}/[D^{\times},D^{\times}].\] The Dieudonn\'e determinant realizes this quotient; every group homomorphism from the stable general linear group to an abelian group factors uniquely through it. In characteristic zero this also holds in each finite matrix size at least two.}
\GSfield{Reconstruction}{The notes do not specify matrix size. The stable formulation removes the small finite-field exception \(\mathrm{GL}_2(\mathbb F_2)\); the Weyl fields here have characteristic zero.}
\GSfield{Fate}{The division-ring determinant and its universal abelian quotient were established by \citet[\S\S5--7]{dieudonne1943determinants}.}
\GSfield{Present status}{As of 9 November 2026: This is the established Dieudonn\'e determinant, not a remaining open determinant problem.}
\GSfield{Assessment}{literal statement: correct with missing hypotheses; publication match: exact; verification: read the result; development: none; confidence: high; problem \hyperref[prob:GP-1965-07]{GP-1965-07}}
\end{GSitem}
\begin{GSitem}{1965-80.10}{C}{Determine all automorphisms of the Weyl division ring}
\GSfield{Kind; source}{problem; \GSpage{P000124}}
\GSfield{The notes}{\ru{Найти все автоморфизмы этого тела (тр. зад.)}; below it: \ru{(Пр. Фурье + умн. на ф-ю)}.}
\GSfield{Statement}{For the previously defined \(D_{n,k}\), determine its automorphisms over the coefficient field. A precise minimal case is \(\operatorname{Aut}_{\mathbb C}D_{1,0}\), where \(D_{1,0}=Q(\mathbb C\langle x,p\mid px-xp=1\rangle)\). The question concerns the division ring, not just automorphisms or endomorphisms of its polynomial Weyl subalgebra.}
\GSfield{Reconstruction}{The base field and the case \(n=1,k=0\) are explicit restrictions for development. Nothing on the page asserts that Fourier transforms and multiplication operators generate every automorphism. No novel general solution is claimed below.}
\GSfield{Fate}{Polynomial automorphisms are understood \citep{makarlimanov1984automorphisms}; rational subfield stabilizers are classified by \citet[Proposition 2.2]{alevdumas1998invariants}. These are partial answers to the notebook\textquotesingle s unrestricted question.}
\GSfield{Present status}{As of 9 November 2026: The full group is not classified in the verified literature: \citet[\S5.3.3]{dumas2011invariants} states the gap; searches through this date found no later full classification. \citet[Theorem 0.12]{schwarz2024generalizations} concerns induced finite-group invariant fields, not all division-ring automorphisms. The development proves a restricted classification and an exact rational-pole criterion; the unrestricted problem is left open. The November 2026 index-symmetry theorem of \citet[Corollary 8.12]{lauritzenthomsen2026duality} is a further related result, not an automorphism classification.}
\GSfield{Assessment}{literal statement: correct with missing hypotheses; publication match: partial; verification: read the result; development: developed in this chapter; confidence: medium; problem \hyperref[prob:GP-1965-08]{GP-1965-08}}
\GSfield{Worth developing}{Worth a bounded development of rational subfield stabilizers and inner-conjugation obstructions, because these answer a precise part of the handwritten question without presuming a classification of the full group.}
\end{GSitem}
\begin{GSitem}{1965-80.11}{E}{Fourier transforms and multiplication as automorphism constructions}
\GSfield{Kind; source}{method; \GSpage{P000124}}
\GSfield{The notes}{The parenthesis to the automorphism problem suggests Fourier transform plus multiplication by a function.}
\GSfield{Statement}{For \(D_{1,0}\) with \([p,x]=1\), the algebraic Fourier map \(x\mapsto p\), \(p\mapsto-x\) is an automorphism. For \(u\in\mathbb C(x)^{\times}\), conjugation by multiplication by \(u\) gives \(x\mapsto x\), \(p\mapsto p-u^{\prime}/u\). Multiplication alone is not a unital algebra automorphism. More generally every rational shift \(p\mapsto p+h(x)\) is an automorphism, but need not be inner.}
\GSfield{Reconstruction}{The notebook\textquotesingle s operation on functions is interpreted as conjugation on the operator algebra; invertibility of \(u\) and the rational coefficient restriction are supplied. The last distinction is proved in the development.}
\GSfield{Fate}{Fourier transformations already motivate \citet{gelfandkirillov1966corps}; the rational stabilizer context is made precise in \citet{alevdumas1998invariants}.}
\GSfield{Present status}{As of 9 November 2026: The stated result is established with the hypotheses supplied here.}
\GSfield{Assessment}{literal statement: correct with missing hypotheses; publication match: partial; verification: read the result; development: annotated; confidence: high; problem \hyperref[prob:GP-1965-08]{GP-1965-08}}
\end{GSitem}
\begin{GSitem}{1965-80.12}{A}{Growth dimension and the missing logarithm}
\GSfield{Kind; source}{definition; \GSpage{P000124}}
\GSfield{The notes}{For generators \(\alpha=(a_1,\ldots,a_s)\), the displayed numerator is \(d(\alpha,N)\), described as the number of linearly independent monomials of degree at most \(N\); the denominator is \(\ln N\).}
\GSfield{Statement}{The correct algebra growth invariant is \[\operatorname{GKdim}_L A=\sup_V\limsup_{N\to\infty}\frac{\log\dim_L V^N}{\log N},\] where \(V\) runs over finite-dimensional \(L\)-subspaces containing \(1\), and \(V^N\) is the span of products of at most \(N\) factors. For finitely generated \(A\), any finite-dimensional generating \(V\ni1\) suffices. Without the numerator logarithm, even \(L[x]\) has an infinite displayed value.}
\GSfield{Reconstruction}{The normalized logarithmic growth is a mathematical emendation, not a new reading of the ink. Independence of generators follows from \(V\subseteq W^c\) and \(W\subseteq V^d\). For quotient division rings the 1966 paper uses a different invariant with a denominator-clearing infimum; replacing that invariant by naive algebra growth is unjustified.}
\GSfield{Fate}{The logarithm is present in formula (6) of \citet[\S4]{gelfandkirillov1966corps}; the scan confirms its omission in the notebook.}
\GSfield{Present status}{As of 9 November 2026: The corrected invariant is Gelfand--Kirillov dimension. The handwritten omission is substantive and is recorded in the correction apparatus.}
\GSfield{Assessment}{literal statement: false as written; publication match: exact; verification: read the result; development: annotated; confidence: high}
\end{GSitem}
\begin{GSitem}{1965-80.13}{A}{Growth dimension of the parameterized Weyl algebra}
\GSfield{Kind; source}{theorem; \GSpage{P000124}}
\GSfield{The notes}{The last line is \(\dim R_{n,k}=2n+k\).}
\GSfield{Statement}{For \(R_{n,k}=A_n(L)\otimes L[z_1,\ldots,z_k]\) with the total-degree generating filtration and \(\operatorname{char}L=0\), ordered PBW monomials form a basis and \[\dim_L F_N R_{n,k}=\binom{N+2n+k}{2n+k}.\] Consequently \(\operatorname{GKdim}_L R_{n,k}=2n+k\), including the scalar case \(n=k=0\).}
\GSfield{Reconstruction}{This conclusion uses the corrected logarithmic definition. The ordered-monomial count includes the identity and central parameters and is not asserted to count all noncommuting words separately. \oursComputed\ The independent finite checks, including \(n=k=0\) and \(N=0\), are saved in \texttt{code/check\_weyl.py} and its two output files.}
\GSfield{Fate}{This is Lemma 5 of \citet[\S4]{gelfandkirillov1966corps}; the count is checked independently in the supplied computation.}
\GSfield{Present status}{As of 9 November 2026: The stated result is established with the hypotheses supplied here.}
\GSfield{Assessment}{literal statement: correct with missing hypotheses; publication match: exact; verification: read the result; development: annotated; confidence: high}
\end{GSitem}

\GStalk{81}{1966-02-14}{M. S. Agranovich}{P000125-P000128}{Pseudodifferential and singular integral operators}
The first Agranovich lecture introduces homogeneous symbols, spherical harmonic kernels, rough-symbol boundedness and the first rules of symbolic calculus. It is primarily an exposition of the already published 1965 theory, not a preview whose first publication was necessarily after the seminar. Several handwritten slips in the boundedness argument are separated from transcription mistakes. The lecture ends in the upper part of P000128, before Gelfand's integral-geometry heading.

\begin{GSitem}{1965-81.1}{E}{Homogeneous symbol and Fourier quantization}
\GSfield{Kind; source}{definition; \GSpage{P000125}}
\GSfield{The notes}{The notes write \(Au=F^{-1}[\sigma(x,\xi)Fu]\), with \(\sigma\) positively homogeneous of degree \(s\geq0\).}
\GSfield{Statement}{For a smooth symbol away from \(\xi=0\), satisfying \(\sigma(x,t\xi)=t^s\sigma(x,\xi)\) for \(t>0\), left quantization is \(Au(x)=(2\pi)^{-n}\int e^{ix\cdot\xi}\sigma(x,\xi)\widehat u(\xi)\,d\xi\), initially on test functions, with an oscillatory interpretation when necessary.}
\GSfield{Reconstruction}{The Fourier normalization and the meaning of the inverse transform with the external parameter \(x\) are supplied. A cutoff near zero gives the standard smooth symbol class.}
\GSfield{Fate}{Already treated in \citet{agranovich1965elliptic}, especially \S\S6--8; this is expository by the date of the talk.}
\GSfield{Present status}{As of November 2026: established under the stated hypotheses; no unresolved problem is asserted by this occurrence.}
\GSfield{Assessment}{literal statement: correct with missing hypotheses; publication match: exact; verification: read the result; development: none; confidence: high}
\end{GSitem}
\begin{GSitem}{1965-81.2}{E}{Differential operators, multiplication and the square-root Laplacian}
\GSfield{Kind; source}{construction; \GSpage{P000125}}
\GSfield{The notes}{Examples are \(\sum_{|\alpha|=m}a_\alpha D^\alpha\), \(s=0\), \(\sigma(x)\), and \(\Lambda=\sqrt{-\Delta}\).}
\GSfield{Statement}{With \(D_j=-i\partial_{x_j}\), a homogeneous differential operator has symbol \(\sum_{|\alpha|=m}a_\alpha(x)\xi^\alpha\). A symbol independent of \(\xi\) quantizes as multiplication. On \(\mathbb R^n\), \(\sqrt{-\Delta}\) is the Fourier multiplier \(|\xi|\) on its natural Sobolev domain.}
\GSfield{Reconstruction}{These examples specify principal homogeneous symbols; a general differential operator also has lower-degree terms.}
\GSfield{Fate}{Already treated in \citet{agranovich1965elliptic}, especially \S\S6--8; this is expository by the date of the talk.}
\GSfield{Present status}{As of November 2026: established under the stated hypotheses; no unresolved problem is asserted by this occurrence.}
\GSfield{Assessment}{literal statement: correct; publication match: exact; verification: read the result; development: none; confidence: high}
\end{GSitem}
\begin{GSitem}{1965-81.3}{E}{Low-frequency regularization of negative orders}
\GSfield{Kind; source}{method; \GSpage{P000125}}
\GSfield{The notes}{For \(s<0\), multiply by \(\varphi(|\xi|)\), zero near zero and one away from zero; the difference is called inessential.}
\GSfield{Statement}{Two smooth cutoffs agreeing for large \(|\xi|\) give symbols differing in a compact frequency annulus. Their properly localized operators differ by an operator with a smooth kernel and hence by an operator of order \(-\infty\).}
\GSfield{Reconstruction}{One cannot subtract an undefined unregularized singular multiplier at zero. The assertion compares two admissible regularizations and is local in the base variable.}
\GSfield{Fate}{Already treated in \citet{agranovich1965elliptic}, especially \S\S6--8; this is expository by the date of the talk.}
\GSfield{Present status}{As of November 2026: established under the stated hypotheses; no unresolved problem is asserted by this occurrence.}
\GSfield{Assessment}{literal statement: correct with missing hypotheses; publication match: exact; verification: read the result; development: none; confidence: high}
\end{GSitem}
\begin{GSitem}{1965-81.4}{E}{Spherical harmonic expansion and homogeneous kernels}
\GSfield{Kind; source}{construction; \GSpage{P000125-P000126}}
\GSfield{The notes}{\(\sigma(x,\xi)=\sum_{m,k}a_{mk}(x)Y_{mk}(\xi/|\xi|)\). For \(m>0\), the inverse multiplier kernel is a constant times \(Y_{mk}(z/|z|)|z|^{-n}\).}
\GSfield{Statement}{Choose orthonormal spherical harmonics of degree \(m\). In the distributional principal-value sense, the Fourier transform preserves each nonconstant harmonic type and maps \(|z|^{-n}Y_{mk}(z/|z|)\) to \(\gamma_mY_{mk}(\xi/|\xi|)\), with \(|\gamma_m|\asymp m^{-n/2}\) for the conventional transform normalization.}
\GSfield{Reconstruction}{The exact complex phase and numerical constant depend on the Fourier convention; only the reciprocal relation and asymptotic order are used below. The degree-zero term is excluded from this principal-value kernel.}
\GSfield{Fate}{\citet{agranovich1965elliptic}, \S6, supplies this harmonic-kernel correspondence.}
\GSfield{Present status}{As of November 2026: established under the stated hypotheses; no unresolved problem is asserted by this occurrence.}
\GSfield{Assessment}{literal statement: correct with missing hypotheses; publication match: exact; verification: read the result; development: none; confidence: high}
\end{GSitem}
\begin{GSitem}{1965-81.5}{E}{Cancellation and the missing multiplication term}
\GSfield{Kind; source}{definition; \GSpage{P000125-P000126}}
\GSfield{The notes}{The notes collect the harmonic kernels into a homogeneous \(K\), state a sphere mean-zero condition, and display a principal-value integral for \(A\).}
\GSfield{Statement}{An order-zero operator of this form is \(Au=a_{00}(x)u+\operatorname{pv}\int K(x,x-y)u(y)\,dy\), where \(K(x,tz)=t^{-n}K(x,z)\) and \(\int_{S^{n-1}}K(x,\omega)\,d\omega=0\).}
\GSfield{Reconstruction}{The multiplication term must be restored unless the degree-zero spherical coefficient vanishes. The fixed first argument of the cancellation condition is the base point \(x\), not the integration variable.}
\GSfield{Fate}{Already treated in \citet{agranovich1965elliptic}, especially \S\S6--8; this is expository by the date of the talk.}
\GSfield{Present status}{As of November 2026: established under the stated hypotheses; no unresolved problem is asserted by this occurrence.}
\GSfield{Assessment}{literal statement: false as written; publication match: exact; verification: read the result; development: annotated; confidence: high}
\end{GSitem}
\begin{GSitem}{1965-81.6}{E}{Fourier-side kernel and stabilization at infinity}
\GSfield{Kind; source}{construction; \GSpage{P000126}}
\GSfield{The notes}{The notes set \(\sigma=\sigma_\infty+\sigma^{\prime}\), call \(\sigma^{\prime}\) finite in \(\xi\), and give a Fourier-side integral with \(\widehat{\sigma^{\prime}}(\xi-\eta,\eta)\).}
\GSfield{Statement}{For a symbol equal to \(\sigma_\infty(\xi)\) outside a compact set in \(x\), \(\widehat{Au}(\xi)=\sigma_\infty(\xi)\widehat u(\xi)+(2\pi)^{-n}\int\widehat{\sigma^{\prime}}_x(\xi-\eta,\eta)\widehat u(\eta)\,d\eta\), with the corresponding normalization adjustment for a unitary transform.}
\GSfield{Reconstruction}{Compact support must refer to \(x\) in this argument. A nonzero homogeneous function cannot be compactly supported in \(\xi\). The scan itself says \(\xi\); this is a notebook slip, not a transcription error.}
\GSfield{Fate}{This representation is part of the 1965 calculus; compare \citet{agranovich1965elliptic} and the independently published \citet{kohnnirenberg1965algebra}.}
\GSfield{Present status}{As of November 2026: established under the stated hypotheses; no unresolved problem is asserted by this occurrence.}
\GSfield{Assessment}{literal statement: false as written; publication match: exact; verification: read the result; development: annotated; confidence: high}
\end{GSitem}
\begin{GSitem}{1965-81.7}{E}{Smooth equivalence and rough-symbol inequivalence}
\GSfield{Kind; source}{remark; \GSpage{P000126}}
\GSfield{The notes}{The two definitions are said to agree for smooth symbols but to lead to different classes after smoothness is relaxed.}
\GSfield{Statement}{Left/right quantizations have the same principal symbol in a smooth classical calculus. For rough symbols, boundedness and the ability to convert quantizations require explicit regularity; no equivalence follows merely from the formal integral expressions.}
\GSfield{Reconstruction}{The notes do not specify a pair of rough symbol spaces for a sharp inequivalence theorem. This occurrence records the methodological warning, not an unrecorded counterexample.}
\GSfield{Fate}{Already treated in \citet{agranovich1965elliptic}, especially \S\S6--8; this is expository by the date of the talk.}
\GSfield{Present status}{As of November 2026: established under the stated hypotheses; no unresolved problem is asserted by this occurrence.}
\GSfield{Assessment}{literal statement: correct with missing hypotheses; publication match: partial; verification: read the result; development: none; confidence: medium}
\end{GSitem}
\begin{GSitem}{1965-81.8}{E}{Sobolev scale and order convention}
\GSfield{Kind; source}{definition; \GSpage{P000126}}
\GSfield{The notes}{\(\|u\|_\ell^2=\int(1+|\xi|^2)^\ell|\widehat u|^2\); an order-\(s\) operator maps \(H^\ell\) to \(H^{\ell-s}\).}
\GSfield{Statement}{For symbols in the smooth uniform class \(S^s_{1,0}\), or for properly supported local operators, \(A:H^\ell_{\mathrm{comp}}\to H^{\ell-s}_{\mathrm{loc}}\) continuously. Uniform global symbol estimates give the global bounded map on \(\mathbb R^n\).}
\GSfield{Reconstruction}{Uniform coefficient bounds or localization are needed on a noncompact base.}
\GSfield{Fate}{Already treated in \citet{agranovich1965elliptic}, especially \S\S6--8; this is expository by the date of the talk.}
\GSfield{Present status}{As of November 2026: established under the stated hypotheses; no unresolved problem is asserted by this occurrence.}
\GSfield{Assessment}{literal statement: correct with missing hypotheses; publication match: exact; verification: read the result; development: none; confidence: high}
\end{GSitem}
\begin{GSitem}{1965-81.9}{E}{Angular Sobolev regularity sufficient for L2 boundedness}
\GSfield{Kind; source}{theorem; \GSpage{P000126-P000127}}
\GSfield{The notes}{The displayed criterion assumes \(q>(n-1)/2\) and \(\sup_x\|\sigma(x,\cdot)\|_{H^q(S^{n-1})}<\infty\).}
\GSfield{Statement}{For \(n\geq2\), a degree-zero symbol measurable in \(x\) with \[M=\operatorname*{ess\,sup}_x\sum_{m,k}(1+m)^{2q}|a_{mk}(x)|^2<\infty\], \(q>(n-1)/2\), defines an \(L^2\)-bounded left-quantized operator. Its norm is at most a dimensional constant times \(M^{1/2}\).}
\GSfield{Reconstruction}{The weight \(1+m\) retains the constant harmonic; the printed use of \(m\) alone is harmless only after separating multiplication.}
\GSfield{Fate}{\citet{agranovich1965elliptic}, Proposition 7.1 and Remark 7.1, already give the criterion and its rough-base-variable interpretation.}
\GSfield{Present status}{As of November 2026: established under the stated hypotheses; no unresolved problem is asserted by this occurrence.}
\GSfield{Assessment}{literal statement: correct with missing hypotheses; publication match: exact; verification: read the result; development: none; confidence: high}
\end{GSitem}
\begin{GSitem}{1965-81.10}{E}{Addition theorem and the boundedness proof}
\GSfield{Kind; source}{method; \GSpage{P000126-P000127}}
\GSfield{The notes}{The proof uses Schwarz, the sum \(\sum_k|Y_{mk}|^2\), and a series with exponent \(n-2\). Several factors and signs in the handwritten estimates are lost or reversed.}
\GSfield{Statement}{Pointwise Schwarz gives \(|Au(x)|^2\leq M\sum_{m,k}(1+m)^{-2q}|Y_{mk}(D)u(x)|^2\). Plancherel and the addition theorem give \(\|Au\|_2^2\leq C M\|u\|_2^2\sum_m(1+m)^{n-2-2q}\). The last series converges exactly when \(2q>n-1\).}
\GSfield{Reconstruction}{\ours{} The explanatory verification here is supplied by us for the established result or correction. The two Schwarz sums are multiplied, not equated; \(|\widehat u|^2\) remains under the Fourier integral, and the weight in the second sum has exponent \(-2q\). These are mathematical repairs to handwriting verified in the scans.}
\GSfield{Fate}{The complete estimate is in \citet{agranovich1965elliptic}, proof of Proposition 7.1.}
\GSfield{Present status}{As of November 2026: established under the stated hypotheses; no unresolved problem is asserted by this occurrence.}
\GSfield{Assessment}{literal statement: false as written; publication match: exact; verification: read the result; development: annotated; confidence: high}
\end{GSitem}
\begin{GSitem}{1965-81.11}{E}{Endpoint and absence of base-variable smoothness}
\GSfield{Kind; source}{remark; \GSpage{P000127}}
\GSfield{The notes}{The notes say \ru{По $\xi$ гладк. не нужно} near the angular argument and discuss divergence at the limiting exponent.}
\GSfield{Statement}{The angular criterion needs no regularity in \(x\). At \(q=(n-1)/2\), its proof gives a harmonic series and no uniform bound; individual operators, such as the identity, remain bounded. The endpoint is not a claim that every symbol there is unbounded.}
\GSfield{Reconstruction}{In this local context the handwritten \(\xi\) must be \(x\); otherwise it contradicts the angular hypothesis. A later, different criterion genuinely removes angular regularity.}
\GSfield{Fate}{\citet{agranovich1965elliptic}, Remark 7.1, discusses the sharp failure of the general convergence assertion.}
\GSfield{Present status}{As of November 2026: established under the stated hypotheses; no unresolved problem is asserted by this occurrence.}
\GSfield{Assessment}{literal statement: false as written; publication match: exact; verification: read the result; development: annotated; confidence: high}
\end{GSitem}
\begin{GSitem}{1965-81.12}{E}{Kernel version of the angular criterion}
\GSfield{Kind; source}{theorem; \GSpage{P000127}}
\GSfield{The notes}{For \(K(x,z)=k(x,z/|z|)|z|^{-n}\), the threshold becomes \(q>-1/2\).}
\GSfield{Statement}{If \(k(x,\cdot)\) has zero spherical mean, is measurable in \(x\), and is uniformly in \(H^q(S^{n-1})\), \(q>-1/2\), then the associated distributional singular integral is \(L^2\)-bounded. For negative \(q\), the angular kernel is interpreted as a distribution.}
\GSfield{Reconstruction}{The harmonic transform raises the angular Sobolev index by \(n/2\). Cancellation and the distributional interpretation are essential.}
\GSfield{Fate}{\citet{agranovich1965elliptic}, \S\S6--7, contains the symbol-to-kernel transfer.}
\GSfield{Present status}{As of November 2026: established under the stated hypotheses; no unresolved problem is asserted by this occurrence.}
\GSfield{Assessment}{literal statement: correct with missing hypotheses; publication match: exact; verification: read the result; development: none; confidence: high}
\end{GSitem}
\begin{GSitem}{1965-81.13}{E}{Planar Riesz-transform example}
\GSfield{Kind; source}{construction; \GSpage{P000127}}
\GSfield{The notes}{The example is \(z_1/(z_1^2+z_2^2)^{3/2}\).}
\GSfield{Statement}{In \(\mathbb R^2\), \(\operatorname{pv}z_1/|z|^3\) is, up to the Fourier normalization and sign, the kernel of the first Riesz transform, whose multiplier is a constant times \(i\xi_1/|\xi|\); it is \(L^2\)-bounded.}
\GSfield{Fate}{Already treated in \citet{agranovich1965elliptic}, especially \S\S6--8; this is expository by the date of the talk.}
\GSfield{Present status}{As of November 2026: established under the stated hypotheses; no unresolved problem is asserted by this occurrence.}
\GSfield{Assessment}{literal statement: correct; publication match: exact; verification: read the result; development: none; confidence: high}
\end{GSitem}
\begin{GSitem}{1965-81.14}{E}{A complementary criterion using base-variable derivatives}
\GSfield{Kind; source}{theorem; \GSpage{P000127}}
\GSfield{The notes}{Compactly supported dependence on \(x\) with \(n+1\) derivatives gives Fourier decay \((1+|\xi-\eta|^2)^{-(n+1)/2}\).}
\GSfield{Statement}{If \(\sigma^{\prime}(\cdot,\eta)\) is supported in one fixed compact set and its \(x\)-derivatives through order \(n+1\) are uniformly bounded, then \(|\widehat{\sigma^{\prime}}_x(\zeta,\eta)|\leq C(1+|\zeta|)^{-n-1}\). The Fourier kernel is dominated by an \(L^1\) convolution kernel, hence defines a bounded \(L^2\) operator, without angular smoothness.}
\GSfield{Reconstruction}{Uniformity in \(\eta\) and fixed spatial support are supplied; boundedness of the separate multiplier \(\sigma_\infty\) is also needed.}
\GSfield{Fate}{Already treated in \citet{agranovich1965elliptic}, especially \S\S6--8; this is expository by the date of the talk.}
\GSfield{Present status}{As of November 2026: established under the stated hypotheses; no unresolved problem is asserted by this occurrence.}
\GSfield{Assessment}{literal statement: correct with missing hypotheses; publication match: exact; verification: read the result; development: none; confidence: high}
\end{GSitem}
\begin{GSitem}{1965-81.15}{E}{First symbolic composition rule}
\GSfield{Kind; source}{theorem; \GSpage{P000127}}
\GSfield{The notes}{The notes announce that \(AB-\sigma_1\sigma_2(x,D)\) is subordinate.}
\GSfield{Statement}{For smooth classical operators of orders \(s_1,s_2\), \(AB-\operatorname{Op}(\sigma_{s_1}(A)\sigma_{s_2}(B))\) has order at most \(s_1+s_2-1\), provided the representative operators have only the stated principal homogeneous term plus lower-order remainders.}
\GSfield{Reconstruction}{Proper support is understood locally. Rough-symbol boundedness alone does not imply this calculus.}
\GSfield{Fate}{The result precedes the seminar in \citet{agranovich1965elliptic} and \citet{kohnnirenberg1965algebra}; record 84 continues its proof.}
\GSfield{Present status}{As of November 2026: established under the stated hypotheses; no unresolved problem is asserted by this occurrence.}
\GSfield{Assessment}{literal statement: correct with missing hypotheses; publication match: exact; verification: read the result; development: none; confidence: high}
\end{GSitem}
\begin{GSitem}{1965-81.16}{E}{Pseudolocality and localization}
\GSfield{Kind; source}{theorem; \GSpage{P000127}}
\GSfield{The notes}{Pseudolocality and extension to manifolds are announced.}
\GSfield{Statement}{For a smooth pseudodifferential operator, \(\operatorname{sing\,supp}(Au)\subseteq\operatorname{sing\,supp}(u)\). If smooth compact cutoffs have disjoint supports, \(\varphi A\psi\) has a smooth kernel.}
\GSfield{Reconstruction}{The full smoothing assertion uses the smooth calculus, stronger than the one-order estimates elsewhere in the notes.}
\GSfield{Fate}{Already treated in \citet{agranovich1965elliptic}, especially \S\S6--8; this is expository by the date of the talk.}
\GSfield{Present status}{As of November 2026: established under the stated hypotheses; no unresolved problem is asserted by this occurrence.}
\GSfield{Assessment}{literal statement: correct with missing hypotheses; publication match: exact; verification: read the result; development: none; confidence: high}
\end{GSitem}
\begin{GSitem}{1965-81.17}{E}{Left and right quantizations differ below principal order}
\GSfield{Kind; source}{theorem; \GSpage{P000128}}
\GSfield{The notes}{The Fourier kernel uses either \(\sigma(\xi-\eta,\eta)\) or \(\sigma(\xi-\eta,\xi)\); Taylor expansion makes their difference subordinate.}
\GSfield{Statement}{For a smooth symbol of order \(s\), left and right quantizations using the same symbol agree modulo order \(s-1\). The full conversion is an asymptotic series involving mixed base and frequency derivatives.}
\GSfield{Reconstruction}{Only the first-order consequence, not a complete asymptotic formula, is asserted by the notes.}
\GSfield{Fate}{Already treated in \citet{agranovich1965elliptic}, especially \S\S6--8; this is expository by the date of the talk.}
\GSfield{Present status}{As of November 2026: established under the stated hypotheses; no unresolved problem is asserted by this occurrence.}
\GSfield{Assessment}{literal statement: correct with missing hypotheses; publication match: exact; verification: read the result; development: none; confidence: high}
\end{GSitem}
\begin{GSitem}{1965-81.18}{E}{Adjoint principal symbol}
\GSfield{Kind; source}{theorem; \GSpage{P000128}}
\GSfield{The notes}{The adjoint reverses the two quantizations and conjugates the symbol.}
\GSfield{Statement}{With a fixed smooth positive density and Hermitian metrics, the formal adjoint of an order-\(s\) operator has principal symbol \(\sigma_s(A)^*\), the fiberwise Hermitian transpose. The remainder is of order \(s-1\).}
\GSfield{Reconstruction}{For scalar operators the star is complex conjugation. Density and bundle metrics matter only below the leading symbol in a compatible local trivialization.}
\GSfield{Fate}{Already treated in \citet{agranovich1965elliptic}, especially \S\S6--8; this is expository by the date of the talk.}
\GSfield{Present status}{As of November 2026: established under the stated hypotheses; no unresolved problem is asserted by this occurrence.}
\GSfield{Assessment}{literal statement: correct with missing hypotheses; publication match: exact; verification: read the result; development: none; confidence: high}
\end{GSitem}
\begin{GSitem}{1965-81.19}{F}{Gelfand asks about replacing polynomials in the algebra}
\GSfield{Kind; source}{question; \GSpage{P000128}}
\GSfield{The notes}{Gelfand asks \ru{Можно ли зам. многочл. на другие, чтобы получ. та же алгебра.}}
\GSfield{Statement}{No precise question can be recovered: the polynomials to be replaced, the admissible replacement functions and the sense of equality of algebras are not specified.}
\GSfield{Reconstruction}{Possible readings include generating symbols, harmonic expansions, or an alternative quantization. None is imposed as the historical statement.}
\GSfield{Fate}{The nearby 1965 calculus papers clarify context but do not identify this wording with a definite published problem.}
\GSfield{Present status}{As of November 2026: unidentifiable from the surviving sentence; no claim of an open mathematical problem is warranted.}
\GSfield{Assessment}{literal statement: not determinable; publication match: none found; verification: unresolved lead; development: none; confidence: low}
\end{GSitem}
\begin{GSitem}{1965-81.20}{E}{Cotangent location of the principal symbol}
\GSfield{Kind; source}{remark; \GSpage{P000128}}
\GSfield{The notes}{The lecture ends by placing the symbol on the cotangent bundle.}
\GSfield{Statement}{The scalar principal symbol of an order-\(s\) operator is a degree-\(s\) homogeneous function on \(T^*M\setminus0\); for operators between bundles it is a section of \(\operatorname{Hom}(\pi^*E,\pi^*F)\).}
\GSfield{Reconstruction}{The bundle-valued statement is the modern completion of the scalar shorthand; record 84 supplies the coordinate argument.}
\GSfield{Fate}{Already treated in \citet{agranovich1965elliptic}, especially \S\S6--8; this is expository by the date of the talk.}
\GSfield{Present status}{As of November 2026: established under the stated hypotheses; no unresolved problem is asserted by this occurrence.}
\GSfield{Assessment}{literal statement: correct; publication match: exact; verification: read the result; development: none; confidence: high}
\end{GSitem}

\GStalk{82}{1966-02-21}{I. M. Gelfand}{P000128-P000132}{Integral geometry on affine k-planes}
The lecture develops differential forms attached to the affine \(k\)-plane transform: a generalized determinant, a closedness criterion for the range, and inversion by integrating over cycles. The close publication is \citet{gelfand1967planes}, with Graev and Shapiro. Its hypotheses and quantifiers correct two potentially serious overstatements in the notebook. P000128 starts this record after Agranovich; P000132 ends it before Karpelevich.

\begin{GSitem}{1965-82.1}{A}{Reconstruction from k-plane integrals and economical families}
\GSfield{Kind; source}{problem; \GSpage{P000128}}
\GSfield{The notes}{Reconstruct \(f\) from \(\varphi(h)=\int_hf\); for lines in \(\mathbb R^3\), can the four-parameter data be reduced to three parameters?}
\GSfield{Statement}{For a rapidly decreasing smooth function on affine space, recover \(f\) from its integrals over affine \(k\)-planes. For a prescribed smaller family, an inversion question additionally requires its geometry and a nonvanishing inversion coefficient; mere equality of parameter counts does not imply injectivity.}
\GSfield{Reconstruction}{The first question is precise. The three-parameter subfamily is not specified; later cycle formulas answer an appropriate family-dependent version, not every subfamily.}
\GSfield{Fate}{Published soon after in \citet{gelfand1967planes}; the relevant definitions and Theorems 1--2 were checked in the Russian full text.}
\GSfield{Present status}{As of November 2026: the full-data problem and admissible cycle inversions are established; the unspecified economical-family sentence does not define a further open problem.}
\GSfield{Assessment}{literal statement: correct with missing hypotheses; publication match: partial; verification: read the result; development: none; confidence: high; problem \hyperref[prob:GP-1965-09]{GP-1965-09}}
\end{GSitem}
\begin{GSitem}{1965-82.2}{A}{Affine-plane bundle and translation-incidence section}
\GSfield{Kind; source}{construction; \GSpage{P000129}}
\GSfield{The notes}{\(H_{n,k}\) consists of affine \(k\)-planes with volume data, \(G\) records parallel classes, and \(S_\beta\) selects the translate through \(\beta\).}
\GSfield{Statement}{An affine plane is \(\alpha+b\), where \(\alpha\in\operatorname{Gr}(k,\mathbb F^n)\) and \(b\in\mathbb F^n/\alpha\), with compatible real or complex volume normalization. The pullback \(\varphi(\alpha,\beta)\) depends only on \(\beta\bmod\alpha\).}
\GSfield{Reconstruction}{Volume weights must transform with a change of frame. The real and complex versions use different integration densities; they cannot be interchanged without a normalization convention.}
\GSfield{Fate}{Published soon after in \citet{gelfand1967planes}; the relevant definitions and Theorems 1--2 were checked in the Russian full text.}
\GSfield{Present status}{As of November 2026: established under the stated hypotheses; no unresolved problem is asserted by this occurrence.}
\GSfield{Assessment}{literal statement: correct with missing hypotheses; publication match: exact; verification: read the result; development: none; confidence: high}
\end{GSitem}
\begin{GSitem}{1965-82.3}{A}{Determinant of vector entries paired with a symmetric form}
\GSfield{Kind; source}{construction; \GSpage{P000129-P000130}}
\GSfield{The notes}{For a symmetric \(k\)-linear form \(B\), the notes define a signed sum over permutations of entries \(\xi_{ij}\).}
\GSfield{Statement}{For vectors \(\xi_{ij}\in W\), set \(\det_B(\xi)=\sum_{\pi\in S_k}\operatorname{sgn}(\pi)B(\xi_{1\pi(1)},\ldots,\xi_{k\pi(k)})\). Symmetry of \(B\) makes this alternating in rows and columns. If \(B\) vanishes whenever an argument belongs to \(\alpha\), the expression descends to \(W/\alpha\).}
\GSfield{Reconstruction}{This is a polarized determinant, not a determinant with an unspecified multiplication on \(W\).}
\GSfield{Fate}{Published soon after in \citet{gelfand1967planes}; the relevant definitions and Theorems 1--2 were checked in the Russian full text.}
\GSfield{Present status}{As of November 2026: established under the stated hypotheses; no unresolved problem is asserted by this occurrence.}
\GSfield{Assessment}{literal statement: correct; publication match: exact; verification: read the result; development: none; confidence: high}
\end{GSitem}
\begin{GSitem}{1965-82.4}{A}{Translation derivatives define a differential form on directions}
\GSfield{Kind; source}{construction; \GSpage{P000130}}
\GSfield{The notes}{The symmetric derivative \(\partial_\beta^k\varphi\) and nearby frames \(e_i+\xi_{ij}\) produce a \(k\)-form; frame-independence is Theorem 1.}
\GSfield{Statement}{The identification \(T_\alpha\operatorname{Gr}(k,\mathbb F^n)=\operatorname{Hom}(\alpha,\mathbb F^n/\alpha)\) lets the polarized determinant pair the translation derivative with \(k\) tangent directions. The accompanying frame-volume factor makes the result intrinsic. In the complex formulation the corresponding object has bidegree \((k,k)\), with holomorphic and antiholomorphic derivative factors.}
\GSfield{Reconstruction}{The real \(k\)-form and complex \((k,k)\)-form are distinguished explicitly. Indexing a basis of \(\alpha\) uses \(k\) entries, even where the notebook writes \(n\).}
\GSfield{Fate}{Published soon after in \citet{gelfand1967planes}; the relevant definitions and Theorems 1--2 were checked in the Russian full text.}
\GSfield{Present status}{As of November 2026: established under the stated hypotheses; no unresolved problem is asserted by this occurrence.}
\GSfield{Assessment}{literal statement: correct with missing hypotheses; publication match: exact; verification: read the result; development: none; confidence: high}
\end{GSitem}
\begin{GSitem}{1965-82.5}{A}{Closedness characterizes the range with restricted k and decay}
\GSfield{Kind; source}{theorem; \GSpage{P000130-P000131}}
\GSfield{The notes}{Theorem 2 says \(\varphi\) is in the plane-transform range if and only if its associated form is closed.}
\GSfield{Statement}{In the smooth rapidly decreasing class \(\Phi\) of the published theorem, with rapid decay in the translation variable together with all derivatives, and \(1\leq k<n-1\), the Gelfand--Graev--Shapiro closedness equations characterize the transform range. The paper gives the complex theorem and notes its real analogue. The case \(k=n-1\) is excluded.}
\GSfield{Reconstruction}{These hypotheses are restored from the published Theorem 1, not inferred from the short notebook theorem. For hyperplanes, closedness in top degree cannot substitute for the missing global moment conditions.}
\GSfield{Fate}{\citet{gelfand1967planes}, Theorem 1, p. 23 of the Russian edition, and its two footnotes explicitly impose the restriction and discuss the real version.}
\GSfield{Present status}{As of November 2026: established for the published range class; the unrestricted notebook wording is not a valid universal characterization.}
\GSfield{Assessment}{literal statement: correct with missing hypotheses; publication match: exact; verification: read the result; development: annotated; confidence: high; problem \hyperref[prob:GP-1965-09]{GP-1965-09}}
\end{GSitem}
\begin{GSitem}{1965-82.6}{A}{Cycle inversion and its homological coefficient}
\GSfield{Kind; source}{theorem; \GSpage{P000131}}
\GSfield{The notes}{Theorem 3 states \(\int_\gamma\kappa\varphi=C_\gamma f(\beta)\); \(C_\gamma\) is topological.}
\GSfield{Statement}{For transform data \(\varphi=R_k f\) and an admissible closed cycle in the direction Grassmannian, integration of the associated form through \(\beta\) equals a fixed homological pairing times \(f(\beta)\). The cycle has real dimension \(k\) in the real construction and \(2k\) in the complex construction. It yields inversion only if its pairing is nonzero.}
\GSfield{Reconstruction}{Orientations, density normalization and the relevant homology class determine \(C_\gamma\). A merely parameter-counting choice of cycle is insufficient.}
\GSfield{Fate}{\citet{gelfand1967planes}, Theorem 2 and the subsequent cycle calculation, pp. 29--31, establish this formula.}
\GSfield{Present status}{As of November 2026: established under the stated hypotheses; no unresolved problem is asserted by this occurrence.}
\GSfield{Assessment}{literal statement: correct with missing hypotheses; publication match: exact; verification: read the result; development: none; confidence: high; problem \hyperref[prob:GP-1965-09]{GP-1965-09}}
\end{GSitem}
\begin{GSitem}{1965-82.7}{A}{Nonzero complex coefficient is an existence statement}
\GSfield{Kind; source}{claim; \GSpage{P000131}}
\GSfield{The notes}{The notes suggest that in the complex case the coefficient is always nonzero, with a contrast to real parity.}
\GSfield{Statement}{For complex affine \(k\)-planes there exist \(2k\)-cycles with nonzero inversion coefficient. An arbitrary cycle can have zero coefficient, including a null-homologous cycle or the zero class. The real local differential inversion has the familiar parity obstruction for odd \(k\).}
\GSfield{Reconstruction}{\ours{} The explanatory verification here is supplied by us for the established result or correction. The published theorem has an existential quantifier. This is a substantive correction to the literal universal wording, not a new open case.}
\GSfield{Fate}{\citet{gelfand1967planes}, Theorem 2, states existence of a suitable nonzero class; it does not assert nonvanishing on every cycle.}
\GSfield{Present status}{As of November 2026: established under the stated hypotheses; no unresolved problem is asserted by this occurrence.}
\GSfield{Assessment}{literal statement: false as written; publication match: exact; verification: read the result; development: annotated; confidence: high}
\end{GSitem}
\begin{GSitem}{1965-82.8}{A}{Delta-current form of the inversion identity}
\GSfield{Kind; source}{construction; \GSpage{P000131}}
\GSfield{The notes}{The notes give the distributional identity \(C_\gamma\delta_\beta=\int_\gamma\kappa\delta_h\).}
\GSfield{Statement}{Interpret plane integration as pairing with the plane delta distribution. Applying the derivative-form construction and then integrating over the admissible cycle gives the evaluation current \(C_\gamma\delta_\beta\), with the same conventions as the function inversion theorem.}
\GSfield{Reconstruction}{The formula is an equality of distributions, after pairing with test functions; it is not a pointwise product of delta functions.}
\GSfield{Fate}{Published soon after in \citet{gelfand1967planes}; the relevant definitions and Theorems 1--2 were checked in the Russian full text.}
\GSfield{Present status}{As of November 2026: established under the stated hypotheses; no unresolved problem is asserted by this occurrence.}
\GSfield{Assessment}{literal statement: correct with missing hypotheses; publication match: exact; verification: read the result; development: none; confidence: high}
\end{GSitem}
\begin{GSitem}{1965-82.9}{A}{Complex-line coordinates and the (1,1)-form}
\GSfield{Kind; source}{construction; \GSpage{P000131-P000132}}
\GSfield{The notes}{Lines are written \(x=\alpha t+\beta\), with integration \(dt\,d\bar t\); a form involving \(\partial_{\beta_i}d\alpha_i\) and its conjugate is displayed.}
\GSfield{Statement}{On the graph chart \(x_j=\alpha_jx_n+\beta_j\), \(j=1,\ldots,n-1\), the line space has complex dimension \(2(n-1)\). Up to the volume factor, the form is \(\kappa\varphi=\sum_{i,j}\partial_{\beta_i}\partial_{\bar\beta_j}\varphi\,d\alpha_i\wedge d\bar\alpha_j\).}
\GSfield{Reconstruction}{Writing arbitrary \(n\)-vectors \(\alpha,\beta\) introduces the reparametrizations \((\alpha,\beta)\sim(\lambda\alpha,\beta+t_0\alpha)\). The graph chart removes this redundancy and fixes the derivative order.}
\GSfield{Fate}{Published soon after in \citet{gelfand1967planes}; the relevant definitions and Theorems 1--2 were checked in the Russian full text.}
\GSfield{Present status}{As of November 2026: established under the stated hypotheses; no unresolved problem is asserted by this occurrence.}
\GSfield{Assessment}{literal statement: correct with missing hypotheses; publication match: exact; verification: read the result; development: none; confidence: high}
\end{GSitem}
\begin{GSitem}{1965-82.10}{F}{Incomplete delta-kernel calculation at the lecture boundary}
\GSfield{Kind; source}{remark; \GSpage{P000132}}
\GSfield{The notes}{The last formula represents line constraints by delta factors; part of the final expression is cancelled or unfinished before Karpelevich begins.}
\GSfield{Statement}{The readable graph constraints define a plane delta kernel. No further definite equality can be reconstructed from the cancelled ending without specifying its coefficient, densities and differentiation conventions.}
\GSfield{Reconstruction}{Only the valid kernel construction is retained. The erased remainder is not promoted to a theorem.}
\GSfield{Fate}{\citet{gelfand1967planes} and \citet{gelfand1969forms} provide the surrounding method, but not a recoverable identity for the cancelled line.}
\GSfield{Present status}{As of November 2026: the incomplete expression remains textually indeterminate; the surrounding inversion method is established.}
\GSfield{Assessment}{literal statement: not determinable; publication match: thematic only; verification: unresolved lead; development: none; confidence: low}
\end{GSitem}

\GStalk{83}{1966-02-21}{F. I. Karpelevich}{P000132-P000133}{Klein on the spinning top}
A short exposition of Klein's Princeton lectures describes rotations by fractional linear transformations and reduces a top problem to an elliptic integral. The lecture crosses the page boundary: its final quadrature is at the top of P000133, before the dated Agranovich continuation. The two cyclic-momentum formulas contain a factor-of-two slip already present in the handwriting.

\begin{GSitem}{1965-83.1}{E}{Euler angles and a curve in SO(3)}
\GSfield{Kind; source}{construction; \GSpage{P000132}}
\GSfield{The notes}{A motion of the top is a curve in \(SO(3)\), factored as successive rotations with angles \(\varphi,\theta,\psi\).}
\GSfield{Statement}{A fixed-point rigid-body attitude is \(R(t)\in SO(3)\). Away from the polar coordinate singularities it admits Euler-angle coordinates, with the middle angle \(\theta\) measuring the inclination of the body axis.}
\GSfield{Reconstruction}{The factor ordering and sign convention fix the names of the conserved momenta; the subsequent equations are read consistently with the displayed kinetic energy.}
\GSfield{Fate}{The subject is explicitly Klein; \citet{klein1897top}, the published Princeton lectures, predates this exposition by nearly seventy years.}
\GSfield{Present status}{As of November 2026: established under the stated hypotheses; no unresolved problem is asserted by this occurrence.}
\GSfield{Assessment}{literal statement: correct; publication match: exact; verification: read the result; development: none; confidence: high}
\end{GSitem}
\begin{GSitem}{1965-83.2}{E}{Rotations as fractional linear transformations}
\GSfield{Kind; source}{construction; \GSpage{P000132}}
\GSfield{The notes}{The notes identify rotations with fractional linear transformations and draw the unitary two-by-two matrix.}
\GSfield{Statement}{The action of \(SU(2)\) on the Riemann sphere by \(z\mapsto(az+b)/(-\bar bz+\bar a)\), \(|a|^2+|b|^2=1\), has kernel \(\{\pm I\}\) and gives \(SO(3)\cong PSU(2)\).}
\GSfield{Reconstruction}{The isomorphism is with the projective transformations, not with the matrix group \(SU(2)\) itself.}
\GSfield{Fate}{The subject is explicitly Klein; \citet{klein1897top}, the published Princeton lectures, predates this exposition by nearly seventy years.}
\GSfield{Present status}{As of November 2026: established under the stated hypotheses; no unresolved problem is asserted by this occurrence.}
\GSfield{Assessment}{literal statement: correct with missing hypotheses; publication match: exact; verification: read the result; development: none; confidence: high}
\end{GSitem}
\begin{GSitem}{1965-83.3}{E}{Normalized spherical-top Lagrangian}
\GSfield{Kind; source}{construction; \GSpage{P000132}}
\GSfield{The notes}{\(T=\tfrac12(\dot\varphi^2+\dot\psi^2+2\dot\varphi\dot\psi\cos\theta+\dot\theta^2)\), \(V=p\cos\theta\).}
\GSfield{Statement}{For equal normalized principal moments of inertia and potential \(V=p\cos\theta\), set \(L=T-V\). The coordinates \(\varphi,\psi\) are cyclic and energy \(h=T+p\cos\theta\) is conserved.}
\GSfield{Reconstruction}{The printed kinetic energy is the spherical-inertia specialization; the general symmetric top has two moments of inertia and a different coefficient pattern.}
\GSfield{Fate}{The subject is explicitly Klein; \citet{klein1897top}, the published Princeton lectures, predates this exposition by nearly seventy years.}
\GSfield{Present status}{As of November 2026: established under the stated hypotheses; no unresolved problem is asserted by this occurrence.}
\GSfield{Assessment}{literal statement: correct with missing hypotheses; publication match: exact; verification: read the result; development: none; confidence: high}
\end{GSitem}
\begin{GSitem}{1965-83.4}{E}{Correct cyclic momenta}
\GSfield{Kind; source}{claim; \GSpage{P000132}}
\GSfield{The notes}{The handwritten first integrals contain \(\dot\varphi+2\dot\psi\cos\theta=n\) and \(\dot\psi+2\dot\varphi\cos\theta=e\).}
\GSfield{Statement}{Differentiation of the recorded \(L\) gives \(n=\dot\varphi+\dot\psi\cos\theta\), \(e=\dot\psi+\dot\varphi\cos\theta\). Thus for \(u=\cos\theta\), \(\dot\varphi=(n-eu)/(1-u^2)\), \(\dot\psi=(e-nu)/(1-u^2)\), when \(|u|<1\).}
\GSfield{Reconstruction}{\oursComputed{} The two factors 2 are inconsistent with the preceding one-half in \(T\). The corrections follow independently from differentiating that displayed kinetic energy, not from changing the physical model.}
\GSfield{Fate}{The subject is explicitly Klein; \citet{klein1897top}, the published Princeton lectures, predates this exposition by nearly seventy years.}
\GSfield{Present status}{As of November 2026: established under the stated hypotheses; no unresolved problem is asserted by this occurrence.}
\GSfield{Assessment}{literal statement: false as written; publication match: exact; verification: read the result; development: annotated; confidence: high}
\end{GSitem}
\begin{GSitem}{1965-83.5}{E}{The cubic governing nutation}
\GSfield{Kind; source}{method; \GSpage{P000132-P000133}}
\GSfield{The notes}{Elimination leads to \(dt=du/\sqrt{U(u)}\), where \(U\) is indicated as cubic but not completely expanded.}
\GSfield{Statement}{With the corrected momenta, \(\dot u^2=U(u)=2(h-pu)(1-u^2)-n^2-e^2+2neu\). Equivalently, \(U(u)=2pu^3-2hu^2+2(ne-p)u+2h-n^2-e^2\). Physical intervals also require \(|u|\leq1\) and \(U(u)\geq0\).}
\GSfield{Reconstruction}{\oursComputed{} This full polynomial is an independent reconstruction; it is not a completed transcription of the truncated notebook expression.}
\GSfield{Fate}{The subject is explicitly Klein; \citet{klein1897top}, the published Princeton lectures, predates this exposition by nearly seventy years.}
\GSfield{Present status}{As of November 2026: established under the stated hypotheses; no unresolved problem is asserted by this occurrence.}
\GSfield{Assessment}{literal statement: correct with missing hypotheses; publication match: exact; verification: read the result; development: annotated; confidence: high}
\end{GSitem}
\begin{GSitem}{1965-83.6}{E}{Elliptic quadrature and degenerate limits}
\GSfield{Kind; source}{theorem; \GSpage{P000133}}
\GSfield{The notes}{The top of P000133 identifies the integral as elliptic.}
\GSfield{Statement}{For a cubic \(U\) with distinct roots, the curve \(v^2=U(u)\) has genus one, and \(t-t_0=\pm\int_{u_0}^u ds/\sqrt{U(s)}\) gives the nutation locally; its inversion is elliptic. Multiple-root and degree-drop cases degenerate and require separate elementary limiting formulas.}
\GSfield{Reconstruction}{The sign chooses a motion branch. The assertion does not exclude stationary solutions or repeated-root trajectories.}
\GSfield{Fate}{The subject is explicitly Klein; \citet{klein1897top}, the published Princeton lectures, predates this exposition by nearly seventy years.}
\GSfield{Present status}{As of November 2026: established under the stated hypotheses; no unresolved problem is asserted by this occurrence.}
\GSfield{Assessment}{literal statement: correct with missing hypotheses; publication match: exact; verification: read the result; development: none; confidence: high}
\end{GSitem}

\GStalk{84}{1966-02-28}{M. S. Agranovich}{P000133-P000136}{Symbolic calculus, elliptic operators and index}
The continuation assembles the local calculus on a manifold and derives the Fredholm and regularity theorems from a parametrix. It then states Seeley's special degree formula for the index and ends with Hörmander's oscillatory characterization. The exact Seeley theorem, including its equivalence with vanishing of lower-rank indices, was checked in the full paper's Russian translation; the 1965 English original is the historical publication.

\begin{GSitem}{1965-84.1}{E}{Right-quantized symbol and order}
\GSfield{Kind; source}{definition; \GSpage{P000133}}
\GSfield{The notes}{This lecture displays the kernel using \(\sigma(y,\xi)\), and repeats the low-frequency cutoff and \(H^{\ell+s}\to H^\ell\) convention.}
\GSfield{Statement}{Right quantization is \(Au(x)=(2\pi)^{-n}\iint e^{i(x-y)\cdot\xi}\sigma(y,\xi)u(y)\,dy\,d\xi\). Smooth order-\(s\) symbols give the indicated Sobolev loss under proper-support or uniform global hypotheses.}
\GSfield{Reconstruction}{The change from the earlier left convention affects lower terms, not the principal symbol.}
\GSfield{Fate}{Already in \citet{seeley1965bundles}, \S7 and the elliptic theory, and \citet{agranovich1965elliptic}; the seminar is expository.}
\GSfield{Present status}{As of November 2026: established under the stated hypotheses; no unresolved problem is asserted by this occurrence.}
\GSfield{Assessment}{literal statement: correct with missing hypotheses; publication match: exact; verification: read the result; development: none; confidence: high}
\end{GSitem}
\begin{GSitem}{1965-84.2}{E}{Product symbol and scalar commutator order}
\GSfield{Kind; source}{theorem; \GSpage{P000133}}
\GSfield{The notes}{\(AB\) has symbol \(\sigma_1\sigma_2\) modulo a subordinate operator.}
\GSfield{Statement}{For smooth classical operators of orders \(s_1,s_2\), the product principal symbol is \(\sigma_{s_1}(A)\sigma_{s_2}(B)\). If the leading symbols commute, in particular for scalar operators, \([A,B]\) has order at most \(s_1+s_2-1\).}
\GSfield{Reconstruction}{The lower-order commutator conclusion does not hold for arbitrary noncommuting matrix symbols.}
\GSfield{Fate}{Already in \citet{seeley1965bundles}, \S7 and the elliptic theory, and \citet{agranovich1965elliptic}; the seminar is expository.}
\GSfield{Present status}{As of November 2026: established under the stated hypotheses; no unresolved problem is asserted by this occurrence.}
\GSfield{Assessment}{literal statement: correct with missing hypotheses; publication match: exact; verification: read the result; development: none; confidence: high}
\end{GSitem}
\begin{GSitem}{1965-84.3}{E}{Multiplication cutoffs and the ambiguous dash}
\GSfield{Kind; source}{theorem; \GSpage{P000133}}
\GSfield{The notes}{The line \(\varphi A\)--\(A\varphi\) is followed by the statement that each has symbol \(\varphi\sigma\) to subordinate terms.}
\GSfield{Statement}{Both \(\varphi A\) and \(A\varphi\) have principal symbol \(\varphi\sigma_s(A)\); their difference has order at most \(s-1\).}
\GSfield{Reconstruction}{The horizontal mark is read as separating two operators, in accordance with the Russian word meaning each. Reading it as a subtraction with principal symbol \(\varphi\sigma\) would be false.}
\GSfield{Fate}{Already in \citet{seeley1965bundles}, \S7 and the elliptic theory, and \citet{agranovich1965elliptic}; the seminar is expository.}
\GSfield{Present status}{As of November 2026: established under the stated hypotheses; no unresolved problem is asserted by this occurrence.}
\GSfield{Assessment}{literal statement: correct with missing hypotheses; publication match: exact; verification: read the result; development: annotated; confidence: high}
\end{GSitem}
\begin{GSitem}{1965-84.4}{E}{Separated cutoffs eliminate the principal symbol}
\GSfield{Kind; source}{theorem; \GSpage{P000133-P000134}}
\GSfield{The notes}{For separated \(\varphi,\psi\), \(\varphi A\psi\) is subordinate, with a one-order Sobolev gain displayed.}
\GSfield{Statement}{If \(\varphi\psi=0\), the order-\(s\) principal symbol of \(\varphi A\psi\) vanishes. If their supports are disjoint, the smooth calculus gives an infinitely smoothing operator, a stronger assertion than the recorded one-order gain.}
\GSfield{Reconstruction}{Pointwise vanishing product and disjoint closed supports are distinguished; only the latter immediately gives the full off-diagonal smoothing statement.}
\GSfield{Fate}{Already in \citet{seeley1965bundles}, \S7 and the elliptic theory, and \citet{agranovich1965elliptic}; the seminar is expository.}
\GSfield{Present status}{As of November 2026: established under the stated hypotheses; no unresolved problem is asserted by this occurrence.}
\GSfield{Assessment}{literal statement: correct with missing hypotheses; publication match: exact; verification: read the result; development: none; confidence: high}
\end{GSitem}
\begin{GSitem}{1965-84.5}{E}{Vanishing principal symbol does not annihilate every operator}
\GSfield{Kind; source}{claim; \GSpage{P000133}}
\GSfield{The notes}{A smoothing order-zero operator is said to have zero symbol and therefore to vanish.}
\GSfield{Statement}{An operator of order zero with improved Sobolev order has zero principal order-zero symbol. In the full calculus this means that it belongs to a lower-order class; it need not be zero. Only the restricted model with a sole exactly homogeneous symbol and no independent remainder has the asserted injectivity.}
\GSfield{Reconstruction}{\ours{} A nonzero operator with a smooth compactly supported kernel is a direct counterexample to the unrestricted conclusion. The recorded proof uses a positive symbol obtained from the adjoint product and angular changes of variables; positivity after finitely many angular localizations detects a nonzero principal symbol. This justifies symbol vanishing, not the disappearance of an arbitrary lower-order remainder.}
\GSfield{Fate}{The kernel of the principal-symbol map is explicitly the lower-order class in \citet{seeley1965bundles}, Theorem 7.1.}
\GSfield{Present status}{As of November 2026: established under the stated hypotheses; no unresolved problem is asserted by this occurrence.}
\GSfield{Assessment}{literal statement: false as written; publication match: exact; verification: read the result; development: annotated; confidence: high}
\end{GSitem}
\begin{GSitem}{1965-84.6}{E}{Change of coordinates and cotangent transformation}
\GSfield{Kind; source}{theorem; \GSpage{P000133-P000134}}
\GSfield{The notes}{The symbol transforms like the highest characteristic polynomial of a differential operator.}
\GSfield{Statement}{Under \(y=\kappa(x)\), the transformed principal symbol is \[\widetilde\sigma(y,\eta)=\sigma(\kappa^{-1}(y),D\kappa(\kappa^{-1}(y))^T\eta)\]. Bundle frames transform by the corresponding source and target changes of basis. Quantization corrections have order at most \(s-1\).}
\GSfield{Reconstruction}{This covector rule is stated explicitly to avoid confusing the tangent and cotangent Jacobians.}
\GSfield{Fate}{Already in \citet{seeley1965bundles}, \S7 and the elliptic theory, and \citet{agranovich1965elliptic}; the seminar is expository.}
\GSfield{Present status}{As of November 2026: established under the stated hypotheses; no unresolved problem is asserted by this occurrence.}
\GSfield{Assessment}{literal statement: correct with missing hypotheses; publication match: exact; verification: read the result; development: none; confidence: high}
\end{GSitem}
\begin{GSitem}{1965-84.7}{E}{Historical attributions in the calculus}
\GSfield{Kind; source}{remark; \GSpage{P000134}}
\GSfield{The notes}{The notes list Calderón, Zygmund, Mikhlin, Kohn, Nirenberg, Vishik, Eskin, Seeley and a surname retained with Dynin/Dynkin uncertainty.}
\GSfield{Statement}{The manifold singular-operator reference intended in this context is A. S. Dynin, whose 1961 paper is cited in the contemporaneous literature; it is not a mathematical result of E. B. Dynkin.}
\GSfield{Reconstruction}{The scan has a Dynin/Dynkin-like handwritten form. This is an identification supported by the bibliographies, not a license to silently alter the diplomatic text.}
\GSfield{Fate}{\citet{seeley1965bundles}, bibliography item 5, identifies Dynin and the 1961 paper; \citet{kohnnirenberg1965algebra} confirms the 1965 calculus chronology.}
\GSfield{Present status}{As of November 2026: established under the stated hypotheses; no unresolved problem is asserted by this occurrence.}
\GSfield{Assessment}{literal statement: correct with missing hypotheses; publication match: partial; verification: read the result; development: annotated; confidence: high}
\end{GSitem}
\begin{GSitem}{1965-84.8}{E}{Global operators from local representatives}
\GSfield{Kind; source}{construction; \GSpage{P000134}}
\GSfield{The notes}{The notes impose separated-support and local-coordinate conditions, then write \(A=\sum\varphi_\nu A_\nu\psi_\nu+T\).}
\GSfield{Statement}{A global classical operator can be defined by compatible local operators with the cotangent transformation rule and a smooth off-diagonal kernel. A finite partition of unity on a closed manifold yields local quantizations patched by cutoffs; changing patching changes only lower-order terms at the principal-symbol level.}
\GSfield{Reconstruction}{The notebook employs an order-minus-one remainder in its coarser calculus; a full classical construction refines this to asymptotic symbols modulo smoothing.}
\GSfield{Fate}{\citet{seeley1965bundles}, Definition 7.1 and Theorem 7.1, give the precise construction.}
\GSfield{Present status}{As of November 2026: established under the stated hypotheses; no unresolved problem is asserted by this occurrence.}
\GSfield{Assessment}{literal statement: correct with missing hypotheses; publication match: exact; verification: read the result; development: none; confidence: high}
\end{GSitem}
\begin{GSitem}{1965-84.9}{E}{Quantizing a global homogeneous symbol}
\GSfield{Kind; source}{construction; \GSpage{P000134}}
\GSfield{The notes}{The concluding remark says an operator can be constructed from a symbol on the cotangent bundle.}
\GSfield{Statement}{Every smooth degree-\(s\) homogeneous bundle homomorphism on \(T^*M\setminus0\) is the principal symbol of a classical pseudodifferential operator, obtained by local quantization and a partition of unity. The choice is unique only modulo order \(s-1\).}
\GSfield{Reconstruction}{Existence does not specify a canonical quantization, and symbol equality does not mean equality of operators.}
\GSfield{Fate}{\citet{seeley1965bundles}, Theorem 7.1.}
\GSfield{Present status}{As of November 2026: established under the stated hypotheses; no unresolved problem is asserted by this occurrence.}
\GSfield{Assessment}{literal statement: correct with missing hypotheses; publication match: exact; verification: read the result; development: none; confidence: high}
\end{GSitem}
\begin{GSitem}{1965-84.10}{E}{The Sobolev mapping line loses an order shift}
\GSfield{Kind; source}{claim; \GSpage{P000134}}
\GSfield{The notes}{The handwritten line says an operator of order \(\ell\) is bounded \(H_\ell\to H_\ell\).}
\GSfield{Statement}{For an operator of order \(s\), the correct mapping is \(A:H^{t+s}\to H^t\), equivalently \(H^t\to H^{t-s}\), for every real Sobolev index \(t\) in the smooth closed-manifold calculus.}
\GSfield{Reconstruction}{The two identical subscripts are visibly present in the scan. The preceding and subsequent lecture conventions determine the repair; it is a mathematical annotation, not an OCR correction.}
\GSfield{Fate}{Already in \citet{seeley1965bundles}, \S7 and the elliptic theory, and \citet{agranovich1965elliptic}; the seminar is expository.}
\GSfield{Present status}{As of November 2026: established under the stated hypotheses; no unresolved problem is asserted by this occurrence.}
\GSfield{Assessment}{literal statement: false as written; publication match: exact; verification: read the result; development: annotated; confidence: high}
\end{GSitem}
\begin{GSitem}{1965-84.11}{E}{Local proof of global composition}
\GSfield{Kind; source}{method; \GSpage{P000134-P000135}}
\GSfield{The notes}{A double partition-of-unity sum for \(AB\) is rearranged by commuting cutoffs modulo subordinate terms.}
\GSfield{Statement}{Insert cutoffs equal to one on relevant supports; terms away from a common chart are lower order, while terms in a common chart have product principal symbols. Summing gives \(\sigma(AB)=\sigma(A)\sigma(B)\) globally.}
\GSfield{Reconstruction}{The rearrangement is modulo the correct lower-order class; it is not an exact equality of the displayed localized operators.}
\GSfield{Fate}{\citet{seeley1965bundles}, proof of Theorem 7.2, contains the corresponding localization argument.}
\GSfield{Present status}{As of November 2026: established under the stated hypotheses; no unresolved problem is asserted by this occurrence.}
\GSfield{Assessment}{literal statement: correct with missing hypotheses; publication match: exact; verification: read the result; development: none; confidence: high}
\end{GSitem}
\begin{GSitem}{1965-84.12}{E}{Ellipticity and the Fredholm property}
\GSfield{Kind; source}{theorem; \GSpage{P000135}}
\GSfield{The notes}{\(\sigma(P,\xi_P)\neq0\), or nonzero determinant, is declared equivalent to being Noetherian.}
\GSfield{Statement}{On a closed smooth manifold, a classical operator \(A\) of order \(s\) between equal-rank complex bundles is elliptic exactly when \(A:H^{t+s}(E)\to H^t(F)\) is Fredholm, for one and then every \(t\). Ellipticity means invertibility of the principal symbol for every nonzero covector.}
\GSfield{Reconstruction}{Compactness and absence of boundary are essential. On a noncompact manifold or a manifold with boundary, this statement needs additional global or boundary conditions.}
\GSfield{Fate}{Already in \citet{seeley1965bundles}, \S7 and the elliptic theory, and \citet{agranovich1965elliptic}; the seminar is expository.}
\GSfield{Present status}{As of November 2026: established under the stated hypotheses; no unresolved problem is asserted by this occurrence.}
\GSfield{Assessment}{literal statement: correct with missing hypotheses; publication match: exact; verification: read the result; development: none; confidence: high}
\end{GSitem}
\begin{GSitem}{1965-84.13}{E}{Parametrix and compact remainders}
\GSfield{Kind; source}{method; \GSpage{P000135}}
\GSfield{The notes}{Quantizing \(\sigma^{-1}\) gives \(BA=I+T_1\), \(AB=I+T_2\), and finite defects.}
\GSfield{Statement}{An elliptic principal symbol gives an order-\(-s\) parametrix with remainders of negative order; on a closed manifold these are compact on the relevant Sobolev spaces. Consequently kernel and cokernel are finite-dimensional and the range is closed. In the full calculus the remainders can be made smoothing.}
\GSfield{Reconstruction}{Compactness uses the compact Sobolev embedding. Finite-dimensional kernel alone would not establish the Fredholm property.}
\GSfield{Fate}{Already in \citet{seeley1965bundles}, \S7 and the elliptic theory, and \citet{agranovich1965elliptic}; the seminar is expository.}
\GSfield{Present status}{As of November 2026: established under the stated hypotheses; no unresolved problem is asserted by this occurrence.}
\GSfield{Assessment}{literal statement: correct with missing hypotheses; publication match: exact; verification: read the result; development: none; confidence: high}
\end{GSitem}
\begin{GSitem}{1965-84.14}{E}{Fredholm estimate forces ellipticity}
\GSfield{Kind; source}{method; \GSpage{P000135}}
\GSfield{The notes}{For \(s=1\), a Fredholm estimate \(\|u\|_1\leq C(\|Au\|_0+\|u\|_0)\) and localized plane waves give the converse.}
\GSfield{Statement}{If a principal symbol has a null vector at \((x_0,\xi_0)\), localized high-frequency sections concentrated near \(x_0\), with phase \(\lambda x\cdot\xi_0\), contradict the Fredholm estimate after normalization. A shrinking support or a two-stage localization limit controls variation of the symbol in \(x\).}
\GSfield{Reconstruction}{\ours{} The explanatory verification here is supplied by us for the established result or correction. A single fixed bump does not literally freeze a variable symbol at \(x_0\); the localization limit is the missing step.}
\GSfield{Fate}{\citet{seeley1965bundles}, Theorem 7.3 and the subsequent Fredholm discussion, verify the high-frequency argument.}
\GSfield{Present status}{As of November 2026: established under the stated hypotheses; no unresolved problem is asserted by this occurrence.}
\GSfield{Assessment}{literal statement: correct with missing hypotheses; publication match: exact; verification: read the result; development: none; confidence: high}
\end{GSitem}
\begin{GSitem}{1965-84.15}{E}{Elliptic regularity and the incorrect remainder argument}
\GSfield{Kind; source}{theorem; \GSpage{P000135}}
\GSfield{The notes}{For \(Au=f\), order zero, \(u\in H^0\), \(f\in H^1\), the notes conclude \(u\in H^1\) using \(u+T_2f=Bf\).}
\GSfield{Statement}{From \(BA=I+T_1\) one obtains \(u+T_1u=Bf\). The remainder sends \(H^0\) into \(H^1\), giving the asserted gain. Iteration yields \(Au\in H^t\Rightarrow u\in H^{t+s}\) for elliptic order \(s\), locally for distributions.}
\GSfield{Reconstruction}{The source term \(T_2f\) is the handwritten error; the remainder acts on \(u\), and the left parametrix determines its index.}
\GSfield{Fate}{Already in \citet{seeley1965bundles}, \S7 and the elliptic theory, and \citet{agranovich1965elliptic}; the seminar is expository.}
\GSfield{Present status}{As of November 2026: established under the stated hypotheses; no unresolved problem is asserted by this occurrence.}
\GSfield{Assessment}{literal statement: false as written; publication match: exact; verification: read the result; development: annotated; confidence: high}
\end{GSitem}
\begin{GSitem}{1965-84.16}{E}{Homotopy invariance of the index}
\GSfield{Kind; source}{theorem; \GSpage{P000135-P000136}}
\GSfield{The notes}{The index \(\chi(A)=\dim\ker A-\dim\operatorname{coker}A\) does not change under small symbol deformations.}
\GSfield{Statement}{The Fredholm index is locally constant in operator norm. A continuous path of elliptic principal symbols, quantized continuously between fixed Sobolev spaces, therefore gives the same index at both ends.}
\GSfield{Reconstruction}{The path must remain elliptic; an arbitrary continuous pointwise deformation without uniform control is not the asserted operator-norm statement.}
\GSfield{Fate}{Already in \citet{seeley1965bundles}, \S7 and the elliptic theory, and \citet{agranovich1965elliptic}; the seminar is expository.}
\GSfield{Present status}{As of November 2026: established under the stated hypotheses; no unresolved problem is asserted by this occurrence.}
\GSfield{Assessment}{literal statement: correct with missing hypotheses; publication match: exact; verification: read the result; development: none; confidence: high}
\end{GSitem}
\begin{GSitem}{1965-84.17}{E}{Reducing the order without changing the index}
\GSfield{Kind; source}{construction; \GSpage{P000136}}
\GSfield{The notes}{\(B=A\Lambda^s\) and \(\chi(B)=\chi(A)+\chi(\Lambda^s)\), with the last index zero.}
\GSfield{Statement}{Choose a positive invertible elliptic order-one operator, for example \((I+\Delta)^{1/2}\) for the nonnegative Laplacian on a closed manifold. Its real powers shift Sobolev order isomorphically. Multiplying by the appropriate power reduces any order to zero without changing the index.}
\GSfield{Reconstruction}{Using \(\sqrt{-\Delta}\) literally requires handling its kernel and sign convention. A self-adjoint Fredholm operator has index zero but need not be invertible.}
\GSfield{Fate}{Already in \citet{seeley1965bundles}, \S7 and the elliptic theory, and \citet{agranovich1965elliptic}; the seminar is expository.}
\GSfield{Present status}{As of November 2026: established under the stated hypotheses; no unresolved problem is asserted by this occurrence.}
\GSfield{Assessment}{literal statement: correct with missing hypotheses; publication match: exact; verification: read the result; development: none; confidence: high}
\end{GSitem}
\begin{GSitem}{1965-84.18}{E}{Seeley degree formula and the lower-rank obstruction}
\GSfield{Kind; source}{theorem; \GSpage{P000136}}
\GSfield{The notes}{The notes give \(\operatorname{ind}A=\deg(\rho\sigma)/(n-1)!\) if and only if every \((n-1)\)-square system has zero index; boundaries in \(\mathbb R^{n+1}\) are mentioned.}
\GSfield{Statement}{For elliptic \(n\times n\) systems on a closed \(n\)-manifold with unitary principal symbol \(\sigma:S^*M\to U(n)\), suitable fixed orientations make \((n-1)!\operatorname{ind}A=\deg(\rho\circ\sigma)\) for every such system if and only if every elliptic \((n-1)\times(n-1)\) system on \(M\) has index zero. Here \(\rho:U(n)\to S^{2n-1}\) selects a column.}
\GSfield{Reconstruction}{This is a theorem about systems on trivial bundles with a global matrix symbol. The quantified equivalence, orientation choice and manifold restriction prevent treating it as a universal degree formula for all bundles.}
\GSfield{Fate}{Exact match: \citet{seeley1965bundles}, Appendix A.3, Theorem 2 (Russian translation p. 95), including the boundary example.}
\GSfield{Present status}{As of November 2026: established under the stated hypotheses; no unresolved problem is asserted by this occurrence.}
\GSfield{Assessment}{literal statement: correct with missing hypotheses; publication match: exact; verification: read the result; development: none; confidence: high}
\end{GSitem}
\begin{GSitem}{1965-84.19}{E}{External product of elliptic operators}
\GSfield{Kind; source}{theorem; \GSpage{P000136}}
\GSfield{The notes}{The operation \(A\#B\) satisfies \(\operatorname{ind}(A\#B)=\operatorname{ind}A\operatorname{ind}B\).}
\GSfield{Statement}{The elliptic external product, formed with the standard graded tensor-product symbol and realized on the product manifold, has index equal to the product of the two indices. It is distinct from ordinary composition, whose index is additive.}
\GSfield{Reconstruction}{The notebook does not define \(\#\); its identification is fixed by the Seeley reference and the displayed multiplicativity.}
\GSfield{Fate}{\citet{seeley1965bundles}, Theorem 13.2 and Appendix A.2, give the operation used here.}
\GSfield{Present status}{As of November 2026: established under the stated hypotheses; no unresolved problem is asserted by this occurrence.}
\GSfield{Assessment}{literal statement: correct with missing hypotheses; publication match: exact; verification: read the result; development: none; confidence: high}
\end{GSitem}
\begin{GSitem}{1965-84.20}{E}{Hörmander oscillatory characterization}
\GSfield{Kind; source}{method; \GSpage{P000136}}
\GSfield{The notes}{\(e^{-i\lambda g}A(fe^{i\lambda g})\sim\sum_jP_j(f,g)\lambda^{s_j}\), \(s_j\downarrow-\infty\), with \(dg\neq0\) on \(\operatorname{supp}f\), bounded remainders and continuous \(A:C_c^\infty\to C^\infty\).}
\GSfield{Statement}{A classical smooth operator can be characterized through its complete high-frequency action on localized nonstationary phases. In the usual classical order-\(s\) case the exponents are \(s-j\); each normalized remainder is bounded in every prescribed \(C^q\) seminorm, uniformly for compact families of admissible phases. Linear phases recover the symbol.}
\GSfield{Reconstruction}{The notebook allows a more general decreasing sequence of exponents and leaves the precise converse class implicit. This entry records the characterization method, not an unverified full equivalence for every such sequence.}
\GSfield{Fate}{\citet{hormander1965pseudo} predates the talk. Publisher metadata and the 1965 publication were checked; the exact generalized-exponent formulation has not been verified against its full theorem.}
\GSfield{Present status}{As of November 2026: established under the stated hypotheses; no unresolved problem is asserted by this occurrence.}
\GSfield{Assessment}{literal statement: correct with missing hypotheses; publication match: partial; verification: publisher metadata checked; development: none; confidence: medium}
\end{GSitem}

\GStalk{85}{1966-02-28}{A. A. Kirillov}{P000137-P000138}{Borel and Tits: algebraic and reductive groups}
This is an introductory report on Borel and Tits, not an announcement of a new classification. It passes from schemes and Chevalley's structure theorem to radicals, tori and root systems. Connectedness and the ground field are repeatedly tacit. The example calling an algebraic torus an abelian variety is incompatible with the standard algebraic-group meaning of torus and is retained as a source-level error.

\begin{GSitem}{1965-85.1}{E}{Schemes and Spec as the algebraic-geometric language}
\GSfield{Kind; source}{definition; \GSpage{P000137}}
\GSfield{The notes}{An algebraic variety is called a locally affine scheme; \(\operatorname{Spec}A\) is built from the prime ideals of \(A\).}
\GSfield{Statement}{An affine scheme is the locally ringed space \((\operatorname{Spec}A,\mathcal O_A)\), with \(\mathcal O_A(D(f))=A_f\). Schemes are locally affine. To describe an algebraic variety over a field one additionally specifies finite type, separatedness and the chosen reducedness/irreducibility convention.}
\GSfield{Reconstruction}{The source conflates a scheme with a variety. These extra conditions are a modern terminology clarification, not a new theorem.}
\GSfield{Fate}{This is background for the explicitly cited \citet{borel1965reductifs}, published in 1965; no later priority claim is made.}
\GSfield{Present status}{As of November 2026: established under the stated hypotheses; no unresolved problem is asserted by this occurrence.}
\GSfield{Assessment}{literal statement: correct with missing hypotheses; publication match: partial; verification: read the result; development: none; confidence: high}
\end{GSitem}
\begin{GSitem}{1965-85.2}{E}{Prime ideals of the affine plane: the omitted generic point}
\GSfield{Kind; source}{claim; \GSpage{P000137}}
\GSfield{The notes}{For \(A=\mathbb C[x_1,x_2]\), primes are said to be point ideals or principal ideals generated by irreducible polynomials.}
\GSfield{Statement}{The primes are \((0)\), the height-one ideals \((P)\) for irreducible \(P\), and maximal point ideals \((x_1-a,x_2-b)\). The zero ideal is the generic point of the irreducible affine plane.}
\GSfield{Reconstruction}{\ours{} The explanatory verification here is supplied by us for the established result or correction. The missing zero ideal makes the literal exhaustive list false; the repair follows from unique factorization, dimension two and the algebraically closed base field.}
\GSfield{Fate}{This is background for the explicitly cited \citet{borel1965reductifs}, published in 1965; no later priority claim is made.}
\GSfield{Present status}{As of November 2026: established under the stated hypotheses; no unresolved problem is asserted by this occurrence.}
\GSfield{Assessment}{literal statement: false as written; publication match: partial; verification: read the result; development: annotated; confidence: high}
\end{GSitem}
\begin{GSitem}{1965-85.3}{E}{Algebraic groups and linear algebraic groups}
\GSfield{Kind; source}{definition; \GSpage{P000137}}
\GSfield{The notes}{The notes require a regular everywhere-defined multiplication and describe linear groups by algebraic equations inside \(GL(n,K)\).}
\GSfield{Statement}{An algebraic group is a finite-type group scheme over the base field, with multiplication, identity and inversion morphisms satisfying the group axioms. In the smooth-variety convention these maps are regular. An affine algebraic group admits a faithful finite-dimensional representation and thus a closed embedding into some \(GL_n\).}
\GSfield{Reconstruction}{The inverse and identity morphisms and the finite-type requirement complete the abbreviated definition.}
\GSfield{Fate}{This is background for the explicitly cited \citet{borel1965reductifs}, published in 1965; no later priority claim is made.}
\GSfield{Present status}{As of November 2026: established under the stated hypotheses; no unresolved problem is asserted by this occurrence.}
\GSfield{Assessment}{literal statement: correct with missing hypotheses; publication match: exact; verification: read the result; development: none; confidence: high}
\end{GSitem}
\begin{GSitem}{1965-85.4}{E}{Chevalley structure theorem}
\GSfield{Kind; source}{theorem; \GSpage{P000137}}
\GSfield{The notes}{A maximal linear subgroup \(G_0\) is asserted with abelian-variety quotient \(G/G_0\).}
\GSfield{Statement}{For a smooth connected algebraic group over a perfect field, there is a unique smooth connected affine normal subgroup \(G_{\mathrm{aff}}\) such that \(G/G_{\mathrm{aff}}\) is an abelian variety.}
\GSfield{Reconstruction}{Connectedness, normality, smoothness and the field hypothesis are made explicit. Stronger scheme formulations over imperfect fields require a different smoothness statement.}
\GSfield{Fate}{Known before the talk; the dated Barsotti--Chevalley theorem and a proof are checked in \citet{milne2017groups}, Theorem 8.27 and Notes 8.30--8.33.}
\GSfield{Present status}{As of November 2026: established under the stated hypotheses; no unresolved problem is asserted by this occurrence.}
\GSfield{Assessment}{literal statement: correct with missing hypotheses; publication match: exact; verification: read the result; development: none; confidence: high}
\end{GSitem}
\begin{GSitem}{1965-85.5}{E}{Complete connected group varieties are abelian varieties}
\GSfield{Kind; source}{theorem; \GSpage{P000137}}
\GSfield{The notes}{An abelian variety is a complete group, and is then commutative.}
\GSfield{Statement}{A smooth connected complete group variety over a field is commutative and projective; this is the abelian-variety notion used here. A nonconnected finite group scheme does not become commutative merely because it is complete.}
\GSfield{Reconstruction}{Connectedness is necessary for the stated implication, and completeness is the algebraic properness condition rather than compactness of rational points.}
\GSfield{Fate}{The pre-seminar result is checked in \citet{milne2017groups}, 8.20 and Theorem 8.45, with historical attribution.}
\GSfield{Present status}{As of November 2026: established under the stated hypotheses; no unresolved problem is asserted by this occurrence.}
\GSfield{Assessment}{literal statement: correct with missing hypotheses; publication match: exact; verification: read the result; development: none; confidence: high}
\end{GSitem}
\begin{GSitem}{1965-85.6}{E}{Homomorphisms between affine groups and abelian varieties}
\GSfield{Kind; source}{theorem; \GSpage{P000137}}
\GSfield{The notes}{Every homomorphism from an abelian variety to a linear group and conversely is declared trivial.}
\GSfield{Statement}{A homomorphism from an abelian variety to an affine algebraic group is trivial. A homomorphism from a smooth connected affine algebraic group to an abelian variety is also trivial. Without connectedness, a finite constant group can map nontrivially into torsion points.}
\GSfield{Reconstruction}{\ours{} The explanatory verification here is supplied by us for the established result or correction. The first direction follows from constancy of regular functions on a proper connected variety. In the reverse direction the connected affine image is simultaneously complete, hence trivial.}
\GSfield{Fate}{This uses the classical affine-quotient and completeness results; modern precise forms are checked in \citet{milne2017groups}, Proposition 8.1 and \S8e.}
\GSfield{Present status}{As of November 2026: established under the stated hypotheses; no unresolved problem is asserted by this occurrence.}
\GSfield{Assessment}{literal statement: correct with missing hypotheses; publication match: exact; verification: read the result; development: none; confidence: high}
\end{GSitem}
\begin{GSitem}{1965-85.7}{E}{An algebraic torus is not an abelian variety}
\GSfield{Kind; source}{claim; \GSpage{P000137}}
\GSfield{The notes}{The scan explicitly says \ru{Например, алгебраический тор --- аб. мн-зие}.}
\GSfield{Statement}{A positive-dimensional algebraic torus, geometrically \((\mathbb G_m)^r\), is affine and not complete, so is not an abelian variety. A compact complex analytic torus is a different object; it is an abelian variety only with the appropriate projectivity condition.}
\GSfield{Reconstruction}{No evidence establishes that an adjective or a negation was omitted by the transcriber. The literal algebraic statement is false; a possible analytic intention remains conjectural.}
\GSfield{Fate}{The distinction is already present in the algebraic-group framework of \citet{borel1965reductifs}.}
\GSfield{Present status}{As of November 2026: established under the stated hypotheses; no unresolved problem is asserted by this occurrence.}
\GSfield{Assessment}{literal statement: false as written; publication match: exact; verification: read the result; development: annotated; confidence: high}
\end{GSitem}
\begin{GSitem}{1965-85.8}{E}{The additive and multiplicative groups}
\GSfield{Kind; source}{construction; \GSpage{P000137}}
\GSfield{The notes}{\(G_a,G_m\) are the additive and multiplicative groups of the field; a unipotent two-by-two matrix is drawn.}
\GSfield{Statement}{\(\mathbb G_a(R)=(R,+)\) is represented by \(x\mapsto\begin{pmatrix}1&x\\0&1\end{pmatrix}\); \(\mathbb G_m(R)=R^\times\) is represented by invertible one-by-one matrices. These are affine algebraic groups.}
\GSfield{Reconstruction}{The functorial formulation supplies the scheme meaning without changing the elementary examples.}
\GSfield{Fate}{This is background for the explicitly cited \citet{borel1965reductifs}, published in 1965; no later priority claim is made.}
\GSfield{Present status}{As of November 2026: established under the stated hypotheses; no unresolved problem is asserted by this occurrence.}
\GSfield{Assessment}{literal statement: correct; publication match: exact; verification: read the result; development: none; confidence: high}
\end{GSitem}
\begin{GSitem}{1965-85.9}{E}{Solvable groups and one-dimensional factors}
\GSfield{Kind; source}{theorem; \GSpage{P000137}}
\GSfield{The notes}{Solvability is described by a normal chain with factors \(G_a\) or \(G_m\), with a parenthetical warning that the definition is imprecise.}
\GSfield{Statement}{For a smooth connected solvable affine algebraic group over an algebraically closed field, a subnormal series can be chosen with successive quotients isomorphic to \(\mathbb G_a\) or \(\mathbb G_m\). The abstract derived-series definition of solvability is the invariant starting point.}
\GSfield{Reconstruction}{The field and connectedness assumptions explain the notebook warning. Subnormal means normal in the preceding member; all members need not be normal in the whole group.}
\GSfield{Fate}{This is background for the explicitly cited \citet{borel1965reductifs}, published in 1965; no later priority claim is made.}
\GSfield{Present status}{As of November 2026: established under the stated hypotheses; no unresolved problem is asserted by this occurrence.}
\GSfield{Assessment}{literal statement: correct with missing hypotheses; publication match: exact; verification: read the result; development: none; confidence: high}
\end{GSitem}
\begin{GSitem}{1965-85.10}{E}{Radical and unipotent radical}
\GSfield{Kind; source}{definition; \GSpage{P000137-P000138}}
\GSfield{The notes}{The radical is the maximal connected solvable normal subgroup; the unipotent radical is described using nilpotent matrices and \((u-E)^k=0\).}
\GSfield{Statement}{For a smooth connected affine algebraic group over an algebraically closed field, \(R(G)\) is its largest connected solvable normal subgroup and \(R_u(G)\) its largest connected unipotent normal subgroup. In a faithful representation, unipotent elements satisfy \((u-I)^N=0\); the matrix \(u\) itself is invertible, not nilpotent.}
\GSfield{Reconstruction}{The nilpotent adjective applies to \(u-I\), as the formula on P000138 confirms.}
\GSfield{Fate}{This is background for the explicitly cited \citet{borel1965reductifs}, published in 1965; no later priority claim is made.}
\GSfield{Present status}{As of November 2026: established under the stated hypotheses; no unresolved problem is asserted by this occurrence.}
\GSfield{Assessment}{literal statement: correct with missing hypotheses; publication match: exact; verification: read the result; development: none; confidence: high}
\end{GSitem}
\begin{GSitem}{1965-85.11}{E}{Reductive and semisimple groups}
\GSfield{Kind; source}{definition; \GSpage{P000138}}
\GSfield{The notes}{Reductive means the radical is a torus; semisimple means \(R(G)=1\).}
\GSfield{Statement}{For smooth connected affine groups over an algebraically closed field, reductive means \(R_u(G)=1\), equivalently that \(R(G)\) is a torus; semisimple means \(R(G)=1\). Over a general field these conditions are geometric, after base extension to an algebraic closure.}
\GSfield{Reconstruction}{The word torus means an algebraic torus, not an abelian variety. General-field forms need not be split.}
\GSfield{Fate}{This is background for the explicitly cited \citet{borel1965reductifs}, published in 1965; no later priority claim is made.}
\GSfield{Present status}{As of November 2026: established under the stated hypotheses; no unresolved problem is asserted by this occurrence.}
\GSfield{Assessment}{literal statement: correct with missing hypotheses; publication match: exact; verification: read the result; development: none; confidence: high}
\end{GSitem}
\begin{GSitem}{1965-85.12}{E}{GLn and SLn as the basic examples}
\GSfield{Kind; source}{construction; \GSpage{P000138}}
\GSfield{The notes}{\(GL(n,K)\) is reductive and \(SGL(n,K)\) is semisimple.}
\GSfield{Statement}{\(GL_n\) is reductive, with connected solvable radical its scalar torus; \(SL_n\) is semisimple for \(n\geq2\).}
\GSfield{Reconstruction}{The written \(SGL\) is interpreted as \(SL\); no separate standard group \(SGL\) is defined.}
\GSfield{Fate}{This is background for the explicitly cited \citet{borel1965reductifs}, published in 1965; no later priority claim is made.}
\GSfield{Present status}{As of November 2026: established under the stated hypotheses; no unresolved problem is asserted by this occurrence.}
\GSfield{Assessment}{literal statement: correct with missing hypotheses; publication match: exact; verification: read the result; development: annotated; confidence: high}
\end{GSitem}
\begin{GSitem}{1965-85.13}{E}{Root-system classification and its scope}
\GSfield{Kind; source}{theorem; \GSpage{P000138}}
\GSfield{The notes}{Four infinite series and five exceptional cases are listed, with maximal tori and root systems as the basic notions.}
\GSfield{Statement}{The irreducible reduced crystallographic root systems are \[A_n,B_n,C_n,D_n,E_6,E_7,E_8,F_4,G_2\], with the usual low-rank coincidences. Over an algebraically closed field, connected semisimple groups require a product of such types together with the isogeny/root-datum information; arbitrary-field forms require additional descent data.}
\GSfield{Reconstruction}{The nine-family list classifies simple root-system types, not literally every semisimple algebraic group up to isomorphism.}
\GSfield{Fate}{This is background for the explicitly cited \citet{borel1965reductifs}, published in 1965; no later priority claim is made.}
\GSfield{Present status}{As of November 2026: established under the stated hypotheses; no unresolved problem is asserted by this occurrence.}
\GSfield{Assessment}{literal statement: correct with missing hypotheses; publication match: exact; verification: read the result; development: none; confidence: high}
\end{GSitem}
\begin{GSitem}{1965-85.14}{E}{A real nonsplit torus}
\GSfield{Kind; source}{construction; \GSpage{P000138}}
\GSfield{The notes}{The real matrices \(\begin{pmatrix}a&b\\-b&a\end{pmatrix}\) are identified with nonzero complex numbers under multiplication.}
\GSfield{Statement}{With \(a^2+b^2\neq0\), this group is \(\operatorname{Res}_{\mathbb C/\mathbb R}\mathbb G_m\), a two-dimensional real algebraic torus. It becomes \(\mathbb G_m^2\) over \(\mathbb C\) but is not a split real torus.}
\GSfield{Reconstruction}{The invertibility condition is required. The example reinforces the need to distinguish geometric tori from split tori over the given field.}
\GSfield{Fate}{This is background for the explicitly cited \citet{borel1965reductifs}, published in 1965; no later priority claim is made.}
\GSfield{Present status}{As of November 2026: established under the stated hypotheses; no unresolved problem is asserted by this occurrence.}
\GSfield{Assessment}{literal statement: correct with missing hypotheses; publication match: exact; verification: read the result; development: none; confidence: high}
\end{GSitem}
\begin{GSitem}{1965-85.15}{E}{Conjugacy of maximal tori and the rational-field caveat}
\GSfield{Kind; source}{theorem; \GSpage{P000138}}
\GSfield{The notes}{All maximal tori are said to be conjugate.}
\GSfield{Statement}{Maximal tori of a connected linear algebraic group over an algebraically closed field are conjugate. For a connected reductive group over a general field, maximal split tori defined over that field are conjugate under its rational points. Arbitrary maximal real tori need not be conjugate under \(G(\mathbb R)\).}
\GSfield{Reconstruction}{\ours{} The explanatory verification here is supplied by us for the established result or correction. The nearby real example makes the distinction material. In \(SL_2(\mathbb R)\), split and compact maximal tori already show why the unqualified rational reading fails.}
\GSfield{Fate}{The general-field split-torus theorem is a central result of \citet{borel1965reductifs}; the algebraically closed result is its standard background.}
\GSfield{Present status}{As of November 2026: established under the stated hypotheses; no unresolved problem is asserted by this occurrence.}
\GSfield{Assessment}{literal statement: correct with missing hypotheses; publication match: exact; verification: read the result; development: none; confidence: high}
\end{GSitem}

\GStalk{86}{1966-03-21}{V. A. Borovikov; comment by I. M. Gelfand}{P000138}{Boundary integrals and area}
After the dated March heading, Borovikov states a double boundary cosine integral for area. Gelfand follows with a very compressed slice-area/volume formula. The former is valid with convexity and consistent inward chord-angle conventions; it has no missing distance denominator. The latter is false for a single slicing direction and remains uncertain as a historical assertion because an angular average may have been omitted.

\begin{GSitem}{1965-86.1}{E}{Double boundary cosine integral for planar area}
\GSfield{Kind; source}{theorem; \GSpage{P000138}}
\GSfield{The notes}{With the normals drawn at \(M,N\), \(\iint\cos\varphi_M\cos\varphi_N\,dM\,dN=2\pi S\), where \(S\) is area.}
\GSfield{Statement}{Let \(K\subset\mathbb R^2\) be a bounded strictly convex domain with \(C^2\) boundary. At each ordered pair of distinct boundary points, take the inward normal and the chord directed toward the other endpoint; let \(c_M,c_N\) be their nonnegative cosines. Then \(\int_{\partial K}\int_{\partial K}c_Mc_N\,ds_Mds_N=2\pi|K|\).}
\GSfield{Reconstruction}{\ours{} Convexity, arc-length measures and the two endpoint angle conventions are supplied. Chord-space coordinates give \(|dp\,d\theta|=c_Mc_N|M-N|^{-1}ds_Mds_N\); the pair integral is twice the integral of chord length over unoriented lines. A complete independent derivation is included in the audit notes. A second independent proof uses the divergence theorem and the flux of the logarithmic potential. The circle normalization check is an independent computation, \oursComputed{}.}
\GSfield{Fate}{This is a direct consequence of classical line-space integration: \citet{santalo1953introduction} predates the talk; the relevant measure is checked in \citet{santalo1989integral}, \S1.2. No exact Borovikov publication or priority claim was verified.}
\GSfield{Present status}{As of November 2026: the stated convex-domain identity is verified by the included derivation; attribution of this exact presentation remains unresolved.}
\GSfield{Assessment}{literal statement: correct with missing hypotheses; publication match: partial; verification: read the result; development: annotated; confidence: medium}
\end{GSitem}
\begin{GSitem}{1965-86.2}{F}{Gelfand slice-area formula and the missing directional average}
\GSfield{Kind; source}{claim; \GSpage{P000138}}
\GSfield{The notes}{The scan shows \(M(\int S^{\prime\,2}(x)\,dx)=V\) and explicitly calls \(M\) a universal constant. The prime reading is marked uncertain in the transcription.}
\GSfield{Statement}{If \(S(t)\) is the area of planar sections of a three-dimensional body in one fixed direction, no universal scalar \(M\) makes \(M\int|S^{\prime}(t)|^2dt=V\) for all bodies. For an ellipsoid with semiaxes \(a,b,c\), sliced perpendicular to the \(c\)-axis, the ratio of the integral to volume is \(2\pi ab/c^2\).}
\GSfield{Reconstruction}{\oursComputed{} The ellipsoid counterexample is independently computed here. \ours{} An angularly integrated Radon-Plancherel identity does have a universal constant: with ordinary area on \(S^2\), \(\int_{S^2}\int_{\mathbb R}|\partial_tR f(\omega,t)|^2dt\,d\omega=8\pi^2\|f\|_2^2\). Taking \(f=1_K\) gives volume. The notes do not record this angular integral, so it is a possible repair, not the recovered historical statement. The complete Fourier-slice and Plancherel derivation is included in the audit notes.}
\GSfield{Fate}{The nearby Gelfand integral-geometry literature supplies thematic context, but no exact match to this one-line formula with its stated constant was verified.}
\GSfield{Present status}{As of November 2026: the fixed-direction literal interpretation is disproved by the ellipsoid computation; the intended averaged statement remains textually uncertain.}
\GSfield{Assessment}{literal statement: false as written; publication match: thematic only; verification: unresolved lead; development: annotated; confidence: low}
\end{GSitem}

\GStalk{87}{1966-03-21}{V. B. Lidskii}{P000138}{Zeta function of all zeros of an exponential-type entire function}
This compact announcement has an unusually exact later match: Lidskii and Sadovnichii's 1967 paper, especially Lemmas 2--3. The entire continuation concerns a branch-dependent sum over all nonzero zeros, with complete sectorial asymptotics. The comparison with the ordinary zeta function refers to the difference between all signed sine zeros and only the positive half.

\begin{GSitem}{1965-87.1}{A}{The exponential-asymptotic class of entire functions}
\GSfield{Kind; source}{definition; \GSpage{P000138}}
\GSfield{The notes}{\(f(z)=\sum_{k=1}^Ne^{\alpha_kz}\mathcal P_k(z)\), where \(\mathcal P_k\) are polynomial-plus-inverse-power asymptotic series.}
\GSfield{Statement}{The matching class consists of entire functions represented sectorially by finitely many exponential factors with analytic coefficients admitting complete differentiable asymptotic expansions \(P_k(z)\sim z^{\mu_k}\sum_{j\geq0}p_{kj}z^{-j}\), \(\mu_k\in\mathbb Z\), \(p_{k0}\neq0\), on a finite covering by open sectors outside a disk.}
\GSfield{Reconstruction}{Complete asymptotics, sectorial analyticity and termwise differentiability are supplied from the published class \(K\). A formal divergent Laurent series by itself does not define an entire function.}
\GSfield{Fate}{Published soon after in \citet{lidskii1967zeros}; the class definition and Lemmas 2--3 were read in full text.}
\GSfield{Present status}{As of November 2026: established under the stated hypotheses; no unresolved problem is asserted by this occurrence.}
\GSfield{Assessment}{literal statement: correct with missing hypotheses; publication match: exact; verification: read the result; development: none; confidence: high}
\end{GSitem}
\begin{GSitem}{1965-87.2}{A}{All-zero zeta sum and its logarithm branch}
\GSfield{Kind; source}{construction; \GSpage{P000138}}
\GSfield{The notes}{\(\zeta(\sigma)=\sum z_n^{-\sigma}\), with \(z_n\) the zeros of \(f\).}
\GSfield{Statement}{Remove a possible power of \(z\) so that \(f(0)\neq0\), count all zeros with multiplicity, and choose a logarithm branch cut along a ray in a noncritical zero-free sector, avoiding the finite remaining zeros. Then \(Z(s)=\sum_n\exp(-s\operatorname{Log}z_n)\) initially converges for \(\operatorname{Re}s>1\).}
\GSfield{Reconstruction}{The branch and exclusion of zero are indispensable. Changing the admissible cut can change the zeta function even though the root set is unchanged.}
\GSfield{Fate}{\citet{lidskii1967zeros}, \S1, definition (9)--(10) and Lemma 2.}
\GSfield{Present status}{As of November 2026: established under the stated hypotheses; no unresolved problem is asserted by this occurrence.}
\GSfield{Assessment}{literal statement: correct with missing hypotheses; publication match: exact; verification: read the result; development: none; confidence: high}
\end{GSitem}
\begin{GSitem}{1965-87.3}{A}{Entire continuation by a Cauchy contour}
\GSfield{Kind; source}{theorem; \GSpage{P000138}}
\GSfield{The notes}{The zeta function is said to be always entire, with proof by a Cauchy integral.}
\GSfield{Statement}{For the stated complete-asymptotic class and an admissible branch, the all-zero zeta function extends to an entire function of \(s\). A keyhole integral of \(z^{-s}f^{\prime}(z)/f(z)\) first represents it in a right half-plane; subtraction of arbitrarily many asymptotic terms continues it globally.}
\GSfield{Reconstruction}{The unrestricted word always is valid only within the restored class. With a finite-order asymptotic expansion the paper guarantees only a corresponding half-plane continuation.}
\GSfield{Fate}{Exact result: \citet{lidskii1967zeros}, Lemma 3 and its footnote, Russian p. 55.}
\GSfield{Present status}{As of November 2026: established under the stated hypotheses; no unresolved problem is asserted by this occurrence.}
\GSfield{Assessment}{literal statement: correct with missing hypotheses; publication match: exact; verification: read the result; development: none; confidence: high}
\end{GSitem}
\begin{GSitem}{1965-87.4}{A}{Why this does not contradict the Riemann zeta pole}
\GSfield{Kind; source}{remark; \GSpage{P000138}}
\GSfield{The notes}{The notes say ordinary zeta does not use all the zeros of sine, with a reference to not starting at minus infinity.}
\GSfield{Statement}{For \(f(z)=\sin(\pi z)/z\), the nonzero zeros are \(\pm1,\pm2,\ldots\). With a branch assigning argument \(\pi\) to negative reals, the all-zero sum is \((1+e^{-i\pi s})\zeta_{\mathrm R}(s)\); the factor vanishes at \(s=1\), cancelling the sole pole. A sum only over positive zeros retains the pole.}
\GSfield{Reconstruction}{\oursComputed{} The explicit example and branch choice reconstruct the compressed sine remark; they are checked algebraically rather than quoted as a formula from the seminar.}
\GSfield{Fate}{Published soon after in \citet{lidskii1967zeros}; the class definition and Lemmas 2--3 were read in full text.}
\GSfield{Present status}{As of November 2026: established under the stated hypotheses; no unresolved problem is asserted by this occurrence.}
\GSfield{Assessment}{literal statement: correct with missing hypotheses; publication match: exact; verification: read the result; development: annotated; confidence: high}
\end{GSitem}

\GStalk{89}{1966-03-21}{S. G. Gindikin}{P000138-P000139}{Generalized parabolic equations and cone-valued ellipticity}
Gindikin's announcement precedes Minlos in the physical notebook, although the catalogue numbers it 89. The 1967 multidimensional-time paper identifies both the hyperbolic replacement of the time derivative and the \(V\)-elliptic mapping. A reinspection corrects the transcription's pseudo-operator to hyperbolic operator. The random-walk/heat sentence before the boxed Minlos heading is recorded here with attribution uncertainty.

\begin{GSitem}{1965-89.1}{A}{Hyperbolic time operator generalizing the heat equation}
\GSfield{Kind; source}{construction; \GSpage{P000138}}
\GSfield{The notes}{The supplied transcription says any pseudo-operator in place of \(\partial_t\); the scan reads \ru{любой гиперб. оператор} \(P(\partial/\partial t)\).}
\GSfield{Statement}{For a homogeneous polynomial \(p\) hyperbolic with respect to a direction in a proper cone \(V\subset\mathbb R^r\), a compatible \(V\)-elliptic spatial symbol produces a pluriparabolic operator with \(r\) time variables. In a consistent Fourier convention it has symbol \(q(\tau,\xi)=p(\tau+F(\xi))\). For one time variable and \(p(\tau)=\tau\), this recovers the ordinary parabolic construction.}
\GSfield{Reconstruction}{Hyperbolicity is recovered from the scan; homogeneity, the cone and spatial positivity are restored from the precise published theorem. No claim is made for an arbitrary pseudodifferential time operator.}
\GSfield{Fate}{\citet{gindikin1967multitime}, \S3, Theorem 4, submitted 27 April 1966; the matching construction was read in the original four-page paper.}
\GSfield{Present status}{As of November 2026: established under the stated hypotheses; no unresolved problem is asserted by this occurrence.}
\GSfield{Assessment}{literal statement: correct with missing hypotheses; publication match: exact; verification: read the result; development: annotated; confidence: high}
\end{GSitem}
\begin{GSitem}{1965-89.2}{A}{V-elliptic mappings and the hyperbolicity cone}
\GSfield{Kind; source}{definition; \GSpage{P000139}}
\GSfield{The notes}{The first line replaces ellipticity by \(V\)-ellipticity and refers to a cone corresponding to a hyperbolic polynomial.}
\GSfield{Statement}{In the Fourier-symbol convention where \(p(\tau+i\eta)\neq0\) for \(\eta\in V\), let \(F:\mathbb R^n\to\mathbb C^r\) have homogeneous polynomial components of common even degree. The required cone positivity is \(\operatorname{Im}F(\xi)\in V\) for real \(\xi\neq0\). Compactness of the unit sphere gives uniform interior positivity after normalization.}
\GSfield{Reconstruction}{Factors of \(i\) move positivity between real and imaginary parts when differential-operator conventions change. The condition is stated at the symbol level to avoid silently mixing the notebook convention with the article convention.}
\GSfield{Fate}{Exact terminology and construction in \citet{gindikin1967multitime}, \S3, immediately before Theorem 4.}
\GSfield{Present status}{As of November 2026: established under the stated hypotheses; no unresolved problem is asserted by this occurrence.}
\GSfield{Assessment}{literal statement: correct with missing hypotheses; publication match: exact; verification: read the result; development: none; confidence: high}
\end{GSitem}
\begin{GSitem}{1965-89.3}{A}{The associated family of operators is a cone semigroup}
\GSfield{Kind; source}{construction; \GSpage{P000138-P000139}}
\GSfield{The notes}{A family of operators is associated with \(P\), reducing to the parabolic family when \(P=\partial_t\).}
\GSfield{Statement}{Under the cone-positivity assumptions, commuting spatial Fourier multipliers generate a family \(T(t)\), \(t\) in the dual time cone, satisfying \(T(t)T(s)=T(t+s)\). Along a fixed interior ray the family is an ordinary parabolic evolution. The multidimensional-time fundamental solution combines this family with that of the hyperbolic time operator.}
\GSfield{Reconstruction}{The source does not state all generator-domain conditions. This is the constant-coefficient Fourier-multiplier realization, not a theorem for arbitrary unbounded commuting operators.}
\GSfield{Fate}{\citet{gindikin1967multitime}, Theorem 5 and \S4, develop the family; \citet{gindikin1967cauchy} treats the related strongly homogeneous setting. The variable-coefficient development is documented by \citet{volevich1969pluriparabolic}, introduction and Theorem 1.1.}
\GSfield{Present status}{As of November 2026: established under the stated hypotheses; no unresolved problem is asserted by this occurrence.}
\GSfield{Assessment}{literal statement: correct with missing hypotheses; publication match: partial; verification: read the result; development: none; confidence: high}
\end{GSitem}
\begin{GSitem}{1965-89.4}{E}{A random walk on a line and the heat equation}
\GSfield{Kind; source}{remark; \GSpage{P000139}}
\GSfield{The notes}{Before the boxed Minlos name, the source says \ru{Блуждание по прямой --- ур-е теплопроводности}.}
\GSfield{Statement}{A centered nearest-neighbor walk with step length \(\varepsilon\) and step time \(\varepsilon^2\), under diffusive scaling, converges to Brownian motion with variance rate one; its limiting density solves \(\partial_tu=\tfrac12\partial_x^2u\). Different variance/time normalizations change the diffusion constant.}
\GSfield{Reconstruction}{This is the minimal standard mathematical interpretation; the walk law and scaling are absent. Page position favors continuity with Gindikin, but the isolated sentence may be an interjection.}
\GSfield{Fate}{The diffusion-limit connection predates the seminar; \citet{donsker1951invariance} is a rigorous functional-limit reference, checked at bibliographic level.}
\GSfield{Present status}{As of November 2026: established under the stated hypotheses; no unresolved problem is asserted by this occurrence.}
\GSfield{Assessment}{literal statement: correct with missing hypotheses; publication match: partial; verification: publisher metadata checked; development: none; confidence: medium}
\end{GSitem}

\GStalk{88}{1966-03-21}{R. A. Minlos}{P000139}{Ising lattice gas and thermodynamic limit}
The boxed Minlos heading introduces a lattice-gas configuration sum and the limit of its logarithm per volume. This is a brief substantive reminder, not a content-free heading, but no low-temperature contour theorem or phase-transition statement is actually recorded. It recurs to the elementary material indexed as 1964-3.1 and 1964-3.2.

\begin{GSitem}{1965-88.1}{E}{Lattice configurations, occupied-neighbor energy and partition sum}
\GSfield{Kind; source}{definition; \GSpage{P000139}}
\GSfield{The notes}{\(L\) is a subset of the lattice, \(E(L)\) counts neighbors, \(N(L)\) counts points, and \(\Xi=\sum_Le^{-\beta E(L)-\mu N(L)}\).}
\GSfield{Statement}{For finite \(\Lambda\subset\mathbb Z^d\), a configuration is \(L\subseteq\Lambda\). With \(n_x=1_{x\in L}\), define \(N(L)=\sum_xn_x\) and, with a stated convention, \(E(L)=\sum_{\langle x,y\rangle}n_xn_y\) over unordered nearest-neighbor bonds. Then \(Z_\Lambda(\beta,\mu)=\sum_{L\subseteq\Lambda}e^{-\beta E(L)-\mu N(L)}\).}
\GSfield{Reconstruction}{The finite region and bond-counting convention are supplied. The displayed sign penalizes occupied-neighbor pairs for positive \(\beta\); calling it ferromagnetic would require changing the interaction sign. The uncertain parenthesis is not expanded.}
\GSfield{Fate}{An established lattice-gas formulation, recurring from 1964-3.1. \citet{ising1925ferromagnetismus} is the original model reference, checked bibliographically; no exact claim of equivalence to its notation is made.}
\GSfield{Present status}{As of November 2026: established under the stated hypotheses; no unresolved problem is asserted by this occurrence.}
\GSfield{Assessment}{literal statement: correct with missing hypotheses; publication match: partial; verification: publisher metadata checked; development: none; confidence: high}
\end{GSitem}
\begin{GSitem}{1965-88.2}{E}{Thermodynamic pressure and the normalization ambiguity}
\GSfield{Kind; source}{theorem; \GSpage{P000139}}
\GSfield{The notes}{\(\lim(\ln Z)/N=f(\theta,\beta)\) is written; the first argument and the role of \(N\) are unclear.}
\GSfield{Statement}{For a fixed finite-range translation-invariant interaction on a finite-state lattice, along a van Hove sequence \(\Lambda_j\), \(\lim_j|\Lambda_j|^{-1}\log Z_{\Lambda_j}(\beta,\mu)\) exists and is independent of standard boundary conditions. A canonical free energy instead fixes particle density and uses a different partition sum.}
\GSfield{Reconstruction}{In the grand-canonical sum above, \(N(L)\) varies across summands and cannot be the external volume denominator. The limit is reconstructed as pressure per site. No analyticity, coexistence or critical-temperature conclusion follows from this one line.}
\GSfield{Fate}{This recurs from 1964-3.2. Pre-seminar thermodynamic-limit results are documented by \citet{fisher1964freeenergy}; its publisher abstract and metadata, rather than its full proof, were checked.}
\GSfield{Present status}{As of November 2026: established under the stated hypotheses; no unresolved problem is asserted by this occurrence.}
\GSfield{Assessment}{literal statement: correct with missing hypotheses; publication match: partial; verification: publisher metadata checked; development: none; confidence: high}
\end{GSitem}

\section{Developments}\label{dev:1965}
\subsection{A bounded answer to the Weyl-field automorphism problem}
\label{dev:weyl}

\paragraph{Question, scope and evidence.}
Item 1965-80.10 asks, literally, \ru{Найти все автоморфизмы этого тела}.
The strongest formulation supported by its preceding notation is to determine
\(\operatorname{Aut}_{L}D_{n,k}\), once a characteristic-zero coefficient field
\(L\) and its fixedness are specified. The page supplies neither proposed
generators nor relations for that group. In particular, its parenthesis about
Fourier transformation and multiplication is a suggestion, not a generation
conjecture. The following work fixes
\[
 K=\mathbb C(x),\qquad R=K[p;d/dx],\qquad
 D=Q(R)=D_{1,0},\qquad [p,x]=1.
\]
It determines the subgroup preserving this particular \(K\), and distinguishes
its rational translations that are inner from those that are outer. All
automorphisms fix \(\mathbb C\) pointwise. These hypotheses are restrictions
made here; they are not attributed to the notebook.

The full polynomial-algebra automorphism theorem is different:
\citet[Theorem 2]{makarlimanov1984automorphisms} treats \(A_1\), not its entire
fraction field. The rational-subfield stabilizer below is already
\citet[Proposition 2.2(i)]{alevdumas1998invariants}; their Theorem 2.5 also gives
\(D^G\simeq D\) for finite groups stabilizing \(\mathbb C(p)\), \(\mathbb C(x)\)
or \(\mathbb C(px)\). We re-prove the stabilizer and derive an explicit
inner-translation criterion. These are bounded, self-contained partial
answers, with no claim of new literature results. The notation \ours\ means
that a proof is supplied in this companion, not that priority is claimed.

\citet[\S5.3.3]{dumas2011invariants} explicitly records the missing full
automorphism-group description. The current search, dated 9 November 2026,
found no later full classification. A recent related result,
\citet[Theorem 0.12]{schwarz2024generalizations}, proves rationality of invariant
Weyl fields for finite groups induced from rational function fields when the
commutative invariant field is rational. It does not classify arbitrary
automorphisms of \(D\). A still more recent preprint,
\citet[Corollary 8.12]{lauritzenthomsen2026duality}, version 2 dated
2 November 2026, proves equality and finiteness of the left and right
indices for every noncommutative division \(\mathbb C\)-subalgebra of
\(D_{1,0}\); commutative subalgebras have infinite indices. This concerns
extension dimensions, not a description of all automorphisms.
These distinctions identify the remaining mathematical obstruction,
but do not establish why particular researchers did or did not pursue the
1966 suggestion; no documentary claim of that kind is made.

\paragraph{A centralizer calculation.}
We use formal pseudodifferential series only as an algebraic device. Let
\(\mathcal P\) consist of sums \(\sum_{j\le m}a_j(x)p^j\), with
\(a_j\in K\), \(m\in\mathbb Z\), and multiplication determined by
\[
 p^j a=\sum_{r\ge0}\binom{j}{r}a^{(r)}p^{j-r}
 \quad(j\in\mathbb Z).
\]
Every coefficient of a product receives finitely many contributions. For
completeness, associativity reduces, after the Leibniz rule, to the identity
\[
 \binom{i}{r}\binom{i+j-r}{s}
 =\sum_{t=0}^{s}\binom{j}{t}\binom{i}{r+s-t}
                         \binom{r+s-t}{r}
 \quad(i,j\in\mathbb Z,\ r,s\ge0).
\]
Indeed, first take out \(\binom{i}{r}\) and then apply the generalized
Vandermonde identity to the remaining sum. The two sides are exactly the
coefficients of \(a b^{(r)}c^{(s)}p^{i+j+k-r-s}\) in the two bracketings
of \((ap^i)(bp^j)(cp^k)\). Coefficientwise finiteness extends this calculation
from monomials to all the stated series.
A nonzero series can be factored as \(a_mp^m(1+Q)\), where \(Q\) has order
at most \(-1\); the formal geometric series in \(-Q\) constructs its inverse.
Thus \(\mathcal P\) is a division ring containing \(R\), and the Ore
universal property embeds \(D\) into it. No analytic convergence is asserted.

\begin{lemma}[The rational centralizer \ours]
The centralizer of \(x\) in \(D\) is precisely \(K\).
\end{lemma}
\begin{proof}
The multiplication rule gives \([p^j,x]=jp^{j-1}\) for every integer \(j\),
including negative integers, since all derivatives of \(x\) above the first
vanish. If \(A=\sum_{j\le m}a_jp^j\) commutes with \(x\), uniqueness of its
series expansion gives \(ja_j=0\) for every \(j\ne0\). Characteristic zero
therefore gives \(A=a_0\in K\). Conversely, every element of \(K\) commutes
with \(x\).
\end{proof}

\begin{proposition}[The rational-subfield stabilizer \ours]
An automorphism \(\phi\in\operatorname{Aut}_{\mathbb C}D\) satisfies
\(\phi(K)=K\) if and only if
\[
 \phi(x)=f(x)=\frac{ax+b}{cx+d},\qquad
 \phi(p)=\frac{1}{f'(x)}p+h(x),
 \quad ad-bc\ne0,\quad h\in K.
\]
Scalar multiples of the matrix represent the same \(f\). Every displayed pair
defines an automorphism, including when \(c\ne0\).
\end{proposition}
\begin{proof}
The restriction to \(K\) is a rational function field automorphism. Its image
\(f\) has a rational compositional inverse. Degrees of nonconstant rational
maps of the projective line multiply under composition, so \(\deg f=1\):
\(f\) is the displayed fractional linear map. In particular \(f'\ne0\).
The defining relation gives \([\phi(p),f]=1\), while
\([f'^{-1}p,f]=1\). Their difference commutes with \(f\), hence with
\(\mathbb C(f)=K\). The lemma puts that difference in \(K\).

Conversely, substitution \(a(x)\mapsto a(f(x))\) is an automorphism of \(K\).
For \(P=f'^{-1}p+h\), the chain rule gives
\([P,a(f)]=a'(f)\); thus the prescription respects the Ore relation and
defines a map on \(R\). Its images generate \(R\), since \(K\) is already
in the image and \(p=f'(P-h)\). More explicitly, if \(u=f^{-1}\), its inverse
has \(x\mapsto u(x)\) and
\[
 p\mapsto f'(u(x))p-f'(u(x))h(u(x)).
\]
The chain rule verifies both compositions on the generators. The maps extend
to inverse automorphisms of \(D\).
\end{proof}

The equality \(\phi(K)=K\) is indispensable: the algebraic Fourier
automorphism \(x\mapsto p\), \(p\mapsto-x\) does not preserve \(K\).
Conversely, preserving \(K\) need not mean preserving \(A_1\):
\(x\mapsto x^{-1}\), \(p\mapsto-x^2p\) is an involution of \(D\) whose image
of \(x\) is not polynomial. This is a statement about the chosen polynomial
subalgebra, not about whether the involution is conjugate to another action.

\paragraph{Which rational translations are inner?}
Write \(T_h(x)=x\) and \(T_h(p)=p+h(x)\).
The maps \(T_h\) form a copy of the additive group of \(K\).

\begin{proposition}[Exact rational-pole criterion \ours]
For \(h\in\mathbb C(x)\), the following are equivalent:
\begin{enumerate}[label=(\roman*)]
\item \(T_h\) is an inner automorphism of \(D\).
\item \(h=-u'/u\) for some \(u\in\mathbb C(x)^{\times}\).
\item The polynomial part of \(h\) is zero, every finite pole is simple,
and every finite-pole residue is an integer.
\end{enumerate}
If \(h=\sum_a r_a/(x-a)\) has these properties, one implementing element is
\(u=\prod_a(x-a)^{-r_a}\).
\end{proposition}
\begin{proof}
If \(T_h=\operatorname{Ad}_u\) for some \(u\in D^{\times}\), its fixing
\(x\) implies \(ux=xu\). The centralizer lemma forces \(u\in K^{\times}\).
The Ore relation yields
\[
 upu^{-1}=p-\frac{u'}u.
\]
This proves (i) implies (ii), and the same identity proves the converse.
Factor a nonzero rational function over \(\mathbb C\) as
\(u=C\prod_a(x-a)^{m_a}\), with integral \(m_a\) and finitely many nonzero
exponents. Then \(-u'/u=-\sum_a m_a/(x-a)\), proving (ii) implies (iii).
Conversely, unique partial fraction decomposition and (iii) give exactly
\(h=\sum_a r_a/(x-a)\), \(r_a\in\mathbb Z\); the displayed product has
\(-u'/u=h\). No extra condition at infinity is missing: the differential
\(h(x)\,dx\) has residue there equal to the negative sum of the finite
residues.
\end{proof}

\begin{corollary}[Translation classes in the outer automorphism group \ours]
The image of \(\{T_h:h\in K\}\) in
\(\operatorname{Out}_{\mathbb C}D=\operatorname{Aut}_{\mathbb C}D/\operatorname{Inn}D\)
is, as an abelian group,
\[
 K\Big/\left\{-\frac{u'}u:u\in K^{\times}\right\}
 \simeq
 \mathbb C[x]\ \oplus\!\bigoplus_{a\in\mathbb C,\ j\ge2}\!\mathbb C
 \ \oplus\!\bigoplus_{a\in\mathbb C}\!(\mathbb C/\mathbb Z).
\]
All direct-sum elements have finite support. The first term records the
polynomial part, the second higher-order principal parts, and the third
simple-pole residues modulo integers.
\end{corollary}
\begin{proof}
The kernel of the homomorphism \(h\mapsto[T_h]\) is exactly the subgroup in
the denominator by the proposition. Partial fractions give an additive direct
sum decomposition of \(K\). The denominator is precisely the integral lattice
in its simple-pole summands, so quotienting gives the formula. This is an
abelian-group quotient; the \(\mathbb C/\mathbb Z\) terms are not complex
vector spaces.
\end{proof}

\begin{example}[Outer torsion and actual infinite order \ours]
For an integer \(m\ge2\), \(T_{1/(mx)}\) is not inner, but
\(T_{1/(mx)}^m=T_{1/x}=\operatorname{Ad}_{x^{-1}}\). Its class in the outer
group has exact order \(m\), because \(r/m\notin\mathbb Z\) for
\(1\le r<m\). The automorphism itself has infinite order:
\(T_{1/(mx)}^r=T_{r/(mx)}\ne\mathrm{id}\) for every positive \(r\).
Thus statements about finite order in \(\operatorname{Out}D\) must not be
substituted for statements about finite order in \(\operatorname{Aut}D\).
\end{example}

\paragraph{Finite subgroups of this stabilizer.}
For \(\sigma\in\operatorname{Aut}_{\mathbb C}K\) put \(f=\sigma(x)\) and
let \(L_\sigma\) be the lift with \(p\mapsto f'^{-1}p\).
Then
\[
 L_\sigma T_b L_\sigma^{-1}=T_{R_\sigma b},\qquad
 R_\sigma b=f'(x)b(f(x)).
\]
With automorphisms composed as maps of algebras, \(f_{\sigma\tau}=
f_\tau\circ f_\sigma\), and the chain rule gives
\(R_{\sigma\tau}=R_\sigma R_\tau\). Thus every stabilizer element has a
unique form \(T_bL_\sigma\), with multiplication
\[
 (T_bL_\sigma)(T_cL_\tau)=T_{b+R_\sigma c}L_{\sigma\tau}.
\]

\begin{proposition}[Removing a finite-group translation cocycle \ours]
Every finite subgroup \(H\) of the stabilizer of \(K\) is conjugate, by a
rational translation \(T_a\), to the canonical lift of its image in
\(\operatorname{Aut}_{\mathbb C}K\).
\end{proposition}
\begin{proof}
The kernel of restriction to \(K\) is the torsion-free additive group
\(\{T_b\}\), so restriction is injective on \(H\). Write its image as
\(Q\), and its elements uniquely as \(T_{b_\sigma}L_\sigma\),
\(\sigma\in Q\). The multiplication law gives
\(b_{\sigma\tau}=b_\sigma+R_\sigma b_\tau\).
Let \(a=|Q|^{-1}\sum_{\tau\in Q}b_\tau\). Sum this identity over \(\tau\)
to obtain \(b_\sigma=a-R_\sigma a\). Hence
\[
 T_{-a}(T_{b_\sigma}L_\sigma)T_a=L_\sigma
\]
for every \(\sigma\). Averaging uses characteristic zero; the assertion is
not made in positive characteristic.
\end{proof}

This proves a normal form for finite actions within the chosen stabilizer.
It says nothing about whether an arbitrary finite subgroup of
\(\operatorname{Aut}D\) stabilizes a conjugate rational subfield. Nor does
it make the conjugating translation inner: the pole criterion decides that
additional question separately.

\paragraph{Computations and independent checks.}
The accompanying \texttt{code/check\_weyl.py} implements multiplication of
finite differential operators from the Leibniz rule, with exact rational
coefficients. It checks the commutator, inverse and composition identities
above for fractional linear changes including inversion, higher-order poles,
nonintegral residues and polynomial shifts. A separate evaluation on
monomials checks the operator commutator without using the multiplication
routine. The same script enumerates PBW exponent tuples by a recurrence and
compares them with \(\binom{N+2n+k}{2n+k}\), including zero generators and
\(N=0\). The complete output is retained; finite tests are used to detect
algebraic or coding mistakes, not to establish the universally quantified
results.

\begin{example}[Rational shifts checked by symbolic differentiation \oursComputed]
The exact examples give
\[
\begin{array}{c|c|c}
h(x)&\text{inner?}&\text{implementer, when it exists}\\\hline
3/(x-1)-2/(x+2)&\text{yes}&(x-1)^{-3}(x+2)^2\\
1/(2x)&\text{no}&\text{its square is implemented by }x^{-1}\\
1/x^2&\text{no}&\text{higher-order pole}\\
x&\text{no}&\text{nonzero polynomial part}\\
0&\text{yes}&1
\end{array}
\]
Each positive case is checked by differentiating the implementer; negative
cases are excluded by the proved criterion, not by a failed search.
\end{example}

The second proof check re-derives the conjugation sign from
\(pu=up+u'\), and the composition law from substitution before applying the
abstract semidirect-product notation. Deliberately reversing the sign of
\(u'/u\), or omitting \(f'\) in the group action, fails the symbolic checks.
The proof also distinguishes the fixed subfield condition from the weaker
condition that an image is rational in some other generator. A separate
review of the written proof and a second parameter run are documented in
\texttt{evidence/development\_review.md} and the computation output.

\paragraph{The exact obstruction and a further bounded run.}
The unresolved step is global: an arbitrary automorphism can move
\(\mathbb C(x)\) to a different commutative subfield of \(D\). The
centralizer argument classifies maps only after the source and target
rational subfield have been fixed. Fourier interchange already escapes that
stabilizer. No reduction of every Weyl generating pair to the normal form
above has been proved here, and no generation theorem is conjectured on the
strength of the computations.

A further bounded run could explicitly compare the stabilizers of
\(\mathbb C(x)\) and \(\mathbb C(xp)\), using the second normal form in
\citet[Proposition 2.2(ii)]{alevdumas1998invariants}, and compute innerness
and conjugacy obstructions for rational functions with a fixed finite pole
set. That would refine the notebook's concrete suggestions while remaining
far short of classifying the full automorphism group. The present contribution
is the explicit, checked restricted answer and the precise location of the
remaining gap.

\section{Problem index}\label{problems:1965}
Earlier-year occurrence identifiers are retained as text; current-year identifiers link to their entries.

\subsection{GP-1965-01: Solving d-bar globally to construct a holomorphic boundary singularity}\label{prob:GP-1965-01}
For a bounded \(C^\infty\) strongly pseudoconvex domain \(\Omega\subset\mathbb C^n\), solve \(\bar\partial u=\alpha\) with smooth dependence and regularity up to the boundary for every smooth closed \((0,1)\)-form, and use this to construct, at each boundary point, a holomorphic function that does not extend across that point.
\par\noindent\textit{Occurrences:} \hyperref[item:1965-53.1]{1965-53.1}, \hyperref[item:1965-53.4]{1965-53.4}.
\par\noindent\textit{Status (2026-11):} As of 9 November 2026, this precise strongly pseudoconvex version was already solved before the seminar \citep{kohn1963potential,kohn1963harmonic,kohn1964harmonic}. No claim of a new unresolved Levi problem is made.

\subsection{GP-1965-02: Multiplication of distributions conormal to transverse submanifolds}\label{prob:GP-1965-02}
For transverse smooth submanifolds \(S,T\subset X\), determine a sufficient class of distributions singular on them whose product is canonically defined, including \(P\delta_S\) and \(Q\delta_T\) for properly supported classical pseudodifferential operators. Mere information about singular supports is not a sufficient formulation.
\par\noindent\textit{Occurrences:} \hyperref[item:1965-69.6]{1965-69.6}, \hyperref[item:1965-69.7]{1965-69.7}.
\par\noindent\textit{Status (2026-11):} As of 9 November 2026, settled for the stated conormal class by the wavefront multiplication criterion. H\"ormander's original Theorem 2.5.10, pp. 126--127, is read; pseudodifferential operators do not enlarge wavefront sets, and transverse conormal bundles have no opposite nonzero covectors. Bott's flat-trace theorem is the contemporary graph/diagonal case. The later terminology is a reconstruction, not a quotation from the notes. \citep{bott1966report,hormander1971fourier}

\subsection{GP-1965-03: The folds and cusps of z squared plus epsilon conjugate z}\label{prob:GP-1965-03}
Determine the real singularities and critical-value curve of \(F_\varepsilon(z)=z^2+\varepsilon\overline z\), \(\varepsilon\ne0\).
\par\noindent\textit{Occurrences:} \hyperref[item:1965-72.5]{1965-72.5}.
\par\noindent\textit{Status (2026-11):} As of 9 November 2026, solved: the critical locus is \(|z|=|\varepsilon|/2\), with three ordinary cusp points and folds at all other critical points; its image is a three-cusped hypocycloid. The same example is published in Arnold's 1968 exposition, Example 6. The determinant, caustic and ordinary-cusp tests are independently checked here. \citep{arnold1968singularities}

\subsection{GP-1965-04: Algebraic realization of finite curvature jets}\label{prob:GP-1965-04}
For a fixed dimension and order \(k\), characterize the pointwise arrays \((R,\nabla R,\ldots,\nabla^kR)\) realized by a Riemannian metric with prescribed zeroth-order inner product. The notebook asks for the algebraic relations among these tensors.
\par\noindent\textit{Occurrences:} \hyperref[item:1965-73.5]{1965-73.5}.
\par\noindent\textit{Status (2026-11):} As of 9 November 2026, settled in the finite-order realization formulation by the classical metric-jet isomorphism theorem. The exact modern algebraic-jet conditions and sufficiency are read in Jentsch, Definition 1 and Theorem 3; this checked 2026 revision is not the date of the theorem's discovery. Pre-1965 normal-tensor and curvature-reduction theory is documented through the Thomas/Veblen/Schouten tradition, with the earliest complete modern identity list not independently dated here. No minimal invariant-ring-presentation problem is inferred. \citep{jentsch2026jet,janyska2011reduction,thomas1926identities}

\subsection{GP-1965-05: Intersections of the zero torus with its Hamiltonian image}\label{prob:GP-1965-05}
For the zero section \(T^n\subset T^*T^n\), bound the number of intersections with its image under a compactly supported Hamiltonian diffeomorphism. The transverse formulation asks for at least \(2^n\) intersections; in the possibly degenerate formulation the cup-length lower bound is \(n+1\). The notebook writes only exact periods and an unqualified \(2^n\), so these additional hypotheses are explicitly reconstructed.
\par\noindent\textit{Occurrences:} \hyperref[item:1965-74.14]{1965-74.14}, \hyperref[item:1965-74.15]{1965-74.15}.
\par\noindent\textit{Status (2026-11):} As of 9 November 2026, both stated Hamiltonian bounds are established. The unqualified \(2^n\) version is false: a graph of the differential of \(\sin x+\sin y-\sin(x+y)\) on \(T^2\) has three intersections. No open general Arnold-conjecture variant is inferred from this record. \citep{chaperon1983questions,auroux2026floer,gong2023cuplength}

\subsection{GP-1965-06: Distinguishing Weyl rank and central parameters}\label{prob:GP-1965-06}
Can \(R_{n,k}\) or \(D_{n,k}\), over a fixed characteristic-zero field, be isomorphic for different pairs \((n,k)\)? The earlier occurrence 1964-38.1 concerns the ring version without additional central parameters; the present occurrence states the stronger division-ring answer.
\par\noindent\textit{Occurrences:} 1964-38.1, \hyperref[item:1965-80.6]{1965-80.6}.
\par\noindent\textit{Status (2026-11):} As of 9 November 2026: solved in the negative, exactly as formulated over the fixed base field, by \citet[\S4]{gelfandkirillov1966corps}. This is independent of the fixed-rank Dixmier endomorphism problem. No earlier GP identifier was supplied in the previous index.

\subsection{GP-1965-07: Universal determinant for a division ring}\label{prob:GP-1965-07}
Determine the universal abelian quotient of \(\mathrm{GL}_{\infty}(D)\) for a division ring \(D\).
\par\noindent\textit{Occurrences:} \hyperref[item:1965-80.9]{1965-80.9}.
\par\noindent\textit{Status (2026-11):} As of 9 November 2026: already solved before the talk. The Dieudonn\'e determinant identifies this quotient with \(D^{\times}_{\mathrm{ab}}\) \citep{dieudonne1943determinants}.

\subsection{GP-1965-08: Automorphisms of Weyl division rings}\label{prob:GP-1965-08}
Determine \(\operatorname{Aut}_{L}D_{n,k}\). The developed subproblem fixes \(L=\mathbb C\), \(n=1\), \(k=0\), and a rational subfield \(\mathbb C(x)\).
\par\noindent\textit{Occurrences:} \hyperref[item:1965-80.10]{1965-80.10}, \hyperref[item:1965-80.11]{1965-80.11}.
\par\noindent\textit{Status (2026-11):} As of 9 November 2026: partial, with no full classification found in the documented searches. The rational-subfield stabilizer is known \citep{alevdumas1998invariants}; polynomial automorphism classifications and results on finite-group invariant fields \citep{makarlimanov1984automorphisms,schwarz2024generalizations} do not determine the full group. This chapter proves the restricted normal form, finite-group splitting, and an exact criterion for which rational translations are inner; no novelty claim is attached to these elementary consequences.

\subsection{GP-1965-09: Range and cycle inversion for the affine k-plane transform}\label{prob:GP-1965-09}
Characterize the rapidly decreasing affine \(k\)-plane transform data by differential conditions, and recover the original function by integration over a suitable cycle of directions; specify which cycle classes give a nonzero coefficient.
\par\noindent\textit{Occurrences:} \hyperref[item:1965-82.1]{1965-82.1}, \hyperref[item:1965-82.5]{1965-82.5}, \hyperref[item:1965-82.6]{1965-82.6}.
\par\noindent\textit{Status (2026-11):} As of November 2026: the precise range theorem for \(1\leq k<n-1\) and the complex cycle inversion theorem are established in \citet{gelfand1967planes}. The notebook does not specify an additional open family-selection problem.

\section{Corrections to the transcription and translation}\label{corr:1965}
The evidence distinguishes a transcription correction from a mathematical emendation of wording faithfully preserved in the notebook.
\subsection*{P000088--P000089 (both): Record attribution}
\noindent\textit{Transcription reads.} Included in the attached catalogue range for record 53.\par
\noindent\textit{Critical reading.} Carry-over from the invariant-subspace discussion already indexed as 1964-50.11--24; Kohn starts on P000090.\par
\noindent\textit{Evidence.} P000090 has the new date \(20/\mathrm{IX}\) and Kohn heading. The preceding two pages match the supplied earlier occurrence titles. The catalogue ranges overlap adjacent records.\par
\subsection*{P000090 (both): display following the compatibility equations}
\noindent\textit{Transcription reads.} \(\sum u_{z_j}\,d\bar z_j\)\par
\noindent\textit{Critical reading.} \(\bar\partial u=\sum u_{\bar z_j}\,d\bar z_j\)\par
\noindent\textit{Evidence.} The scan itself omits the bar in this display, although the preceding derivative definition has it. This is a mathematical emendation of a notebook slip, not a demonstrated OCR error.\par
\subsection*{P000090 (both): last sentence of the cut-off argument}
\noindent\textit{Transcription reads.} u does not continue holomorphically\par
\noindent\textit{Critical reading.} h does not continue holomorphically across the chosen boundary point\par
\noindent\textit{Evidence.} The scan reads u, while the immediately preceding construction defines h=u-g/W. The singularity is in h; preserve u diplomatically and annotate the mathematical correction.\par
\subsection*{P000093 (both): de Rham product-rule display}
\noindent\textit{Transcription reads.} \(d(s f)(x_0)=s(x_0)\wedge df_{x_0}\) in positive form degree\par
\noindent\textit{Critical reading.} for a \(q\)-form \(s\), \(d(fs)(x_0)=df_{x_0}\wedge s(x_0)=(-1)^q s(x_0)\wedge df_{x_0}\) when \(f(x_0)=0\)\par
\noindent\textit{Evidence.} The scan genuinely puts the old form before df with no degree sign. The standard graded product rule supplies the sign. Symbol exactness is unchanged, so this is a convention-sensitive notebook correction, not a change in the ellipticity conclusion.\par
\subsection*{P000094 (both): fibre map in the geometric-lift construction}
\noindent\textit{Transcription reads.} \(E_{i,f^{-1}x}\to E_{i,x}\) beside a map \(f:X\to X\)\par
\noindent\textit{Critical reading.} Choose consistently either \(\phi:f^*E\to E,\ \phi_x:E_{f(x)}\to E_x\), or an inverse-map convention requiring invertibility.\par
\noindent\textit{Evidence.} The inverse superscript is visible in the scan. This is a convention mismatch in the notebook, not an established transcription error; there is no single diplomatic substitution valid for every later formula.\par
\subsection*{P000095 (both): spectral description}
\noindent\textit{Transcription reads.} positive eigenvalues, with zero eigenspace then used\par
\noindent\textit{Critical reading.} nonnegative spectrum; strictly positive eigenvalues considered separately from zero\par
\noindent\textit{Evidence.} The handwritten positive wording is genuine. The immediately following harmonic discussion and the identity \((\Box u,u)=\|du\|^2+\|d^*u\|^2\) require nonnegative.\par
\subsection*{P000096 (both): coordinate-ring description}
\noindent\textit{Transcription reads.} quotient by the ideal spanned by the displayed polynomials\par
\noindent\textit{Critical reading.} for the classical zero set, quotient by the radical of that ideal; alternatively specify a possibly nonreduced scheme\par
\noindent\textit{Evidence.} The scan confirms the ideal-of-equations wording. This is a missing mathematical qualification: \(x^2=0\) supplies an immediate counterexample to the unrestricted classical identification.\par
\subsection*{P000099 (both): aside after the closed-diagonal criterion}
\noindent\textit{Transcription reads.} a presheaf is called a sheaf if the diagonal is closed\par
\noindent\textit{Critical reading.} retain as a notebook error; the closed algebraic diagonal characterizes separation of the prevariety, whereas the sheaf axiom is unique gluing\par
\noindent\textit{Evidence.} Both transcription and scan genuinely have the presheaf/sheaf words. No transcription repair is justified.\par
\subsection*{P000100 (both): Chern-class heading and final map}
\noindent\textit{Transcription reads.} the map \(K(X)\to A(X)\) is linear\par
\noindent\textit{Critical reading.} leave the map unidentified; an additive Chern character needs rational coefficients, and total Chern classes obey a multiplicative Whitney formula\par
\noindent\textit{Evidence.} The linearity word and cancelled first axiom are visible in the scan. This is an unresolved mathematical identification, not evidence that the scribe meant one uniquely recoverable formula.\par
\subsection*{P000102 (both): Dynin's boxed identification}
\noindent\textit{Transcription reads.} \(H^i(\Gamma_\lambda(E))=\ker(\Box-\lambda)\)\par
\noindent\textit{Critical reading.} for \(\lambda>0\), \(H^i(\Gamma_\lambda(E))=0\), proved by \(h=d^*/\lambda\); the displayed kernel is the eigenspace, not its cohomology\par
\noindent\textit{Evidence.} The boxed equality is genuinely present in the scan. It is a notebook mathematical error or incomplete failed argument, not a mistranscribed valid use of the Hodge theorem.\par
\subsection*{P000103 (both): circle example}
\noindent\textit{Transcription reads.} \(\Box=d^2/dt^2\)\par
\noindent\textit{Critical reading.} \(\Box=-d^2/dt^2\) for \(d=d/dt\) and the standard periodic \(L^2\) metric\par
\noindent\textit{Evidence.} The positive sign is visible in the scan. The same notes define \(\Box=dd^*+d^*d\), whereas integration by parts gives \((d/dt)^*=-d/dt\); this is a notebook mathematical correction.\par
\subsection*{P000104 (both): second displayed spectral cancellation}
\noindent\textit{Transcription reads.} \((-1)^i\operatorname{tr}(T_\lambda|E_i)=0\), \(\lambda\ne0\), without a sum\par
\noindent\textit{Critical reading.} \(\sum_i(-1)^i\operatorname{tr}(T_\lambda|E_i)=0\)\par
\noindent\textit{Evidence.} The scan has no summation sign on this line. Acyclicity cancels alternating traces over all degrees; it does not force each chain-space trace to vanish.\par
\subsection*{P000105 (both): flat-trace description at the bottom of Gelfand's talk}
\noindent\textit{Transcription reads.} \(\int\delta(x,f(x))\,dx\) is the number of fixed points\par
\noindent\textit{Critical reading.} \(\sum_{p=f(p)}|\det(I-df_p)|^{-1}\), assuming nondegenerate fixed points\par
\noindent\textit{Evidence.} The unweighted description is present in the scan; the multivariable delta change-of-variables formula and the explicit weighted formula on P000106 require the determinant. This is a notebook mathematical omission.\par
\subsection*{P000105 (both): individual regularized trace series}
\noindent\textit{Transcription reads.} \(\sum_\lambda(1+\lambda)^{-z}\operatorname{tr}(T_\lambda^*|E_i)\)\par
\noindent\textit{Critical reading.} use the chain-space compression \(T_\lambda\) for the individual trace; the cohomological replacement applies to the alternating sum\par
\noindent\textit{Evidence.} The star is visible but its role is unstable. Earlier T* denotes induced cohomology. If read that way here, every positive-eigenvalue contribution would disappear degreewise and would not regularize the original chain operator. Retain the mark diplomatically and explain the type distinction.\par
\subsection*{P000106 (both): definition of K0}
\noindent\textit{Transcription reads.} projective modules, with no finiteness qualification\par
\noindent\textit{Critical reading.} isomorphism classes of finitely generated projective modules\par
\noindent\textit{Evidence.} The finiteness word is absent in the scan. The standard K0 definition needs it; otherwise countable direct sums give the displayed Eilenberg-swindle collapse. This is an omitted hypothesis in the notes.\par
\subsection*{P000106 (en): Smirnov's parenthesis}
\noindent\textit{Transcription reads.} cont. representations\par
\noindent\textit{Critical reading.} retain \ru{непр. предст.}; continuous representations / irreducible representations are both possible expansions\par
\noindent\textit{Evidence.} The scan has only the abbreviated Russian phrase. The English expansion selects one meaning without sufficient evidence; the nearby extreme-point question makes irreducibility a live alternative. The later central-character paper does not uniquely disambiguate the notebook.\par
\subsection*{P000106 (both): operator, kernel, and local-trace formula in record 65}
\noindent\textit{Transcription reads.} \(\phi(x)s(f^{-1}(x))\), then \(\phi(x)\delta(y-f(x))\), then denominator \(|1-J_\rho(f)|\)\par
\noindent\textit{Critical reading.} make the same choice of \(f\) or \(f^{-1}\) in both operator and kernel, and write \(|\det_{\mathbb R}(I-df_\rho)|\) for the forward convention\par
\noindent\textit{Evidence.} All three notations are visible in the scan. The inverse-map mismatch is genuine; the missing determinant notation is shorthand at best. This is an annotated mathematical normalization, not a unique claim about the scribe's intended convention.\par
\subsection*{P000107 (both): de Rham local numerator}
\noindent\textit{Transcription reads.} \(\operatorname{tr}\varphi_i(\rho)=\operatorname{det}-1\)\par
\noindent\textit{Critical reading.} \(\sum_i(-1)^i\operatorname{tr}(\Lambda^i(df_\rho)^*)=\det(I-df_\rho)\)\par
\noindent\textit{Evidence.} The scan leaves the determinant argument incomplete and the alternating sum is absent in this line. The following eigenvalue-products line and the final signs determine the standard supertrace identity, not an individual-degree equality.\par
\subsection*{P000107 (both): conjugate-average formula}
\noindent\textit{Transcription reads.} \(2\operatorname{Re}(1/(1-J_f))\) followed by an uncertain factor \(1/[?]\) times the conjugate sum\par
\noindent\textit{Critical reading.} retain the uncertain handwritten denominator; distinguish \(\operatorname{Re}z=(z+\overline z)/2\) from the realified trace \(z+\overline z\)\par
\noindent\textit{Evidence.} The active numerator 2 is clear, but the later prefactor denominator remains visually uncertain. The scan supports an ambiguity report, not a confident one-character transcription correction.\par
\subsection*{P000107 (both): last visible arrow in the long exact sequence}
\noindent\textit{Transcription reads.} \(F^1(B)\xrightarrow{\delta}\cdots\)\par
\noindent\textit{Critical reading.} \(R^1F(B)\to R^1F(C)\xrightarrow{\delta^1}R^2F(A)\)\par
\noindent\textit{Evidence.} The scan visibly places delta over the arrow immediately after \(F^1(B)\), at the edge of the image. Its next term is not visible. The correction gives the standard continuation and locates the genuine connecting arrow, without claiming the missing term was read.\par
\subsection*{P000108 (both): naturality square}
\noindent\textit{Transcription reads.} upper-left \(F^0(A)\); vertical arrows point upward\par
\noindent\textit{Critical reading.} upper-left \(F^0(C)\); for the displayed downward morphism of short exact sequences, both induced vertical arrows point downward\par
\noindent\textit{Evidence.} The two sloping strokes and crossbar of handwritten A are clear, and the arrowheads are genuinely upward. The correction is mathematical: covariant functoriality and the source of the connecting map determine the proper square.\par
\subsection*{P000113 (en): Arnold opening discussion}
\noindent\textit{Transcription reads.} Study the roughness[?] of singularities\par
\noindent\textit{Critical reading.} Study the structural stability [reading of the Russian word uncertain] of singularities\par
\noindent\textit{Evidence.} Russian grubost is the standard structural-stability term; the following parenthesis and equivalence diagram specify persistence under small perturbations. Preserve the source-reading uncertainty, but avoid literal roughness as the mathematical translation.\par
\subsection*{P000124 (both): Growth dimension; notebook error}
\noindent\textit{Transcription reads.} \(\limsup d(\alpha,N)/\ln N\)\par
\noindent\textit{Critical reading.} For the mathematical definition: \(\limsup\log d(\alpha,N)/\log N\).\par
\noindent\textit{Evidence.} The scan really omits the numerator logarithm; both transcriptions preserve that error. Formula (6) in \citet{gelfandkirillov1966corps} includes it. This is an annotation, not a diplomatic replacement.\par
\subsection*{P000124 (both): Ore proof; omitted argument}
\noindent\textit{Transcription reads.} \(x_0a=(\quad)b\)\par
\noindent\textit{Critical reading.} Choose the first nonzero \(x_j\), cancel \(b^j\), then factor the remaining right-hand side by \(b\).\par
\noindent\textit{Evidence.} The blank is visible in the scan. The \(x_0\ne0\) step cannot be assumed; the full proof in Lemma 2 of \citet{gelfandkirillov1966corps} supplies it. Notebook incompleteness, not a mistranscription.\par
\subsection*{P000124 (both): Ore fractions; handedness}
\noindent\textit{Transcription reads.} \(xa=yb\), then fractions \(ab^{-1}\).\par
\noindent\textit{Critical reading.} Left Ore corresponds to left fractions; use right Ore for the displayed right fractions, or assume both.\par
\noindent\textit{Evidence.} The two displays are legible in the scan. The rings used later are two-sided Noetherian domains, so the mathematics can be repaired without changing the ink.\par
\subsection*{P000138 (both): Gindikin, penultimate line}
\noindent\textit{Transcription reads.} Russian: любой псевдооператор; English: any pseudo-operator\par
\noindent\textit{Critical reading.} Russian: любой гиперб. оператор; English: any hyperbolic operator\par
\noindent\textit{Evidence.} The scan reads гиперб., with the same abbreviation repeated at the top of P000139. Independently rechecked by the parent reviewer. Gindikin 1967, Theorem 4, corroborates the reading but does not replace the visual evidence.\par
